% main.tex --- The Maya Book
%
% Build:   make            (in this directory; uses pdflatex twice)
% Output:  maya-book.pdf
%
% Chapter sources live under chapters/. To rearrange, edit \include lines.

\documentclass[11pt,openany,oneside]{book}

% preamble.tex --- shared setup for the maya book.
% Targets pdflatex (no fontspec / unicode-math). Safe on every TeX Live.

\usepackage[utf8]{inputenc}
\usepackage[T1]{fontenc}
\usepackage{lmodern}
\usepackage{microtype}

\usepackage[a4paper,margin=1in,headheight=14pt]{geometry}
\usepackage{xcolor}
\usepackage{graphicx}
\usepackage{amsmath,amssymb,amsthm}
\usepackage{enumitem}
\usepackage{hyperref}
\usepackage{fancyhdr}
\usepackage{titlesec}
\usepackage{tcolorbox}
\usepackage{listings}
\usepackage{caption}

\tcbuselibrary{breakable,skins}

% ---- Colour palette (terminal-themed) -----------------------------------
\definecolor{mayablue}{HTML}{2E5C8A}
\definecolor{mayalavender}{HTML}{6E5A8A}
\definecolor{mayagreen}{HTML}{3F7B5A}
\definecolor{mayarose}{HTML}{8A3F5A}
\definecolor{mayaink}{HTML}{1A1A22}
\definecolor{mayagrey}{HTML}{5E5E68}
\definecolor{mayalight}{HTML}{F2F2F4}
\definecolor{mayaamber}{HTML}{B07A1A}

% ---- Hyperref -----------------------------------------------------------
\hypersetup{
  colorlinks=true,
  linkcolor=mayablue,
  citecolor=mayablue,
  urlcolor=mayalavender,
  pdftitle={The Maya Book},
  pdfauthor={agentty docs},
  pdfsubject={A complete tour of the maya TUI library},
  bookmarksnumbered=true,
  bookmarksopen=true
}

% ---- Headers / footers --------------------------------------------------
\pagestyle{fancy}
\fancyhf{}
\fancyhead[LE]{\small\textit{\nouppercase{\leftmark}}}
\fancyhead[RO]{\small\textit{\nouppercase{\rightmark}}}
\fancyfoot[C]{\thepage}
\renewcommand{\headrulewidth}{0.4pt}
\renewcommand{\footrulewidth}{0pt}

% ---- Chapter / section titles ------------------------------------------
\titleformat{\chapter}[display]
  {\normalfont\sffamily\Huge\bfseries\color{mayablue}}
  {\filleft\Large\textit{Chapter \thechapter}}
  {1ex}
  {\titlerule\vspace{1ex}\filright}
\titlespacing*{\chapter}{0pt}{-30pt}{20pt}

\titleformat{\section}{\normalfont\sffamily\Large\bfseries\color{mayalavender}}{\thesection}{1em}{}
\titleformat{\subsection}{\normalfont\sffamily\large\bfseries\color{mayagreen}}{\thesubsection}{1em}{}

% ---- Listings: a C++ flavour we control ---------------------------------
% pdflatex's listings package can't render arbitrary Unicode; we keep code
% strictly ASCII inside lstlistings. For pretty terminal glyphs we drop
% into tcolorbox bodies.
\lstdefinelanguage{maya-cpp}{
  language=C++,
  morekeywords={
    constexpr,consteval,concept,requires,co_await,co_yield,co_return,
    char32_t,char8_t,string_view,expected,variant,visit,move_only_function,
    Signal,Computed,Effect,Batch,Cmd,Sub,Element,Model,Msg,Program,
    BoxNode,TextNode,DynNode,RuntimeTextNode,MapNode,Style,CTStyle,
    Color,Border,BorderStyle,FlexDirection,Align,Justify,Ctx,RunConfig,
    Canvas,StylePool,LayoutNode,Event,Key,Mouse,Resize,Paste,
    overload,Str,FixedName,Refined,EffectSet,Profile,Decision
  }
}
\lstdefinestyle{maya}{
  language=maya-cpp,
  basicstyle=\ttfamily\footnotesize,
  keywordstyle=\color{mayablue}\bfseries,
  stringstyle=\color{mayarose},
  commentstyle=\color{mayagrey}\itshape,
  showstringspaces=false,
  breaklines=true,
  breakatwhitespace=false,
  tabsize=2,
  columns=fullflexible,
  keepspaces=true,
  upquote=true,
  numbers=left,
  numberstyle=\ttfamily\scriptsize\color{mayaink},
  numbersep=12pt,
  frame=leftline,
  framerule=2pt,
  rulecolor=\color{mayablue},
  xleftmargin=3em,
  framexleftmargin=2.5em,
  framesep=8pt,
  literate=%
    {->}{{$\rightarrow$}}2
    {=>}{{$\Rightarrow$}}2
}
\lstset{style=maya}

\lstdefinestyle{shell}{
  language=bash,
  basicstyle=\ttfamily\footnotesize,
  keywordstyle=\color{mayablue},
  commentstyle=\color{mayagrey}\itshape,
  showstringspaces=false,
  breaklines=true,
  frame=leftline,
  framerule=2pt,
  rulecolor=\color{mayagreen},
  xleftmargin=1.5em,
  framesep=8pt,
  numbers=none
}

% ---- Boxes: idea, warning, terminal, theory ----------------------------
\newtcolorbox{idea}[1][]{
  enhanced, breakable,
  colback=mayablue!5, colframe=mayablue, fonttitle=\bfseries,
  title={Idea}, #1
}
\newtcolorbox{theory}[1][]{
  enhanced, breakable,
  colback=mayalavender!7, colframe=mayalavender, fonttitle=\bfseries,
  title={Theory}, #1
}
\newtcolorbox{warning}[1][]{
  enhanced, breakable,
  colback=mayaamber!10, colframe=mayaamber, fonttitle=\bfseries,
  title={Pitfall}, #1
}
\newtcolorbox{recap}[1][]{
  enhanced, breakable,
  colback=mayagreen!8, colframe=mayagreen, fonttitle=\bfseries,
  title={Recap}, #1
}
\newtcolorbox{terminalbox}[1][]{
  enhanced, breakable,
  colback=mayaink, colframe=mayaink, coltext=mayalight,
  fonttitle=\bfseries\color{mayalight},
  title={Terminal}, fontupper=\ttfamily\small, #1
}

% ---- Misc convenience ---------------------------------------------------
\newcommand{\maya}{\textsf{maya}}
\newcommand{\agentty}{\textsf{agentty}}
\newcommand{\cpp}{C\nolinebreak\hspace{-.05em}\raisebox{.2ex}{\small{++}}}
\newcommand{\code}[1]{\texttt{#1}}
\newcommand{\file}[1]{\textsf{\textit{#1}}}
\newcommand{\TODO}[1]{\textcolor{mayarose}{\textbf{TODO:} #1}}

\setlength{\parskip}{0.5\baselineskip}
\setlength{\parindent}{0pt}


\date{}

\begin{document}

\frontmatter

% ----------------------------------------------------------------------
% Cover page --- one word, one line
% ----------------------------------------------------------------------
\begin{titlepage}
\thispagestyle{empty}

\vspace*{\fill}

\begin{center}
  {\fontsize{96}{100}\selectfont\sffamily\bfseries\color{mayablue} maya}
\end{center}

\vspace{0.6cm}

\begin{center}
  {\color{mayablue}\rule{4cm}{0.6pt}}
\end{center}

\vspace{0.6cm}

\begin{center}
  {\fontsize{12}{16}\selectfont\sffamily\color{mayagrey}
    a c\nolinebreak\hspace{-.05em}\raisebox{.2ex}{\small{++}}26 terminal ui library, from the byte up\par}
\end{center}

\vspace*{\fill}

\begin{center}
  {\fontsize{9}{12}\selectfont\sffamily\color{mayagrey}
    companion to the \textbf{agentty} sources\par}
\end{center}

\vspace*{1.5cm}
\end{titlepage}

% ----------------------------------------------------------------------
% Half-title / abstract
% ----------------------------------------------------------------------
\thispagestyle{empty}
\vspace*{\fill}
\begin{center}
\begin{minipage}{0.78\textwidth}
\small\itshape
This book is a deep guide to \maya{}, the header-mostly \cpp{}26
terminal-UI library that powers \agentty. We start with the terminal
as a machine, work up through the \cpp{} language features \maya{}
relies on, and then walk every layer of the library --- the DSL,
elements, the layout engine, the SIMD-accelerated renderer, the
signal graph, the Elm-style runtime, and the widget catalogue. Each
chapter ends with a working program.

\medskip

Everything in this book is grounded in the actual source under
\file{maya/}. Where the code disagrees with the prose, the code wins
--- report the mismatch as a bug.
\end{minipage}
\end{center}
\vspace*{\fill}
\clearpage

% ----------------------------------------------------------------------
% Preface
% ----------------------------------------------------------------------
\chapter*{Preface}
\addcontentsline{toc}{chapter}{Preface}

\section*{Who this book is for}

You have probably opened a terminal today. You have probably also
used a program that draws into it --- \code{htop}, \code{vim}, a git
client, a music player, a chat tool. This book is for the reader who
wants to write such programs themselves, in modern \cpp{}, and who
wants to understand every layer between the keystroke and the
flickering glyph.

It assumes nothing. You do not need to have written a TUI before. You
do not need to know \cpp{}26. You do not need to know what a pty is,
or what an ANSI escape sequence is, or why anyone would put a
compile-time string inside a type. Each of those is introduced from
zero, with the smallest example that makes the idea concrete, and
then connected to the place in \maya{}'s source where you will see
it again.

If you already know \cpp{}26 fluently and just want the \maya{}
tour, you can skim Chapter~\ref{ch:cpp}. Everyone else: read in
order. The book is a rope ladder; each rung sits on the one below
it.

\section*{What you will be able to do by the end}

Four concrete things, in this order:

\begin{enumerate}[leftmargin=*]
  \item Read the bytes flowing between a program and a terminal,
    and write a TUI ``Hello, red world'' in plain \cpp{} with no
    library at all (Chapter~\ref{ch:terminal}).
  \item Read every line of every \maya{} header --- templates,
    NTTPs, concepts, variants, \code{constexpr} chains --- and know
    what each construct is doing (Chapter~\ref{ch:cpp}).
  \item Build a non-trivial \maya{} program of your own, from a
    \code{Program} struct through a \code{view} function to a
    \code{Cmd<Msg>} effect, and trust that the compile-time checks
    catch the impossible states (Parts II and IV).
  \item Open the SIMD diff loop, the signal graph, the layout
    engine, or the inline-mode witness chain and follow what each
    one is doing on its own terms (Part III).
\end{enumerate}

That is the contract. If a chapter does not move you closer to one
of those four outcomes, file an issue.

\section*{How the book is shaped}

Four parts, mirroring four levels of zoom.

\paragraph{Part I --- Foundations.} The terminal as a state machine
speaking bytes; the \cpp{} language features \maya{} relies on. The
output of Part I is a reader who can run the examples and read the
headers.

\paragraph{Part II --- Architecture.} The Elm-style runtime
(\code{init}/\code{update}/\code{view}/\code{subscribe}), the
compile-time DSL that produces UI descriptions, the in-memory
\code{Element} tree, and the flexbox layout engine that places
them. The output of Part II is a reader who can build apps.

\paragraph{Part III --- Rendering.} Cells and canvases, SIMD-accelerated
frame diffing, the ANSI byte serialiser, and \maya{}'s unusual inline
mode that lives inside your scrollback instead of taking over the
terminal. The output of Part III is a reader who understands why
\maya{} feels fast.

\paragraph{Part IV --- Reactivity, runtime, and widgets.} The signal
graph that powers fine-grained reactivity, the input parser, the
full \code{Cmd}/\code{Sub} effect algebra, the catalogue of widgets,
and how to package and ship a binary built on \maya{}. The output
of Part IV is a reader who can ship.

\section*{How to read the boxes}

Five coloured boxes recur throughout. They are signposts; you can
read the book without them, but they tell you what the prose around
them is doing.

\begin{idea}
\textbf{Idea} (blue). One sentence summarising the key intuition of
the surrounding section. If you skim, read these.
\end{idea}

\begin{theory}
\textbf{Theory} (purple). The background concept from computer
science or programming-language theory behind what we just
explained. Skippable on first read; useful on second.
\end{theory}

\begin{warning}
\textbf{Pitfall} (amber). A mistake people new to the material
commonly make, and how to avoid it.
\end{warning}

\begin{recap}
\textbf{Recap} (green). A two-or-three-sentence summary of the
section just finished.
\end{recap}

\begin{terminalbox}
\textbf{Terminal} (black). Literal bytes or output, the way the
terminal sees or shows them.
\end{terminalbox}

Code in monospace with a coloured rule on the left is \cpp{} (or
occasionally shell). Lines are numbered. Long examples are split
into pieces and each piece is decoded after.

\section*{A note on accuracy}

Every code fragment in this book is either copied from \maya{}'s
headers or paraphrased from them. Where we simplify --- to keep an
example short, or to elide an implementation detail not yet
introduced --- we say so and name the file where the real version
lives. Read the book with the source open in another window;
that is the intended posture.

\bigskip

\noindent The chapters that follow assume nothing and explain
everything. Let us begin with the terminal.

\clearpage

\tableofcontents

\mainmatter

% ---- Part I: Foundations -----------------------------------------------
\part{Foundations}
\chapter{Hello, maya}
\label{ch:intro}

This book teaches you \maya{} --- a library for building programs that
draw their interface inside a terminal window --- assuming you are new
to almost everything involved. We will not skip steps. By the end of
this first chapter you will know what kind of program \maya{} is, why
anyone bothers building such a thing in 2025, what is happening on
your screen \emph{right now} as you read this in a terminal, and how
to compile and run your first \maya{} program. None of that requires
prior knowledge of \cpp{}; we will treat the example as a thing to
look at, and Chapter~\ref{ch:cpp} will teach you the language behind
it.

\section{What is a terminal? (the honest answer)}

You have almost certainly used a terminal: it is the black-or-dark
window that says things like \code{\$} or \code{>} and lets you type
commands. It might be called Terminal, iTerm, kitty, Alacritty, GNOME
Terminal, Konsole, Windows Terminal, or just \emph{the console}.

Here is what the terminal actually is, in plain words:

\begin{enumerate}[leftmargin=*]
  \item A \textbf{window on your screen} that shows a rectangular grid
    of text. Every position in the grid holds \emph{one character} ---
    a letter, a digit, a symbol --- possibly with a colour.
  \item A \textbf{program} (the \emph{terminal emulator}) that owns
    that window. The terminal emulator's job is twofold:
    \begin{itemize}[leftmargin=*]
      \item When the user types on the keyboard, send those keystrokes
        to whatever program is running inside.
      \item When that program prints text, draw the text into the
        window.
    \end{itemize}
  \item A \textbf{program running inside it}. Usually this is a
    \emph{shell} like \code{bash} or \code{zsh} --- the thing that
    shows the \code{\$} prompt --- but it can be \emph{any} program.
    When you run \code{ls}, the shell starts \code{ls}, which prints
    file names, then exits, and you are back at the shell.
\end{enumerate}

So the terminal emulator is a kind of two-way pipe with a screen
attached: text goes in (from the keyboard), text comes out (from the
program), and the rectangle of cells on screen is updated to reflect
what the program has printed so far.

\begin{idea}
A terminal is a window that displays a grid of characters, and a
program that pumps text between your keyboard and whatever else is
running.
\end{idea}

\subsection{Why does it look so old?}

Because it is. The protocol the terminal speaks --- the rules for
``move the cursor to row 5, column 12'' and ``print this in red'' ---
was finalised on a physical device called the DEC VT100 in 1978, and
every modern terminal still understands it. Programs written today
talk the same byte language as programs written for an actual VT100.
This is unusual in computing and turns out to be very valuable: any
program that follows the rules will work on every machine you can
\code{ssh} into, today, forever.

\section{What is a TUI?}

A \textbf{TUI} is a \emph{Text User Interface}: a program that does
not just dump output line by line, but \emph{takes over} the terminal
window and draws a proper interface inside it. Menus, panels,
scrollable lists, progress bars, even charts --- all rendered out of
characters.

You have used TUIs without thinking of them as such:

\begin{itemize}[leftmargin=*]
  \item \code{vim} or \code{nano} when you edit a file.
  \item \code{htop} when you watch processes.
  \item \code{git log} (the pager part), \code{less}, \code{man}.
  \item Installer screens on a Linux ISO.
\end{itemize}

A \textbf{GUI} (Graphical User Interface) program asks the operating
system to paint individual pixels: \emph{draw a blue rectangle from
$(100, 80)$ to $(300, 120)$}. A TUI program does something more
modest: it asks the terminal emulator to put specific characters into
specific grid cells, possibly with colour. The terminal emulator
itself decides how each character looks (the font, the cell size); the
TUI program only chooses what character goes where.

\begin{idea}
GUI = pixels you place anywhere. TUI = characters you place inside a
fixed grid of cells.
\end{idea}

That restriction --- one character per cell --- is the entire reason
TUIs are smaller, faster, and easier to write than GUIs. There is no
font engine, no anti-aliasing, no input method editor; just a grid and
some bytes.

\subsection{Why bother with TUIs in 2025?}

You might reasonably ask: why not just write a web page? Several
reasons:

\begin{itemize}[leftmargin=*]
  \item \textbf{They run anywhere there is a terminal.} A developer
    on a server they reached by \code{ssh} has no browser, no
    desktop, sometimes not even a graphical session. They do have a
    terminal.
  \item \textbf{They are very fast.} A typical TUI redraw is a few
    kilobytes of bytes. A web page is megabytes of HTML, CSS, JS,
    images, and a browser.
  \item \textbf{They are scriptable.} A TUI program's output is
    still text, so it composes with pipes, \code{grep}, \code{ssh},
    \code{tmux}, and the rest of the Unix toolbox.
  \item \textbf{They are pleasant to use for developers.} Keyboard
    first, no mouse round-trips, low latency, look good in any
    monospace font.
\end{itemize}

The trade-off is the visual restriction: no images (well, mostly), no
freeform layouts, no animations beyond what a grid of glyphs can do.
For a large class of programs, that is exactly the right trade.

\section{One frame of a TUI: what literally happens}
\label{sec:intro:one-frame}

To make the rest of this book concrete, let us walk what happens in
one ``frame'' of a typical TUI --- the unit of work between two
updates of the screen. We will use a counter app as the example: it
shows a number, and the keys \code{+} and \code{-} change it.

Suppose the screen currently shows the number $7$. The user presses
\code{+}.

\paragraph{1. The keyboard sends a byte to the kernel.}
When you press \code{+}, the OS kernel receives the byte $0x2B$
(ASCII for \code{+}) from your keyboard driver and places it into
the input buffer of whichever terminal program currently has focus
--- your terminal emulator.

\paragraph{2. The terminal emulator forwards the byte.}
Your terminal emulator (kitty, iTerm, GNOME Terminal, \dots) writes
the byte into one end of the \emph{pseudo-terminal} --- the kernel
byte-pipe between the emulator and the program running inside it.
Chapter~\ref{ch:terminal} covers the pseudo-terminal in detail.

\paragraph{3. The maya program reads the byte.}
The \maya{} program is sitting in a loop, asking the OS ``is there
any input ready?''. Now there is. It reads $0x2B$ from its standard
input.

\paragraph{4. Maya turns the byte into an Event.}
A single byte is not yet a useful concept. \maya{}'s input parser
decides ``this is a printable character, namely the character
\code{+}'' and constructs a structured value: \code{Key\{.kind =
KeyKind::Char, .character = '+'\}}. That value is then wrapped in
the broader \code{Event} type and handed to your application code.

\paragraph{5. Your code computes a new state.}
Your application receives the \code{Event}, decides it means
``increment'', and produces the new state: the counter is now $8$.
In an Elm-style \maya{} program this is one pure function:
\code{update(state=7, event=Inc) = 8}.

\paragraph{6. Maya turns the new state into an Element tree.}
Your \code{view} function is called with the new state ($8$) and
returns an \code{Element}: a description of the UI, ``a box with a
rounded border containing the text \code{Count: 8}''. This is a
value in memory, not yet bytes on the wire.

\paragraph{7. The layout engine assigns positions.}
Given the current terminal size --- say 80 columns by 24 rows ---
\maya{} walks the element tree and decides where every piece goes:
the outer box at $(0, 0)$ of size $14 \times 3$, the inner text at
$(1, 1)$, and so on. The output is a list of styled cells, one per
position.

\paragraph{8. The renderer builds a virtual screen.}
Those cells are written into a 2-D grid in memory --- the
\emph{canvas}. Every position on the screen now has a known
character, foreground colour, background colour, and attribute set.

\paragraph{9. The differ finds what changed.}
A copy of last frame's canvas is still around. \maya{} compares the
two canvases cell by cell. For our counter, almost everything is
identical; only one cell changed --- the \code{7} became an \code{8}.

\paragraph{10. Maya emits ANSI bytes for only that cell.}
A short string of bytes is produced: ``move cursor to (col, row)
of the changed cell; set the right colour; print \code{8}''. Even
though there might be hundreds of cells on screen, only a dozen
bytes are sent.

\paragraph{11. Those bytes flow back through the pty to the emulator.}
The emulator reads them, parses them, and updates its window. The
user sees the digit change. From keypress to pixel update is
typically a couple of milliseconds.

That is one frame. The rest of the book is, in a sense, just
detailed coverage of each of those eleven steps. The DSL
(Chapter~\ref{ch:dsl}) and Element tree
(Chapter~\ref{ch:elements}) cover steps 5--6; the layout engine
(Chapter~\ref{ch:layout}) is step 7; the canvas
(Chapter~\ref{ch:canvas}) is step 8; the differ
(Chapter~\ref{ch:simd}) is step 9; ANSI
(Chapter~\ref{ch:ansi}) is step 10; the runtime that orchestrates
the whole loop (Chapter~\ref{ch:elm}) is what sits around steps
3--6.

\begin{idea}
One frame: a byte arrives, it becomes an event, you compute new
state, maya builds the new screen, compares it to the old one, and
sends bytes for just the differences. That tight loop is what every
TUI library is, including \maya{}.
\end{idea}

\section{What is \maya{}?}

\maya{} is a \cpp{} \textbf{library} --- a bundle of code you include
in your own program --- whose job is to make writing a TUI pleasant.
Without a library, writing a TUI means manually printing escape
sequences (we will look at the bytes in Chapter~\ref{ch:terminal}),
manually tracking what is on screen, manually parsing keyboard input,
and manually re-laying-out everything when the window resizes.

With \maya{}, you describe what you \emph{want} on screen as a tree of
boxes and text, and \maya{} works out the bytes to send. Concretely,
the work the library does for you is:

\begin{enumerate}[leftmargin=*]
  \item Take your description (\emph{``a box containing the text
    `Hello' with a rounded border''}) and figure out the position
    and size of every piece given the current window size.
  \item Build a virtual grid of cells representing what the screen
    should look like.
  \item Compare that grid to what is \emph{already} on screen, and
    emit just enough escape sequences to update the cells that
    changed.
  \item Read keyboard, mouse, paste, and resize events and feed them
    back to your program.
  \item Restore the terminal to a clean state when your program
    exits, even on crashes.
\end{enumerate}

This book covers each of those jobs in detail. For now, treat
\maya{} as a black box that turns ``description of a UI'' into ``the
correct bytes on the wire''.

\section{What is a tree, and why a tree?}
\label{sec:intro:tree}

In computing, a \textbf{tree} is a data structure that looks like a
family tree turned upside down: one thing at the top, called the
\textbf{root}, which has zero or more \textbf{children}, each of
which may have its own children, and so on. The end of every branch
--- a node with no children --- is called a \textbf{leaf}.

A terminal UI is naturally tree-shaped, and once you see it you
cannot unsee it. Consider a chat application:

\begin{center}
\small
\begin{verbatim}
  +-------------------------------------------+
  | Title: my chat                           |
  +-------------------------------------------+
  | sidebar  |  message list                  |
  |  - alice |   alice: hi                    |
  |  - bob   |   me   : hello                 |
  |          |   bob  : evening               |
  +----------+--------------------------------+
  | type here...                             |
  +-------------------------------------------+
\end{verbatim}
\end{center}

Described as a tree, this UI is:

\begin{center}
\small
\begin{verbatim}
  Root: vertical stack
  |-- Title bar (text)
  |-- Body: horizontal split
  |   |-- Sidebar: vertical list
  |   |   |-- "alice"  (text)
  |   |   `-- "bob"    (text)
  |   `-- Messages: vertical list
  |       |-- "alice: hi"      (text)
  |       |-- "me   : hello"   (text)
  |       `-- "bob  : evening" (text)
  `-- Input: bordered text box
\end{verbatim}
\end{center}

Four kinds of node appear: ``vertical stack'', ``horizontal split'',
``list'', ``text''. The first three are \emph{container} nodes ---
they have children that they arrange. Text nodes are leaves: they
have no children, just a string.

A tree is the right structure because UIs are recursive: a sidebar
\emph{is} a UI; a message list \emph{is} a UI; both are pieces of a
bigger UI. Trees compose. You can reuse the message-list subtree in
a completely different program by grafting it onto a different root.

\maya{}'s \code{Element} type is exactly the tree representation of
a UI. Container nodes (\code{BoxElement}, \code{ElementList}) hold
\code{std::vector<Element>} of children; leaf nodes
(\code{TextElement}) hold the bytes of a string. The DSL is the
surface syntax that lets you spell that tree out as readable \cpp{}
code. Chapter~\ref{ch:elements} dissects the \code{Element} type in
full.

\begin{idea}
A UI is a tree of containers and text. Containers arrange their
children; text leaves carry the actual characters. \maya{}'s
\code{Element} type \emph{is} that tree.
\end{idea}

\section{Why does \maya{} exist when there are already others?}

There are many TUI libraries. \code{ncurses} (C, ancient and
respected), \code{ratatui} (Rust), \code{bubbletea} (Go),
\code{textual} (Python), \code{ink} (Node.js). Each is good in its
own way. \maya{} picks three positions that distinguish it.

\paragraph{1. Wrong UI code does not compile.}
Most libraries accept your code happily and let you discover the bug
when you run the program: a border colour set without a border becomes
a no-op, a negative padding silently becomes zero, a layout meant to
be a row accidentally becomes a column. \maya{} is built on \cpp{}'s
\emph{type system} so that a lot of these mistakes refuse to compile.
The compiler tells you, at build time, exactly which line is wrong.

As a tiny taste: in \maya{} you might write

\begin{lstlisting}
auto ui = t<"Hello"> | bcol<60, 65, 80>;   // border colour, no border yet
\end{lstlisting}

A library written in Python or Go would accept that and silently do
nothing with the border colour at run-time. \maya{} refuses to
compile it, because the \code{bcol} (``border colour'') modifier is
defined to only apply to nodes that have already been given a
border. The compiler error names the constraint that failed; you fix
the code by adding \code{| border\_<Round>} first.

You will not see why this is possible until Chapter~\ref{ch:dsl}; the
short version is that the description you write is stored in the
\emph{types} of the \cpp{} values, and \maya{}'s API states the rules
between those types.

\paragraph{2. Re-drawing is wired to be fast.}
Once \maya{} has decided what the next frame should look like, it
compares that virtual grid of cells with the grid that is currently on
screen. Only the cells that changed are sent to the terminal. That
comparison is done with the CPU's \textbf{SIMD} instructions ---
special instructions that operate on many bytes at once. The effect
is that even on a slow \code{ssh} link, a \maya{} program redraws
smoothly. We will dissect the SIMD code in Chapter~\ref{ch:simd}; you
do not need to know what SIMD is yet.

\paragraph{3. The runtime is borrowed from a language called Elm.}
\maya{} organises a whole program around one idea: your program's
state is a single value, your program's events are a single closed
list of cases, and the function that updates the state given an event
is \emph{pure} --- meaning it computes a new value and does nothing
else (no printing, no opening files, no network calls). Anything that
needs to ``happen in the world'' is described as data and handed to
the library to perform. This is unusual for \cpp{}, and the payoff is
programs that are easy to test, easy to replay, and easy to reason
about. Chapter~\ref{ch:elm} is dedicated to this idea.

If those three sentences mean nothing to you yet, that is fine. The
book is structured so each one is earned, slowly, before you have to
use it.

\section{The layers of a \maya{} program}

It will help to have a picture of how the work flows, even if many of
the words are unfamiliar now. Read it once; you do not need to
memorise it.

\begin{center}
\small
\begin{tabular}{r|l}
\textbf{Your code} & describe a UI; say what should happen on each event \\
\hline
\textbf{DSL} & \cpp{} syntax (\code{|}, \code{t<"...">}) that builds your description \\
\textbf{Element} & in-memory tree of boxes and text \\
\textbf{Layout} & gives every box a position and size \\
\textbf{Canvas} & a grid of cells with characters, colours, attributes \\
\textbf{Diff} & finds which cells changed since last frame \\
\textbf{ANSI} & escape-sequence bytes for the changed cells \\
\textbf{Writer} & sends those bytes to the terminal \\
\hline
\textbf{Terminal} & draws characters into the window \\
\end{tabular}
\end{center}

A few words on the unfamiliar ones:

\begin{itemize}[leftmargin=*]
  \item \textbf{DSL} stands for \emph{Domain-Specific Language}. It is
    not a separate language --- it is just normal \cpp{} ---  but the
    surface looks specialised. You write
    \code{t<"Hi"> | Bold | border\_<Round>} and that is valid \cpp{}
    that produces a UI description. The DSL is what makes \maya{}
    code short.
  \item \textbf{Element} is the name we give to the in-memory shape
    of your UI. It is a tree because boxes contain other boxes (and
    text). The element tree is the central data structure of
    \maya{}.
  \item \textbf{ANSI} is the byte language the terminal understands.
    \code{ANSI} is short for \emph{American National Standards
    Institute}, who blessed the standard in 1976. We will look at
    individual ANSI bytes in Chapter~\ref{ch:terminal}.
\end{itemize}

Each layer knows only about the one immediately below it. Your code
talks to the DSL, the DSL produces an Element tree, the layout engine
walks the tree, and so on. This separation is what makes the library
both testable and replaceable: someone could swap out the SIMD diff
for a slower scalar one without anything above noticing.

\section{Your first \maya{} program}

Here is a complete, real \maya{} program. It compiles, runs, prints a
small banner, and exits.

\begin{lstlisting}
#include <maya/maya.hpp>
using namespace maya;
using namespace maya::dsl;

int main() {
    constexpr auto ui = v(
        t<"Hello, maya"> | Bold | Fg<100, 180, 255>,
        t<"a tiny banner"> | Dim
    ) | border_<Round> | pad<1>;

    print(ui.build());
}
\end{lstlisting}

Before we decode it line by line, here is roughly what you should
see in your terminal after running this:

\begin{terminalbox}
\textcolor{mayablue}{\textbf{.----------------.}}\\
\textcolor{mayablue}{\textbf{|}}\, \textbf{Hello, maya}~~\, \textcolor{mayablue}{\textbf{|}}\\
\textcolor{mayablue}{\textbf{|}}\, a tiny banner \textcolor{mayablue}{\textbf{|}}\\
\textcolor{mayablue}{\textbf{`----------------'}}
\end{terminalbox}

A two-line text stack, the first line bold and pale blue, the second
dimmed, surrounded by a rounded border with one cell of padding
inside it. The corners are drawn with the Unicode box-drawing
glyphs the terminal renders as smooth round corners.

You are not expected to understand every token of the program. You
\emph{are} expected to be able to point at each piece and at least
name what it is. Let us go line by line.

\subsection{Line by line}

\paragraph{\code{\#include <maya/maya.hpp>}.}
\code{\#include} is the \cpp{} way of saying ``copy and paste the
contents of this file here before compiling''. \file{maya/maya.hpp} is
a single header that pulls in the public parts of \maya{}. After this
line, the names \code{v}, \code{t}, \code{Bold}, \code{print}, and so
on are available.

\paragraph{\code{using namespace maya;} and \code{using namespace maya::dsl;}.}
A \textbf{namespace} is a labelled container of names. \maya{}'s
public API lives inside the namespace \code{maya}, and the DSL pieces
live in \code{maya::dsl}. The two \code{using namespace} lines say
``for the rest of this file, names from those namespaces are visible
without their prefix''. Without them you would have to write
\code{maya::print(maya::dsl::v(...))}. Verbose; we elide it for
examples.

\paragraph{\code{int main()}.}
Every \cpp{} program starts at a function called \code{main}. The
\code{int} says \code{main} returns an integer (the exit status); we
do not write a \code{return} statement, and \cpp{} treats that as
\code{return 0} which means ``success''.

\paragraph{\code{constexpr auto ui = ...}.}
Three things to unpack:
\begin{itemize}[leftmargin=*]
  \item \code{auto} means ``figure out the type of this for me, based
    on the right-hand side''. \cpp{} normally requires you to spell
    the type; \code{auto} lets the compiler infer it. The actual type
    of \code{ui} here is something elaborate that we do not care to
    write out.
  \item \code{constexpr} means ``this value is computed at
    \emph{compile time}, not at runtime''. The whole tree --- the
    border, the colour, the text --- is baked into the binary. When
    the program starts, the tree is already there.
  \item \code{ui} is just the name of the variable. We could have
    called it anything.
\end{itemize}

\paragraph{\code{v(...)}.}
\code{v} is \maya{}'s function for ``stack these things vertically''.
The arguments inside the parentheses are the things to stack ---
in this case, two pieces of text. There is also \code{h(...)} for
horizontal.

\paragraph{\code{t<"Hello, maya">}.}
\code{t} is a \textbf{template} parameterised by a compile-time
string. The angle brackets \code{<\ldots>} are the template-argument
syntax. \code{t<"Hello, maya">} is a value representing the literal
text ``Hello, maya''. The peculiar thing here is that the text itself
lives inside the type --- swapping \code{"Hello"} for \code{"World"}
gives a different type. This is one of the tricks that lets the
compiler check things. Chapter~\ref{ch:cpp} explains how it works in
detail.

\paragraph{\code{| Bold | Fg<100, 180, 255>}.}
The vertical bar \code{|} is the \textbf{pipe operator} in \cpp{}. By
default it means bitwise-or on integers, but you can teach it to mean
anything for your own types. \maya{} has taught it to mean
\emph{compose this style onto the thing on the left}. \code{Bold} is
a tag value meaning ``apply bold''; \code{Fg<R, G, B>} is a template
that takes three numbers in the range $0$--$255$ and represents a
foreground colour.

So \code{t<"Hello, maya"> | Bold | Fg<100, 180, 255>} reads
left-to-right: ``the text `Hello, maya', made bold, given a foreground
colour of $(100, 180, 255)$''. A pleasant pale blue.

\paragraph{\code{| Dim}.}
Same idea: applies the \emph{dim} attribute (terminals usually render
this as a faded foreground). One of the few attributes whose
appearance varies most across terminals.

\paragraph{\code{| border\_<Round> | pad<1>}.}
After the inner \code{v(...)}, two more pipes wrap the result. The
first puts a rounded border around the whole stack; the second adds
one cell of padding (space) inside the border. \code{Round} is a
choice from a small set (\code{Round}, \code{Square}, \code{Double},
\code{Thick}, etc.); \code{pad<1>} carries the number $1$ in its type.

\paragraph{\code{print(ui.build())}.}
Finally, the call that does something. \code{ui.build()} converts the
compile-time description \code{ui} into a runtime \code{Element} ---
the in-memory tree we mentioned earlier. \code{print(...)} is the
simplest of four output modes \maya{} provides; it paints the element
once, to standard output, and returns. No event loop, no alternate
screen, no clearing the terminal. For a static banner, that is what
we want. Bigger programs use \code{run<P>(...)}, which we will meet in
Chapter~\ref{ch:elm}.

\subsection{What you should take away}

Three observations are worth carrying forward:

\begin{enumerate}[leftmargin=*]
  \item The structure of the UI lives in the \cpp{} source as a
    \emph{value}: \code{ui} is a thing you can pass to a function,
    return from a function, store in an array. It is not strings of
    markup, not a separate file format, not XML.
  \item Because the value is \code{constexpr}, the compiler builds
    it. If you wrote \code{Fg<300, 0, 0>} (out of range), the
    compiler would refuse. If you wrote \code{border\_<Squiggle>}
    (not a real style), the compiler would refuse. These checks are
    free at runtime --- they have already happened.
  \item Reading code with \code{|} is like reading English left to
    right: subject first, modifiers afterwards. Once you get used
    to it the syntax is quite pleasant to scan.
\end{enumerate}

\section{Building and running}

The \maya{} repository builds with \textbf{CMake}, a tool that
generates build files for your platform. You need a recent \cpp{}
compiler --- GCC 15 or newer, or a recent Clang --- because \maya{}
uses \cpp{}26 features. On Linux:

\begin{lstlisting}[style=shell]
git clone --recursive https://github.com/1ay1/agentty
cd agentty/maya
cmake -B build
cmake --build build -j$(nproc)
\end{lstlisting}

What each command does, briefly:

\begin{itemize}[leftmargin=*]
  \item \code{git clone --recursive ...} downloads the project and
    its sub-projects (called \emph{submodules}).
  \item \code{cd agentty/maya} moves into the maya directory.
  \item \code{cmake -B build} reads \maya{}'s \file{CMakeLists.txt}
    file and produces build files in a new directory called
    \file{build}.
  \item \code{cmake --build build -j\$(nproc)} actually compiles
    things, using as many parallel jobs as you have CPU cores.
\end{itemize}

If it succeeds, \file{build/} contains around two dozen example
programs --- \file{build/counter}, \file{build/dashboard},
\file{build/doom\_fire} --- one for each file in \file{examples/}.
Run any of them by typing its path:

\begin{lstlisting}[style=shell]
./build/counter
\end{lstlisting}

A small UI appears. Press the keys it documents (often \code{q} to
quit). Resize the terminal window while it runs; watch the UI
re-flow.

\subsection{Adding your own example}

The easiest way to try the program above is to drop it into
\file{maya/examples/hello.cpp}, re-run CMake, and rebuild:

\begin{lstlisting}[style=shell]
# from inside maya/
$EDITOR examples/hello.cpp        # paste the program
cmake -B build                    # picks up the new file
cmake --build build -j$(nproc)
./build/hello
\end{lstlisting}

This loop --- edit, build, run --- is the loop you will spend the
rest of the book in.

\begin{warning}
If the build complains about \code{constexpr}, \code{concept}, or
unknown standard-library headers, your compiler is too old. \code{g++
--version} should be 15 or newer. On Ubuntu/Debian:
\code{sudo apt install g++-15}. On macOS, install a recent Homebrew
\code{llvm}. On Windows, use WSL with a current Ubuntu.
\end{warning}

\section{What the rest of the book looks like}

The book has four parts.

\textbf{Part I, Foundations}, is the part you are reading. After this
chapter, Chapter~\ref{ch:terminal} explains what a terminal really is
at the byte level --- you will be able to write ``Hello'' in red
without any library at all by the end of it. Chapter~\ref{ch:cpp} is
a careful, from-the-ground-up tour of the \cpp{} features \maya{}
uses; if you are new to \cpp{}, this is the chapter that pays for
everything that follows.

\textbf{Part II, Architecture}, is about the way a \maya{} program is
structured. The Elm runtime (\ref{ch:elm}), the DSL that produces
UI descriptions (\ref{ch:dsl}), the in-memory tree that is the result
(\ref{ch:elements}), and the layout engine that decides where things
go (\ref{ch:layout}).

\textbf{Part III, Rendering}, is the bytes-on-the-wire half. Cells and
canvases (\ref{ch:canvas}), SIMD comparison (\ref{ch:simd}), the ANSI
language (\ref{ch:ansi}), and the special mode where \maya{} renders
\emph{inline} without taking over the terminal (\ref{ch:inline}).

\textbf{Part IV, Runtime and widgets}, is everything that happens
while a program runs. Fine-grained reactivity with signals
(\ref{ch:signals}), input parsing (\ref{ch:events}), the full effect
algebra (\ref{ch:cmdsub}), the widget catalogue (\ref{ch:widgets}),
and how to ship a program built on \maya{} (\ref{ch:build}).

You can read straight through. You can also skip Chapter~\ref{ch:cpp}
on first pass if you already know modern \cpp{} well, and come back
to it as a reference.

\section{How \maya{} compares to other TUI libraries}
\label{sec:intro:compare}

It helps to know what other people are doing. Here is the same
static banner --- text, bold, colour, border --- in four other
popular TUI libraries, alongside the \maya{} version, so you have a
felt sense of where this book is taking you.

\paragraph{ncurses (C, 1990s).}
\begin{lstlisting}[language=C]
#include <ncurses.h>
int main(void) {
    initscr();
    start_color();
    init_pair(1, COLOR_CYAN, COLOR_BLACK);
    box(stdscr, 0, 0);
    attron(A_BOLD | COLOR_PAIR(1));
    mvprintw(1, 2, "Hello, maya");
    attroff(A_BOLD | COLOR_PAIR(1));
    mvprintw(2, 2, "a tiny banner");
    refresh();
    getch();
    endwin();
}
\end{lstlisting}
Imperative. You position every character by row and column. There
is no layout; if the terminal is smaller than your hard-coded
numbers, the UI breaks.

\paragraph{textual (Python).}
\begin{lstlisting}[language=Python]
from textual.app import App
from textual.widgets import Static

class Hello(App):
    def compose(self):
        yield Static("[bold cyan]Hello, maya[/]\n[dim]a tiny banner[/]",
                     classes="box")
    CSS = ".box { border: round cyan; padding: 1; }"

Hello().run()
\end{lstlisting}
Declarative, pleasant, but markup is strings (the bold/colour
spans), and the CSS is a separate sub-language the Python compiler
cannot check.

\paragraph{ratatui (Rust).}
\begin{lstlisting}[language=C]
let block = Block::default()
    .borders(Borders::ALL)
    .border_type(BorderType::Rounded)
    .border_style(Style::default().fg(Color::Cyan));
let text = Text::from(vec![
    Line::from(Span::styled("Hello, maya",
        Style::default().fg(Color::Cyan).add_modifier(Modifier::BOLD))),
    Line::from(Span::styled("a tiny banner",
        Style::default().add_modifier(Modifier::DIM))),
]);
frame.render_widget(Paragraph::new(text).block(block), frame.size());
\end{lstlisting}
Close to \maya{} in spirit --- a tree of widgets, styles as values
--- but more verbose, and styles are runtime values, so e.g. an
out-of-range colour is caught only at runtime (or in your tests).

\paragraph{\maya{} (\cpp{}26).}
\begin{lstlisting}
constexpr auto ui = v(
    t<"Hello, maya"> | Bold | Fg<100, 180, 255>,
    t<"a tiny banner"> | Dim
) | border_<Round> | pad<1>;
print(ui.build());
\end{lstlisting}
The whole description is a \code{constexpr} value. The compiler
builds the tree at translation time; bad styles, bad borders, and
out-of-range colours are caught before the program runs.

The trade is roughly: more verbose than \maya{} in dynamic
languages; more strictly checked in \maya{} than anywhere else;
broadly similar mental model across the declarative tier (textual,
ratatui, \maya{}). Once you can read one, you can read all of them.

\section{What is a ``binary''?}
\label{sec:intro:binary}

We keep mentioning ``the binary'' --- the thing your compiler
produces. One sentence on what it is: a \textbf{binary} (or
\textbf{executable}) is a single file on disk full of CPU
instructions plus the data your program needs. When you run
\code{./hello}, the operating system loads that file into memory
and begins executing the instructions at a fixed entry point
(typically a function the compiler links in, which then calls your
\code{main}).

The binary is built once and run many times. Anything that can be
computed at compile time can live inside the binary directly ---
text strings, constant tables, style values --- without any cost at
startup. That is what \maya{} exploits with \code{constexpr}: the
UI tree of a static banner is, literally, encoded into the bytes of
the executable file.

\section{How this book is laid out}

A few visual conventions you will see in the chapters:

\begin{idea}
Blue boxes are \textbf{ideas} --- a single sentence summarising the
key intuition. If you only remember one thing from a section, this is
what it should be.
\end{idea}

\begin{theory}
Purple boxes are \textbf{theory} --- the background concept from
computer science or programming-language theory behind what the
section just explained. Skippable on first read; useful on second.
\end{theory}

\begin{warning}
Amber boxes are \textbf{pitfalls} --- a mistake people new to the
material commonly make, and how to avoid it.
\end{warning}

\begin{recap}
Green boxes at the end of a section summarise it in two or three
sentences.
\end{recap}

\begin{terminalbox}
Black boxes contain literal terminal bytes or output, the way the
terminal sees or shows them.
\end{terminalbox}

Code in monospace with a coloured rule on the left is \cpp{} (or
occasionally shell). Lines are numbered. Long examples are broken into
pieces and explained.

\begin{recap}
A TUI is a program that draws into the grid of cells inside a terminal
window. \maya{} is a \cpp{}26 library that lets you describe a TUI as
a tree of styled boxes and text and renders it for you. Its three
distinctive ideas --- compile-time correctness, SIMD-fast diffing, and
a pure-functional runtime --- are each given a chapter of their own
later. The first \maya{} program is a five-line static banner; the
loop you spent the rest of the book in is edit, build, run.
\end{recap}

\section*{Workshop}

You should not yet expect to do these without looking things up. They
are calibrated so that doing them takes you exactly through the
machinery the chapter introduced.

\begin{enumerate}[leftmargin=*]
  \item \textbf{Build the repository.} Get a clean
    \code{cmake --build build -j\$(nproc)} succeeding on your
    machine. If the build fails, the error message is the assignment
    --- read it, look up the missing tool, install it.
  \item \textbf{Run three examples} from \file{build/}. Suggested:
    \file{counter} (smallest), \file{dashboard} (mid-sized),
    \file{doom\_fire} (visually striking). Resize the window while
    each is running.
  \item \textbf{Compile the Hello-banner.} Drop the program from
    earlier into \file{maya/examples/hello.cpp}, re-run CMake,
    rebuild, run.
  \item \textbf{Mutate the banner.} Change the foreground colour to
    something else. Swap \code{Round} for \code{Double} and then
    \code{Thick}. Change \code{pad<1>} to \code{pad<3>}. Each
    change is one number or one identifier; each requires a rebuild
    to take effect.
  \item \textbf{Provoke a compile error.} Change \code{Fg<100, 180,
    255>} to \code{Fg<100, 180, 999>}. Read the compiler's error
    message. You do not need to fix it elegantly --- just observe
    that it caught the mistake.
\end{enumerate}

\section*{Glossary of terms introduced in this chapter}

Keep this list handy; every term here is referenced freely in the
rest of the book.

\begin{description}[leftmargin=*, style=nextline]
  \item[\textbf{ANSI}] The byte language a program writes to a
    terminal to control colour, cursor position, and screen modes.
    Standardised in 1976. Chapter~\ref{ch:terminal} covers it.
  \item[\textbf{Binary, executable}] A file on disk holding compiled
    CPU instructions for your program. \code{./hello} runs the
    binary called \file{hello}.
  \item[\textbf{Canvas}] \maya{}'s in-memory 2-D grid of styled
    cells; the working representation of the next screen.
    Chapter~\ref{ch:canvas}.
  \item[\textbf{Cell}] One position in the terminal grid. Holds one
    character (sometimes two-cell-wide), foreground colour,
    background colour, and attribute bits.
  \item[\textbf{\code{constexpr}}] \cpp{} keyword meaning ``compute
    this at compile time''. \maya{}'s DSL leans on it heavily.
  \item[\textbf{Diff}] The pass that compares the new canvas to the
    previous one and decides which cells changed.
    Chapter~\ref{ch:simd}.
  \item[\textbf{DSL}] Domain-Specific Language. The pleasant
    \code{t<"..."> | Bold | ...} syntax in \maya{} is a DSL
    embedded inside \cpp{}. Chapter~\ref{ch:dsl}.
  \item[\textbf{Element}] \maya{}'s in-memory UI tree. The output of
    your \code{view} function; the input to the layout engine.
    Chapter~\ref{ch:elements}.
  \item[\textbf{Elm runtime}] The Elm-inspired loop
    (\code{init}--\code{update}--\code{view}--\code{subscribe}) that
    \maya{} programs are organised around. Chapter~\ref{ch:elm}.
  \item[\textbf{Escape sequence}] A short string of bytes beginning
    with \code{ESC} (\code{0x1B}) that tells the terminal to do
    something. The grammar of ANSI.
  \item[\textbf{Frame}] One complete update cycle: read events,
    compute new state, build new canvas, diff, emit bytes.
  \item[\textbf{GUI}] Graphical User Interface --- pixel-based.
    Contrast with TUI.
  \item[\textbf{Layout}] The step that assigns each Element a
    position and size given the terminal dimensions.
    Chapter~\ref{ch:layout}.
  \item[\textbf{Leaf, node, root}] Tree vocabulary. The root is the
    top of a tree; leaves are nodes with no children; everything
    else is an internal node.
  \item[\textbf{NTTP}] Non-Type Template Parameter --- a template
    argument that is a value, not a type. The trick that lets
    \code{t<"Hello">} encode a string into a type.
    Chapter~\ref{ch:cpp}.
  \item[\textbf{Pseudo-terminal, pty}] The kernel-managed byte pipe
    between a terminal emulator and the program running inside it.
    Chapter~\ref{ch:terminal}.
  \item[\textbf{Raw mode}] A configuration of the terminal where the
    kernel passes every keystroke through immediately, with no line
    editing or echo. The mode every TUI runs in.
  \item[\textbf{SIMD}] Single Instruction Multiple Data: CPU
    instructions that operate on many bytes at once. \maya{}'s diff
    pass uses SIMD. Chapter~\ref{ch:simd}.
  \item[\textbf{stdin / stdout / stderr}] The three default file
    descriptors --- 0, 1, 2 --- a process inherits. In a TUI all
    three are wired to the pty.
  \item[\textbf{Tree}] A data structure: a root with children, each
    of which may have children. A UI is naturally a tree.
  \item[\textbf{TUI}] Text User Interface. A program that draws
    into the grid of cells inside a terminal window.
  \item[\textbf{UTF-8}] The standard variable-length encoding of
    Unicode characters into bytes. Modern terminals expect UTF-8.
\end{description}

\chapter{The terminal, byte by byte}
\label{ch:terminal}

Chapter~\ref{ch:intro} described a terminal in plain words. This
chapter looks underneath. We will see exactly what bytes travel
between your program and the terminal emulator, what those bytes mean,
and why the protocol is the way it is. By the end you will be able to
read raw terminal output, write ``Hello'' in red without any library,
and trust your mental model of what \maya{} is doing one floor below
the API.

There is nothing magical here. The terminal is one of the simplest
abstract machines in computing. It has a small alphabet and a small
list of rules; we will learn both.

\section{What is a ``state machine''?}
\label{sec:term:fsm}

The chapter is titled ``the terminal as a state machine'', so it is
worth saying what that phrase means before we use it.

A \textbf{state machine} (formally, a \emph{finite-state automaton})
is a model of computation with three parts:

\begin{enumerate}[leftmargin=*]
  \item A finite set of \textbf{states} --- things the system can
    \emph{be in} right now.
  \item An \textbf{alphabet} --- the kinds of \emph{input}
    the system reacts to.
  \item A \textbf{transition table} --- for each (state, input)
    pair, the state the machine moves to next, and any output it
    produces along the way.
\end{enumerate}

The machine sits in a state, reads an input, looks up the transition,
and moves. That is the entire model. A vending machine is one: state
is ``how much money has been inserted''; alphabet is ``coin'',
``button press''; transitions either accept the coin (raising the
balance) or, if the button is pressed and balance is sufficient,
dispense a snack and return to zero.

A terminal emulator is a state machine. Its states include:

\begin{itemize}[leftmargin=*]
  \item \emph{Ground} --- the default, awaiting the next byte.
  \item \emph{Escape} --- the previous byte was \code{ESC};
    expecting a follow-up.
  \item \emph{CSI} --- inside a Control Sequence; reading parameter
    digits.
  \item \emph{OSC} --- inside an Operating System Command; reading
    a string up to a terminator.
  \item \emph{DCS} --- inside a Device Control String (rare).
\end{itemize}

Its alphabet is ``one byte''. Its transitions read like:

\begin{center}
\small
\begin{tabular}{lll}
\textbf{State} & \textbf{Input byte} & \textbf{Action / next state} \\
\hline
Ground  & printable & emit the glyph; stay in Ground \\
Ground  & \code{0x1B} (ESC) & move to Escape \\
Ground  & \code{0x0A} (LF) & move cursor down; stay in Ground \\
Escape  & \code{[}       & move to CSI \\
Escape  & other         & dispatch a one-byte escape; back to Ground \\
CSI     & digit         & accumulate into current parameter \\
CSI     & \code{;}       & start the next parameter \\
CSI     & letter         & dispatch the sequence; back to Ground \\
\end{tabular}
\end{center}

This is a tiny excerpt of what a real terminal emulator implements
(Paul Williams's well-known VT500 state diagram has a few more rows),
but it is the whole essence. Every byte the emulator receives moves
it through these states; the side effects --- drawing glyphs,
changing colours, moving the cursor --- happen on the transitions.

\maya{}'s input parser is also a state machine, but in reverse: it
reads bytes from the keyboard and assembles them into structured
\code{Key} events. Its states are ``Ground'', ``after ESC'', ``inside
CSI'', ``maybe Alt+key''. The grammar is the same; only the
direction of the byte flow is reversed.

\begin{idea}
The terminal is a small state machine that reads bytes and updates a
grid. Your program is a state machine in the other direction:
structured intent in, bytes out.
\end{idea}

\section{Three things called ``terminal''}

The word \emph{terminal} is overloaded. Pinning down which one we mean
is half the battle. There are three distinct things, often confused.

\paragraph{The hardware terminal.}
A physical device with a keyboard and a screen that connected to a
mainframe over a serial cable. Digital Equipment Corporation's
\textbf{VT100} (released 1978) is the famous one. Nobody owns one
today, but the byte protocol it spoke survives essentially unchanged
on every modern system. When we say a modern emulator ``speaks
ANSI'', we mean it speaks VT100's grandchildren.

\paragraph{The terminal emulator.}
A program on your machine --- iTerm2, kitty, Alacritty, GNOME
Terminal, Konsole, Windows Terminal, xterm --- that pretends to be a
VT100 (or a richer descendant). It owns a window on your desktop,
draws characters into it, and forwards your keystrokes to whatever
program is running ``inside'' it.

\paragraph{The pseudo-terminal, or \emph{pty}.}
A pair of file descriptors in the operating-system kernel that hooks
the emulator up to the program. We will look at it in a moment. The
pty is the part most people never hear named, and yet it is the part
doing the actual work of being a terminal.

When this book says ``\maya{} writes ANSI to the terminal'', it means:
\maya{} calls the operating-system primitive \code{write} on the file
descriptor numbered $1$ (called \code{stdout}); the bytes flow through
the pty; the emulator on the other end reads them and updates the
window.

\subsection{What is a file descriptor?}

If you have not run into the word before: in Unix, every open ``thing
you can read from or write to'' is named by a small integer called a
\textbf{file descriptor}. Three are conventional:

\begin{center}
\small
\begin{tabular}{rl}
\textbf{0} & \code{stdin}  -- bytes the program reads as input. \\
\textbf{1} & \code{stdout} -- bytes the program writes as ordinary output. \\
\textbf{2} & \code{stderr} -- bytes the program writes as error output. \\
\end{tabular}
\end{center}

When you run a program inside a terminal, all three are wired to the
pty. When you write \code{ls > out.txt}, the shell rewires the
program's file descriptor $1$ to a freshly opened file before starting
it. The program is none the wiser: it calls \code{write(1, ...)} as
always, and bytes happen to land in a file. This is the famous Unix
``everything is a file'' design, and it is what makes terminal
programs so flexible.

\section{The pseudo-terminal: a kitchen-pass-through}

A useful mental model for the pty: imagine a kitchen with a
pass-through window. The cook (your program) puts plates of bytes on
the counter; the server (the terminal emulator) picks them up and
delivers them to the diner (the window on your screen). When the
diner says something to the server, the server sets a slip of paper
(keystrokes) on the counter for the cook to pick up. The cook and the
server never see each other; they only exchange bytes via the
counter.

In reality the ``counter'' is a chunk of kernel memory --- a buffer
with two ends. The cook (your program) sees one end as a normal file:
\code{write} to send bytes into it, \code{read} to take bytes out of
it. The server (the emulator) sees the other end the same way. The
kernel shuffles bytes between the two ends, with a few useful
behaviours layered on (input modes, signal handling --- we will get to
those).

\begin{center}
\small
\begin{verbatim}
   +----------+        write(1)         +-----+      read     +----------+
   | program  | ----------------------> |     | ------------> | emulator |
   |          | <---------------------- | pty | <------------ |          |
   +----------+        read(0)          +-----+      write    +----------+
\end{verbatim}
\end{center}

This picture is enough for the rest of the chapter. We do not need to
look at the implementation of the pty itself.

\begin{idea}
The pty is a two-way byte pipe between your program and the terminal
emulator. \code{stdout} goes \emph{to} the emulator; \code{stdin}
comes \emph{from} it.
\end{idea}

\section{The byte language: ASCII plus ANSI}

Now the central question: what bytes does the program send, and what
do they do?

\subsection{Ordinary bytes: ASCII}

Most bytes you write are interpreted literally. The byte \code{0x48}
is the letter \code{H}; \code{0x69} is \code{i}; writing the two-byte
sequence \code{0x48 0x69} prints \code{Hi}. This mapping is
\textbf{ASCII}, a standard from 1963 covering the bytes from $0$ to
$127$. Every English letter, every digit, every common punctuation
mark has a fixed ASCII byte.

Characters outside ASCII (accented letters, CJK, emoji) are sent as
\textbf{UTF-8} sequences of two-to-four bytes per character. Modern
terminals all understand UTF-8; \maya{} produces UTF-8 throughout. We
will return to the question of \emph{cell width} for non-ASCII
characters in the canvas chapter.

\subsubsection*{A quick word on UTF-8}

UTF-8 is the variable-length encoding of Unicode that the world
standardised on. The rule, in one paragraph: code points 0--127 (the
old ASCII range) are encoded as a single byte, exactly as in ASCII;
code points 128--2047 take two bytes; 2048--65535 take three bytes;
the remainder (including emoji and obscure scripts) take four. The
first byte of a multi-byte sequence has a specific high-bit pattern
identifying the length (\code{110xxxxx} for two-byte, \code{1110xxxx}
for three-byte, \code{11110xxx} for four-byte), and subsequent bytes
all begin with \code{10xxxxxx}.

Example. The letter \code{H} is \code{0x48} (one byte). The Greek
letter mu, $\mu$ (U+03BC), takes two: \code{0xCE 0xBC}. The Chinese
character zh\=ong (U+4E2D) takes three: \code{0xE4 0xB8 0xAD}. A
grinning emoji (U+1F600) takes four: \code{0xF0 0x9F 0x98 0x80}.

Why this matters for a TUI: the number of \emph{bytes} in a string
is not the number of \emph{characters}, and the number of
characters is not the number of \emph{cells} the terminal will use
to display it. Three different quantities. \maya{}'s text layer
distinguishes them carefully; you usually do not have to.

\subsection{Control bytes}

A handful of bytes in the range $0$--$31$ have \emph{control}
meanings: they do not print a glyph, they tell the terminal to do
something. The interesting ones:

\begin{center}
\small
\begin{tabular}{rll}
\textbf{Byte (hex)} & \textbf{Name} & \textbf{Effect} \\
\hline
\code{0x07} & BEL & beep / bell \\
\code{0x08} & BS  & backspace --- cursor left one cell \\
\code{0x09} & HT  & horizontal tab --- advance to next tab stop \\
\code{0x0A} & LF  & line feed --- cursor down one row \\
\code{0x0D} & CR  & carriage return --- cursor to column $1$ \\
\code{0x1B} & ESC & escape --- start of an escape sequence \\
\code{0x7F} & DEL & delete --- treated as backspace on input \\
\end{tabular}
\end{center}

You have met most of these. \code{LF} is the newline. \code{CR} is
why old-school progress bars work --- print \code{CR}, then a fresh
line of text, and you overwrite the previous one in place.

The interesting byte is \code{ESC}, hex \code{0x1B}, decimal $27$. It
is the door into the ANSI escape sequence language.

\subsection{Anatomy of an escape sequence}

An ANSI escape sequence has a strict shape:

\begin{center}
\small
\begin{tabular}{l}
\code{ESC} \, \code{[} \, \emph{parameters} \, \emph{final-byte} \\
\end{tabular}
\end{center}

\begin{itemize}[leftmargin=*]
  \item \code{ESC} is the single byte \code{0x1B}.
  \item \code{[} is the literal ASCII byte \code{0x5B}. Together,
    \code{ESC [} (two bytes) form the \textbf{Control Sequence
    Introducer}, abbreviated \textbf{CSI}.
  \item \emph{parameters} is a list of decimal numbers separated by
    semicolons. Each number is written as its ASCII digits ---
    so the number $31$ is two bytes, \code{0x33 0x31}, not the binary
    value $31$.
  \item \emph{final-byte} is a single ASCII letter that says which
    command this sequence is. The letter chooses which family the
    command belongs to.
\end{itemize}

Here is a complete example, byte by byte. Suppose we want to apply
the colour ``bold red'':

\begin{terminalbox}
ESC  [   1   ;   3   1   m
0x1b 0x5b 0x31 0x3b 0x33 0x31 0x6d
\end{terminalbox}

Seven bytes. The first two are CSI. \code{1 ; 3 1} is the parameter
list ``$1$, $31$''. The final byte \code{m} says ``this is an
\textbf{SGR} command --- Select Graphic Rendition''. Parameter $1$
within SGR means ``bold''; $31$ means ``red foreground''. The terminal
emulator parses the seven bytes, updates its internal style for
\emph{future} text, and waits. The next ordinary bytes you send are
drawn bold and red.

If we now write \code{Hi}, then send \code{ESC [ 0 m} to reset, the
full transcript is:

\begin{terminalbox}
\textbackslash x1b[1;31mHi\textbackslash x1b[0m
\end{terminalbox}

This is what \maya{} ultimately produces. Every layer of the library
exists to spare you from writing those bytes by hand.

\begin{idea}
An escape sequence is \code{ESC [}, then a few numbers with
semicolons, then one letter. The letter picks the family; the numbers
parameterise it.
\end{idea}

\subsection{The families}

The final byte categorises the command. The families you will see in
\maya{}'s output, in rough order of frequency:

\begin{center}
\small
\begin{tabular}{lll}
\textbf{Family} & \textbf{Final byte} & \textbf{Purpose} \\
\hline
SGR             & \code{m} & graphic rendition: colour, bold, italic, etc. \\
CUP             & \code{H} or \code{f} & move cursor to (row, column) \\
CUU / CUD       & \code{A} / \code{B} & cursor up / down by $n$ \\
CUF / CUB       & \code{C} / \code{D} & cursor forward / back by $n$ \\
CHA             & \code{G} & cursor to absolute column \\
ED              & \code{J} & erase part of the display \\
EL              & \code{K} & erase part of the line \\
DECSET / DECRST & \code{h} / \code{l} & set / reset private modes (with \code{?}) \\
\end{tabular}
\end{center}

Three additional things to know.

\paragraph{ANSI columns and rows are 1-based.}
\code{ESC [ 5 ; 12 H} moves the cursor to row $5$, column $12$, where
the top-left cell of the screen is $(1, 1)$, \emph{not} $(0, 0)$. This
is a constant source of off-by-one bugs; \maya{} converts to/from
$0$-based internally and \code{move\_to}'s signature is documented as
1-based at the boundary.

\paragraph{Private modes use \code{?}.}
DEC's extensions to ANSI use a question mark inside the parameter list
to distinguish ``private'' parameters that only DEC-style terminals
understand. \code{ESC [ ? 25 h} (show cursor) and
\code{ESC [ ? 1049 h} (enter alternate screen) are private; the plain
SGR \code{ESC [ 1 m} is not.

\paragraph{Numbers default to $1$.}
If you write \code{ESC [ A} (just the family with no parameter), the
parameter is implicitly $1$ --- ``cursor up one row''. Most families
follow this convention; \maya{} explicitly emits the number anyway,
because being explicit avoids surprises across emulators.

\subsection{SGR: colours and attributes}

SGR is the family you will see most. It composes: \code{ESC [ 1 ; 31
; 4 m} applies bold, red, and underline together. Then ordinary text
is drawn with all three until another SGR command changes them.

The parameters you will care about:

\begin{lstlisting}
ESC [ 0 m              // reset every attribute
ESC [ 1 m              // bold
ESC [ 2 m              // dim
ESC [ 3 m              // italic
ESC [ 4 m              // underline
ESC [ 7 m              // inverse (swap fg / bg)
ESC [ 9 m              // strikethrough

ESC [ 22 m             // bold + dim off
ESC [ 23 m             // italic off
ESC [ 24 m             // underline off
ESC [ 27 m             // inverse off
ESC [ 29 m             // strikethrough off

ESC [ 30..37 m         // 3-bit foreground (black..white)
ESC [ 40..47 m         // 3-bit background
ESC [ 90..97 m         // bright fg
ESC [ 100..107 m       // bright bg
ESC [ 39 m             // default fg
ESC [ 49 m             // default bg

ESC [ 38 ; 5 ; N m     // 256-colour fg, palette index N
ESC [ 38 ; 2 ; R ; G ; B m   // 24-bit truecolour fg
ESC [ 48 ; 2 ; R ; G ; B m   // 24-bit truecolour bg
\end{lstlisting}

Modern terminals all support truecolour (24 bits = 16,777,216 colours).
\maya{}'s \code{Fg<R, G, B>} compiles to exactly \code{ESC [ 38 ; 2 ;
R ; G ; B m}. If you read \file{maya/include/maya/terminal/ansi.hpp}
you will find these as \code{constexpr std::string\_view} constants:

\begin{lstlisting}
inline constexpr std::string_view bold      = "\x1b[1m";
inline constexpr std::string_view italic    = "\x1b[3m";
inline constexpr std::string_view underline = "\x1b[4m";
inline constexpr std::string_view reset     = "\x1b[0m";
\end{lstlisting}

\code{constexpr} means the string literally lives in the binary; no
allocation, no construction at startup. We will explain
\code{string\_view} and \code{constexpr} properly in
Chapter~\ref{ch:cpp}.

\subsection{Cursor movement and erasing}

A short worked example: print ``A'' at row $3$, column $5$, then
``B'' at row $3$, column $20$:

\begin{terminalbox}
ESC [ 3 ; 5 H A ESC [ 3 ; 20 H B
\end{terminalbox}

Two CUP sequences sandwich the printable letters. Each is six bytes
(seven, with the two-digit column): the cost of cursor positioning is
exactly the length of the decimal numbers plus four (\code{ESC},
\code{[}, \code{;}, \code{H}). This is why diff-based rendering is so
effective: most cursor moves are tiny.

For erasing, the three useful ED forms:

\begin{lstlisting}
ESC [ 0 J     // erase from cursor to end of screen
ESC [ 1 J     // erase from start of screen to cursor
ESC [ 2 J     // erase entire screen
\end{lstlisting}

And EL, line-erasing:

\begin{lstlisting}
ESC [ 0 K     // erase from cursor to end of line
ESC [ 1 K     // erase from start of line to cursor
ESC [ 2 K     // erase entire line
\end{lstlisting}

The cursor itself does not move when you erase --- the cells under the
cursor change to blank but the position stays. If you want the
classic \emph{clear-screen-and-home} effect, you send both:
\code{ESC [ 2 J ESC [ H} (the \code{H} with no parameters means ``move
to row $1$, column $1$'').

\begin{warning}
Erasing does not reset the colour. A cell ``erased'' to blank keeps
its background colour. If you erased after setting a red background,
you get a red rectangle. To truly blank the screen to the default,
send \code{ESC [ 0 m} (reset colours) \emph{before} erasing.
\end{warning}

\subsection{Modes: cursor, alt-screen, mouse, paste}

A different style of command toggles \emph{modes}: terminal-wide flags
the emulator carries between sequences. These use the DEC private form
with a \code{?}:

\begin{lstlisting}
ESC [ ? 25 l           // hide the cursor
ESC [ ? 25 h           // show the cursor
ESC [ ? 1049 h         // enter alternate screen
ESC [ ? 1049 l         // leave alternate screen
ESC [ ? 1000 h         // mouse press/release events
ESC [ ? 1002 h         // mouse + drag motion events
ESC [ ? 1006 h         // SGR-encoded mouse coordinates (no 223-col limit)
ESC [ ? 2004 h         // bracketed paste (paste arrives wrapped)
ESC [ ? 2026 h         // begin synchronised update
ESC [ ? 2026 l         // end synchronised update
\end{lstlisting}

A few of these deserve names of their own.

\paragraph{The alternate screen.}
When you run \code{vim} or \code{less}, you notice that quitting
restores whatever you had on screen \emph{before} the editor started
--- the editor's contents do not pollute your scrollback. This is the
alternate screen. The emulator keeps two screen buffers; \code{?1049
h} switches to the secondary one (saving the primary), \code{?1049 l}
switches back. \maya{} programs in fullscreen mode emit the enter
sequence at startup and the leave sequence at exit, which is why they
behave the same way.

\paragraph{Bracketed paste.}
When the user pastes a chunk of text, the emulator wraps it with the
sentinels \code{ESC [ 200 \~{}} (before) and \code{ESC [ 201 \~{}}
(after). Without this, a pasted multi-line block looks identical to
the user typing very fast, and a TUI editor cannot tell the
difference. With it, the editor sees: ``a paste begins, here are the
bytes, the paste ends'' --- and can treat the whole block atomically
(no auto-indent per line, for example). \maya{} surfaces this as a
\code{Paste} variant of its \code{Event} type.

\paragraph{Synchronised update.}
A new and very useful mode. When \code{?2026 h} is sent, the terminal
buffers all output until \code{?2026 l}. The frame is then drawn
atomically --- no tearing, no half-painted UI even on a slow connection.
Modern terminals (kitty, WezTerm, foot, Ghostty, recent xterm)
support this; older ones ignore the sequence harmlessly. \maya{}'s
\code{env\_supports\_synchronized\_output} in
\file{terminal/ansi.hpp} detects this from environment variables and
decides whether to wrap each frame in sync brackets.

\subsection{Encoding non-ASCII keys: the Kitty Keyboard Protocol}

Plain ANSI cannot tell \code{Ctrl+M} from \code{Enter} (both arrive as
byte $0x0D$), or \code{Shift+Tab} from a regular escape sequence. This
has been a chronic source of bugs in TUI editors for thirty years.
The \textbf{Kitty Keyboard Protocol} (KKP) and the older
\code{modifyOtherKeys=2} extension fix this by sending fully
disambiguated key reports.

\maya{} opts in on startup with two sequences:

\begin{lstlisting}
ESC [ > 1 u            // push KKP flag 1 ("disambiguate")
ESC [ > 4 ; 2 m        // xterm modifyOtherKeys=2
\end{lstlisting}

and opts out on exit with \code{ESC [ < u} and \code{ESC [ > 4 m}.
Terminals that do not understand the sequences ignore them silently;
the worst case is the older behaviour. We discuss the input side
of this in Chapter~\ref{ch:events}.

\section{Cooked mode, raw mode, and \code{termios}}

Now we turn from output to input. In a normal shell, what happens
when you type a line of text?

The kernel sits between your keyboard and any program reading
\code{stdin}. By default it is in \textbf{canonical} (cooked) mode:

\begin{itemize}[leftmargin=*]
  \item Keystrokes are buffered until you press \code{Enter}, then
    delivered as a whole line.
  \item Backspace edits the buffer in place (the program never sees
    the deletion).
  \item \code{Ctrl-C} is intercepted and sent as a \emph{signal}
    (\code{SIGINT}) rather than as data.
  \item \code{Ctrl-D} on an empty line is interpreted as
    end-of-input.
  \item Each typed character is also \emph{echoed} back to the
    screen automatically.
\end{itemize}

This is exactly the behaviour you want at a shell prompt. It is
exactly the wrong behaviour for a TUI. A TUI wants every keystroke
delivered immediately, byte for byte, with no automatic echo and no
special handling of control characters.

The opposite mode is \textbf{raw mode}, configured via the
\code{termios} structure --- a C \code{struct} the kernel exposes
through \code{tcgetattr} and \code{tcsetattr}.

\subsection{The \code{termios} struct}

\code{termios} has four flag fields plus a small array of special
characters. Each flag field is a bitmask: bits enabled mean different
behaviours active.

\begin{center}
\small
\begin{tabular}{ll}
\textbf{Field} & \textbf{Meaning} \\
\hline
\code{c\_iflag} & input-processing flags \\
\code{c\_oflag} & output-processing flags \\
\code{c\_cflag} & control flags (baud rate, parity --- mostly serial relics) \\
\code{c\_lflag} & local flags (canonical mode, echo, signals) \\
\code{c\_cc[]}  & special characters (which byte is INTR? EOF? \dots) \\
\end{tabular}
\end{center}

To enter raw mode, you turn off a specific list of bits:

\begin{lstlisting}[language=C]
#include <termios.h>
#include <unistd.h>

termios cooked, raw;
tcgetattr(STDIN_FILENO, &cooked);   // save the current settings
raw = cooked;                       // copy, mutate the copy

raw.c_lflag &= ~(ICANON | ECHO | ISIG | IEXTEN);
raw.c_iflag &= ~(IXON | ICRNL | INPCK | ISTRIP);
raw.c_oflag &= ~(OPOST);
raw.c_cc[VMIN]  = 0;     // non-blocking reads
raw.c_cc[VTIME] = 0;
tcsetattr(STDIN_FILENO, TCSAFLUSH, &raw);

// ... TUI runs here ...

tcsetattr(STDIN_FILENO, TCSAFLUSH, &cooked);   // restore
\end{lstlisting}

The bits, decoded:

\begin{itemize}[leftmargin=*]
  \item \code{ICANON} --- canonical mode (line buffering). Off means
    every keystroke comes through immediately.
  \item \code{ECHO} --- echo input. Off means typed characters do not
    appear on screen automatically; the program decides.
  \item \code{ISIG} --- signal generation on \code{Ctrl-C} /
    \code{Ctrl-Z}. Off means those become normal bytes ($0x03$,
    $0x1A$).
  \item \code{IEXTEN} --- implementation-defined extensions
    (\code{Ctrl-V} as ``literal next'', etc.). Off, just to be
    safe.
  \item \code{IXON} --- flow control (\code{Ctrl-S} freezes the
    terminal, \code{Ctrl-Q} unfreezes). Off; you want those bytes
    yourself.
  \item \code{ICRNL} --- translate \code{CR} ($0x0D$) into \code{LF}
    ($0x0A$) on input. Off, so you see exactly what was typed.
  \item \code{INPCK}, \code{ISTRIP} --- parity checking and 7-bit
    stripping. Both relics from serial-line days.
  \item \code{OPOST} --- output processing, mostly translating
    \code{LF} into \code{CR LF}. Off, because you will be writing
    full escape sequences yourself.
  \item \code{VMIN} = 0, \code{VTIME} = 0 --- ``return from
    \code{read} immediately if there is no input''. The TUI loop
    polls; it does not block.
\end{itemize}

\begin{warning}
You must restore \code{cooked} on exit. If the program crashes before
restoring, the user's shell will appear broken --- typed characters
will not echo, \code{Enter} will not commit, \code{Ctrl-C} will not
work. The remedy is the magic command \code{stty sane}, which resets
\code{termios} to sensible defaults. \maya{} uses a RAII guard (see
Chapter~\ref{ch:cpp}) so the destructor restores cooked mode even on
exceptions.
\end{warning}

\subsection{What keystrokes look like in raw mode}

In raw mode, a single keypress arrives as one or more bytes.

\begin{itemize}[leftmargin=*]
  \item Printable keys: one byte, their ASCII value. \code{a} arrives
    as $0x61$.
  \item Control keys: \code{Ctrl-A} is $0x01$, \code{Ctrl-B} is
    $0x02$, \dots, \code{Ctrl-Z} is $0x1A$. (The pattern: clear the
    top three bits of the letter.)
  \item \code{Enter}: $0x0D$ ($CR$).
  \item \code{Tab}: $0x09$ ($HT$).
  \item \code{Backspace}: $0x7F$ ($DEL$).
  \item \code{Escape}: $0x1B$. By itself \emph{or} as the start of a
    sequence.
  \item Function and arrow keys: an escape sequence. \code{Up}
    arrives as the three bytes \code{ESC [ A}; \code{Down} as
    \code{ESC [ B}; \code{Right} \code{ESC [ C}; \code{Left}
    \code{ESC [ D}.
  \item \code{Alt-}\emph{x}: the bytes \code{ESC} then \emph{x}.
\end{itemize}

The clash is obvious: \code{ESC} alone is a real keypress (some
people bind it), but \code{ESC} is \emph{also} the first byte of every
arrow-key sequence. The standard trick is to read \code{ESC},
\emph{wait a few milliseconds}, and if no more bytes arrive, treat it
as a bare \code{Escape}. If they do, parse the sequence. \maya{}'s
input parser (\file{src/terminal/input.cpp}) implements exactly this;
the timeout is around 50 ms by default.

\section{The screen as a grid}

We have looked at bytes; now look at what those bytes update.

The emulator maintains, in its own memory, a 2-D grid of \textbf{cells}.
Each cell holds:

\begin{itemize}[leftmargin=*]
  \item One Unicode \emph{grapheme} (for most languages, one code
    point --- a letter or a digit; for emoji and some scripts, a
    short cluster of code points that displays as one glyph).
  \item A foreground colour and a background colour.
  \item A set of attribute bits: bold, italic, underline, inverse,
    strikethrough, dim.
\end{itemize}

The cursor has a position (row, column) and a visibility flag. A
\textbf{frame} of UI is a complete specification of every cell that
should be on screen.

When the user resizes the window, the emulator sends the program a
notification: on Linux/macOS, the \code{SIGWINCH} signal; on Windows,
a console-resize event. The program re-queries the new size with
\code{ioctl(TIOCGWINSZ)} and re-lays-out its UI. \maya{} does this
inside the runtime; you do not need to call anything.

\subsection{Wide characters}

A subtle wrinkle: not every Unicode character takes \emph{one} cell.

\begin{itemize}[leftmargin=*]
  \item CJK ideographs (U+4E2D \emph{zh\=ong}, U+3042 \emph{hiragana
    a}) take \textbf{two} cells. Their glyph is roughly square; in a
    monospace font each one is twice as wide as a Latin letter.
  \item Many emoji also take two cells (U+1F600 smiley face).
  \item Combining marks --- a \code{COMBINING ACUTE ACCENT} on top
    of an \code{e} to render \emph{\'e} --- take \textbf{zero}
    cells; they decorate the previous cell.
\end{itemize}

If the program writes a wide character but counts it as one cell, the
emulator and the program disagree about the cursor position after,
and the layout drifts. The Unicode Consortium publishes
\file{EastAsianWidth.txt} listing the width of every code point;
\maya{} ships a generated table at
\file{include/maya/text/unicode\_width\_table.hpp} and consults it
when laying out text. You do not call into the table directly ---
\maya{}'s text widget does it for you.

\section{Signals: \code{SIGWINCH}, \code{SIGINT}, and friends}
\label{sec:term:signals}

A \textbf{signal} on Unix is a tiny asynchronous notification the
kernel delivers to a process. The process can install a handler ---
a function that runs when the signal arrives --- or it can use the
default behaviour (often: ``terminate''). Two signals matter to a
TUI.

\paragraph{\code{SIGWINCH} --- window-size change.}
When the user resizes the terminal window, the kernel delivers
\code{SIGWINCH} to whichever process owns the foreground of the pty.
The signal carries no payload --- the process must re-ask the kernel
for the new size. The mechanism is the system call \code{ioctl},
broadly ``do a non-standard thing on a file descriptor''. The
request \code{TIOCGWINSZ} (``terminal IO control: get window size'')
fills a small struct with row and column counts:

\begin{lstlisting}[language=C]
#include <sys/ioctl.h>
struct winsize ws;
ioctl(STDOUT_FILENO, TIOCGWINSZ, &ws);
// ws.ws_row, ws.ws_col, ws.ws_xpixel, ws.ws_ypixel
\end{lstlisting}

\maya{}'s runtime installs a \code{SIGWINCH} handler that simply
sets an atomic flag; the main loop checks the flag, re-queries the
size, and re-runs layout. Doing real work inside a signal handler
is dangerous (most library functions are not safe to call from one);
flagging-then-handling is the standard idiom.

\paragraph{\code{SIGINT} --- interrupt.}
Delivered by the kernel when the user presses \code{Ctrl-C} \emph{in
cooked mode}. In raw mode (which \maya{} sets), \code{Ctrl-C}
becomes the byte \code{0x03}, which the program sees as ordinary
input and may interpret however it likes. \maya{} also installs a
backup \code{SIGINT} handler so that if something forces cooked
mode back on mid-run (a misbehaving sub-process, for example), a
Ctrl-C still cleans up the terminal state before exiting.

\paragraph{\code{SIGTERM}, \code{SIGHUP}, \code{SIGSEGV}.}
\maya{}'s terminal driver installs handlers for all three so that
catastrophic exits still restore the cooked-mode terminal
before the process dies. Without this, killing a TUI with
\code{kill <pid>} or letting it crash would leave the user's shell
in raw mode. RAII is not enough alone; the destructor only runs if
the stack unwinds, and \code{SIGSEGV} does not unwind. The signal
handler is a belt-and-braces backup.

\subsection{What \code{ioctl} actually is}

The ``standard'' file operations are \code{open}, \code{read},
\code{write}, \code{close}. \code{ioctl} (``IO control'') is the
escape hatch for everything else: ``perform request \code{X} on file
descriptor \code{fd}, with these arguments''. Hundreds of requests
exist; each one is essentially a separate API hiding behind a
common name. You meet \code{ioctl} mostly for terminals
(\code{TIOCGWINSZ}), sockets, and device drivers.

In portable code you would prefer dedicated APIs; in practice no
such dedicated API exists for ``get terminal size'', so \code{ioctl}
with \code{TIOCGWINSZ} is the only path.

\section{Mouse reporting, on the wire}
\label{sec:term:mouse}

Mouse input is also delivered as escape sequences. When mouse
reporting is enabled (\code{ESC [ ?1000 h} or one of its richer
variants), every mouse event arrives as a CSI sequence the program
must parse.

The modern format is \textbf{SGR mouse reporting} (\code{?1006 h}):

\begin{center}
\small
\code{ESC [ < button ; column ; row M}~~(press / drag) \\
\code{ESC [ < button ; column ; row m}~~(release)
\end{center}

\code{button} is a small bitfield: low bits encode left/middle/right
(0/1/2); higher bits flag \code{Shift}, \code{Alt}, \code{Ctrl}
modifiers (4, 8, 16); a bit at position 32 marks ``motion event''
(this is a drag, not a click); a bit at position 64 marks wheel
scroll. \code{column} and \code{row} are 1-based.

Example. A left-click at (col 12, row 5):

\begin{terminalbox}
ESC [ < 0 ; 12 ; 5 M
\end{terminalbox}

A Shift-drag of the same button to (15, 5):

\begin{terminalbox}
ESC [ < 36 ; 15 ; 5 M
\end{terminalbox}

(Button code $36 = 0 + 4 \text{ (Shift)} + 32 \text{ (motion)}$.)

A wheel-scroll down at (40, 10):

\begin{terminalbox}
ESC [ < 65 ; 40 ; 10 M
\end{terminalbox}

(Button code $65 = 1 + 64 \text{ (wheel)}$; the low bit on wheel
means ``down''. Wheel up would be 64; left/right wheel use 66/67.)

\maya{}'s event parser decodes these into structured
\code{Mouse} events. The widget layer then routes them by position
to the widget under the cursor.

\subsubsection*{Why the SGR form?}

The old X10 mouse-reporting format (\code{?1000 h}) encoded the
column and row as \emph{single bytes}, with an offset of 32 (so the
byte for column 1 was 33, the byte for column 12 was 44). It broke
on column or row numbers above $223 - 32 = 191$, because the bytes
were only 7-bit-safe and the offsetting overflowed. Modern wide
terminals routinely exceed 200 columns; the SGR form, with decimal
coordinates, has no such limit. \maya{} enables both \code{?1000 h}
and \code{?1006 h} on startup, the latter overrides; emulators that
do not understand \code{?1006 h} fall back gracefully to the older
encoding.

\section{Bracketed paste, on the wire}
\label{sec:term:paste}

When \code{?2004 h} is enabled and the user pastes a block of text
into the terminal, the emulator wraps the pasted bytes in two
sentinel sequences:

\begin{terminalbox}
ESC [ 200 \textasciitilde~ ... pasted bytes ... ESC [ 201 \textasciitilde
\end{terminalbox}

The content between can be arbitrary, including newlines and other
control characters --- the program should treat the whole block
as one atomic event.

Example. Pasting the two lines

\begin{terminalbox}
foo \\
bar
\end{terminalbox}

produces on the wire:

\begin{terminalbox}
ESC [ 200 \textasciitilde~f o o LF b a r ESC [ 201 \textasciitilde
\end{terminalbox}

Without bracketed paste, that arrival is indistinguishable from a
very fast user typing \code{f}, \code{o}, \code{o}, \code{Enter},
\code{b}, \code{a}, \code{r}. With it, an editor knows to disable
auto-indent, treat the whole block as a single undo unit, and so
on. \maya{}'s input parser produces a single \code{Paste} event
carrying the entire content as a \code{std::string}.

\section{The input parser as a state machine}
\label{sec:term:input-fsm}

We sketched the input parser earlier as ``a state machine in the
other direction''. Concretely: it reads bytes from \code{stdin} and
builds structured \code{Event} values. Here is its core loop in
pseudo-code:

\begin{lstlisting}
enum State { Ground, AfterEsc, InCsi };
State s = Ground;
std::string params;          // accumulator for CSI parameter bytes
for (;;) {
    byte b = read_one();
    switch (s) {
      case Ground:
        if (b == 0x1B) { s = AfterEsc; start_esc_timeout(50ms); }
        else            emit_key(b);
        break;
      case AfterEsc:
        if (b == '[')  { s = InCsi; params.clear(); }
        else if (timed_out) { emit_key(Escape); s = Ground; pushback(b); }
        else            { emit_key_alt(b); s = Ground; }
        break;
      case InCsi:
        if (is_param_byte(b)) params.push_back(b);
        else if (is_final_byte(b)) {
            dispatch_csi(params, b);   // -> emit Arrow, Mouse, Paste, ...
            s = Ground;
        }
        break;
    }
}
\end{lstlisting}

A single keystroke can therefore become a single \code{Event}
(\code{A} arrives, \code{KeyKind::Char 'A'} emitted) or take a few
bytes (\code{ESC [ A} becomes \code{KeyKind::Up}) or many
(a paste of a hundred bytes is one \code{Paste} event). The 50 ms
timeout on \code{ESC} is the only piece that is not a pure function
of the byte stream --- it disambiguates a bare \code{Escape} from
an incomplete escape sequence.

\maya{}'s real parser is more elaborate (it handles modified
arrows, the Kitty Keyboard Protocol's CSI-u format,
\code{modifyOtherKeys=2}, function keys, focus events, OSC
responses), but every additional case follows the same state-machine
shape.

\section{Putting it together: a no-library Hello world}

We now know enough to write a complete TUI ``Hello, red world'' with
no library at all. It is a useful exercise: you will see exactly what
\maya{} is doing for you.

\begin{lstlisting}[language=C]
#include <stdio.h>
#include <termios.h>
#include <unistd.h>

int main(void) {
    termios cooked, raw;
    tcgetattr(STDIN_FILENO, &cooked);
    raw = cooked;
    raw.c_lflag &= ~(ICANON | ECHO);
    tcsetattr(STDIN_FILENO, TCSAFLUSH, &raw);

    // 1. Hide the cursor
    // 2. Clear the screen
    // 3. Move to row 1, column 1
    // 4. Set bold red foreground
    // 5. Write "Hi"
    // 6. Reset attributes
    // 7. Show the cursor again
    write(STDOUT_FILENO,
          "\x1b[?25l"          // hide cursor
          "\x1b[2J"            // erase entire screen
          "\x1b[1;1H"          // cursor to (row 1, col 1)
          "\x1b[1;31m"         // bold red foreground
          "Hi"
          "\x1b[0m"            // reset attributes
          "\x1b[?25h",         // show cursor
          /*len=*/ 33);

    char c;
    read(STDIN_FILENO, &c, 1);   // wait for any keypress

    tcsetattr(STDIN_FILENO, TCSAFLUSH, &cooked);
    return 0;
}
\end{lstlisting}

Save it as \file{hello.c}, build with \code{cc hello.c -o hello}, run
\code{./hello}. The screen clears; ``Hi'' appears in bold red at the
top-left; press any key; the screen returns. You have just written a
TUI by hand.

This is also a faithful caricature of what \maya{} does \emph{every
frame}. The whole library is in service of generating exactly this
style of byte stream, only longer, more carefully diffed, and
described to you as a tree of styled elements instead of a literal
string.

\section{A real maya frame, traced byte by byte}
\label{sec:term:trace}

To close the chapter, here is a hand-traced version of one
incremental redraw in a maya program. The screen currently shows
the word \code{Count: 7} on row 5, starting at column 3, in the
default white foreground. The user pressed \code{+}; maya wants the
screen to now read \code{Count: 8}. Only one cell changed --- column
10 of row 5 --- so maya emits only what is needed to update it.

The bytes \maya{} writes:

\begin{terminalbox}
\textbackslash x1b[?2026h   ~ESC [ ?2026 h ~  begin synchronised update \\
\textbackslash x1b[5;10H   ~ESC [ 5;10 H  ~  move cursor to row 5, col 10 \\
\textbackslash x1b[39;49m  ~ESC [ 39;49 m ~  default fg \& bg (style reset) \\
8                       ~~~~~~~~~~~~~~~~~~  the new digit \\
\textbackslash x1b[?2026l  ~ESC [ ?2026 l ~  end synchronised update
\end{terminalbox}

Byte-by-byte that is:

\begin{terminalbox}
1B 5B 3F 32 30 32 36 68 1B 5B 35 3B 31 30 48 1B \\
5B 33 39 3B 34 39 6D 38 1B 5B 3F 32 30 32 36 6C
\end{terminalbox}

32 bytes in total to redraw the change. Compare to what a naive
renderer would do: clear the screen, redraw every cell. For an 80
by 24 terminal that is roughly $80 \times 24 \times 10$ bytes (most
cells take around ten bytes including SGR and the character), or
close to 20\,kB. Two orders of magnitude saved by the diff.

The sync wrapper (\code{?2026 h} \dots \code{?2026 l}) is optional
and omitted by terminals that do not advertise it; \maya{}'s
renderer checks an environment-derived hint
(\code{env\_supports\_synchronized\_output} in
\file{terminal/ansi.hpp}) and elides the wrapper if the terminal
will not honour it.

The cursor positioning (\code{ESC [ 5;10 H}) and the style transition
(\code{ESC [ 39;49 m}) are produced by \maya{}'s
\code{StyleApplier::transition\_to} helper, which writes the minimal
SGR change from the previous cell's style to the new one (in this
case, the previous cell on the row was already styled the same, so
the SGR transition would actually be empty --- we include it here
for illustration).

\begin{idea}
The whole \maya{} rendering pipeline exists to make this incremental
byte stream short and correct. Everything from the canvas to the
SIMD diff to the ANSI builder is in service of: ``write the minimum
bytes needed to update only the cells that actually changed''.
\end{idea}

\begin{recap}
A terminal is a state machine on the far end of a pty. Bytes flowing
into it are mostly literal text, with a small number of control
characters and the ANSI escape language (\code{ESC [} \emph{params}
\emph{letter}) for everything else. Cooked mode lets the kernel
pre-process input; raw mode hands bytes through unchanged. The screen
is a 2-D grid of styled Unicode cells. Every layer above this in
\maya{} speaks in \cpp{} values; this one alone speaks in bytes.
\end{recap}

\section*{Workshop}

\begin{enumerate}[leftmargin=*]
  \item \textbf{Inspect your terminal's flags.} Run \code{stty -a}
    in a normal shell. Find the lines beginning with \code{iflag}
    and \code{lflag}. Identify \code{icanon} and \code{echo} (with
    or without a leading \code{-}, which means ``negated''). Those
    are the bits a TUI clears.

  \item \textbf{Identify maya's constants.} Open
    \file{maya/include/maya/terminal/ansi.hpp} (we read it above).
    Find:
    \begin{itemize}[leftmargin=*]
      \item the constant for entering the alternate screen;
      \item the constant for enabling SGR mouse reporting (1006);
      \item the constant for the Kitty Keyboard Protocol push.
    \end{itemize}
    Cross-check each one against the table earlier in this chapter.

  \item \textbf{Drive a terminal by hand.} In a scratch \cpp{} file:
    \begin{lstlisting}
#include <cstdio>
#include <unistd.h>
int main() {
    std::fputs("\x1b[?1049h", stdout);     // enter alt-screen
    std::fputs("\x1b[1;1HHello", stdout);  // print at top-left
    sleep(2);
    std::fputs("\x1b[?1049l", stdout);     // leave alt-screen
}
    \end{lstlisting}
    Build and run. Observe that your shell is exactly as you left it
    afterwards.

  \item \textbf{Type the no-library example.} Compile and run the C
    program in this chapter. Then mutate it: change \code{31}
    (red) to \code{32} (green), \code{1} (bold) to \code{4}
    (underline). Each change is one digit.

  \item \textbf{Provoke a raw-mode mistake.} Comment out the final
    \code{tcsetattr} that restores cooked mode. Recompile and run.
    Your shell will behave strangely after. Recover by typing
    \code{stty sane}, then \code{Enter} (you may need to type
    blind). You now appreciate the RAII guard.
\end{enumerate}

\chapter{Modern \cpp{} for \maya{}}
\label{ch:cpp}

This chapter is a self-contained course in the slice of \cpp{} that
\maya{} uses. It assumes no \cpp{} background. If you have written
\cpp{}11 or \cpp{}14 it will still be useful, because \maya{} leans
on features that arrived in \cpp{}20, \cpp{}23, and \cpp{}26; some of
the patterns will be new even to experienced \cpp{} programmers.

Read it linearly the first time. The sections build on each other:
later ones (templates, concepts, variants) lean on earlier ones
(values, references, classes). When you finish, you will be able to
read every line of \maya{} source code.

\section{What a \cpp{} program is}
\label{sec:cpp:program}

A \cpp{} program is one or more \textbf{source files} compiled
together into an \textbf{executable}. The smallest possible one fits
on one line, although by convention we spread it out:

\begin{lstlisting}
#include <print>

int main() {
    std::println("Hello, world");
}
\end{lstlisting}

Four things are happening, and each deserves a paragraph.

\paragraph{The preprocessor runs first.}
Before the compiler proper looks at your code, a small program called
the \textbf{preprocessor} does textual substitution. Lines beginning
with \code{\#} are its instructions. \code{\#include <print>}
literally means ``find the file named \file{print} in the standard
library's directories and paste its contents here''. After the
preprocessor finishes, your source file has grown to several
thousand lines, all the standard-library declarations now visible.

\paragraph{The compiler reads what is left.}
The expanded text is parsed, type-checked, and translated into
\textbf{object code} --- machine instructions in a not-yet-runnable
format. Each source file becomes one \emph{object file}, also called
a \textbf{translation unit}. A small project has one; \maya{} has
several hundred.

\paragraph{The linker stitches object files together.}
The \textbf{linker} takes all object files plus any libraries you
depend on and produces a single executable. Its main job is
resolving \textbf{symbols}: when one file uses a function defined in
another file, the linker matches the use to the definition.

\paragraph{The operating system runs the result.}
When you type \code{./hello}, the OS loads the executable into
memory and jumps to a specific named function: \code{main}. The
integer \code{main} returns is the process's \emph{exit status}
(zero = success). The whole program ends when \code{main} returns.

\begin{idea}
The preprocessor pastes text. The compiler turns text into object
files. The linker joins object files into an executable. The OS runs
the executable from \code{main}.
\end{idea}

\subsection{Headers, source files, and the one-definition rule}

A \textbf{header file} (suffix \file{.hpp} or \file{.h}) is a file
you \code{\#include} into other files. It typically contains
\emph{declarations} --- statements about what exists, but not how it
works. A \textbf{source file} (suffix \file{.cpp}) typically contains
\emph{definitions} --- the actual function bodies and data.

\begin{lstlisting}
// math.hpp --- declaration
int add(int a, int b);
\end{lstlisting}

\begin{lstlisting}
// math.cpp --- definition
#include "math.hpp"
int add(int a, int b) { return a + b; }
\end{lstlisting}

\begin{lstlisting}
// main.cpp --- user
#include "math.hpp"
int main() { return add(2, 3); }
\end{lstlisting}

Every use of \code{add} in any source file sees the same declaration
(from \file{math.hpp}), and there is exactly one definition in the
whole program (in \file{math.cpp}). That is the \textbf{One Definition
Rule}, ODR: every entity must have exactly one definition across all
translation units linked together.

\maya{} is a \emph{header-mostly} library. Almost all its code lives
in headers, with most functions and classes defined \code{inline} or
as templates. This is unusual and has trade-offs (compile times can
be long), but it lets the compiler aggressively inline and constant-
fold across the library boundary. We discuss the trade-off in
Chapter~\ref{ch:build}.

\subsection{\code{\#pragma once} and include guards}

If two headers include the same third header, the preprocessor would
paste it twice, producing duplicate declarations. The traditional
fix is an \textbf{include guard}:

\begin{lstlisting}
// math.hpp
#ifndef MATH_HPP
#define MATH_HPP
int add(int a, int b);
#endif
\end{lstlisting}

The second inclusion finds \code{MATH\_HPP} already defined and skips
the contents. Modern code uses the shorter and equivalent:

\begin{lstlisting}
// math.hpp
#pragma once
int add(int a, int b);
\end{lstlisting}

Every file in \maya{} starts with \code{\#pragma once}.

\section{Values, references, pointers, \code{const}}
\label{sec:cpp:memory}

\cpp{} forces you to be precise about \emph{where} a piece of data
lives and \emph{who can change it}. This precision is what makes
\cpp{} fast and what makes beginners stumble. We will walk it slowly.

\subsection{Values}

A \textbf{value} is just data: \code{int x = 42;} creates a variable
\code{x} of type \code{int} holding the value $42$. The variable
owns its storage. When \code{x} goes out of scope (the block ends),
the storage is gone.

\begin{lstlisting}
int x = 42;          // x owns 4 bytes holding 42
int y = x;           // y owns its own 4 bytes, copy of x
y = 100;             // x is still 42
\end{lstlisting}

This is the \textbf{value semantics} model: every assignment is a
copy, every variable is its own thing. The same applies to classes:

\begin{lstlisting}
std::string a = "hello";
std::string b = a;   // b is its own string, copy of a\'s characters
b += " world";
// a == "hello", b == "hello world"
\end{lstlisting}

For a \code{std::string} this is a real copy of every character.
That sounds wasteful and sometimes is; we will see \code{std::move}
shortly as the way to opt out of copies when you no longer need the
original.

\subsection{References}

A \textbf{reference} is an alias for an existing value. It is \emph{not}
a separate variable; it is another name for the same storage.

\begin{lstlisting}
int x = 42;
int& r = x;        // r is an alias for x
r = 100;
// x is now 100
\end{lstlisting}

References must be bound when created and cannot be re-seated.
\code{int\& r;} alone is not legal --- a reference always names
something.

References are the standard way to pass large values into a function
without copying:

\begin{lstlisting}
void print_long(const std::string& s) {
    std::println("{}", s);     // s is the caller\'s string
}                                  // no copy, no destruction
\end{lstlisting}

The \code{const} qualifier on a reference promises ``I will not
modify the thing you let me look at''. \code{const T\&} is the
idiomatic way to receive read-only inputs of any non-trivial type.

\subsection{Pointers}

A \textbf{pointer} is a value that holds an address --- the location
of another value in memory. The syntax \code{T*} means ``pointer to a
\code{T}''. You dereference with \code{*}.

\begin{lstlisting}
int x = 42;
int* p = &x;         // p holds the address of x
*p = 100;            // write through the pointer
// x is now 100
\end{lstlisting}

Unlike references, pointers \emph{can} be re-seated, can be null
(\code{nullptr}), and can point at nothing meaningful (a real source
of bugs). \maya{} avoids raw pointers in user-facing APIs.
Internally, it uses them for things that genuinely model ``maybe
points to something'': linked lists, lookup tables, parent links
in a tree.

When you do see a pointer in \maya{}, prefer the modern alternatives:
\code{std::unique\_ptr<T>} (single owner, automatically destroyed),
\code{std::shared\_ptr<T>} (reference-counted shared owner),
\code{std::optional<T>} (a value that may or may not be present, no
heap allocation), \code{std::span<T>} (a non-owning view of a
contiguous range), \code{std::string\_view} (the same for character
data).

\subsection{\code{const} on values and references}

A \code{const} variable cannot be reassigned:

\begin{lstlisting}
const int n = 5;
// n = 10;          // error: cannot assign to const
\end{lstlisting}

A \code{const} reference promises read-only access:

\begin{lstlisting}
void f(const std::string& s) {
    // s.push_back('!');  // error: s is const here
    std::println("{}", s);
}
\end{lstlisting}

A \code{const} method on a class promises not to modify the object:

\begin{lstlisting}
struct Counter {
    int n = 0;
    int get() const { return n; }   // const method, ok to call on const objects
    void inc()       { ++n; }
};

const Counter c;
c.get();           // ok
// c.inc();        // error: cannot call non-const method on const object
\end{lstlisting}

\code{const}-correctness is a deceptively powerful discipline: it
lets the compiler check at compile time that your read-only paths
really are read-only.

\begin{recap}
A value owns its storage. A reference is an alias; a pointer holds
an address. \code{const} marks read-only at every level: variable,
reference, method.
\end{recap}

\section{Functions, overloads, default arguments}
\label{sec:cpp:functions}

A function takes inputs, may return an output, and may have side
effects:

\begin{lstlisting}
int square(int n) { return n * n; }
void greet(std::string_view name) {
    std::println("hi, {}", name);
}
\end{lstlisting}

The return type goes first, then the name, then the parameters in
parentheses. \code{void} means ``returns nothing''.

\subsection{Parameter passing}

A parameter declared by value (\code{int n}) is a fresh copy. A
parameter declared by reference (\code{const T\&}) shares storage
with the caller's value. The rule of thumb:

\begin{itemize}[leftmargin=*]
  \item Trivial small types (\code{int}, \code{double},
    \code{std::string\_view}, \code{std::span<T>}) --- by value.
    Copying is cheap.
  \item Larger types you only read --- by \code{const} reference.
    No copy.
  \item Types you intend to modify in place --- by non-\code{const}
    reference. The caller sees the change.
  \item Types you want to take ownership of --- by value (and let
    the caller \code{std::move} into the call).
\end{itemize}

\subsection{Overloading}

Multiple functions can share a name if their parameter lists differ.
This is \textbf{overloading}:

\begin{lstlisting}
void log(int n);
void log(double d);
void log(std::string_view s);

log(42);          // calls log(int)
log(3.14);        // calls log(double)
log("hello");     // calls log(string_view)
\end{lstlisting}

The compiler picks the best match. If two are equally good, you get
a compile error and you spell out the choice.

\subsection{Default arguments}

A parameter can have a default:

\begin{lstlisting}
void greet(std::string_view name, std::string_view greeting = "hi") {
    std::println("{}, {}", greeting, name);
}

greet("world");            // "hi, world"
greet("world", "hello");   // "hello, world"
\end{lstlisting}

Defaults are part of the declaration; later overloads of the same
name cannot redeclare them.

\section{Classes, constructors, destructors}
\label{sec:cpp:classes}

A \textbf{class} (introduced with \code{class} or \code{struct} ---
the only difference is the default access level) is a user-defined
type that bundles data and behaviour.

\begin{lstlisting}
struct Point {
    double x;
    double y;

    double length() const {
        return std::sqrt(x*x + y*y);
    }
};

Point p{3.0, 4.0};
std::println("{}", p.length());   // 5
\end{lstlisting}

\code{struct} members are public by default; \code{class} members
are private by default. Most \maya{} types are \code{struct}.

\subsection{Constructors}

A \textbf{constructor} is a special function that runs when an object
is created. It has the same name as the class, no return type, and
may take parameters.

\begin{lstlisting}
struct File {
    FILE* fp;

    explicit File(const char* path) {
        fp = std::fopen(path, "r");
    }
};

File log{"app.log"};   // calls the constructor with "app.log"
\end{lstlisting}

\code{explicit} prevents implicit conversion: without it, you could
write \code{File f = "app.log";} and \cpp{} would silently call the
constructor. Mark single-argument constructors \code{explicit} unless
you genuinely want the implicit form.

\subsection{Member initialiser lists}

The preferred way to set members from a constructor is the
\textbf{initialiser list}, which appears between a colon and the
body:

\begin{lstlisting}
struct File {
    FILE* fp;

    explicit File(const char* path) : fp(std::fopen(path, "r")) {}
};
\end{lstlisting}

Members are initialised in the order they are declared in the class,
\emph{not} in the order they appear in the initialiser list. Modern
compilers warn if you write them in the wrong order; pay attention
to the warning.

\subsection{Destructors}

A \textbf{destructor} is the constructor's counterpart: it runs when
the object is about to be destroyed. Its name is the class name
prefixed with \code{\~{}}.

\begin{lstlisting}
struct File {
    FILE* fp;

    explicit File(const char* path) : fp(std::fopen(path, "r")) {}
    ~File() { if (fp) std::fclose(fp); }
};

void load() {
    File log{"app.log"};   // fopen runs here
    // ... use log.fp ...
}                              // ~File runs here, fclose runs automatically
\end{lstlisting}

The destructor runs at the end of the scope where the object was
created, even if the function returns early, even if an exception is
thrown. This is the foundation of the next idea.

\section{RAII --- the central pattern}
\label{sec:cpp:raii}

\textbf{Resource Acquisition Is Initialization}, or \textbf{RAII}, is
the single most important idiom in \cpp{}. It is the principle that
every resource your program holds --- memory, file descriptor,
terminal mode, mutex lock, network connection --- should be tied to
the lifetime of an object. The constructor acquires; the destructor
releases. You cannot forget to release because the language runs the
destructor for you.

The \code{File} class above is RAII: constructing a \code{File} opens
the file; the destructor closes it. There is no \code{close()}
method to forget to call.

\subsection{A \code{TerminalGuard} for raw mode}

Let us write the RAII type that fixes the bug from
Chapter~\ref{ch:terminal}, where forgetting to restore cooked mode
left the user's shell broken.

\begin{lstlisting}
class TerminalGuard {
    termios cooked_;
public:
    TerminalGuard() {
        tcgetattr(STDIN_FILENO, &cooked_);
        termios raw = cooked_;
        raw.c_lflag &= ~(ICANON | ECHO | ISIG | IEXTEN);
        raw.c_iflag &= ~(IXON | ICRNL | INPCK | ISTRIP);
        raw.c_oflag &= ~(OPOST);
        raw.c_cc[VMIN]  = 0;
        raw.c_cc[VTIME] = 0;
        tcsetattr(STDIN_FILENO, TCSAFLUSH, &raw);
    }

    ~TerminalGuard() {
        tcsetattr(STDIN_FILENO, TCSAFLUSH, &cooked_);
    }

    TerminalGuard(const TerminalGuard&)            = delete;
    TerminalGuard& operator=(const TerminalGuard&) = delete;
};

int main() {
    TerminalGuard guard;   // raw mode on
    // ... run a TUI ...
}                          // raw mode off, automatically
\end{lstlisting}

The two \code{= delete} lines forbid copying. Copying would let two
\code{TerminalGuard} objects fight over the terminal state. We will
see \code{= delete} again shortly.

\maya{}'s actual terminal driver is more elaborate (it also handles
the alt-screen, mouse mode, KKP push, signal restoration on crash),
but the shape is exactly this: a class whose lifetime is the
lifetime of the terminal being in TUI mode.

\begin{idea}
RAII: every resource is owned by an object. Acquire in the
constructor, release in the destructor. Forgetting becomes a
compile-time impossibility, not a runtime risk.
\end{idea}

\subsection{The rule of zero, three, five}

If a class manages a resource, you usually need to define five
special member functions: destructor, copy constructor, copy
assignment, move constructor, move assignment. This is the
\textbf{rule of five}.

If you do not manage any resource yourself --- you only hold other
RAII types as members (\code{std::string}, \code{std::vector},
\code{std::unique\_ptr}) --- you should define \emph{none} of them.
The compiler generates correct versions automatically. This is the
\textbf{rule of zero}, and it is the rule you should follow first.

The \textbf{rule of three} (\cpp{}98) is the pre-move-semantics
version: define destructor, copy constructor, copy assignment
together.

Most \maya{} classes follow the rule of zero. The few that hold raw
resources (the terminal guard, signal handlers, file descriptor
wrappers) follow the rule of five.

\section{\code{auto} and type deduction}
\label{sec:cpp:auto}

Writing out the type of every variable is fine for \code{int} and
\code{double}; it is exhausting for things like
\code{std::vector<std::pair<std::string, Element>>::const\_iterator}.
\cpp{}11 introduced \code{auto}: ``compiler, figure out the type from
the initialiser''.

\begin{lstlisting}
auto x = 42;                  // int
auto pi = 3.14;               // double
auto name = std::string{"hi"};// std::string
auto v = std::vector<int>{1, 2, 3};
for (auto& item : v) ++item;  // item is int&
\end{lstlisting}

Three rules of thumb:

\begin{itemize}[leftmargin=*]
  \item Plain \code{auto} drops references and \code{const}. If you
    write \code{auto x = some\_const\_ref();} you get a fresh copy.
  \item \code{auto\&} keeps references and \code{const}. Use this in
    range-\code{for} loops to avoid copying each element.
  \item \code{const auto\&} is the most common form for read-only
    iteration: \code{for (const auto\& cell : row) \{ ... \}}.
\end{itemize}

In \maya{}, \code{auto} is the default everywhere a type would be
verbose. Spell types out only when the spelling is short or when it
clarifies intent.

\section{Move semantics}
\label{sec:cpp:move}

We saw that assignment copies. For a \code{std::vector<int>} with a
million elements, copying is a million-element allocation plus a
million-element copy. If you do not need the original after, the
copy is pure waste.

Move semantics is the \cpp{}11 feature that lets you say ``transfer
the contents, do not copy''. A move is usually a pointer swap: the
new vector inherits the buffer; the old vector is left empty.

\begin{lstlisting}
std::vector<int> a = {1, 2, 3};
std::vector<int> b = std::move(a);   // O(1) --- b owns the buffer
// a is in a "valid but unspecified" state; do not read it.
\end{lstlisting}

Key points:

\begin{itemize}[leftmargin=*]
  \item \code{std::move(x)} does not actually move anything; it is a
    \emph{cast} that says ``treat \code{x} as a thing whose contents
    may be stolen''. The move actually happens in the constructor
    or assignment operator that receives the rvalue.
  \item After moving from \code{x}, you must not assume \code{x}
    holds anything specific. You can still destroy it (safely) and
    you can still assign a fresh value to it.
  \item Move is an \emph{optimisation}, not a semantic change. A
    correctly written program works whether \code{std::move} is
    present or not; \code{std::move} only changes how fast it runs.
\end{itemize}

The \textbf{rvalue reference} type \code{T\&\&} is how a function
opts in to receiving a move:

\begin{lstlisting}
void take(std::string s);              // by value: caller can move into it
void take_ref(const std::string& s);   // by const ref: read-only
void take_rref(std::string&& s);       // rvalue ref: must be moved in

std::string name = "alice";
take(name);              // copies name into the parameter
take(std::move(name));   // moves name into the parameter
// after this, name is moved-from
\end{lstlisting}

In \maya{} you see \code{std::move} most often in the runtime:
\code{Model} flows through \code{update(Model, Msg)} without being
copied, because \code{update} takes its \code{Model} by value and
the runtime moves the current model into the call.

\begin{warning}
Do not \code{std::move} a local on the return statement (\code{return
std::move(x);}). The compiler already moves return values
automatically; writing \code{std::move} explicitly disables a
separate optimisation called RVO (Return Value Optimisation) that is
usually \emph{better} than move.
\end{warning}

\section{Templates}
\label{sec:cpp:templates}

A \textbf{template} is a recipe the compiler runs at compile time.
You write the recipe once; the compiler bakes a fresh function or
class for each set of arguments you give it.

\subsection{Function templates}

\begin{lstlisting}
template <typename T>
T add(T a, T b) { return a + b; }

int    x = add(1, 2);          // T = int
double y = add(1.5, 2.5);      // T = double
std::string s = add<std::string>("foo", "bar");  // T = std::string
\end{lstlisting}

\code{typename T} introduces a type parameter named \code{T}. The
function body uses \code{T} as if it were a type. The compiler
\emph{instantiates} the template once per distinct \code{T}: there
are separate \code{add<int>}, \code{add<double>}, and
\code{add<std::string>} functions in the final binary.

The template only compiles correctly if \code{T} actually supports
\code{operator+}. Try \code{add(File, File)} and you will get a long
error. Concepts (Section~\ref{sec:cpp:concepts}) fix this.

\subsection{Class templates}

\begin{lstlisting}
template <typename T>
struct Box {
    T value;
    T& get() { return value; }
};

Box<int>    a{42};
Box<double> b{3.14};
Box<std::string> c{"hi"};
\end{lstlisting}

Each instantiation is a different type. \code{Box<int>} and
\code{Box<double>} share no relationship --- you cannot assign one
to the other.

\code{std::vector<T>}, \code{std::optional<T>},
\code{std::unique\_ptr<T>}, \code{std::variant<T, U, ...>} are all
class templates.

\subsection{Multiple template parameters}

\begin{lstlisting}
template <typename K, typename V>
struct Pair {
    K first;
    V second;
};

Pair<std::string, int> ages[] = {{"alice", 30}, {"bob", 25}};
\end{lstlisting}

\subsection{Non-type template parameters (NTTPs)}

The lesser-known but crucial flavour. A template parameter does not
have to be a type --- it can be a \emph{value}, computed at compile
time:

\begin{lstlisting}
template <int N>
struct Array {
    int data[N];
    int size() const { return N; }
};

Array<8>  a;   // data is int[8]
Array<16> b;   // data is int[16]
// a and b are different types
\end{lstlisting}

\code{std::array<T, N>} works this way. The size is part of the
type, so \code{std::array<int, 4>} and \code{std::array<int, 8>} are
distinct types and the size is known at compile time.

The types of NTTPs were limited for a long time --- integers,
pointers, enums. Since \cpp{}20, NTTPs can be \textbf{structural
class types}: any literal type with all-public members and no
user-defined comparison oddities. This is the trick that powers
\maya{}'s DSL.

\subsection{NTTPs of structural type --- the climax}

We want to write \code{t<"Hello">} and have the string become part of
the type. The straight path is to define a small ``string of fixed
length'' type and use it as an NTTP:

\begin{lstlisting}
template <std::size_t N>
struct Str {
    char data[N]{};

    consteval Str(const char (&s)[N]) {
        for (std::size_t i = 0; i < N; ++i) data[i] = s[i];
    }
};
\end{lstlisting}

Line by line:

\begin{itemize}[leftmargin=*]
  \item \code{template <std::size\_t N>}: a class template
    parameterised by an integer.
  \item \code{char data[N]\{\}}: a fixed-size array of \code{N}
    characters, zero-initialised.
  \item \code{consteval Str(const char (\&s)[N])}: a constructor that
    \emph{must} run at compile time (\code{consteval}) and takes a
    \emph{reference to a literal array of \code{N} chars}. The
    array-reference syntax \code{(\&s)[N]} captures the length of
    the literal --- the compiler picks \code{N} for you.
  \item The loop copies the string into \code{data}.
\end{itemize}

Now we can use \code{Str} as an NTTP:

\begin{lstlisting}
template <Str S>
struct Text {
    static constexpr std::string_view text() {
        return {S.data};
    }
};

Text<"Hello"> a;     // S = Str{"Hello\\0"}, baked into the type
Text<"World"> b;     // distinct type
\end{lstlisting}

\code{Text<"Hello">} and \code{Text<"World">} are different types.
The \emph{bytes} of the string are part of the type identity. This
is exactly how \maya{}'s \code{t<"...">} works, and it is the
foundation of the whole compile-time DSL we tour in
Chapter~\ref{ch:dsl}.

\begin{idea}
NTTPs let you encode \emph{values} into the type system. \maya{}
encodes text, colours, paddings, borders, and style flags this way.
Once a value is in a type, the compiler can check rules between
values at compile time.
\end{idea}

\subsection{Templates and the \code{<>} syntax}

When you see \code{Foo<X, Y>} in a \cpp{} expression, the angle
brackets are template-argument syntax, not less-than / greater-than.
The ambiguity is one of the language's papercuts; the compiler
usually figures it out from context, and where it cannot, you write
\code{template} explicitly: \code{x.template foo<int>()}.

You will see \code{t<"Hi">}, \code{Fg<255, 100, 100>}, \code{pad<1>},
\code{border\_<Round>} throughout \maya{}. They are all template
instantiations.

\section{\code{constexpr}, \code{consteval}, \code{constinit}}
\label{sec:cpp:constexpr}

A \textbf{constant expression} is one the compiler can evaluate at
translation time, before the program runs. \cpp{} has three
keywords governing this:

\begin{itemize}[leftmargin=*]
  \item \code{constexpr} on a function or variable: \emph{may} be
    constant-evaluated; otherwise behaves like a normal function or
    variable.
  \item \code{consteval} on a function: \emph{must} be
    constant-evaluated. Calling it at runtime is a compile error.
  \item \code{constinit} on a variable: must be initialised by a
    constant expression, but the variable itself is otherwise
    mutable. Used to forbid the ``static initialisation order
    fiasco''.
\end{itemize}

Example of \code{constexpr}:

\begin{lstlisting}
constexpr int square(int n) { return n * n; }

constexpr int sixteen = square(4);   // evaluated at compile time
int runtime_n = 7;
int forty_nine = square(runtime_n);  // also legal, evaluated at runtime
\end{lstlisting}

Example of \code{consteval}:

\begin{lstlisting}
consteval int square_only_ct(int n) { return n * n; }

constexpr int a = square_only_ct(4);  // ok
int runtime_n = 7;
// int b = square_only_ct(runtime_n); // error: must be constant-evaluated
\end{lstlisting}

\maya{}'s DSL is heavy on \code{constexpr} and \code{consteval}. The
entire chain \code{t<"Hi"> | Bold | Fg<255, 80, 80>} produces a
\code{constexpr} value --- every \code{operator|} along the way is
\code{constexpr}. Conversion to the runtime \code{Element} (which
may allocate) happens only when you call \code{.build()}.

\section{Concepts and \code{requires}}
\label{sec:cpp:concepts}

A template constrains its arguments implicitly: writing \code{a + b}
in a template body means \code{T} must support \code{operator+}, but
the constraint is invisible until you try to instantiate the
template with a wrong type and get a wall of compiler errors.

\cpp{}20 \textbf{concepts} make those constraints explicit and
named:

\begin{lstlisting}
template <typename T>
concept Addable = requires(T a, T b) {
    { a + b } -> std::convertible_to<T>;
};

template <Addable T>
T add(T a, T b) { return a + b; }

add(1, 2);                // ok, int is Addable
// add(File{...}, File{...});  // clean error: File does not satisfy Addable
\end{lstlisting}

The \code{requires} expression lists statements that must be valid
for \code{T}. \code{\{ a + b \} -> std::convertible\_to<T>} means
``\code{a + b} must compile, and its result must convert to
\code{T}''. If a caller passes a \code{T} that does not satisfy
\code{Addable}, the compiler points at the constraint, not at
something seventeen layers deep inside the template body.

\subsection{Maya's concepts}

Maya defines several concepts in \file{core/concepts.hpp}. Three of
the important ones:

\begin{lstlisting}
// (paraphrased from maya source)

template <typename T>
concept Node = requires(const T& t) {
    { t.build() } -> std::convertible_to<Element>;
};

template <typename P>
concept Program = requires(typename P::Model m, typename P::Msg msg) {
    typename P::Model;
    typename P::Msg;
    { P::init() };
    { P::update(m, msg) };
    { P::view(m) }       -> std::convertible_to<Element>;
    { P::subscribe(m) };
};
\end{lstlisting}

\code{Node} says ``has a \code{.build()} that returns something
convertible to \code{Element}''. Every DSL piece satisfies it.
\code{Program} says ``has the four Elm functions with the right
signatures''. \code{run<P>()} requires \code{Program<P>}; if your
program is missing a function, the compiler tells you which one.

\begin{idea}
Concepts are documentation that the compiler checks. They turn
``this template requires the type to do X'' from a comment into a
compile-time predicate.
\end{idea}

\section{\code{std::variant} and \code{std::visit}}
\label{sec:cpp:variant}

A \code{std::variant<A, B, C>} is a value that holds \emph{exactly
one} of \code{A}, \code{B}, or \code{C}. It is the \cpp{} closed sum
type --- the same construct functional languages call an algebraic
data type.

\begin{lstlisting}
#include <variant>
std::variant<int, std::string, double> v;
v = 42;          // holds an int
v = "hello";     // now holds a string
v = 3.14;        // now holds a double
\end{lstlisting}

You inspect a variant with \code{std::visit}, which dispatches to
the overload matching the current alternative. The idiomatic helper
is \code{overload}, a small struct that combines multiple lambdas
into a single callable:

\begin{lstlisting}
template <typename... Fs> struct overload : Fs... { using Fs::operator()...; };
template <typename... Fs> overload(Fs...) -> overload<Fs...>;

std::visit(overload{
    [](int n)               { std::println("int  : {}", n); },
    [](const std::string& s){ std::println("str  : {}", s); },
    [](double d)            { std::println("real : {}", d); },
}, v);
\end{lstlisting}

The \code{overload} template uses two features:

\begin{itemize}[leftmargin=*]
  \item \code{Fs...} is a \textbf{parameter pack} --- a variadic
    template, accepting any number of types.
  \item Inheriting from each one and using \code{using
    Fs::operator()...} brings every lambda's \code{operator()} into
    the combined class. \code{overload\{f, g, h\}} is then a single
    object with three call-operator overloads.
\end{itemize}

This is enough \cpp{} sorcery to write off as a one-liner you copy.
It is in \file{maya/core/overload.hpp}.

\subsection{Maya's variants}

Maya uses \code{std::variant} relentlessly:

\begin{itemize}[leftmargin=*]
  \item \code{Element} is
    \code{std::variant<BoxElement, TextElement, ElementList>}.
  \item \code{Event} is
    \code{std::variant<Key, Mouse, Resize, Paste, ...>}.
  \item \code{Cmd<Msg>} is a variant over every effect the runtime
    knows: \code{None}, \code{Quit}, \code{Batch}, \code{After},
    \code{SetTitle}, \code{WriteClipboard}, \code{Task}, ...
  \item \code{Sub<Msg>} is a variant over event sources.
\end{itemize}

Whenever you see a \code{std::visit} in \maya{}, it is the dispatch
hub for one of these types.

\begin{warning}
If you add a new alternative to a variant and forget to handle it in
a visitor, the visitor stops compiling. This is a feature: it makes
refactors safe. Embrace it; do not paper over it with a generic
fallback.
\end{warning}

\section{\code{std::expected<T, E>}}
\label{sec:cpp:expected}

\cpp{}23's \code{std::expected<T, E>} is a sum type for fallible
operations: either a success of type \code{T} or an error of type
\code{E}. It is \maya{}'s preferred alternative to throwing
exceptions for ``failure is expected, just rare'' cases.

\begin{lstlisting}
#include <expected>
#include <charconv>

std::expected<int, std::string> parse_int(std::string_view s) {
    int n;
    auto [end, ec] = std::from_chars(s.data(), s.data() + s.size(), n);
    if (ec != std::errc{}) return std::unexpected("not an int");
    return n;
}

if (auto r = parse_int("42")) {
    std::println("value: {}", *r);
} else {
    std::println("error: {}", r.error());
}
\end{lstlisting}

The API at a glance:

\begin{itemize}[leftmargin=*]
  \item \code{(bool)r} --- true if success.
  \item \code{*r} or \code{r.value()} --- the success value.
  \item \code{r.error()} --- the error value.
  \item \code{std::unexpected(e)} --- factory for the error case.
\end{itemize}

In \maya{} the error type is usually a small struct:

\begin{lstlisting}
enum class HttpErrorKind { Tls, Timeout, Refused, Status, Cancel };
struct HttpError {
    HttpErrorKind kind;
    int           http_status;
    std::string   detail;
};
using HttpResult = std::expected<Response, HttpError>;
\end{lstlisting}

The caller switches on \code{kind} without parsing strings --- a
proper closed sum of failure modes. This is what we call a
``refined'' error type.

\section{Lambdas and capture}
\label{sec:cpp:lambdas}

A \textbf{lambda} is an anonymous function object. Under the hood,
the compiler synthesises a nameless class with an \code{operator()}
that runs the lambda body.

\begin{lstlisting}
auto add = [](int a, int b) { return a + b; };
int n = add(3, 4);    // 7
\end{lstlisting}

The square brackets at the front are the \textbf{capture list}: they
list names from the surrounding scope the lambda should remember.

\begin{lstlisting}
int x = 10;

auto by_value = [x](int n) { return x + n; };   // copy of x
auto by_ref   = [&x](int n){ return x + n; };   // reference to x

x = 99;
by_value(0);   // 10  --- the lambda has the old copy
by_ref(0);     // 99  --- the lambda sees the new x
\end{lstlisting}

Shorthand captures:

\begin{lstlisting}
[=]        // capture every used name by value
[&]        // capture every used name by reference
[=, &out]  // mostly by value, but out by reference
[&, n]     // mostly by reference, but n by value
\end{lstlisting}

Default captures (\code{[=]} and \code{[\&]}) are convenient but
dangerous in long-lived lambdas; you can easily capture a reference
to a local that has gone out of scope. Prefer explicit capture lists.

\subsection{Init-captures and moving into a lambda}

\cpp{}14 added \textbf{init-captures}: \code{[name = expression]}
introduces a new name local to the lambda with any initialiser. This
is how you capture a move-only type into a lambda:

\begin{lstlisting}
auto buf = std::make_unique<Buffer>();
auto cb  = [buf = std::move(buf)] {
    buf->flush();
};
// outer buf is moved-from; only the lambda owns the buffer.
\end{lstlisting}

\subsection{Storing lambdas: \code{std::function} vs.\ \code{std::move\_only\_function}}

A lambda's type is unique to it; you cannot write the type by hand.
To store a lambda in a class member or container, you need a
type-erased wrapper:

\begin{itemize}[leftmargin=*]
  \item \code{std::function<R(Args...)>} (\cpp{}11) holds any
    callable matching the signature, but the callable must be
    \emph{copyable}. That excludes lambdas capturing move-only state
    (\code{std::unique\_ptr}, sockets, file handles).
  \item \code{std::move\_only\_function<R(Args...)>} (\cpp{}23)
    requires only \emph{movable}. \maya{}'s signal callbacks and
    \code{Cmd<Msg>::Task} use this; effect descriptions routinely
    carry move-only state.
\end{itemize}

\section{Standard-library bits you will see}
\label{sec:cpp:stdlib}

The standard library is large. A non-exhaustive list of names that
appear in \maya{} and what they are:

\begin{itemize}[leftmargin=*]
  \item \code{std::string\_view} --- a non-owning view of a
    contiguous \code{char} sequence. Two pointers under the hood.
    Cheap to pass; does not own the bytes; the underlying data must
    outlive the view.
  \item \code{std::span<T>} --- the same idea for any contiguous
    range of \code{T}. \code{std::span<const Cell>} is how the
    renderer receives a row.
  \item \code{std::array<T, N>} --- a fixed-size array; size is a
    template argument. Same memory layout as \code{T[N]} but with
    proper value semantics and methods.
  \item \code{std::vector<T>} --- the dynamic-array workhorse.
    Heap-allocated buffer that grows.
  \item \code{std::optional<T>} --- either a \code{T} or nothing. Use
    when ``absent'' is a real state but you do not need a richer
    error.
  \item \code{std::format} / \code{std::print} / \code{std::println}
    --- \cpp{}20/23 formatted output, like Python's f-strings.
  \item \code{std::source\_location} --- a value carrying the file,
    line, and function name where it was constructed. \maya{} uses
    it for stable cache identifiers.
  \item \code{std::chrono} --- duration and time-point types.
    \code{std::chrono::milliseconds(50)} reads cleanly.
  \item \code{std::atomic<T>} --- lock-free operations on a value
    shared across threads.
  \item \code{std::jthread} --- a thread that joins on destruction
    (no leaked threads on early return). \cpp{}20.
  \item \code{std::ranges} / \code{std::views} --- lazy, composable
    algorithms over sequences (\code{xs | views::transform(...) |
    views::filter(...)}).
\end{itemize}

Learn the views as you meet them; the ranges library is large and you
do not need it all at once.

\section{Reading real \maya{} code}
\label{sec:cpp:reading}

With the toolkit in hand, let us read a fragment of actual \maya{}.
This is from the style header; it declares the tag types behind
\code{Bold}, \code{Fg<R, G, B>}, and friends.

\begin{lstlisting}
template <CTStyle V>
struct StyTag {
    constexpr operator Style() const noexcept { return V.runtime(); }
};

inline constexpr StyTag<CTStyle::bold()>   Bold{};
inline constexpr StyTag<CTStyle::italic()> Italic{};

template <std::uint8_t R, std::uint8_t G, std::uint8_t B>
inline constexpr StyTag<CTStyle::fg(R, G, B)> Fg{};
\end{lstlisting}

Line by line:

\begin{itemize}[leftmargin=*]
  \item \code{template <CTStyle V>}: \code{StyTag} is a class
    template parameterised by an NTTP \code{V} of type
    \code{CTStyle}. (\code{CTStyle} is a structural class type
    --- another struct defined elsewhere --- carrying a compile-time
    description of a style: which attributes are on, the colours,
    and so on.)
  \item \code{struct StyTag \{ ... \}}: the type has no data
    members. It is a \emph{tag} type --- empty at runtime, distinct
    at compile time. Two \code{StyTag<X>}'s with different \code{X}
    are different types.
  \item \code{constexpr operator Style() const noexcept}: an
    implicit conversion from \code{StyTag<V>} to the runtime
    \code{Style} type. \code{constexpr}, so it can run at compile
    time; \code{noexcept} promises it does not throw. \code{const}
    on the trailing position says the conversion does not modify
    the object.
  \item \code{return V.runtime();}: ask the compile-time style for
    its runtime form. \code{V} is the NTTP, available inside the
    class.
  \item \code{inline constexpr StyTag<CTStyle::bold()> Bold\{\};}: a
    \code{constexpr} variable of type
    \code{StyTag<CTStyle::bold()>}, initialised with default braces.
    \code{inline} on a variable in a header lets it appear in
    multiple translation units without violating ODR. \code{Bold} is
    the value you write in user code.
  \item \code{Fg<R, G, B>}: a \emph{variable template} ---  a
    template that yields a variable, not a function or class. The
    expression \code{Fg<255, 80, 80>} instantiates the template with
    those three integers, producing a unique \code{StyTag<...>}
    value.
\end{itemize}

If you can read that fragment now and feel the pieces in place ---
NTTP, class template, \code{constexpr}, \code{operator T()},
variable template, the implicit conversion that lets \code{Bold}
appear in places that want a \code{Style} --- you can read the rest
of \maya{}'s headers. From here on, the rest of the book is mostly
about what each file \emph{does}, not how its syntax decodes.

\begin{recap}
The \cpp{} \maya{} uses, in one paragraph: templates with type and
non-type parameters; \code{constexpr}/\code{consteval} for compile-
time evaluation; concepts and \code{requires} for named constraints;
\code{std::variant} + \code{std::visit} + \code{overload} for closed
sum types; \code{std::expected} for structured errors; lambdas with
explicit capture and \code{std::move\_only\_function} for storage;
move semantics to avoid copies; RAII to make resources self-managing.
Together they support one discipline: encode invariants in the type
system, make impossible states unrepresentable, defer side effects to
the edge.
\end{recap}

\section*{Workshop}

The exercises here are calibrated for someone who has never used the
feature. If a step takes you fifteen minutes and three Google
searches, that is the right speed.

\begin{enumerate}[leftmargin=*]
  \item \textbf{RAII guard.} Write a class \code{Timer} whose
    constructor records \code{std::chrono::steady\_clock::now()} and
    whose destructor prints the elapsed time. Use it at the top of
    a function; verify the message appears on return.
  \item \textbf{Templates and NTTPs.} Implement a \code{FixedArray<T,
    N>} class template with a member \code{T data[N]}, a
    \code{size()} method returning \code{N}, and \code{operator[]}.
    Instantiate \code{FixedArray<int, 4>} and \code{FixedArray<int,
    8>}; convince yourself they are different types by trying to
    assign one to the other.
  \item \textbf{NTTP of structural type.} Type out the \code{Str<N>}
    template from this chapter. Define a class template \code{Tag<Str
    S>} parameterised by it. Instantiate \code{Tag<"hello">} and
    \code{Tag<"world">} and call a method that returns the string.
  \item \textbf{Concept.} Write a concept \code{Drawable} requiring
    a \code{void draw() const} method. Use it on a template function
    \code{render(const T\& t)}. Make two structs satisfying it, one
    not, and observe the error.
  \item \textbf{Variant.} Define
    \code{using Shape = std::variant<Circle, Square, Triangle>}
    with three trivial structs (each with one field). Write a
    \code{double area(const Shape\&)} using \code{std::visit} and the
    \code{overload} helper.
  \item \textbf{Expected.} Write
    \code{std::expected<int, std::string> safe\_div(int a, int b)}.
    Call it with \code{(10, 2)} and \code{(10, 0)}; print each
    result.
  \item \textbf{Lambda with move capture.} Build a
    \code{std::vector<int>}, capture it into a lambda by move, and
    return the lambda from a function. Call the returned lambda;
    confirm the vector outlives the function it came from.
\end{enumerate}


% ---- Part II: Architecture ---------------------------------------------
\part{Architecture}
\chapter{The Elm loop and algebraic effects}
\label{ch:elm}

\maya{} organises a whole application around a single idea borrowed from
the Elm language: every program is a pure function from the current
state and an incoming event to the next state. Anything that has to
\emph{happen} in the outside world --- a timer firing, a key being
read, a clipboard being written, an HTTP request --- is returned as a
\emph{description of work to do}, not performed inside the function.
This chapter explains the idea, why it is worth the discipline, and
how \maya{} encodes it.

\begin{aside}
If you have written GUIs before, you have written event handlers ---
short functions hung off buttons and inputs that mutate fields, fire
requests, and redraw. They work. They also tend to accumulate, over
months, into a tangle nobody can hold in their head. Elm's response
was counter-intuitive: \emph{remove} the ability to do things in
handlers, force handlers to return plain data, and have a separate
layer perform the actual work. The architecture sounds restrictive
and turns out to be liberating: there is suddenly one place where
state changes (your \code{update}) and one place where the world is
touched (the runtime). This chapter is the long version of that
sentence.
\end{aside}

\section{What does ``pure'' actually mean?}
\label{sec:elm:purity}

The word ``pure'' gets thrown around in programming circles with the
vigour of a moral verdict; it has a precise meaning that has nothing
to do with virtue.

A function is \textbf{pure} when both of these hold:

\begin{enumerate}[leftmargin=*]
  \item \textbf{Same inputs $\Rightarrow$ same output.} Call
    \code{f(3)} twice and you get the same value both times. The
    function does not consult the clock, the filesystem, or a random
    number generator. It does not look at any global variable.
  \item \textbf{No side effects.} Calling the function does nothing
    observable to the rest of the program: no field is mutated, no
    file written, no byte sent on the wire. The only thing the
    caller can observe is the return value.
\end{enumerate}

Both halves are needed. \code{std::time(nullptr)} satisfies (2) ---
it writes nothing --- but fails (1) because it returns a different
value every second. \code{std::printf("hi")} returns the same value
every time but fails (2) by writing to a file descriptor.

Here are some functions you have written, classified:

\begin{center}
\small
\begin{tabular}{lll}
\textbf{Function} & \textbf{Pure?} & \textbf{Why} \\
\hline
\code{int square(int n) \{ return n*n; \}}        & yes & input only, no effects \\
\code{void log(\dots)} writing to stderr           & no  & writes to a file descriptor \\
\code{int next\_id()} with a \code{static} counter & no  & reads/writes hidden state \\
\code{std::sort(begin, end)}                       & no  & mutates the range (side effect) \\
\code{Model update(Model m, Msg ev)} returning a new \code{Model} & yes & inputs in, value out \\
\end{tabular}
\end{center}

\subsection{Why anyone wants this}

Purity buys you four properties that imperative code does not have:

\begin{itemize}[leftmargin=*]
  \item \textbf{Referential transparency.} \code{update(m, Inc\{\})}
    can be replaced anywhere by its value. The compiler can cache it,
    you can paste it into a test, you can move it across threads ---
    none of those changes the meaning of the program.
  \item \textbf{Local reasoning.} To understand a pure function you
    read its body. You do not have to ask ``but what if some other
    thread mutates the field between line 3 and line 4?'' --- there
    are no fields.
  \item \textbf{Trivial testing.} The test is the function call.
    \code{ASSERT(update(\{5\}, Inc\{\}).first.count == 6)} needs no
    mock, no fixture, no setup.
  \item \textbf{Deterministic replay.} A pure function fed the same
    sequence of inputs produces the same sequence of outputs. Bug
    reports become reproducible by attaching a log of \code{Msg}s.
\end{itemize}

The price is structure: state changes are explicit, side effects are
described as values instead of done in place, and small programs
become a little more verbose. For programs that grow past a thousand
lines, the trade is plainly worth it. \maya{} bakes it in.

\begin{idea}
Pure = ``same inputs give same output, and nothing observable
happens on the side''. Pure functions can be tested, cached, moved
between threads, and replayed; impure functions can only be run.
\end{idea}

\subsection{Functional programming, briefly}

The family of languages built around pure functions is called
\textbf{functional}. Haskell, Elm, OCaml, F\#, Clojure, Scala, and
the pure subset of every modern language belong to this family. The
two big ideas are:

\begin{itemize}[leftmargin=*]
  \item \textbf{Data is immutable.} A ``modified'' value is a fresh
    value, not an in-place mutation. Old references keep seeing the
    old data.
  \item \textbf{Functions are first-class.} A function can be
    stored in a variable, passed to another function, returned from
    a function. Lambdas in \cpp{} (Chapter~\ref{ch:cpp}) are this
    feature.
\end{itemize}

\cpp{} is not a functional language --- it is a multi-paradigm one,
which is a polite way of saying ``it lets you do anything'' --- but
it has the pieces required to write functional code: \code{const}
for immutability, lambdas for first-class functions, \code{std::variant}
for sum types, \code{std::move} for cheap functional updates of
large values. \maya{} uses those pieces. It does not enforce purity
the way Haskell does; you can write impure \code{update}s if you
insist. The library is structured so that pure ones cost nothing
extra and impure ones gradually corrupt the benefits.

\begin{funfact}
Elm itself is a \emph{purely functional} language designed in 2012
for browser UIs --- Evan Czaplicki's PhD-era response to React. Its
biggest export is not the language but the pattern: the model /
update / view triangle now turns up under different names in Redux
(JavaScript), \code{bubbletea} (Go), \code{iced} (Rust),
\code{re-frame} (ClojureScript), SwiftUI's \code{Store}\dots and now
\maya{}. The idea outgrew the language.
\end{funfact}

\section{Where state usually lives}

The most common way to write a UI program is to keep state in mutable
fields scattered across whatever objects feel relevant, and to mutate
them from inside event handlers:

\begin{lstlisting}
class Counter {
    int count = 0;
public:
    void on_plus()  { ++count; redraw(); }
    void on_minus() { --count; redraw(); }
};
\end{lstlisting}

This works for tiny programs. It breaks down for the kind of program
where \maya{} is interesting --- programs with several concurrent
sources of events (keys, mouse, timers, async tasks), where you want
to undo, replay, snapshot, or test the logic in isolation. The state
is everywhere; the events come from everywhere; there is no single
place to put a breakpoint.

\section{The Elm shape}

The Elm Architecture says: pick one type to be your state (the
\textbf{Model}), pick one closed sum type to be your events (the
\textbf{Msg}), and write three functions:

\begin{itemize}[leftmargin=*]
  \item \code{init() -> Model} --- the starting state.
  \item \code{update(Model, Msg) -> Model} --- given old state and an
    event, return new state.
  \item \code{view(const Model\&) -> Element} --- given state,
    produce the UI to display.
\end{itemize}

Everything outside those three functions is the \emph{runtime}: it
reads keys, calls \code{update}, calls \code{view}, paints the result,
loops.

Here is the world's smallest example, the counter from
Chapter~\ref{ch:intro} but written in this style:

\begin{lstlisting}
struct Counter {
    struct Model { int count = 0; };
    struct Inc {}; struct Dec {}; struct Quit {};
    using Msg = std::variant<Inc, Dec, Quit>;

    static Model init() { return {}; }

    static Model update(Model m, Msg msg) {
        return std::visit(overload{
            [&](Inc)  -> Model { return {m.count + 1}; },
            [&](Dec)  -> Model { return {m.count - 1}; },
            [&](Quit) -> Model { return m; }
        }, msg);
    }

    static Element view(const Model& m) {
        return v(text(m.count) | Bold,
                 t<"[+/-] q quit"> | Dim) | pad<1> | border_<Round>;
    }
};
\end{lstlisting}

Read it twice. Notice:

\begin{itemize}[leftmargin=*]
  \item \code{Model} is a plain value. No pointers, no inheritance.
  \item \code{Msg} is a closed sum: \emph{exactly} \code{Inc},
    \code{Dec}, or \code{Quit}. If you add a fourth event tomorrow,
    every \code{std::visit} call in the program complains until you
    handle it.
  \item \code{update} is pure: same model + same message $\Rightarrow$
    same result. You can call it from a test with hand-rolled
    \code{Model}s and \code{Msg}s and check equality.
  \item \code{view} is pure too: same model $\Rightarrow$ same UI.
\end{itemize}

\begin{idea}
A pure \code{update} function is the smallest unit of testable
behaviour in a maya program. There is no mock terminal, no fake
event loop, no setup code --- just two values in, one value out.
\end{idea}

\section{But programs need to do things}

Real programs have to perform side effects. Reading a file. Setting
the terminal title. Spawning a subprocess. Quitting. If \code{update}
is pure, it can't do any of those.

The Elm answer is brilliant: have \code{update} \emph{return} the side
effects as values. They are descriptions of work, not work. The
runtime, which has permission to do impure things, interprets them.

This is the technique known as \textbf{algebraic effects} or
\emph{effects as data}. In Haskell it is the \code{IO} monad. In Rust
it is libraries like \code{tokio::Command}. In \maya{} it is
\code{Cmd<Msg>}, which is a \code{std::variant} of all the effects the
runtime knows how to interpret:

\begin{lstlisting}
template <typename Msg>
class Cmd {
public:
    struct None {};
    struct Quit {};
    struct Batch { std::vector<Cmd> cmds; };
    struct After { std::chrono::milliseconds delay; Msg msg; };
    struct SetTitle      { std::string title; };
    struct WriteClipboard { std::string content; };
    struct Task         { std::function<void(std::function<void(Msg)>)> run; };
    struct IsolatedTask { std::function<void(std::function<void(Msg)>)> run; };
    // ...
};
\end{lstlisting}

\code{update} now returns a \code{std::pair<Model, Cmd<Msg>>}:

\begin{lstlisting}
static std::pair<Model, Cmd<Msg>> update(Model m, Msg msg) {
    return std::visit(overload{
        [&](Inc)  -> std::pair<Model, Cmd<Msg>> {
            return {{m.count + 1}, Cmd<Msg>{}};
        },
        [&](Save) -> std::pair<Model, Cmd<Msg>> {
            return {m, Cmd<Msg>::write_clipboard(std::to_string(m.count))};
        },
        [&](Quit) -> std::pair<Model, Cmd<Msg>> {
            return {m, Cmd<Msg>::quit()};
        }
    }, msg);
}
\end{lstlisting}

\code{update} still does no I/O. It \emph{describes} I/O. The runtime
sees \code{WriteClipboard}, calls the OS clipboard API, and the next
call to \code{update} sees the world in a new state.

\subsection{Every \code{Cmd<Msg>} alternative, in detail}
\label{sec:elm:cmd-tour}

Let us walk the actual list. The file is
\file{maya/include/maya/core/cmd.hpp}; every alternative below is one
of the eleven cases of \code{Cmd<Msg>::Variant}.

\paragraph{\code{None}.} The empty command. Returned when
\code{update} has nothing for the runtime to do. The runtime visits
it and immediately returns.

\begin{lstlisting}
return {new_model, Cmd<Msg>::none()};
\end{lstlisting}

\paragraph{\code{Quit}.} Politely asks the runtime to exit its main
loop. The current frame finishes painting, the terminal driver tears
down (raw mode off, alt-screen out, KKP popped, cursor restored),
and \code{run<P>()} returns.

\begin{lstlisting}
return {model, Cmd<Msg>::quit()};
\end{lstlisting}

\paragraph{\code{After}.} Schedule a single \code{Msg} to be
delivered after a delay. The runtime sets a timer; when it fires,
the message is fed into \code{update} as if a user had triggered it.

\begin{lstlisting}
return {model, Cmd<Msg>::after(500ms, HideToast{})};
\end{lstlisting}

Use it for tooltips, transient notifications, debouncing. The timer
is cancelled implicitly if the runtime exits first.

\paragraph{\code{SetTitle}.} Sets the terminal window title.
Internally the runtime emits \code{ESC \textbackslash{}"]0;\dots\textbackslash{}"a}
(an OSC sequence). Outside fullscreen mode the title persists after
the program exits; inside, the alt-screen leave restores the prior
title.

\begin{lstlisting}
return {m, Cmd<Msg>::set_title("todo - " + std::to_string(n) + " items")};
\end{lstlisting}

\paragraph{\code{WriteClipboard}.} Writes a string to the OS
clipboard via OSC 52 (an escape sequence the terminal forwards to
the system clipboard). Works locally and over \code{ssh} if the
emulator opts in.

\paragraph{\code{Task}.} The escape hatch for arbitrary asynchronous
work. You hand the runtime a callable; it runs the callable on a
background worker pool, and the callable is free to dispatch any
number of \code{Msg}s back as results arrive.

\begin{lstlisting}
return {m.with(.loading = true),
        Cmd<Msg>::task([url](auto dispatch) {
            auto response = http_get(url);    // blocking, on a worker
            dispatch(Loaded{std::move(response)});
        })};
\end{lstlisting}

The \code{dispatch} parameter is itself a function: every call
funnels one \code{Msg} back through the runtime, which pushes it
through \code{update} on the main thread. Your task does not touch
the model and does not touch the screen; it only produces messages.

\paragraph{\code{IsolatedTask}.} Like \code{Task} but runs on a
dedicated detached \code{std::thread} instead of the shared pool.
Use it for work that might wedge on a blocking syscall --- a slow
filesystem mount, a network call without a timeout, a sandboxed
subprocess that might never return. A wedged \code{IsolatedTask}
leaks one thread; a wedged \code{Task} permanently consumes one of
the pool's slots and starves all subsequent tasks. The trade is
roughly 100--300\,\textmu{}s of \code{std::thread} construction per
call, paid only for genuinely hang-prone work.

\paragraph{\code{Batch}.} Combine multiple commands. The runtime
performs them all, in unspecified order.

\begin{lstlisting}
return {m, Cmd<Msg>::batch(
    Cmd<Msg>::set_title("saving\dots"),
    Cmd<Msg>::after(2s, ClearStatus{}),
    Cmd<Msg>::task([](auto dispatch) {
        write_to_disk();
        dispatch(Saved{});
    })
)};
\end{lstlisting}

Nested batches are flattened automatically and \code{None}s stripped
(see the \code{batch} static in \file{cmd.hpp}). The two convenient
forms are \code{batch(std::vector<Cmd>\{...\})} and the variadic
\code{batch(c1, c2, c3, ...)}.

\paragraph{\code{CommitScrollback} / \code{CommitScrollbackOverflow} /
\code{ForceRedraw}.} Three rendering-control commands relevant only
to inline-mode programs (Chapter~\ref{ch:inline}). Outside that
context, do not touch them. We will return to them when we cover
the inline renderer.

\subsection{Mapping \code{Cmd<Msg>} between message types}
\label{sec:elm:cmd-map}

Large programs split into components, each with their own
\code{Msg} type. A child component returns \code{Cmd<ChildMsg>}; the
parent's \code{update} returns \code{Cmd<ParentMsg>}. The
conversion is the \code{map} method:

\begin{lstlisting}
Cmd<ChildMsg>  c1 = Child::update(child_model, child_msg).second;
Cmd<ParentMsg> c2 = c1.map([](ChildMsg m) -> ParentMsg {
    return ChildEvent{std::move(m)};
});
\end{lstlisting}

\code{map} threads the conversion function through every alternative:
a \code{Cmd<ChildMsg>::Task} becomes a \code{Cmd<ParentMsg>::Task}
that wraps each \code{ChildMsg} it would dispatch with the
conversion. A \code{Quit} stays a \code{Quit}. A \code{Batch}
recurses element-wise. Read the implementation in \file{cmd.hpp};
it is one \code{std::visit} per alternative.

In category-theory terms, \code{Cmd} is a functor parameterised by
the message type, and \code{map} is the functor's lift. You do not
need to know that label; the trick is that you can refactor a
standalone component into a child of a parent program without
rewriting any of its effect-emitting code --- the parent's update
wraps the child's \code{Cmd} with \code{map}.

\section{Subscriptions: the other half}

\code{Cmd} describes effects the program \emph{starts}. There is also
the symmetric notion: events the program \emph{receives}. Keystrokes,
mouse moves, timers, async results --- these are sources of
\code{Msg}s. \maya{} models them with \code{Sub<Msg>}:

\begin{lstlisting}
template <typename Msg>
class Sub {
public:
    struct Keys     { std::function<std::optional<Msg>(Key)>   f; };
    struct Mouse    { std::function<std::optional<Msg>(Mouse)> f; };
    struct Tick     { std::chrono::milliseconds period; std::function<Msg()> f; };
    // ...
};
\end{lstlisting}

\code{subscribe(const Model\&) -> Sub<Msg>} is the fourth program
function. Like the others, it is pure: given a model, declare which
event sources are active right now. The runtime samples the
subscriptions, reads from those sources, and feeds the \code{Msg}s
back into \code{update}.

\subsection{Why subscriptions are declarative}
\label{sec:elm:why-declarative}

In an imperative GUI, you wire event sources in setup code:
``\code{button.on\_click(\dots)}, \code{timer.start(16ms)}, on
shutdown remember to \code{timer.stop()} and \code{button.disconnect()}
and don't leak the handler''. The wiring is stateful, and forgetting
to unwind it leaks event sources --- a category of bug every GUI
framework has seen.

Maya's \code{subscribe} is different: it is a pure function the
runtime calls on every frame. The returned \code{Sub} value
\emph{is} the current declaration of what event sources are active.
The runtime diffs the new \code{Sub} against the previous one and
takes care of starting and stopping the underlying listeners. You
do not call \code{start} or \code{stop}; you describe the desired
state, and the runtime makes the world match.

The model is React's hooks rules in a different language: declare
what you want, the framework reconciles. The trick is that ``what
you want'' is a value of an ordinary closed sum type, which means
the set of legal answers is bounded and visible.

\subsection{Every \code{Sub<Msg>} alternative}
\label{sec:elm:sub-tour}

From \file{maya/include/maya/app/sub.hpp}:

\paragraph{\code{None}.} Subscribe to nothing. The empty sub.

\paragraph{\code{OnKey}.} Subscribe to keyboard events. The filter
returns \code{std::optional<Msg>} --- \code{nullopt} to ignore the
event, a \code{Msg} to forward it to \code{update}. The filter is
called for every key press; only keys the application cares about
become messages.

\begin{lstlisting}
Sub<Msg>::on_key([](const KeyEvent& k) -> std::optional<Msg> {
    if (key_is(k, 'q')) return Quit{};
    if (key_is(k, '+')) return Inc{};
    return std::nullopt;          // ignore
});
\end{lstlisting}

\paragraph{\code{OnMouse}.} Same pattern, for mouse events: press,
release, drag, scroll. The filter is called for each event. See
Chapter~\ref{ch:events} for the structure of \code{MouseEvent}.

\paragraph{\code{OnResize}.} The window changed size. The function
receives the new \code{Size\{w, h\}}; it returns a \code{Msg} (not
optional --- you must produce some message, even a no-op one). The
runtime relays out the screen automatically after a resize; the
message is for application logic that depends on size (collapsing a
sidebar below 60 columns, for example).

\paragraph{\code{OnPaste}.} The user pasted multi-byte text
(bracketed paste; Chapter~\ref{ch:terminal}). The function receives
the full pasted string at once, returns a \code{Msg}. Use it where
you want paste to be one atomic event rather than a flood of
synthetic \code{Char} keys.

\paragraph{\code{Every}.} A timer. Every \code{interval} milliseconds
the runtime delivers the same \code{msg} to \code{update}. Use this
for clocks, polling, status refreshes. The interval is honoured
best-effort; missed ticks are coalesced rather than queued.

\paragraph{\code{AnimationFrame}.} A 60\,Hz tick. The runtime
delivers \code{msg} on every animation frame for as long as the sub
is returned. Drop it from the sub tree (do not include it during
the next \code{subscribe} call) to stop the animation. The cost
when no \code{AnimationFrame} sub is present is zero --- idle frames
are not rendered.

\paragraph{\code{Batch}.} Combine multiple subs. Order of activation
is unspecified; ordering of delivery follows the underlying event
sources.

\subsection{Subscribe is pure --- and that matters}

A subtle benefit of \code{subscribe} being pure: it can react to the
model.

\begin{lstlisting}
static Sub<Msg> subscribe(const Model& m) {
    auto keys = Sub<Msg>::on_key(/* always-on filter */);
    if (m.animating) {
        return Sub<Msg>::batch(keys,
            Sub<Msg>::on_animation_frame(Tick{}));
    }
    return keys;
}
\end{lstlisting}

While \code{m.animating} is true, the runtime starts a 60\,Hz tick.
When the model sets \code{animating} back to \code{false}, the next
\code{subscribe} call returns no \code{AnimationFrame}; the runtime
stops the tick. You never wrote a \code{stop()} call. The diffing
is the framework's job.

This matters for things like spinners, transient progress bars,
short-lived polling: subscribe while the work is pending, drop the
sub when it's done. Imperative code requires a flag, a cleanup path,
and a bug.

\begin{idea}
\code{subscribe(m)} returns a \emph{declaration} of which event
sources are wanted right now. The runtime diffs against the prior
declaration and starts / stops listeners. You never write
\code{stop()}.
\end{idea}

A complete Elm-style \maya{} program is therefore a struct with four
static functions:

\begin{center}
\small
\begin{tabular}{lll}
\code{init}      & $\;\to\;$ & \code{(Model, Cmd<Msg>)} \\
\code{update}    & \code{(Model, Msg)} $\;\to\;$ & \code{(Model, Cmd<Msg>)} \\
\code{view}      & \code{(const Model\&)} $\;\to\;$ & \code{Element} \\
\code{subscribe} & \code{(const Model\&)} $\;\to\;$ & \code{Sub<Msg>} \\
\end{tabular}
\end{center}

This four-tuple is the \code{Program} concept:

\begin{lstlisting}
template <typename P>
concept Program = requires(typename P::Model m, typename P::Msg msg) {
    typename P::Model;
    typename P::Msg;
    { P::init() }            -> std::convertible_to<std::pair<typename P::Model, Cmd<typename P::Msg>>>;
    { P::update(m, msg) }    -> std::convertible_to<std::pair<typename P::Model, Cmd<typename P::Msg>>>;
    { P::view(m) }           -> std::convertible_to<Element>;
    { P::subscribe(m) }      -> std::convertible_to<Sub<typename P::Msg>>;
};
\end{lstlisting}

\code{run<P>()} requires \code{Program<P>}; if your program is missing
a function or has a wrong signature, the compiler tells you which one
and why.

\section{Why this is worth the effort}

A pure-functional core has properties that an imperative one doesn't:

\begin{itemize}[leftmargin=*]
  \item \textbf{Testability without mocks.} You write
    \code{ASSERT(update(\{5\}, Inc\{\}).first == Model\{6\})}. No
    terminal, no event loop, no fixtures.
  \item \textbf{Determinism.} Given a recorded sequence of
    \code{Msg}s, you can reproduce any rendered frame. \maya{}
    leverages this for debugging.
  \item \textbf{Time travel.} Storing a list of past models is
    storing a list of plain values --- snapshot, replay, undo.
  \item \textbf{Local reasoning.} A bug in \code{update} cannot
    corrupt the renderer; a bug in the renderer cannot mutate the
    model. The interfaces between layers are values, not method
    calls.
\end{itemize}

The cost is some up-front structure. The four-function skeleton feels
heavy for a 30-line program; \maya{} provides simpler entry points
(\code{run(event\_fn, render\_fn)}) for those cases. For anything that
will live longer than an afternoon, the Elm shape is the right shape.

\section{The runtime loop, sketched}

The \code{run<P>()} runtime is one loop. Stripped of error handling
and the inline-mode branch, it looks like this:

\begin{lstlisting}
template <Program P>
void run(RunConfig cfg) {
    auto [model, cmd0] = P::init();
    Renderer r{cfg};
    interpret(cmd0);                          // perform initial effects

    while (!quit_requested()) {
        auto subs = P::subscribe(model);       // declare event sources
        Msg msg   = wait_for(subs);            // block until something arrives

        auto [next_model, next_cmd] = P::update(std::move(model), msg);
        model = std::move(next_model);
        interpret(next_cmd);                   // perform effects from update

        r.paint(P::view(model));               // re-render
    }
}
\end{lstlisting}

That is the entire shape. Everything else in \maya{}'s architecture
chapter is an embellishment: the canvas double-buffer, the SIMD diff,
the inline-mode renderer, the BG thread pool that interprets
\code{Task} commands. All in service of this loop.

\begin{theory}
\textbf{Algebraic effects} is a programming-language idea: separate
\emph{describing} what should happen from \emph{deciding how}. The
function returns a value of an effect algebra; an interpreter chooses
how to run it. The same description can be run for real, run in a
test (with a mock interpreter), or replayed from a log. Elm's
\code{Cmd} and \code{Sub} are a small, practical algebra of effects.
\end{theory}

\subsection{The loop, traced one frame at a time}
\label{sec:elm:trace}

Let's walk what happens between two consecutive frames of a running
program. Suppose the model has \code{count = 7} and the user has
just pressed \code{+}.

\paragraph{Step 0.} The runtime has already painted the previous
frame and is sitting in \code{wait\_for(subs)}. Its standard input is
open and watching for bytes.

\paragraph{Step 1.} The byte \code{0x2B} (ASCII \code{+}) arrives.
The runtime's input parser (Chapter~\ref{ch:terminal}) recognises a
plain printable key and produces \code{KeyEvent\{kind = Char, char =
'+'\}}.

\paragraph{Step 2.} The runtime walks the current \code{Sub} tree
looking for an \code{OnKey} filter. It finds one, calls it with the
key event, gets back \code{std::optional<Msg>(Inc\{\})}.

\paragraph{Step 3.} The runtime calls \code{P::update(std::move(model),
Inc\{\})}. The model is moved (not copied) into \code{update}. The
function returns \code{\{Model\{8\}, Cmd<Msg>::none()\}}.

\paragraph{Step 4.} The runtime moves the returned model into its
own \code{model} variable. It interprets the \code{Cmd}: \code{None}
means nothing to do.

\paragraph{Step 5.} The runtime calls \code{P::view(model)}. The
view function returns an \code{Element} tree (Chapter~\ref{ch:elements})
describing the new screen.

\paragraph{Step 6.} The renderer lays the tree out
(Chapter~\ref{ch:layout}), rasterises it into the back canvas
(Chapter~\ref{ch:canvas}), diffs against the front canvas
(Chapter~\ref{ch:simd}), and emits exactly the bytes for the cells
that changed (Chapter~\ref{ch:ansi}). The front and back canvases
swap.

\paragraph{Step 7.} The runtime calls \code{P::subscribe(model)} to
learn what event sources to watch this frame, diffs against the
prior subscription, and re-enters \code{wait\_for(subs)}.

The whole cycle is the loop body. Every keystroke, every timer tick,
every async result re-runs steps 1--7. There is no other entry
point into your program's logic.

\section{Testing the model in isolation}
\label{sec:elm:tests}

The reward for keeping \code{update} pure is that you can test the
business logic without spinning up a terminal at all. There is no
mock event loop, no fake renderer, no setup of fixtures. The test is
a function call.

\begin{lstlisting}
#include <cassert>
#include "counter.hpp"

int main() {
    using M = Counter::Model;
    auto [m0, _0] = Counter::init();
    assert(m0.count == 0);

    auto [m1, _1] = Counter::update(m0, Counter::Inc{});
    auto [m2, _2] = Counter::update(m1, Counter::Inc{});
    auto [m3, _3] = Counter::update(m2, Counter::Inc{});
    assert(m3.count == 3);

    auto [m4, _4] = Counter::update(m3, Counter::Dec{});
    assert(m4.count == 2);

    auto [m5, c5] = Counter::update(m4, Counter::Quit{});
    assert(c5.is_quit());
}
\end{lstlisting}

No terminal is opened. No bytes are written. The program prints
nothing and either succeeds or aborts on a failed assertion. Compile
it with whatever \cpp{} testing framework you like ---
\code{Catch2}, \code{doctest}, GoogleTest, or none at all
(\code{assert} is fine for small programs).

\subsection{Testing effects}

What if \code{update} returns a non-trivial \code{Cmd}? You inspect
the returned variant:

\begin{lstlisting}
auto [m, cmd] = MyApp::update(m_in, Save{});
assert(std::holds_alternative<Cmd<Msg>::WriteClipboard>(cmd.inner));
auto& w = std::get<Cmd<Msg>::WriteClipboard>(cmd.inner);
assert(w.content == "expected text");
\end{lstlisting}

The \code{Cmd} is a value; you do not have to perform it to check it.
This is the deepest practical benefit of effects-as-data: the test
asks ``what would the program ask the runtime to do?'' without
asking the runtime to do anything.

For \code{Task} commands the situation is trickier --- the inside of
the task is a \code{std::function} that you would normally have to
run to find out what messages it produces. Two practical paths:

\begin{enumerate}[leftmargin=*]
  \item Make the task call a small helper function (e.g.
    \code{fetch\_user(int)}) instead of inline code, and test the
    helper separately.
  \item Run the task synchronously in a test: hand it a lambda that
    records every \code{Msg} into a vector, then assert on the
    vector.
\end{enumerate}

\begin{lstlisting}
std::vector<Msg> dispatched;
auto run_task = [&](auto& task_cmd) {
    task_cmd.run([&](Msg m) { dispatched.push_back(std::move(m)); });
};
\end{lstlisting}

In maya's own test suite, both patterns appear. The library's
internal HTTP layer factors out \code{http\_get\_blocking} so the
\code{Cmd::task} wrapping it can be tested independently.

\subsection{Property tests}

Because \code{update} is pure, property-based testing is a natural
fit. \emph{For every model \code{m} and every \code{Inc} message,
the resulting count is one more than \code{m.count}.} That is one
\code{rapidcheck} property:

\begin{lstlisting}
rc::check("inc increments by one", [](Counter::Model m) {
    auto [m1, _] = Counter::update(m, Counter::Inc{});
    RC_ASSERT(m1.count == m.count + 1);
});
\end{lstlisting}

The property generator produces hundreds of random models and
verifies the invariant on each. Imperative GUI code typically cannot
be property-tested at all --- the state is too entangled.

\section{Replay, snapshot, time travel}
\label{sec:elm:replay}

If you record every \code{Msg} that arrives, you can reproduce any
state the program ever reached by replaying the log into
\code{update}, starting from \code{init}. This is not a research
property; it is the natural consequence of \code{update} being pure.

A skeletal recorder, wrapping any \code{Program}:

\begin{lstlisting}
template <Program P>
struct Recorded {
    using Model = std::pair<typename P::Model, std::vector<typename P::Msg>>;
    using Msg   = typename P::Msg;

    static auto init() -> std::pair<Model, Cmd<Msg>> {
        auto [m, c] = P::init();
        return {{std::move(m), {}}, c};
    }
    static auto update(Model state, Msg msg) -> std::pair<Model, Cmd<Msg>> {
        auto& [inner, log] = state;
        log.push_back(msg);                  // record
        auto [m1, c] = P::update(std::move(inner), msg);
        return {{std::move(m1), std::move(log)}, c};
    }
    static Element view(const Model& state) { return P::view(state.first); }
    static Sub<Msg> subscribe(const Model& state) {
        return P::subscribe(state.first);
    }
};
\end{lstlisting}

The recorded log is just \code{std::vector<Msg>}. To replay, fold
\code{update} over the log starting from a fresh \code{init}:

\begin{lstlisting}
auto [m, _] = P::init();
for (auto& msg : recorded_log) {
    auto [next, _] = P::update(std::move(m), msg);
    m = std::move(next);
}
// m is now exactly what the user saw at the end of the recorded session
\end{lstlisting}

This is the heart of every ``undo'' implementation: store enough
messages to rewind, replay the prefix. \emph{Redux DevTools} does the
same for JavaScript apps; Elm's debugger likewise. Once your update
is pure, you get this for free.

\begin{idea}
Log the \code{Msg}s. Replay through \code{update}. Get the same
model back, every time. This is undo, snapshot, debug-record, and
deterministic bug repro --- one fact about pure functions, four
product features.
\end{idea}

\section{Splitting a big program into components}
\label{sec:elm:components}

A single \code{Model} / \code{Msg} pair works fine up to a few
thousand lines. Beyond that, you want to split: a child component
has its own \code{Model} and \code{Msg}, and the parent embeds them.

The machinery is mechanical:

\begin{lstlisting}
struct Search {
    struct Model { std::string query; std::vector<std::string> results; };
    struct Type { char c; };
    struct Backspace {};
    struct ResultsArrived { std::vector<std::string> r; };
    using Msg = std::variant<Type, Backspace, ResultsArrived>;

    static auto update(Model, Msg) -> std::pair<Model, Cmd<Msg>>;
    static Element view(const Model&);
    static Sub<Msg> subscribe(const Model&);
};

struct App {
    struct Model {
        Search::Model search;
        // ... other components ...
    };
    struct SearchEvent { Search::Msg child; };
    struct GlobalQuit {};
    using Msg = std::variant<SearchEvent, GlobalQuit>;

    static auto update(Model m, Msg msg) -> std::pair<Model, Cmd<Msg>> {
        return std::visit(overload{
            [&](GlobalQuit) -> std::pair<Model, Cmd<Msg>> {
                return {m, Cmd<Msg>::quit()};
            },
            [&](SearchEvent ev) -> std::pair<Model, Cmd<Msg>> {
                auto [child_model, child_cmd] =
                    Search::update(std::move(m.search), ev.child);
                m.search = std::move(child_model);
                auto parent_cmd = child_cmd.map([](Search::Msg cm) {
                    return Msg{SearchEvent{std::move(cm)}};
                });
                return {std::move(m), std::move(parent_cmd)};
            },
        }, msg);
    }
};
\end{lstlisting}

What to notice:

\begin{itemize}[leftmargin=*]
  \item The child's \code{Model} is a field of the parent's
    \code{Model}. No pointers, no inheritance.
  \item The child's \code{Msg} is wrapped by a parent message
    constructor (\code{SearchEvent}) into the parent's \code{Msg}.
  \item When the parent's \code{update} sees a \code{SearchEvent},
    it forwards to the child's \code{update}, getting back
    \code{(ChildModel, Cmd<ChildMsg>)}.
  \item The new child model is written back into the parent model.
    The child \code{Cmd} is mapped into a parent \code{Cmd} with
    \code{.map(\dots)}.
\end{itemize}

\code{subscribe} composes the same way: the parent's
\code{subscribe} calls the child's, wraps the result with
\code{Sub<>::map(\dots)}, and returns it. \code{view} embeds the
child's view as a subtree.

This is the entire technique for scaling Elm-style code. There is no
``component framework''; it is plain \cpp{} types composed by
hand. The compiler keeps everyone honest about which messages can
appear where.

\subsection{When to split, when to keep it flat}

Resist the urge to split too early. A flat 2,000-line program with
clear \code{Msg} categories is easier to read than three components
wrapped in two layers of \code{map} indirection. Split when:

\begin{itemize}[leftmargin=*]
  \item A subtree of the UI is genuinely reusable (a date picker,
    an autocomplete) and you want to drop it into multiple
    programs.
  \item The \code{Msg} variant has more than 20 alternatives and the
    \code{update} function visibly groups into clusters.
  \item A subteam owns a feature and you want their changes to land
    in one file, not in a global \code{Msg}.
\end{itemize}

Otherwise, flat is fine. Maya's own programs in \code{agentty} are
mostly flat with a few well-isolated child components for
specifically reusable pieces.

\section{A worked refactor: from script to component}
\label{sec:elm:worked-refactor}

For concreteness, here is a small but realistic program in both the
quick-tool form and the Program form. The program is a to-do list:
the user types items into a draft buffer, presses \code{Enter} to
commit, presses \code{space} to toggle the selected item's done
state, \code{d} to delete, arrow keys to navigate, \code{q} to quit.

First the imperative sketch, the kind of code you might write in an
afternoon and regret later:

\begin{lstlisting}
struct Item { std::string text; bool done = false; };

int main() {
    std::vector<Item> items;
    std::string draft;
    int selected = 0;

    run({.title = "todo"},
        [&](const Event& ev) {
            if (key(ev, 'q')) return false;
            if (key(ev, '\n')) {
                if (!draft.empty()) {
                    items.push_back({draft, false});
                    draft.clear();
                }
            } else if (key(ev, ' ')) {
                if (!items.empty()) items[selected].done = !items[selected].done;
            } else if (key(ev, 'd')) {
                if (!items.empty()) {
                    items.erase(items.begin() + selected);
                    if (selected >= (int)items.size())
                        selected = std::max(0, (int)items.size() - 1);
                }
            } else if (ev.is_key(KeyKind::Up))   selected = std::max(0, selected - 1);
            else if   (ev.is_key(KeyKind::Down)) selected = std::min((int)items.size() - 1,
                                                                    selected + 1);
            else if   (ev.is_key(KeyKind::Backspace)) {
                if (!draft.empty()) draft.pop_back();
            } else if (ev.is_char()) draft.push_back(ev.character());
            return true;
        },
        [&] {
            // ... build the view from items, draft, selected ...
        });
}
\end{lstlisting}

Works. Two issues: state lives in main's stack frame (no test can
reach it without launching the program), and the event-handler
branches are not enumerable --- adding a feature means adding
another \code{if} and hoping you covered every case.

Now the same program as a \code{Program}:

\paragraph{Step 1: define the \code{Model}.}

\begin{lstlisting}
struct Todos {
    struct Item { std::string text; bool done = false; };
    struct Model {
        std::vector<Item> items;
        std::string       draft;
        int               selected = 0;
    };
};
\end{lstlisting}

A plain value type. Equality is comparison of fields. Copying is a
vector copy plus two strings.

\paragraph{Step 2: enumerate the \code{Msg}.}

\begin{lstlisting}
    struct AppendDraft   { char c; };
    struct DeleteDraft   {};
    struct SubmitDraft   {};
    struct ToggleDone    {};
    struct DeleteCurrent {};
    struct SelectUp      {};
    struct SelectDown    {};
    struct Quit          {};
    using Msg = std::variant<AppendDraft, DeleteDraft, SubmitDraft,
                              ToggleDone, DeleteCurrent, SelectUp,
                              SelectDown, Quit>;
\end{lstlisting}

Eight cases. The compiler will refuse to compile an \code{update}
that does not handle all of them.

\paragraph{Step 3: write \code{update}.}

\begin{lstlisting}
    static auto update(Model m, Msg msg) -> std::pair<Model, Cmd<Msg>> {
        return std::visit(overload{
            [&](AppendDraft a) -> std::pair<Model, Cmd<Msg>> {
                m.draft.push_back(a.c);
                return {std::move(m), Cmd<Msg>::none()};
            },
            [&](DeleteDraft) -> std::pair<Model, Cmd<Msg>> {
                if (!m.draft.empty()) m.draft.pop_back();
                return {std::move(m), Cmd<Msg>::none()};
            },
            [&](SubmitDraft) -> std::pair<Model, Cmd<Msg>> {
                if (!m.draft.empty()) {
                    m.items.push_back({std::move(m.draft), false});
                    m.draft.clear();
                }
                return {std::move(m), Cmd<Msg>::none()};
            },
            [&](ToggleDone) -> std::pair<Model, Cmd<Msg>> {
                if (!m.items.empty())
                    m.items[m.selected].done = !m.items[m.selected].done;
                return {std::move(m), Cmd<Msg>::none()};
            },
            [&](DeleteCurrent) -> std::pair<Model, Cmd<Msg>> {
                if (!m.items.empty()) {
                    m.items.erase(m.items.begin() + m.selected);
                    if (m.selected >= (int)m.items.size())
                        m.selected = std::max(0, (int)m.items.size() - 1);
                }
                return {std::move(m), Cmd<Msg>::none()};
            },
            [&](SelectUp) -> std::pair<Model, Cmd<Msg>> {
                m.selected = std::max(0, m.selected - 1);
                return {std::move(m), Cmd<Msg>::none()};
            },
            [&](SelectDown) -> std::pair<Model, Cmd<Msg>> {
                m.selected = std::min((int)m.items.size() - 1, m.selected + 1);
                return {std::move(m), Cmd<Msg>::none()};
            },
            [&](Quit) -> std::pair<Model, Cmd<Msg>> {
                return {std::move(m), Cmd<Msg>::quit()};
            },
        }, msg);
    }
\end{lstlisting}

Every case is named. Adding a ninth feature is adding a struct to
the variant and a lambda to the visit; the compiler will list every
callsite that needs updating.

\paragraph{Step 4: write \code{view}.}

\begin{lstlisting}
    static Element view(const Model& m) {
        std::vector<Element> rows;
        rows.reserve(m.items.size());
        for (int i = 0; i < (int)m.items.size(); ++i) {
            const auto& it = m.items[i];
            auto check  = it.done ? "[x] " : "[ ] ";
            auto style  = it.done ? (Style{} | Dim | Strike) : Style{};
            auto line   = text(std::string(check) + it.text) | style;
            if (i == m.selected) line = line | Inverse;
            rows.push_back(line);
        }
        auto list  = rows.empty() ? Element{t<"(empty)"> | Dim} : v(rows);
        auto input = h(t<"> "> | Bold,
                       text(m.draft + "_") | Fg<200, 200, 100>);
        auto help  = t<"enter add | space toggle | d delete | q quit"> | Dim;

        return v(list | pad<1> | border_<Round>,
                 input | pad<1>,
                 help) | pad<1>;
    }
\end{lstlisting}

The view depends only on the model. Hand it any \code{Model} value
and inspect the returned \code{Element} tree without rendering it.

\paragraph{Step 5: write \code{subscribe}.}

\begin{lstlisting}
    static Sub<Msg> subscribe(const Model&) {
        return Sub<Msg>::on_key([](const KeyEvent& k) -> std::optional<Msg> {
            if (key_is(k, 'q'))                 return Quit{};
            if (k.kind == KeyKind::Enter)       return SubmitDraft{};
            if (k.kind == KeyKind::Backspace)   return DeleteDraft{};
            if (k.kind == KeyKind::Up)          return SelectUp{};
            if (k.kind == KeyKind::Down)        return SelectDown{};
            if (key_is(k, ' '))                 return ToggleDone{};
            if (key_is(k, 'd'))                 return DeleteCurrent{};
            if (k.kind == KeyKind::Char)        return AppendDraft{k.character};
            return std::nullopt;
        });
    }
\end{lstlisting}

The keymap is one screen of code. Every other piece of input
handling in the program goes through this single filter.

\paragraph{Step 6: \code{init} and \code{main}.}

\begin{lstlisting}
    static auto init() -> std::pair<Model, Cmd<Msg>> {
        return {Model{}, Cmd<Msg>::none()};
    }
}; // end struct Todos

int main() { run<Todos>({.title = "todo"}); }
\end{lstlisting}

The whole application is one file, a few hundred lines, and every
part is nameable. Tests look like:

\begin{lstlisting}
int main() {
    using T = Todos;
    auto [m0, _] = T::init();

    auto [m1, _1] = T::update(std::move(m0), T::AppendDraft{'h'});
    auto [m2, _2] = T::update(std::move(m1), T::AppendDraft{'i'});
    auto [m3, _3] = T::update(std::move(m2), T::SubmitDraft{});
    assert(m3.items.size() == 1 && m3.items[0].text == "hi");
    assert(m3.draft.empty());

    auto [m4, _4] = T::update(std::move(m3), T::ToggleDone{});
    assert(m4.items[0].done);

    auto [m5, _5] = T::update(std::move(m4), T::DeleteCurrent{});
    assert(m5.items.empty());
}
\end{lstlisting}

\section{Performance: where the loop spends time}
\label{sec:elm:perf}

The loop is simple, but ``one frame per event'' could be expensive
if executed naively. A few notes on what \maya{} does to keep it
fast and where the costs actually fall.

\paragraph{The model is moved, not copied.} \code{update(Model,
Msg)} takes the model \emph{by value}; the runtime
\code{std::move}s the current model in. For a \code{Model} of small
fields, this is the cost of a few pointer copies; for one with a
big vector or string, the buffer's ownership swaps and no
element-by-element copy happens.

\paragraph{The view is rebuilt every frame.} The \code{Element}
tree is allocated fresh on each frame. This sounds wasteful but is
cheap: small bump-style allocations, freed when the tree goes out of
scope at the end of the frame. Chapter~\ref{ch:canvas} shows the
rendering pipeline that absorbs this cost into a single linear
sweep.

\paragraph{Diff, not full repaint.} The renderer compares the new
canvas to the previous one and emits bytes only for changed cells.
A frame with no visible changes emits zero bytes. A frame whose only
change is one cell emits five to ten bytes. The SIMD pass in
Chapter~\ref{ch:simd} compares 32 cells per CPU instruction; a 200x60
terminal compares in microseconds.

\paragraph{Idle frames have no cost.} If no event arrives and no
subscription is firing, the runtime is blocked in a \code{poll}
syscall and uses zero CPU. There is no render-on-tick busy loop. A
\maya{} program that nobody is using is invisible in \code{htop}.

\paragraph{Tasks run on a worker pool.} \code{Cmd::task} dispatches
to a small bounded pool of background threads; \code{Cmd::task\_isolated}
spins a fresh thread per call. The main loop never blocks on user
code. A wedged task does not freeze the UI.

\begin{recap}
The loop is cheap because: the model is moved not copied; the
element tree is allocated cheaply and dropped per-frame; the diff
is SIMD-fast; idle frames cost zero; tasks are off the main thread.
A \maya{} program redraws at hundreds of frames per second on a
modern laptop, and at dozens over slow ssh links.
\end{recap}

\section{Pitfalls and anti-patterns}
\label{sec:elm:pitfalls}

If you ignore Elm's discipline, the library still compiles and runs,
but you start losing the benefits. The common ways:

\begin{warning}
\textbf{Hidden mutation in update.} Calling
\code{some\_global.cache.set(\dots)} inside \code{update} works but
makes the function impure; testing becomes ``set up the global,
then call, then check the global''. Push the mutation out as a
\code{Cmd::task} that performs it, or move the state into the
model.
\end{warning}

\begin{warning}
\textbf{Synchronous I/O in update.} \code{std::fopen} inside
\code{update} blocks the main loop. The UI freezes. Always wrap
file/network/subprocess work in a \code{Cmd::task} so it runs on a
worker.
\end{warning}

\begin{warning}
\textbf{Coupling view to side effects.} \code{view} should depend
only on the model. If it reads the clock, the filesystem, or any
mutable global, frames stop being deterministic, and the diff can
produce a flicker between two equally valid renderings of ``now''.
Capture the world's current value into the model via a
subscription (\code{Sub::every}) and read it from the model.
\end{warning}

\begin{warning}
\textbf{Overusing \code{Cmd::task} for trivial work.} A task is
cheap but not free --- pool dispatch, worker wake, mutex hop. If
your ``async'' work is \code{counter + 1}, just compute it inline.
\end{warning}

\begin{warning}
\textbf{Forgetting the \code{std::optional} in \code{OnKey}.} The
filter must return \code{std::nullopt} for keys the application
doesn't care about. Returning a \code{Msg} for every key produces a
flood and may interfere with the runtime's own bindings (\code{Ctrl-C}
in some configurations).
\end{warning}

\section{When to skip the Program scaffold}
\label{sec:elm:simple}

Elm-style discipline pays for itself on programs that grow.
For a one-shot script --- print a banner, ask one question, exit ---
it is overkill. \maya{} provides smaller entry points for those
cases.

\code{maya::print(element)} renders a static \code{Element} once and
returns. No model, no loop, no subscriptions. We used it for the
five-line banner in Chapter~\ref{ch:intro}.

The second-smallest form, useful when you need a single keystroke
but nothing more elaborate:

\begin{lstlisting}
std::println("press y to continue");
auto key = maya::read_key();
if (key == 'y') /* go */;
\end{lstlisting}

For anything with state that changes over time, switch to the
\code{Program} form. The boundary is roughly: if your program has
three distinct keys it responds to, write the \code{Program}.

\section{From single-function to Program: a worked refactor}

Here is the same program in both styles. First, the quick-tool form
(useful for one-off scripts):

\begin{lstlisting}
Signal<int> n{0};
run({.title = "counter"},
    [&](const Event& ev) {
        if (key(ev, '+')) n.update([](int& x){ ++x; });
        if (key(ev, '-')) n.update([](int& x){ --x; });
        return !key(ev, 'q');
    },
    [&] {
        return v(text("Count: " + std::to_string(n.get())) | Bold,
                 t<"[+/-] q quit"> | Dim) | pad<1>;
    });
\end{lstlisting}

This is fine. But the state is captured by reference inside a lambda,
the event handler returns a \code{bool}, and there is no separation
between business logic and presentation.

The Program version:

\begin{lstlisting}
struct Counter {
    struct Model { int count = 0; };
    struct Inc {}; struct Dec {}; struct Quit {};
    using Msg = std::variant<Inc, Dec, Quit>;

    static auto init() -> std::pair<Model, Cmd<Msg>> {
        return {{}, Cmd<Msg>::none()};
    }
    static auto update(Model m, Msg msg) -> std::pair<Model, Cmd<Msg>> {
        return std::visit(overload{
            [&](Inc)  { return std::pair{Model{m.count + 1}, Cmd<Msg>{}}; },
            [&](Dec)  { return std::pair{Model{m.count - 1}, Cmd<Msg>{}}; },
            [](Quit)  { return std::pair{Model{}, Cmd<Msg>::quit()}; }
        }, msg);
    }
    static Element view(const Model& m) {
        return v(text("Count: " + std::to_string(m.count)) | Bold,
                 t<"[+/-] q quit"> | Dim) | pad<1>;
    }
    static Sub<Msg> subscribe(const Model&) {
        return key_map<Msg>({{'+', Inc{}}, {'-', Dec{}}, {'q', Quit{}}});
    }
};
int main() { run<Counter>({.title = "counter"}); }
\end{lstlisting}

It is a few more lines but the parts are now nameable. Tests look
like:

\begin{lstlisting}
TEST_CASE("incrementing the counter") {
    auto [m, _] = Counter::update({.count = 5}, Counter::Inc{});
    CHECK(m.count == 6);
}
\end{lstlisting}

No terminal, no event loop, no render. That is the entire point.

\begin{recap}
A \maya{} application is a \code{Program}: a \code{Model} type, a
\code{Msg} sum, and four pure functions \code{init}/\code{update}/
\code{view}/\code{subscribe}. Effects (\code{Cmd}) and event sources
(\code{Sub}) are described as data; the runtime interprets them. The
result is a program whose business logic is trivially testable.
\end{recap}

\section*{Workshop}

\begin{enumerate}[leftmargin=*]
  \item Add a \code{Reset} message to the counter that sets
    \code{count} back to zero. Add a \code{r} keybinding.
  \item Add a \code{Tick} subscription that increments the counter
    once per second.
  \item Add a \code{Save} message that writes the current count to
    the clipboard via \code{Cmd<Msg>::write\_clipboard(...)}.
  \item Write a unit test (just a \code{static\_assert} or a
    one-line \code{main}) that pumps three \code{Inc}s through
    \code{update} and checks the final count.
\end{enumerate}

\chapter{The compile-time DSL}
\label{ch:dsl}

\maya{}'s most distinctive feature is its DSL --- a tiny embedded
language inside \cpp{} for describing UI trees. The DSL is unusual
because it lives entirely at the type level: an expression like

\begin{lstlisting}
t<"Hello"> | Bold | Fg<100, 180, 255>
\end{lstlisting}

is a \emph{single value} whose type encodes every choice you made
(\emph{the string}, \emph{bold}, \emph{the colour}). This chapter
shows how that works, why the compiler refuses certain expressions,
and how to escape into the runtime when you need to.

\begin{aside}
If you have written HTML or React, you have written a DSL: a small,
specialised vocabulary for the one thing you are doing (describing
UI). \maya{}'s DSL is the same idea, but it sits inside \cpp{}, runs
at compile time, and refuses to compile if you build an illegal
tree. The bottom-up explanation in this chapter is long, but the
payoff is being able to read any expression in any \maya{} program
and know what every token contributes to the type.
\end{aside}

\section{From scratch: what a DSL even is}
\label{sec:dsl:what}

The phrase ``Domain-Specific Language'' has two flavours in
programming.

An \textbf{external DSL} is a language with its own parser, lexer,
and often its own interpreter: SQL, regular expressions, JSON, CSS,
the shell. You write it in a different syntax from the host
language, and a separate engine reads it.

An \textbf{embedded} (or \emph{internal}) DSL is a sub-vocabulary
within the host language. It looks like normal code, but the
operators and identifiers have been overloaded to express a different
domain. JavaScript's \code{describe("thing", () => \dots)} from
Mocha is an embedded DSL for testing. \cpp{}'s \code{std::cout <<
thing} is the world's most ubiquitous embedded DSL --- the \code{<<}
operator does \emph{stream insertion} when applied to streams, not
bit-shift.

The \maya{} DSL is the embedded kind. Every piece of it is legal
\cpp{}; the magic is in operator overloads, variable templates, and
the fact that we have made the result types interesting enough that
the compiler can reason about them.

Four properties define a good embedded DSL:

\begin{enumerate}[leftmargin=*]
  \item \textbf{It reads like the thing it describes.}
    \code{v(t<"a">, t<"b">)} looks vertical because we picked
    \code{v} for that role.
  \item \textbf{It composes with the rest of the language.} You can
    pass a DSL fragment to a function, return one from a function,
    store one in a variable.
  \item \textbf{The host's type checker still works.} A mistake in
    the DSL is a host-language error, not a runtime ``DSL
    interpreter blew up''.
  \item \textbf{There is an escape hatch back to the host
    language.} \maya{}'s \code{dyn([]{...})} is exactly this.
\end{enumerate}

In the rest of the chapter we will see how \maya{} achieves all four.

\section{Non-type template parameters, properly}
\label{sec:dsl:nttp}

The foundation of the entire DSL is the \cpp{}20 / 26 ability to use
structured values as template parameters. We touched on this in
Chapter~\ref{ch:cpp}; here is the longer version, because the DSL
leans on it constantly.

A \textbf{template parameter} can be a type (\code{typename T}), but
it can also be a \emph{value}:

\begin{lstlisting}
template <int N> struct Array { int data[N]; };
Array<8>  a;   // N = 8
Array<16> b;   // N = 16, different type from a
\end{lstlisting}

The parameter \code{N} is an integer, not a type. The template is
instantiated once per distinct value of \code{N}, and each
instantiation is a separate type. \code{Array<8>} and \code{Array<16>}
do not interconvert.

\subsection{Why integers are not enough}

Integers, pointers, and enum values have always been allowed as
non-type template parameters (NTTPs). What you could \emph{not} do,
until \cpp{}20, was use a class type as an NTTP. There was no way to
write

\begin{lstlisting}
struct Point { int x; int y; };
template <Point P> struct Thing { /* ... */ };
Thing<Point{3, 4}> t;     // illegal in C++17
\end{lstlisting}

This matters for us because a UI description involves not just
numbers but \emph{compound data}: a string is a class type, a style
is a class type, a colour is three bytes. If we cannot put those in
template parameters, the DSL is severely limited.

\cpp{}20 lifted the restriction for \textbf{structural} class types:
any literal type whose members are all public, with no user-defined
oddities, can be an NTTP. \cpp{}26 extended it further. The
consequence is that \code{maya::dsl::Str}, \code{maya::dsl::CTStyle},
and \code{maya::dsl::BoxCfg} can all appear as template parameters,
and the DSL works.

\subsection{\code{Str}: the building block}

The simplest structural NTTP in \maya{} is \code{Str} ---
a fixed-size string class:

\begin{lstlisting}
template <std::size_t N>
struct Str {
    char data[N]{};
    consteval Str(const char (&s)[N]) {
        for (std::size_t i = 0; i < N; ++i) data[i] = s[i];
    }
    consteval std::size_t size() const { return N - 1; }
    constexpr operator std::string_view() const { return {data, N - 1}; }
};
\end{lstlisting}

Line by line:

\begin{itemize}[leftmargin=*]
  \item \code{template <std::size\_t N> struct Str}: parameterised
    by the length \code{N}.
  \item \code{char data[N]\{\}}: a fixed array of \code{N}
    characters, zero-initialised.
  \item \code{consteval Str(const char (\&s)[N])}: constructor that
    \emph{must} run at compile time, taking a reference to a
    \code{char} literal of length \code{N}. The compiler deduces
    \code{N} from the literal's size.
  \item The loop copies the literal into \code{data}.
\end{itemize}

Now you can write \code{Str{"Hello"}} as a constant expression. The
resulting value has a \emph{type} that depends on the length and a
\emph{value} that contains the bytes --- and crucially, both are
available to other templates.

\subsection{Putting \code{Str} into the type system}

\begin{lstlisting}
template <Str S>
struct TextNode {
    static constexpr std::string_view text() { return {S.data, S.size()}; }
};

TextNode<"Hello"> a;
TextNode<"World"> b;
// decltype(a) and decltype(b) are different types,
// even though both store no data per-instance.
\end{lstlisting}

\code{TextNode<"Hello">} and \code{TextNode<"World">} are distinct
types. The string contents are part of the \emph{type identity},
not a field. Two instances of \code{TextNode<"Hello">} occupy zero
bytes and are interchangeable.

Why does this matter? Because the compiler can now reason about your
strings: it can choose overloads based on which string was used, it
can refuse to mix incompatible strings, and it can fold the string's
length into compile-time arithmetic --- without you writing any of
that by hand.

\subsection{\code{CTStyle}: a structural style}

The second key NTTP type:

\begin{lstlisting}
struct CTStyle {
    bool has_fg = false, has_bg = false;
    uint8_t fg_r = 0, fg_g = 0, fg_b = 0;
    uint8_t bg_r = 0, bg_g = 0, bg_b = 0;
    bool bold_ = false, dim_ = false, italic_ = false;
    bool underline_ = false, strike_ = false, inverse_ = false;

    consteval CTStyle merge(CTStyle o) const { /* combine fields */ }
    [[nodiscard]] Style runtime() const       { /* materialise */ }
};
\end{lstlisting}

A \code{CTStyle} value is a \textbf{bag of compile-time style
flags}: bold yes/no, foreground colour yes/no plus three RGB bytes,
and so on. The struct is structural (all-public fields, no virtuals,
no custom comparison), so it qualifies as an NTTP. The
\code{merge} method is \code{consteval} --- it can be called inside
template arguments to compute new \code{CTStyle}s from old ones.

A style tag exists for every bool field:

\begin{lstlisting}
template <CTStyle V>
struct StyTag {
    static constexpr CTStyle value = V;
};

inline constexpr StyTag<CTStyle{.bold_ = true}>      Bold{};
inline constexpr StyTag<CTStyle{.dim_ = true}>       Dim{};
inline constexpr StyTag<CTStyle{.italic_ = true}>    Italic{};
inline constexpr StyTag<CTStyle{.underline_ = true}> Underline{};
\end{lstlisting}

\code{Bold} is a \code{StyTag} whose template argument is a
\code{CTStyle} with one field set. The variable \code{Bold} itself
is a zero-size instance --- the \code{CTStyle} lives in the type.

\subsection{\code{Fg<R,G,B>}: colours in the type system}

Colours are encoded the same way: a \code{CTStyle} with the
foreground fields set, exposed through a variable template:

\begin{lstlisting}
namespace detail_rgb {
    template <uint8_t R, uint8_t G, uint8_t B>
    consteval CTStyle make_fg() {
        return CTStyle{.has_fg=true, .fg_r=R, .fg_g=G, .fg_b=B};
    }
}

template <uint8_t R, uint8_t G, uint8_t B>
inline constexpr StyTag<detail_rgb::make_fg<R, G, B>()> Fg{};
\end{lstlisting}

\code{Fg<100, 180, 255>} is a zero-size value whose type is a
\code{StyTag} whose template parameter is a \code{CTStyle} whose
fields are exactly the RGB you asked for. The numbers fed in become
part of the type. Nothing about colour leaks into runtime data
until \code{.build()}.

\begin{funfact}
The MSVC compiler had (and partly still has) a long-standing bug
around \code{constexpr} evaluation of designated-initialised NTTPs.
If you look in \file{dsl.hpp} you will see a small namespace
\code{detail\_rgb} with \code{make\_fg} / \code{make\_bg} helpers ---
that existence-of-helpers is the workaround. The compiler bug ticket
is cited in a comment. Real libraries are full of small monuments
like this; each is a story of someone who hit the same wall before
you.
\end{funfact}

\begin{idea}
NTTPs let you put \emph{values} into the type system. \maya{} uses
that to encode strings (\code{Str}), styles (\code{CTStyle}),
colours (RGB triples), borders, paddings, and growth factors. Once a
value lives in a type, the compiler can check rules between values
at compile time --- which is the whole point.
\end{idea}

\section{The shape}

The DSL has six kinds of node:

\begin{center}
\small
\begin{tabular}{ll}
\code{TextNode<Str, CTStyle>}     & compile-time text \\
\code{RuntimeTextNode<S>}         & runtime text (from a string or number) \\
\code{BoxNode<Dir, Cfg, Cs...>}   & vertical or horizontal container \\
\code{DynNode<F>}                 & lambda escape hatch \\
\code{MapNode<R, P>}              & project a range to elements \\
\code{SpacerNode}                 & a flexible gap \\
\end{tabular}
\end{center}

Every node satisfies the \code{Node} concept:

\begin{lstlisting}
template <typename T>
concept Node = requires(const T& n) {
    { n.build() } -> std::convertible_to<Element>;
};
\end{lstlisting}

A \code{Node} is anything with a \code{.build()} that yields a runtime
\code{Element}. That is the only contract; the rest is how nodes
combine.

\section{\code{t<"Hello">} --- string in a type}

The simplest node is the compile-time text node:

\begin{lstlisting}
template <Str S, CTStyle Sty = CTStyle{}>
struct TextNode {
    static constexpr Str  str   = S;
    static constexpr CTStyle style = Sty;

    Element build() const {
        return TextElement{std::string{S.view()}, Sty.runtime()};
    }
    operator Element() const { return build(); }
};

template <Str S>
inline constexpr TextNode<S> t{};
\end{lstlisting}

The template parameters are: an \code{Str} (the string contents, an
NTTP from Chapter~\ref{ch:cpp}) and a \code{CTStyle} (a structural
type holding the compile-time style). \code{t<"Hello">} is the
variable template: \code{TextNode<"Hello", CTStyle\{\}>\{\}}.

Two facts to internalise:
\begin{itemize}[leftmargin=*]
  \item \code{t<"a">} and \code{t<"b">} have \emph{different types}.
    Their string contents live in the type system, not in a field.
  \item \code{TextNode} has no nonstatic data members. It is a zero-size
    type. Two instances of \code{TextNode<"a">} are
    indistinguishable; the compiler treats them as the same value.
\end{itemize}

\section{Pipes are operator overloads}

The DSL's pipe (\code{|}) is just \code{operator|}. \maya{} defines
overloads that match \code{(Node, StyTag)} and return a new node with
the style merged in:

\begin{lstlisting}
template <Str S, CTStyle Sty, CTStyle V>
constexpr auto operator|(TextNode<S, Sty>, StyTag<V>) {
    return TextNode<S, Sty.merge(V)>{};   // new type, same instance shape
}
\end{lstlisting}

Read the return type: the merged style is computed by \code{Sty.merge(V)},
which is \code{constexpr}, so the result is a brand-new
\code{TextNode<S, Sty'>}. Apply two style tags and you get yet another
type. The pipeline is sequential type rewriting at compile time.

\begin{idea}
Because the result of each pipe is a new type, the compiler can
enforce \emph{type-state} rules: a method is only available on a node
that has reached a particular state.
\end{idea}

This is the trick behind the border colour example. The pipe overload
for \code{BColTag} (a border colour tag) has a \code{requires} clause:

\begin{lstlisting}
template <FlexDirection Dir, BoxCfg Cfg, typename... Cs,
          uint8_t R, uint8_t G, uint8_t B>
    requires (Cfg.has_border)
constexpr auto operator|(BoxNode<Dir, Cfg, Cs...> n, BColTag<R, G, B>) {
    /* return a new BoxNode with border colour set */
}
\end{lstlisting}

If you call \code{v(...) | bcol<60, 65, 80>} without first applying
\code{border\_<Round>}, \code{Cfg.has\_border} is \code{false}, the
\code{requires} clause fails, no overload matches, and the compiler
prints an error. It is the same mechanism a Rust crate would use with
typestate traits.

\section{\code{v()} and \code{h()} --- variadic containers}

A box takes any number of children and a direction:

\begin{lstlisting}
template <FlexDirection Dir, BoxCfg Cfg, typename... Cs>
struct BoxNode {
    std::tuple<Cs...> children;

    Element build() const {
        BoxBuilder b{Dir, Cfg};
        std::apply([&](const auto&... cs) {
            (b.add(cs.build()), ...);   // fold-expression over children
        }, children);
        return b.finalize();
    }
};

template <typename... Cs>
constexpr auto v(Cs... cs) {
    return BoxNode<FlexDirection::Column, BoxCfg{}, Cs...>{
        std::tuple{std::move(cs)...}
    };
}

template <typename... Cs>
constexpr auto h(Cs... cs) {
    return BoxNode<FlexDirection::Row, BoxCfg{}, Cs...>{
        std::tuple{std::move(cs)...}
    };
}
\end{lstlisting}

Two \cpp{} features worth pausing on:

\begin{itemize}[leftmargin=*]
  \item \textbf{Variadic templates} (\code{typename... Cs}): you can
    declare ``zero or more types''. \code{Cs...} is a \emph{parameter
    pack}; you expand it with \code{...}.
  \item \textbf{Fold expressions} (\code{(b.add(cs.build()), ...)}):
    short for ``apply this expression to every element of the
    pack''. The comma is the \emph{fold operator} --- evaluate each,
    discard intermediate results.
\end{itemize}

The result: \code{v(t<"a">, t<"b">, t<"c">)} is a single
\code{BoxNode<Column, BoxCfg\{\}, TextNode<"a">, TextNode<"b">,
TextNode<"c">>}, holding a \code{tuple} of three zero-size text
nodes. Total runtime cost: zero bytes (the tuple of empty types is
empty).

\section{\code{BoxCfg} --- the type-state}

\code{BoxCfg} is a \code{constexpr} struct that records every layout
decision applied to a box:

\begin{lstlisting}
struct BoxCfg {
    int pad_top    = 0, pad_right = 0, pad_bottom = 0, pad_left = 0;
    int gap        = 0;
    int grow_num   = 0;
    bool has_border       = false;
    BorderStyle border    = BorderStyle::None;
    bool has_border_color = false;
    Color border_color{};
    Align align_items    = Align::Auto;
    Justify justify      = Justify::Auto;
    // ... etc.

    constexpr BoxCfg with_padding(int t, int r, int b, int l) const {
        BoxCfg c = *this;
        c.pad_top = t; c.pad_right = r; c.pad_bottom = b; c.pad_left = l;
        return c;
    }
    constexpr BoxCfg with_border(BorderStyle bs) const {
        BoxCfg c = *this; c.has_border = true; c.border = bs; return c;
    }
    // ...
};
\end{lstlisting}

Every layout pipe (\code{pad<1>}, \code{gap\_<1>}, \code{border\_<Round>},
\code{grow\_<2>}) returns a new \code{BoxNode} whose \code{BoxCfg}
template parameter has been replaced via the corresponding
\code{with\_*} method. The type carries the full configuration; the
final \code{.build()} reads the type and produces the runtime
\code{BoxElement}.

\section{Static, dynamic, and the bridge}

A pure-compile-time tree is fine for a banner but useless for a real
app: the count, the file list, the current time are all known only at
runtime. \maya{} provides three bridge nodes.

\subsection{\code{text(...)} --- runtime strings}

\begin{lstlisting}
text("hello")                // wraps a string_view
text(std::to_string(n))      // wraps a string
text(42)                     // wraps an int, formatted at build time
text(3.14)                   // wraps a double
\end{lstlisting}

\code{text} returns a \code{RuntimeTextNode<S>} --- a node whose
content is a runtime value but which still supports the DSL pipes.
Its \code{build()} formats the value and produces a
\code{TextElement}.

\subsection{\code{dyn([](){ ... })} --- arbitrary lambdas}

\begin{lstlisting}
dyn([&] {
    return loading ? text("Loading...") | Dim
                   : text("Ready")     | Bold | Fg<80, 220, 120>;
})
\end{lstlisting}

\code{dyn} wraps a lambda that returns anything convertible to
\code{Element}. \code{build()} just calls the lambda. This is the
universal escape hatch: when you cannot express something in the DSL
itself, write a lambda.

\subsection{\code{map(range, projection)} --- list comprehension}

\begin{lstlisting}
std::vector<File> files = ...;

auto list = v(
    t<"Files:"> | Bold,
    map(files, [](const File& f) {
        return h(text(f.name) | Bold, space, text(f.size) | Dim);
    })
);
\end{lstlisting}

\code{map} produces a \code{MapNode<R, P>}. At \code{build()} time it
walks the range, applies the projection, and produces an
\code{ElementList} that is spliced into the parent's children.

\section{Reading a real pipeline}

Put it all together. Here is a typical maya panel:

\begin{lstlisting}
auto panel = v(
    t<"Dashboard"> | Bold | Fg<100, 180, 255>,
    sep,
    dyn([&] { return text("Active users: " + std::to_string(state.users)); }),
    map(state.items, [](const auto& it) {
        return h(text(it.name) | Bold, space, text(it.value) | Dim);
    }),
    space,
    t<"[q] quit"> | Dim
) | border_<Round> | bcol<60, 65, 80> | pad<1, 2> | gap_<1>;
\end{lstlisting}

What is the type of \code{panel}? Approximately:

\begin{lstlisting}
BoxNode<
    FlexDirection::Column,
    BoxCfg{ /* round border, colour, pad 1/2, gap 1 */ },
    TextNode<"Dashboard", { bold, fg=(100,180,255) }>,
    SepNode,
    DynNode<lambda1>,
    MapNode<vector<Item>, lambda2>,
    SpacerNode,
    TextNode<"[q] quit", { dim }>
>
\end{lstlisting}

The compiler holds every choice you made --- the strings, the styles,
the border, the padding --- in this one large type. \code{panel.build()}
then walks it once, recursively, and produces an \code{Element} tree
that the layout engine and renderer take from there.

\section{Why compile-time?}

You could write this DSL with runtime polymorphism (each node a
\code{shared\_ptr<Node>}, each pipe a method that mutates fields). It
would work. But:

\begin{itemize}[leftmargin=*]
  \item \textbf{Allocations vanish.} A purely compile-time tree of
    zero-size types has no heap traffic.
  \item \textbf{Errors are at compile time.} Border colour without a
    border, negative padding (template parameter constraint
    \code{N >= 0}), wrong type of child --- all rejected before the
    program runs.
  \item \textbf{The compiler can fold.} A \code{constexpr} tree can
    be assigned to a \code{constexpr} variable and rendered at startup
    with no parsing.
\end{itemize}

The trade-off: error messages can be long. The first time you see a
500-line template error, breathe. The bottom of the message --- the
\code{required from here} line and the constraint that failed --- is
usually the actionable bit.

\section{Where the runtime takes over}

\code{.build()} is the boundary. Above it: types, NTTPs, fold
expressions, the world the compiler sees. Below it: \code{Element}, a
runtime value the layout and renderer understand. The next chapter
goes through that runtime side.

\begin{recap}
\maya{}'s DSL is a set of node types that compose with \code{|} into
larger node types. Strings, styles, borders, paddings, and tree
structure all live as template parameters. Type-state constraints
make invalid trees uncompilable. Runtime content enters through
\code{text}, \code{dyn}, and \code{map}. \code{.build()} converts the
compile-time tree into a runtime \code{Element}.
\end{recap}

\section{Every pipe, named and explained}
\label{sec:dsl:pipes-tour}

The DSL has roughly twenty overloads of \code{operator|}. Reading
the source for each takes ten seconds; reading them together is
useful, because the table makes it obvious which kinds of node
accept which kinds of modifier. Here is the full list, grouped by
the left-hand side.

\subsection{Pipes that take a \code{TextNode}}

\begin{center}
\small
\begin{tabular}{lll}
\textbf{Pipe} & \textbf{Effect} & \textbf{Result type} \\
\hline
\code{TextNode | StyTag<V>}    & merge style V       & \code{TextNode<S, Sty.merge(V)>} \\
\code{TextNode | WidthTag<W>}  & wrap in fixed box   & \code{BoxNode<Row, ct\_width=W>} \\
\code{TextNode | HeightTag<H>} & wrap in fixed box   & \code{BoxNode<Row, ct\_height=H>} \\
\end{tabular}
\end{center}

Text nodes accept any style modifier and nothing else. Adding
\code{pad<1>} to a bare text node is a compile error --- padding is
a container property; text has no notion of inset space.

\subsection{Pipes that take a \code{RuntimeTextNode}}

The runtime text node (built by \code{text(...)}) gets the same
style pipes plus two wrap-mode pipes:

\begin{center}
\small
\begin{tabular}{ll}
\textbf{Pipe} & \textbf{Effect} \\
\hline
\code{| StyTag<V>}    & merge style into the runtime node \\
\code{| clip}         & truncate with ellipsis on overflow \\
\code{| nowrap}       & never break the line, may overflow \\
\end{tabular}
\end{center}

The wrap-mode pipes do not exist on compile-time \code{TextNode}
because compile-time text is meant to be short and known; the
ellipsis case only matters when you wrap a runtime value of unknown
length.

\subsection{Pipes that take a \code{BoxNode}}

This is where the real type-state machine lives. A box accepts every
kind of modifier:

\begin{center}
\small
\begin{tabular}{lll}
\textbf{Pipe} & \textbf{Effect} & \textbf{Notes} \\
\hline
\code{| StyTag<V>}        & merge style on the box itself      & all boxes \\
\code{| PadTag<T,R,B,L>}  & set inner padding                  & all boxes \\
\code{| GapTag<G>}        & set gap between children           & all boxes \\
\code{| BorderTag<BS>}    & set border style                   & enables \code{bcol} pipe \\
\code{| BColTag<R,G,B>}   & set border colour                  & \textbf{requires} \code{has\_border} \\
\code{| GrowTag<G>}       & set flex grow factor               & all boxes \\
\code{| WidthTag<W>}      & set fixed width                    & all boxes \\
\code{| HeightTag<H>}     & set fixed height                   & all boxes \\
\end{tabular}
\end{center}

Four of the modifier types have a \code{static\_assert} on the
template parameter to reject obvious nonsense at compile time:

\begin{itemize}[leftmargin=*]
  \item \code{PadTag}: every value must be \code{>= 0}. Negative
    padding is a category error.
  \item \code{GapTag}: \code{G >= 0}.
  \item \code{GrowTag}: \code{G >= 0}.
  \item \code{WidthTag}, \code{HeightTag}: \code{> 0}. A width of
    zero is a hidden box, which we have a different way to express
    (\code{visible(false)}).
\end{itemize}

The one with the named requires-clause is \code{BColTag}:

\begin{lstlisting}
template <FlexDirection Dir, BoxCfg Cfg, typename... Cs,
          uint8_t R, uint8_t G, uint8_t B>
    requires (Cfg.has_border)
constexpr auto operator|(BoxNode<Dir, Cfg, Cs...> n, BColTag<R, G, B>);
\end{lstlisting}

The constraint \code{requires (Cfg.has\_border)} forbids the
overload when the box has no border yet. We will look at the error
message in section~\ref{sec:dsl:errors}.

\subsection{The style-merging pipe}

A quietly powerful one: two style tags combine into one.

\begin{lstlisting}
template <CTStyle A, CTStyle B>
constexpr auto operator|(StyTag<A>, StyTag<B>) {
    return StyTag<A.merge(B)>{};
}
\end{lstlisting}

This is the foundation of building reusable style presets:

\begin{lstlisting}
constexpr auto heading = Bold | fg<0xFFFFFF>;
constexpr auto error   = Bold | fg<0xFF4444>;
constexpr auto subtle  = Dim  | Italic;

t<"Title">     | heading
text(msg)      | error
t<"caption">   | subtle
\end{lstlisting}

The combined tag is itself a \code{StyTag<...>}, so it composes
further: \code{heading | Underline} is legal. Compose your design
system this way and applying it is a single token.

\subsection{The hex-colour shorthand}

\code{Fg<R, G, B>} is verbose for designers, who think in hex. The
DSL provides \code{fg<Hex>} and \code{bg<Hex>}:

\begin{lstlisting}
t<"warn">  | fg<0xFFAA00>
t<"error"> | fg<0xFF4444>
v(...)     | bg<0x1A1A2E>
\end{lstlisting}

The \code{consteval} helper that decodes the hex into RGB also
validates it: \code{fg<0xFFFFFFFF>} (32 bits) throws a string
literal from inside \code{consteval}, which the compiler reports as
a constant-expression failure --- ``\code{fg<>: hex value exceeds
0xFFFFFF}''. Compile-time error, not a silent black box.

\section{The type-state machine, drawn out}
\label{sec:dsl:state-machine}

The phrase \textbf{type-state} comes from research on programming
languages: encode the legal sequence of operations on an object into
its type, so that calling a method at the wrong stage is a type
error. Rust's \code{Builder} idioms are typical: the constructor
returns a \code{Builder<Empty>}, the \code{name(\dots)} method
returns a \code{Builder<Named>}, and \code{build()} is only defined
on \code{Builder<Named>}.

\maya{}'s \code{BoxNode} works the same way, with \code{BoxCfg} as
the state.

\paragraph{Initial state.} \code{v(\dots)} returns a
\code{BoxNode<Column, BoxCfg\{\}, \dots>}: every flag in
\code{BoxCfg} is false / zero. From this state:

\begin{itemize}[leftmargin=*]
  \item Padding, gap, grow, width, height, raw style: all available.
    Each returns a new \code{BoxNode} whose \code{BoxCfg} has the
    corresponding field updated.
  \item Border: available. Returns a \code{BoxNode} with
    \code{has\_border = true}.
  \item Border colour: \emph{not} available; the requires-clause
    fails because \code{has\_border} is false.
\end{itemize}

\paragraph{After \code{| border\_<Round>}.} \code{has\_border}
becomes true. Border colour becomes available.

\paragraph{Drawn as a state diagram.}

\begin{center}
\small
\begin{tabular}{rcl}
 \code{BoxNode<Cfg(has\_border=false)>}
  & $\xrightarrow{\;\;|\;\code{border\_<R>}\;\;}$
  & \code{BoxNode<Cfg(has\_border=true)>} \\
 \code{BoxNode<Cfg(has\_border=false)>}
  & $\xrightarrow{\;\;|\;\code{bcol<\dots>}\;\;}$
  & \textsc{compile error} \\
 \code{BoxNode<Cfg(has\_border=true)>}
  & $\xrightarrow{\;\;|\;\code{bcol<\dots>}\;\;}$
  & \code{BoxNode<\dots, has\_border\_color=true>} \\
\end{tabular}
\end{center}

The state machine is tiny --- two states for the border flag --- but
the pattern scales. Any rule of the form ``\emph{X} is only legal
after \emph{Y}'' can be encoded the same way, by toggling a field in
the \code{BoxCfg} on the \emph{Y} operation and adding a
\code{requires (Cfg.\emph{field})} to the \emph{X} overload.

\begin{idea}
Type-state in a DSL means: the legal sequence of pipes is encoded
into the chain of types. The compiler refuses pipes the current type
does not permit. The error has a name (a \code{requires} clause) and
a single source line.
\end{idea}

\section{Decoding the compile errors}
\label{sec:dsl:errors}

A powerful type system gets you powerful errors. The first time you
see a 200-line template error, the impulse is to scroll up and
panic; the trick is knowing where to look.

The shape of every constraint-failure error from the DSL is:

\begin{enumerate}[leftmargin=*]
  \item A line saying \emph{no viable overload of \code{operator|}}.
  \item One or more notes saying \emph{candidate template ignored:
    constraints not satisfied}.
  \item For each candidate, the constraint that did not hold:
    \code{Cfg.has\_border} was \code{false}, \code{T >= 0} did not
    hold for \code{T = -1}, and so on.
  \item Finally, the line in \emph{your} code that triggered it.
\end{enumerate}

\subsection{The border-colour error}

The source mistake:

\begin{lstlisting}
auto ui = v(t<"hi">) | bcol<60, 65, 80>;     // forgot border first
\end{lstlisting}

The heart of the error from a modern \code{g++}:

\begin{lstlisting}[style=shell]
error: no match for 'operator|' (operand types are
    'maya::dsl::BoxNode<maya::FlexDirection::Column,
      maya::dsl::BoxCfg{...has_border=false...}, ...>'
    and 'const maya::dsl::BColTag<60, 65, 80>')
note: candidate: 'template<...> constexpr auto
    maya::dsl::operator|(BoxNode<Dir, Cfg, Cs...>, BColTag<R, G, B>)
        requires (Cfg.has_border)'
note: constraints not satisfied
note:   the expression 'Cfg.has_border' evaluated to 'false'
\end{lstlisting}

The relevant fact is in the last note: \code{Cfg.has\_border}
evaluated to false. Fix:

\begin{lstlisting}
auto ui = v(t<"hi">) | border_<Round> | bcol<60, 65, 80>;
\end{lstlisting}

\subsection{The negative-padding error}

The source mistake:

\begin{lstlisting}
auto ui = v(t<"hi">) | pad<-1>;
\end{lstlisting}

The error:

\begin{lstlisting}[style=shell]
error: static assertion failed: Padding values must be non-negative
  in instantiation of 'struct maya::dsl::PadTag<-1, -1, -1, -1>'
\end{lstlisting}

A \code{static\_assert} caught it before any pipe overload was
even considered. The error names exactly what is wrong.

\subsection{The wrong-child-type error}

The DSL's \code{v} and \code{h} take children that satisfy the
\code{DslChild} concept --- either a \code{Node} (anything with
\code{.build()}) or a \code{std::ranges::range} of \code{Element}.
Passing something else fails the concept check:

\begin{lstlisting}
auto bad = v(42);                       // int is not a DslChild
\end{lstlisting}

The error roughly:

\begin{lstlisting}[style=shell]
error: no matching function for call to 'v(int)'
note: candidate template ignored: constraints not satisfied [with Cs = {int}]
note: because 'int' does not satisfy 'DslChild'
\end{lstlisting}

The fix is to wrap the value in \code{text(\dots)} or its appropriate
constructor: \code{v(text(42))}.

\subsection{Strategy for reading a template error}

Three rules that hold for every \maya{} compile failure:

\begin{enumerate}[leftmargin=*]
  \item \textbf{Read the bottom first.} The line that says
    \emph{required from here} or \emph{in the expansion of \dots}
    is closer to your source than the wall of template-instantiation
    backtrace above it.
  \item \textbf{Look for ``constraints not satisfied'' / ``static
    assertion failed''.} Those are the human-readable handles. The
    text after them names the rule that broke.
  \item \textbf{Identify the offending NTTP value.} The compiler
    will print the \code{BoxCfg} or \code{CTStyle} with its fields;
    the field whose value is wrong (\code{has\_border=false},
    \code{pad\_t=-1}) is the actionable bit.
\end{enumerate}

Most of the wall of text is irrelevant; you are looking for two or
three lines.

\begin{warning}
Do \emph{not} \code{\#include <maya/maya.hpp>} into a translation
unit that you then refuse to use because the error is intimidating.
The errors are intimidating because they are precise. The fix is
usually one token away; scrolling past the error makes finding it
harder.
\end{warning}

\section{What \code{.build()} actually does}
\label{sec:dsl:build}

The boundary between compile time and runtime is the \code{.build()}
method. Every node provides it; the implementation depends on the
kind of node.

\paragraph{\code{TextNode<S, Sty>::build()}.}

\begin{lstlisting}
Element build() const {
    return Element{TextElement{
        .content = std::string{std::string_view{S}},
        .style   = Sty.runtime(),
    }};
}
\end{lstlisting}

The compile-time string \code{S} and style \code{Sty} are read out
of the type. \code{S} becomes a \code{std::string} (heap-allocated
if the small-string optimisation does not apply; for short banners
it usually does). \code{Sty.runtime()} converts the
\code{CTStyle} bag into a runtime \code{Style} object (Chapter~\ref{ch:elements}).
The result is a \code{TextElement} wrapped in the \code{Element}
variant.

\paragraph{\code{RuntimeTextNode<S>::build()}.}

\begin{lstlisting}
Element build() const {
    if constexpr (std::integral<S> && !std::same_as<S, bool>) {
        // format integer via std::to_chars
    } else if constexpr (std::floating_point<S>) {
        // format float via std::to_chars, 2 fractional digits
    } else {
        // S is a string-like: copy into a std::string
    }
}
\end{lstlisting}

The \code{if constexpr} picks the right path at compile time, so
there is no per-call branching. \code{text(42)} formats with no
heap allocation (uses a 32-byte stack buffer); \code{text("hi")}
allocates a \code{std::string}.

\paragraph{\code{BoxNode<Dir, Cfg, Cs...>::build()}.} The most
involved:

\begin{lstlisting}
Element build() const {
    return std::apply([](const auto&... cs) {
        auto b = maya::detail::box().direction(Dir);
        if constexpr (Cfg.pad_t || Cfg.pad_r || ...)
            b = std::move(b).padding(Cfg.pad_t, ...);
        if constexpr (Cfg.gap > 0)
            b = std::move(b).gap(Cfg.gap);
        if constexpr (Cfg.has_border)
            b = std::move(b).border(Cfg.bstyle);
        if constexpr (Cfg.has_border_color)
            b = std::move(b).border_color(Color::rgb(Cfg.bc_r, Cfg.bc_g, Cfg.bc_b));
        // ... style, dimensions, etc.
        return b(cs.build()...);   // fold-expression: build each child
    }, children);
}
\end{lstlisting}

Notice the \code{if constexpr} on every \code{BoxCfg} field. A box
with no border generates \emph{no} call to \code{.border()}; the
branch does not exist in the resulting binary. The compiler
specialises one \code{build()} per distinct \code{BoxCfg} and
throws out everything irrelevant.

The trailing \code{b(cs.build()...)} expands the children pack:
every child's \code{.build()} is invoked, and the resulting
\code{Element}s are handed to the runtime box builder via its
function-call operator (Chapter~\ref{ch:elements}).

\paragraph{\code{DynNode<F>::build()}.}

\begin{lstlisting}
Element build() const { return fn(); }
\end{lstlisting}

Literal one-liner. The lambda is invoked; whatever it returns is
converted to \code{Element}.

\paragraph{\code{MapNode<R, Proj>::build()}.}

\begin{lstlisting}
Element build() const {
    std::vector<Element> items;
    if constexpr (std::ranges::sized_range<R>) items.reserve(range.size());
    for (auto&& val : range) items.emplace_back(proj(val).build());
    return Element{ElementList{std::move(items)}};
}
\end{lstlisting}

The range is walked, the projection applied, and the results
gathered into an \code{ElementList}. Because \code{MapNode} produces
an \code{ElementList} (not a single \code{Element}), the parent
\code{v}/\code{h} splices its children into the parent's child list
--- the moral equivalent of a JSX fragment.

\subsection{The cost of \code{.build()}}

A pure compile-time tree (\code{t<"\dots">} and \code{v}/\code{h}
with no runtime children) does \code{.build()} in:

\begin{itemize}[leftmargin=*]
  \item One \code{std::string} construction per text node
    (small-string optimisation often makes this free).
  \item One \code{BoxElement} construction per box node, which holds
    a \code{std::vector<Element>} of children. The vector is
    reserved to the right size.
  \item Some \code{Style} field assignments.
\end{itemize}

No virtual dispatch, no template-method indirection. The whole tree
is built in a single recursive sweep that the compiler often inlines
completely for small trees. For a typical panel with twenty nodes,
\code{.build()} runs in tens of microseconds.

The \code{Element} tree is then dropped on the floor at the end of
the frame (after rasterisation into the canvas), and the next frame
builds a fresh one. This is the rebuild-every-frame discipline from
Chapter~\ref{ch:elm}.

\begin{recap}
\code{.build()} is the bridge from the compile-time tree to the
runtime \code{Element}. It is a one-shot recursive walk: every
node's \code{build()} produces an \code{Element}, children are
gathered into the parent's \code{Element}, and a runtime tree comes
out the bottom. Then it is the runtime side's problem.
\end{recap}

\section{Recipe gallery}
\label{sec:dsl:recipes}

A collection of patterns you will use constantly. Each is a small,
composable fragment; the goal is to make them feel familiar before
you need them.

\subsection{A titled panel with a border}

\begin{lstlisting}
constexpr auto title_panel = v(
    t<"Status"> | Bold | Underline,
    sep,
    t<"all systems nominal"> | Fg<80, 220, 120>
) | pad<1, 2> | border_<Round> | bcol<60, 65, 80>;
\end{lstlisting}

\code{sep} is a horizontal separator (a one-row box with a top
border); \code{pad<1, 2>} is one row of vertical padding, two
columns of horizontal.

\subsection{A two-column layout that grows to fit}

\begin{lstlisting}
auto split = h(
    sidebar  | w_<24>,
    content  | grow_<1>
) | grow_<1>;
\end{lstlisting}

The sidebar has fixed width 24; the content gets every remaining
column. The outer box grows to fill whatever vertical container is
holding it.

\subsection{A list of items from a runtime collection}

\begin{lstlisting}
auto list = v(
    t<"Files"> | Bold,
    map(files, [](const File& f) {
        return h(text(f.name) | w_<30>,
                 text(f.size) | Dim);
    })
);
\end{lstlisting}

The lambda is called once per file; each call produces an
\code{Element}. \code{map} gathers them into an \code{ElementList}
that \code{v} splices in.

\subsection{A conditional region}

Two forms, both common:

\begin{lstlisting}
// when(cond, then) --- hidden if cond is false
auto top = v(
    title,
    when(show_advanced, advanced_panel)
);

// when(cond, then, else) --- one or the other
auto status = when(loading,
    text("loading\dots") | Dim,
    text("ready")        | Fg<80, 220, 120>);
\end{lstlisting}

The one-argument form falls back to \code{BlankNode\{\}} (a
zero-cell empty text). The two-argument form picks one of two
branches; both branches must satisfy \code{Node}.

\subsection{A reusable style preset}

\begin{lstlisting}
constexpr auto heading = Bold | fg<0xFFFFFF>;
constexpr auto error   = Bold | fg<0xFF4444>;
constexpr auto subtle  = Dim  | Italic;

auto ui = v(
    t<"Welcome"> | heading,
    text(err)    | error,
    t<"v0.1">    | subtle
);
\end{lstlisting}

The presets are themselves \code{StyTag}s, computed at compile time.
Applying one is a single token; defining your design system this
way is the path to consistency.

\subsection{Inline width hints on text}

\begin{lstlisting}
text(very_long_string) | clip | w_<30>
\end{lstlisting}

\code{clip} switches the text node into truncate-with-ellipsis
mode. \code{w\_<30>} wraps the result in a fixed-width box, so the
text is forced to 30 columns regardless of available space.

\subsection{An animated spinner}

\begin{lstlisting}
dyn([&] {
    constexpr std::array frames = {"top", "upper-right", "lower-right", "lower-left"};
    return text(frames[(tick / 100) % frames.size()]);
}) | Fg<140, 220, 255>
\end{lstlisting}

The lambda picks a frame from the array based on the current tick
count. Combine with \code{Sub<Msg>::on\_animation\_frame(Tick\{\})}
from Chapter~\ref{ch:elm} for a 60\,Hz spinner.

\subsection{A right-aligned status bar}

\begin{lstlisting}
auto status_bar = h(
    text(left_text),
    space,                          // SpacerNode: grows to fill
    text(right_text) | Dim
);
\end{lstlisting}

\code{space} is the spacer node: a zero-content box with
\code{grow=1}. Between two non-growing siblings it stretches to
push them to opposite ends of the row.

\section{Why all of this, again?}
\label{sec:dsl:why-again}

At the end of a long chapter on templates, fold expressions, and
type-state machines, it is worth stepping back. What did all that
machinery buy us?

At build time:

\begin{itemize}[leftmargin=*]
  \item \textbf{No allocation for the tree shape.} The structure of
    your UI lives in types; only the leaf values (runtime strings,
    runtime children) allocate.
  \item \textbf{Errors are catch-able before the program runs.}
    Border colour without border, negative padding, wrong child
    type, hex overflow --- all rejected by the compiler.
  \item \textbf{The DSL composes with the host language.} Helper
    functions, lambdas, ranges --- all work without special
    integration.
\end{itemize}

At runtime:

\begin{itemize}[leftmargin=*]
  \item \textbf{\code{.build()} is one linear walk}, specialised by
    the compiler per distinct configuration. No virtual dispatch.
  \item \textbf{The runtime side is simple.} It receives a plain
    tree of \code{Element} values and lays it out; it does not
    care that the tree was generated by a compile-time DSL.
\end{itemize}

For build time. For correctness. For speed. Three things at once,
because we paid the cost in the DSL design.

\begin{idea}
The \maya{} DSL is what you get when you take ``the compiler should
catch this'' seriously and follow it all the way down. Strings,
styles, sizes, structure: in the types. Mistakes: caught at build.
Runtime: cheap and uniform. The trade is some up-front complexity in
the DSL itself; the payoff is every \maya{} program you ever write.
\end{idea}

\section*{Workshop}

\begin{enumerate}[leftmargin=*]
  \item Write a panel that shows your name in bold, an underlined
    tagline, and a divided list of three hobbies. Use only the
    compile-time pipes.
  \item Add a runtime counter (an \code{int} variable) with
    \code{dyn}. Increment it in a loop with
    \code{maya::live(...)} (covered in Chapter~\ref{ch:cmdsub}) and
    watch the panel refresh.
  \item Try \code{bcol<60, 65, 80>} on a box without a border. Read
    the compiler error. Find the line that says \code{constraints
    not satisfied} and the \code{Cfg.has\_border} that triggered it.
\end{enumerate}

\chapter{Elements, style, colour, border}
\label{ch:elements}

\code{.build()} converts the compile-time tree from
Chapter~\ref{ch:dsl} into an \code{Element}. That \code{Element} is
the input to the layout engine, the renderer, and every widget. This
chapter dissects \code{Element} and the value types it depends on.

\begin{aside}
Chapter~\ref{ch:dsl} was about the compile-time side --- the
elaborate type machinery that makes \code{t<"hi"> | Bold} legal and
\code{v(\dots) | bcol<\dots>} (without a border) illegal. This
chapter is about the runtime side: the small, boring, fast data
structure the renderer actually consumes. The contrast is part of
the trick. The compile-time DSL is intricate because we want the
compiler to do interesting work; the runtime representation is
plain because we want the renderer to be fast.
\end{aside}

\section{The runtime tree, in one sentence}
\label{sec:elem:sentence}

Everything in this chapter rotates around one idea: an
\textbf{element tree} is a tree where every node is one of three
things.

\begin{itemize}[leftmargin=*]
  \item \textbf{Text} --- a leaf, holding bytes and styling.
  \item \textbf{Box} --- a container, holding a layout configuration
    and a list of child elements.
  \item \textbf{List} --- a transparent grouping, ``these elements
    belong to my parent's children'', used for splicing.
\end{itemize}

Nothing more. Every shape you have seen \maya{} produce ---
dashboards, list views, scrolling chats, modal dialogs --- is built
out of those three nodes nested arbitrarily.

\section{What is a tagged union?}
\label{sec:elem:tagged-union}

Before we go into the variant: a small piece of language theory you
will use over and over in this book.

A \textbf{tagged union} (also called a \emph{sum type} or
\emph{discriminated union}) is a value that is \emph{one of} a fixed
set of alternatives, with a tag indicating which one. The phrase
``one of'' matters: at any moment, the value is exactly one
alternative, never two, never zero. The compiler can check that you
have handled every case.

Contrast with the other big composite type:

\begin{itemize}[leftmargin=*]
  \item A \textbf{product type} (a struct, class, tuple) is
    \emph{this and this and this}. A \code{Point\{int x, int y\}}
    has both an \code{x} \emph{and} a \code{y}, always.
  \item A \textbf{sum type} (a tagged union, variant) is
    \emph{this or this or this}. A \code{Result} is either an
    \code{Ok} or an \code{Error}, never both at once.
\end{itemize}

Natural-language analogies break almost immediately when you push
them, but here is a small one: a parking space's state is a sum
type (\emph{Empty} or \emph{Occupied}, never both); a car's state
is a product type (\emph{has colour} \textbf{and} \emph{has
licence plate}, both at once).

In modern \cpp{} the sum-type primitive is \code{std::variant<A, B,
C>}, which holds exactly one of \code{A}, \code{B}, or \code{C} and
remembers which. You inspect it with \code{std::visit}, which calls
your visitor with the actual stored alternative. The compiler
verifies that your visitor handles every case --- forget one and
you get an error.

\begin{lstlisting}
#include <variant>
using Pet = std::variant<Dog, Cat, Hamster>;

std::string describe(const Pet& p) {
    return std::visit(overload{
        [](const Dog& d)     { return "a good boy named " + d.name; },
        [](const Cat& c)     { return c.name + " who is judging you"; },
        [](const Hamster& h) { return "a hamster named " + h.name; },
    }, p);
}
\end{lstlisting}

Three branches, three alternatives. Add a fourth alternative
(\code{Goldfish}) and the visitor stops compiling --- the compiler
lists the missing branch by name.

\begin{idea}
A sum type is the data-structures version of an \code{if}/\code{else}
chain. It is \emph{closed} (the alternatives are listed in the
type), \emph{exhaustive} (the compiler nags you about missing
cases), and \emph{cheap} (one word for the tag, plus enough storage
for the largest alternative).
\end{idea}

\subsection{Why \maya{} uses sum types everywhere}

Look through the source and you will find sum types in almost every
file:

\begin{itemize}[leftmargin=*]
  \item \code{Element = variant<BoxElement, TextElement, ElementList>}
    --- this chapter.
  \item \code{Cmd<Msg>::Variant} --- algebra of effects
    (Chapter~\ref{ch:elm}).
  \item \code{Sub<Msg>::Variant} --- algebra of subscriptions
    (Chapter~\ref{ch:elm}).
  \item \code{Event = variant<KeyEvent, MouseEvent, ResizeEvent, PasteEvent>}
    --- input events (Chapter~\ref{ch:events}).
  \item \code{RenderOp = variant<EmitText, MoveCursor, SetStyle, ...>}
    --- renderer output (Chapter~\ref{ch:ansi}).
  \item \code{expected<T, E>} = \code{variant<T, E>} essentially ---
    success or failure (Chapter~\ref{ch:cpp}).
\end{itemize}

The reason is uniform: each of these is a finite, named set of
possibilities. Encoding them as sum types gives you exhaustiveness
checking and removes whole classes of \code{nullptr}-ish bugs. When
you learn to recognise the pattern, you start seeing it everywhere
in well-designed C++ codebases.

\section{\code{Element}: a tagged union}

\code{Element} is a \code{std::variant} of three alternatives:

\begin{lstlisting}
using Element = std::variant<BoxElement, TextElement, ElementList>;
\end{lstlisting}

\begin{itemize}[leftmargin=*]
  \item \code{TextElement} --- a piece of styled text, possibly with
    a wrap policy.
  \item \code{BoxElement} --- a container with flex layout, padding,
    border, gap, etc., and a vector of child elements.
  \item \code{ElementList} --- a flat sequence of elements that
    should be inlined into the parent's children (used by \code{map}).
\end{itemize}

Three alternatives, a closed set. Every algorithm in \maya{} that
walks an element tree (layout, render, hit-testing, search) ends up
in a \code{std::visit} on this variant. Add a fourth alternative
tomorrow and the compiler tells you all the places to update.

\subsection{\code{TextElement}}

\begin{lstlisting}
struct TextElement {
    std::string text;        // the bytes (UTF-8)
    Style       style;       // colours + attributes
    TextWrap    wrap = TextWrap::Wrap;
    HyperlinkId link{};      // optional OSC-8 link
};
\end{lstlisting}

A leaf in the tree. The string is owned (so the source data can be a
temporary). Styling lives in a \code{Style} struct (next section).
\code{TextWrap} is one of \code{Wrap}, \code{NoWrap}, \code{Truncate},
\code{Ellipsize}.

\subsection{\code{BoxElement}}

\begin{lstlisting}
struct BoxElement {
    layout::FlexStyle    style;     // direction, padding, gap, ...
    std::optional<Border> border;
    std::vector<Element>  children;
    std::optional<std::string> title;  // border title
    // ... a few more renderer hints
};
\end{lstlisting}

The composite. \code{style} carries every layout property as integer
cell counts; \code{border} is set when the DSL chain included
\code{border\_<...>}. \code{title} is set when a border title was
piped on.

\subsection{\code{ElementList}}

\begin{lstlisting}
struct ElementList { std::vector<Element> items; };
\end{lstlisting}

A pseudo-element. The renderer treats it as ``these children belong
to my parent.'' It is the runtime equivalent of a parameter pack
expansion.

\section{\code{Style}}

\code{Style} is a small POD struct describing the appearance of a run
of text:

\begin{lstlisting}
struct Style {
    Color     fg{};
    Color     bg{};
    Attribute attrs{};   // bold | italic | underline | dim | ...

    constexpr Style with_bold()      const noexcept;
    constexpr Style with_italic()    const noexcept;
    constexpr Style with_fg(Color c) const noexcept;
    constexpr Style with_bg(Color c) const noexcept;
    constexpr Style merge(Style other) const noexcept;
    // ...
};
\end{lstlisting}

Two flavours exist:

\begin{itemize}[leftmargin=*]
  \item \code{Style} (runtime), used in \code{TextElement}.
  \item \code{CTStyle} (compile-time), used as an NTTP in
    \code{TextNode}. It is structurally identical but lives in
    template parameters and is constructed with \code{constexpr}
    factory functions (\code{CTStyle::bold()},
    \code{CTStyle::fg(R, G, B)}, etc.).
\end{itemize}

The DSL pipes operate on \code{CTStyle}; \code{.build()} calls
\code{CTStyle.runtime()} to project back to a \code{Style}. The
distinction matters because NTTPs require structural types whose
fields are compared by value, while runtime \code{Style} can have
fancier constructors.

\subsection{\code{Attribute} bitset}

Attributes (bold, dim, italic, underline, strikethrough, inverse,
blink) are a bitset:

\begin{lstlisting}
enum class Attribute : std::uint8_t {
    None       = 0,
    Bold       = 1 << 0,
    Dim        = 1 << 1,
    Italic     = 1 << 2,
    Underline  = 1 << 3,
    Strike     = 1 << 4,
    Inverse    = 1 << 5,
    Blink      = 1 << 6,
};

constexpr Attribute operator|(Attribute a, Attribute b) noexcept;
constexpr bool      has(Attribute set, Attribute bit) noexcept;
\end{lstlisting}

Combining attributes is bitwise OR. Checking is bitwise AND. The
renderer translates each bit to its SGR parameter (1 for bold, 2 for
dim, 3 for italic, etc.) when serialising.

\section{\code{Color}}

A \code{Color} is a 24-bit RGB triple plus a tag for ``default''
(unspecified, inherit from terminal):

\begin{lstlisting}
struct Color {
    std::uint8_t r = 0, g = 0, b = 0;
    bool         is_default = true;   // means "use terminal default"
};

constexpr Color rgb(uint8_t r, uint8_t g, uint8_t b) noexcept {
    return {r, g, b, false};
}
\end{lstlisting}

\maya{} only deals in 24-bit truecolour internally. The 256-colour and
16-colour palettes were popular when terminals had limited support;
modern terminals all do truecolour. If you really need 256-colour
output (because you are targeting a niche emulator), the writer can
quantise on the way out.

\subsection{Themes}

A \textbf{theme} is a 24-slot palette of colours --- one for primary,
secondary, success, warning, error, plus background tiers and text
tiers. \maya{}'s built-in themes (\code{dark}, \code{light},
\code{terminal}) are \code{constexpr} aggregates; you can derive your
own and pass it in \code{RunConfig}.

\begin{lstlisting}
struct Theme {
    Color bg_primary;
    Color bg_secondary;
    Color text_primary;
    Color text_secondary;
    Color accent;
    Color success;
    Color warning;
    Color error;
    // ... 24 slots total
};
\end{lstlisting}

The \code{Ctx} parameter to a render function (in the quick-tool
\code{run(...)}) carries the active theme. Widgets read it instead of
hard-coding colours --- so the same widget tree looks right under
dark and light themes.

\section{\code{Border}}

A border is described by its style (single, double, round, thick,
classic, arrow), its colour (or default), and which sides are drawn:

\begin{lstlisting}
enum class BorderStyle : std::uint8_t {
    None, Single, Double, Round, Bold, Classic, Arrow,
};

struct BorderSides {
    bool top = true, right = true, bottom = true, left = true;

    static constexpr BorderSides all()        { return {true, true, true, true}; }
    static constexpr BorderSides none()       { return {false,false,false,false}; }
    static constexpr BorderSides horizontal() { return {true,false,true,false}; }
    static constexpr BorderSides vertical()   { return {false,true,false,true}; }
};

struct Border {
    BorderStyle style = BorderStyle::Single;
    Color       colour{};
    BorderSides sides = BorderSides::all();
};
\end{lstlisting}

Each \code{BorderStyle} maps to a set of eight glyphs (corners + edges)
that the renderer paints. The mapping is a \code{constexpr}
\code{std::array} in \file{maya/style/border.hpp}.

\subsection{\code{Sep} and \code{VSep}}

The DSL's \code{sep} and \code{vsep} nodes are syntactic sugar for a
box with a one-side-only border --- horizontal or vertical. The same
machinery; the renderer treats them identically.

\section{\code{BoxBuilder}: runtime construction}

The DSL \code{.build()} fans out into a \code{BoxBuilder} which
assembles a \code{BoxElement}:

\begin{lstlisting}
class BoxBuilder {
public:
    BoxBuilder(FlexDirection dir, BoxCfg cfg)
      : dir_(dir), cfg_(cfg) {}

    void add(Element child) { children_.push_back(std::move(child)); }
    void add(ElementList list) {
        children_.insert(children_.end(),
                         std::make_move_iterator(list.items.begin()),
                         std::make_move_iterator(list.items.end()));
    }

    Element finalize() &&;
private:
    FlexDirection         dir_;
    BoxCfg                cfg_;
    std::vector<Element>  children_;
};
\end{lstlisting}

Two overloads of \code{add}: one for ordinary children, one for an
\code{ElementList} which is spliced. \code{finalize()} translates
\code{BoxCfg} into a \code{layout::FlexStyle}, attaches a \code{Border}
if requested, wraps it in a \code{BoxElement}, and returns the
variant.

\section{Walking an element tree}

Every algorithm that processes an element tree is a recursive
\code{std::visit}. Here is the shape of a pretty-printer:

\begin{lstlisting}
void print(const Element& e, int depth = 0) {
    std::string indent(depth * 2, ' ');
    std::visit(overload{
        [&](const TextElement& t) {
            std::println("{}TEXT  : \"{}\"", indent, t.text);
        },
        [&](const BoxElement& b) {
            std::println("{}BOX   : dir={}, {} children",
                         indent, int(b.style.direction), b.children.size());
            for (const auto& c : b.children) print(c, depth + 1);
        },
        [&](const ElementList& l) {
            for (const auto& c : l.items) print(c, depth);
        }
    }, e);
}
\end{lstlisting}

This is the template every renderer-side pass follows: visit, branch
on alternative, recurse on children.

\section{Hyperlinks}

A small but neat feature: text elements can carry a hyperlink. The
renderer emits an OSC-8 sequence around the run:

\begin{lstlisting}
ESC ] 8 ; ; URL ESC \  text  ESC ] 8 ; ; ESC \
\end{lstlisting}

Modern emulators turn that into a clickable link. \maya{} interns the
URL into a \code{HyperlinkId} so the per-cell representation stays
compact (more on that in Chapter~\ref{ch:canvas}).

\section{Why so many small structs?}

A reasonable question: why isn't \code{Element} a class with virtual
methods? \maya{}'s answer: virtual dispatch precludes a packed cell
representation, prevents \code{constexpr} construction, and makes
serialization harder. The variant + visitor pattern gives you
polymorphism without the overhead --- \code{std::visit} compiles to a
jump table or a chain of branches, both faster than a virtual call,
and the data layout stays cache-friendly.

\section{The forward-declaration dance}
\label{sec:elem:forward-decl}

If you open \file{maya/include/maya/element/element.hpp} you will see
an unusual little ritual: \code{Element} is forward-declared,
\code{box.hpp} and \code{text.hpp} are then included, and only after
that is \code{Element} actually defined. The comments in the file
explain why; the explanation is worth understanding because the same
problem will come up in your own code.

The shape is:

\begin{lstlisting}
// Step 1: forward-declare Element
struct Element;
struct ElementList;

// Step 2: include the concrete element headers
#include "box.hpp"      // BoxElement, which has std::vector<Element>
#include "text.hpp"     // TextElement

// Step 3-4: now define Element and ElementList
struct ElementList { std::vector<Element> items; ... };
struct Element     { /* the variant */ };
\end{lstlisting}

Why this order? Because \code{BoxElement} contains
\code{std::vector<Element>}, which means \code{BoxElement}'s
definition needs to know about \code{Element}. But \code{Element} is
the variant of \code{BoxElement}, \code{TextElement}, and
\code{ElementList}, so its definition needs to know about
\code{BoxElement} --- a circular dependency.

\cpp{} resolves this with \textbf{forward declarations}: a
forward-declared type is enough for the compiler to know
``\code{Element} is a type, you can have pointers and references to
it, and \code{std::vector<Element>} works because the vector holds
an indirect pointer to its element type internally''. The full
definition is supplied later, after the dependent pieces are in
place.

The technique generalises: any time you have a tree where a node's
field holds a vector / optional / pointer of the same type, you
use forward declaration to break the cycle.

\begin{idea}
A tree type that contains itself in a field is a circular type
definition. \cpp{} breaks the circle with a forward declaration:
the outer type is named first, the indirect field can then refer to
it, and the full definition follows.
\end{idea}

\section{Constructing an \code{Element} by hand}
\label{sec:elem:by-hand}

The DSL hides almost all element construction. For tests,
diagnostics, and writing widget primitives that work below the DSL
layer, you sometimes need to make an \code{Element} directly.

The direct way:

\begin{lstlisting}
#include <maya/element/element.hpp>

using namespace maya;

// A leaf text element
Element title = TextElement{
    .text  = "Hello",
    .style = Style{}.with_bold().with_fg(Color::rgb(100, 180, 255)),
};

// A box containing two children
Element panel = BoxElement{
    .style    = layout::FlexStyle{ .direction = FlexDirection::Column,
                                   .padding   = {.top=1, .right=2,
                                                 .bottom=1, .left=2} },
    .border   = Border{.style = BorderStyle::Round},
    .children = {
        std::move(title),
        Element{TextElement{.text = "world"}}
    },
};
\end{lstlisting}

The braces in the field initialisers are designated initialisers
(\cpp{}20); they make the field names explicit so the code reads
like data. The whole construction is value-based: no \code{new},
no \code{shared\_ptr}, just nested structs.

The more idiomatic way uses the \code{box()} builder:

\begin{lstlisting}
#include <maya/element/builder.hpp>

Element panel = maya::detail::box()
    .direction(FlexDirection::Column)
    .padding(1, 2, 1, 2)
    .border(BorderStyle::Round)
    .border_color(Color::rgb(60, 65, 80))
    (text_el("Hello").bold().fg(100, 180, 255),
     text_el("world"));
\end{lstlisting}

The builder is a fluent chain that mirrors the DSL pipes; its final
\code{operator()} takes the children. This is what the DSL's
\code{.build()} ultimately calls --- the builder lives in
\file{maya/include/maya/element/builder.hpp}.

\section{The colour model up close}
\label{sec:elem:color}

We established that \code{Color} is a 24-bit RGB triple plus an
``is default'' bit. Worth pausing on what that bit costs and what it
gains.

Without it, you would need a sentinel: \code{Color\{0, 0, 0\}}
meaning either ``actual black'' or ``unset''. The runtime would
have no way to tell ``the user explicitly chose black'' from ``the
user hasn't specified'' --- and the renderer needs to know, because
the terminal's default foreground colour may not be black (it is
white-on-black for dark themes, dark-on-white for light themes).

With \code{is\_default = true}, the renderer can emit
\code{ESC[39m} (\emph{reset foreground to terminal default}) rather
than a specific colour. The cell inherits whatever the user's
terminal preferences say.

\begin{lstlisting}
Color c1 = Color::rgb(0, 0, 0);                // explicit black
Color c2 = Color::default_();                  // terminal default
Color c3{};                                    // also terminal default
// c1.is_default == false, c2.is_default == true
\end{lstlisting}

The difference matters when a \maya{} program runs in someone
else's terminal: their carefully-chosen colour scheme is preserved
for anything you did not explicitly style.

\subsection{The hex shorthand}

Designers think in hex; the runtime has a constructor for that:

\begin{lstlisting}
Color accent = Color::hex(0x64B4FF);    // = rgb(100, 180, 255)
Color danger = Color::hex(0xFF4444);    // = rgb(255, 68, 68)
\end{lstlisting}

The DSL has matching \code{fg<0x\dots>} and \code{bg<0x\dots>}
pipes (Chapter~\ref{ch:dsl}) which call the same constructor at
compile time.

\subsection{Truecolour, 256-colour, 16-colour}

A brief detour. The terminal world has three historical colour
palettes:

\begin{itemize}[leftmargin=*]
  \item \textbf{8 / 16 colours.} The classic ANSI palette. SGR codes
    30--37 (foreground), 40--47 (background), 90--97 (bright
    foreground), 100--107 (bright background). Universally
    supported.
  \item \textbf{256 colours.} An xterm extension; a 6×6×6 colour
    cube plus 24 greys plus the original 16. Encoded as
    \code{ESC[38;5;<n>m}. Most emulators support it.
  \item \textbf{Truecolour (24-bit).} Three bytes per channel.
    Encoded as \code{ESC[38;2;<R>;<G>;<B>m}. Modern emulators ---
    every one made in the last decade --- support this.
\end{itemize}

\maya{} stores colours internally as truecolour and lets the
renderer worry about quantisation if the target terminal does not
support it. In practice, ``the target terminal does not support
it'' is so rare you can ignore it; if you really need to target a
2003 emulator over a hardened serial link, set
\code{RunConfig::color\_mode = ColorMode::ANSI16} and a quantiser
will snap each truecolour to the nearest ANSI colour on the wire.

\begin{funfact}
24-bit truecolour over a terminal was proposed in the early 2000s
and snuck into emulators piecemeal over the next decade. The
official standard governing it (\code{ITU-T T.416}) was published
in 1993 but ignored for fifteen years. Today, your
\code{printf '\textbackslash{}e[38;2;255;128;0mhi'} works in every
mainstream emulator without anyone needing to think about it. The
gap between a feature existing and it being usable is sometimes
decades long.
\end{funfact}

\section{The border gallery}
\label{sec:elem:border-gallery}

Each \code{BorderStyle} maps to a set of glyphs the renderer paints
at the cells immediately inside the box's bounding rectangle. The
styles are not interchangeable; the glyph choice signals affect.

\paragraph{\code{None}.} No border, no glyphs. The box's content
fills its rectangle.

\paragraph{\code{Single}.} The default. Light-weight box-drawing
characters from Unicode's \code{U+2500}--\code{U+257F} block:
horizontal bar (\code{U+2500}), vertical bar (\code{U+2502}),
top-left corner (\code{U+250C}), and so on.

\paragraph{\code{Round}.} Same edges as \code{Single}, but the four
corners use the rounded variants (\code{U+256D}, \code{U+256E},
\code{U+256F}, \code{U+2570}). Looks friendlier on most fonts.

\paragraph{\code{Double}.} Double-line glyphs (\code{U+2550} for
horizontal, \code{U+2551} for vertical, \code{U+2554} for top-left).
The classic \emph{IBM} look from DOS-era applications.

\paragraph{\code{Bold} (alias \code{Thick}).} Heavy box-drawing
glyphs (\code{U+2501}, \code{U+2503}, \code{U+250F}). Useful for
emphasis or selected state.

\paragraph{\code{Classic}.} ASCII-only fallback using \code{+},
\code{-}, \code{|}. Use for legacy terminals that do not render
Unicode box-drawing characters.

\paragraph{\code{Arrow}.} A novelty style with arrow heads pointing
into the corners. Good for ``selected'' or ``focus'' indicators on a
subpanel.

The glyph table for every style lives in
\file{maya/include/maya/style/border.hpp} as a \code{constexpr}
\code{std::array}, one row per style. You can read all of it in one
minute.

\subsection{Partial-side borders}

\code{BorderSides} controls which of the four edges actually paint.
The defaults are all four; the helpers \code{horizontal()} and
\code{vertical()} pick top+bottom or left+right respectively. This
is how \code{sep} and \code{vsep} are implemented --- they are
boxes with one side worth of border, drawn as a thin separator.

\begin{lstlisting}
// A separator under a heading
auto heading = v(
    t<"Section"> | Bold,
    sep,                       // = a row with .border.sides = horizontal()
    body
);
\end{lstlisting}

A separator costs one row of height. Vertical separators (\code{vsep})
cost one column of width and appear inside row layouts.

\section{Patterns for walking the tree}
\label{sec:elem:walk-patterns}

The pretty-printer earlier was the simplest possible tree walk. Real
algorithms in \maya{} have a few additional patterns; learning them
lets you build your own analyses (linters, accessibility checks,
layout debuggers).

\subsection{Counting nodes}

\begin{lstlisting}
int count(const Element& e) {
    return std::visit(overload{
        [](const TextElement&)    { return 1; },
        [&](const BoxElement& b)  {
            int n = 1;
            for (const auto& c : b.children) n += count(c);
            return n;
        },
        [&](const ElementList& l) {
            int n = 0;
            for (const auto& c : l.items) n += count(c);
            return n;
        },
    }, e);
}
\end{lstlisting}

\code{ElementList} does not count itself (it is a transparent
fragment), but its items do. \code{BoxElement} counts itself plus
its children.

\subsection{Finding all text}

Flatten an \code{Element} into a vector of strings (useful for
testing, search, screen readers):

\begin{lstlisting}
void texts(const Element& e, std::vector<std::string>& out) {
    std::visit(overload{
        [&](const TextElement& t) { out.push_back(t.text); },
        [&](const BoxElement& b)  {
            for (const auto& c : b.children) texts(c, out);
        },
        [&](const ElementList& l) {
            for (const auto& c : l.items) texts(c, out);
        },
    }, e);
}
\end{lstlisting}

The same scaffold; the only thing that changes is what you do at
each node.

\subsection{Transforming the tree}

What if you want to return a \emph{new} tree, with every text node
uppercased?

\begin{lstlisting}
Element upper(const Element& e) {
    return std::visit(overload{
        [](const TextElement& t) -> Element {
            auto copy = t;
            std::transform(copy.text.begin(), copy.text.end(),
                           copy.text.begin(), ::toupper);
            return copy;
        },
        [&](const BoxElement& b) -> Element {
            BoxElement nb = b;
            for (auto& c : nb.children) c = upper(c);
            return nb;
        },
        [&](const ElementList& l) -> Element {
            ElementList nl = l;
            for (auto& c : nl.items) c = upper(c);
            return nl;
        },
    }, e);
}
\end{lstlisting}

The shape is the same again --- visit, recurse on children, build
the new node. With trees of plain values, transformation is
non-destructive and trivially correct: the original tree is
untouched.

\subsection{Hit testing}

Given a screen coordinate, which \code{Element} is at that position?
The layout engine (Chapter~\ref{ch:layout}) assigns each node a
rectangle; hit testing walks the tree picking the deepest node
whose rectangle contains the point. The pattern is identical:
\code{std::visit}, recurse, return.

\begin{idea}
Every tree pass in \maya{} is \code{std::visit} + recursion. The
variant is closed, the visitor is exhaustive, and adding a new node
kind would force every visitor in the codebase to add a branch ---
in one compiler error message, with one source location each. This
is why the variant pattern scales.
\end{idea}

\section{The \code{Style} merge rule}
\label{sec:elem:style-merge}

A detail that catches people out: how does \code{Style} composition
actually work? Two styles, one inheriting from a parent and one
applied locally, must combine somehow.

The rule, implemented by \code{Style::merge}:

\begin{itemize}[leftmargin=*]
  \item Colours: the second style's \code{fg} / \code{bg} replace
    the first's \emph{only if they are not default}. The default
    sentinel acts as ``do not override''.
  \item Attributes: a bitwise OR. Bold can be added on top of dim;
    underline can be added on top of bold. There is no ``unbold''
    attribute in the merged form --- styles only accumulate.
\end{itemize}

In practice this means:

\begin{lstlisting}
Style base = Style{}.with_fg(Color::rgb(100, 100, 100));
Style bold = Style{}.with_bold();
Style merged = base.merge(bold);
// merged.fg    == rgb(100, 100, 100)
// merged.attrs == Bold
\end{lstlisting}

The parent's grey foreground is preserved (bold did not specify a
colour), and the bold attribute is added. If you later wanted to
actively reset bold, you would do it by constructing a fresh style,
not by ``negating'' in the merge.

\section{Hyperlinks}

A small but neat feature: text elements can carry a hyperlink. The
renderer emits an OSC-8 sequence around the run:

\begin{lstlisting}
ESC ] 8 ; ; URL ESC \  text  ESC ] 8 ; ; ESC \
\end{lstlisting}

Modern emulators turn that into a clickable link. \maya{} interns the
URL into a \code{HyperlinkId} so the per-cell representation stays
compact (more on that in Chapter~\ref{ch:canvas}).

\begin{lstlisting}
Element docs = TextElement{
    .text = "see docs",
    .style = Style{}.with_underline(),
    .link  = intern_hyperlink("https://example.org/docs"),
};
\end{lstlisting}

The \code{intern\_hyperlink} call returns a small integer; the URL
string lives in a per-frame interning table, freed at end of frame.
This keeps the per-cell representation a small fixed-size struct
rather than a string.

\section{Pitfalls and anti-patterns}
\label{sec:elem:pitfalls}

\begin{warning}
\textbf{Storing an \code{Element} across frames.} The element tree is
built fresh on every frame; storing a node from frame \code{N} for
use in frame \code{N+1} is not catastrophic, but it defeats the
rebuild-from-model discipline and is usually a sign your model is
missing a field. Put the data in the model, build the element from
it.
\end{warning}

\begin{warning}
\textbf{Mutating an \code{Element} after construction.} The runtime
may assume the element you returned from \code{view()} is
immutable for the duration of the frame. Modifying it after the
fact (from a background thread, say) is undefined behaviour. If you
need to ``change the UI'' from a background thread, send a
\code{Msg} through \code{Cmd::task}'s dispatch; \code{view} will be
re-invoked with the new model.
\end{warning}

\begin{warning}
\textbf{Using raw heap allocation.} \code{new BoxElement\{\dots\}}
and returning a pointer breaks the value-semantics invariant the
rest of the library assumes. The \code{Element} variant is meant to
be constructed by value and moved; mixing pointer ownership in
will break copy / move semantics in confusing ways.
\end{warning}

\begin{warning}
\textbf{Building deeply nested boxes when a flat list would do.}
Layout time grows roughly linearly with node count. A list of 100
rows nested inside a vertical box is fine; a list of 100 rows where
each row is a box containing a box containing a box of cells is
slow. Keep the tree as flat as you can.
\end{warning}

\begin{recap}
The runtime tree is \code{Element = variant<BoxElement, TextElement,
ElementList>}. Style, Color, and Border are small POD structs;
attributes are a bitset; themes are 24-slot palettes. Every
algorithm in \maya{} walks the tree with \code{std::visit}.
\code{BoxBuilder} bridges the compile-time DSL to the runtime tree.
\end{recap}

\section*{Workshop}

\begin{enumerate}[leftmargin=*]
  \item Open \file{maya/include/maya/element/element.hpp} and find
    the \code{Element} alias. Note that it is exactly the variant
    described in this chapter.
  \item Write a function
    \code{int count\_text\_nodes(const Element\&)} using
    \code{std::visit}. Test it on a panel from
    Chapter~\ref{ch:dsl}.
  \item Construct an \code{Element} by hand --- not via the DSL ---
    using \code{BoxBuilder}. Render it with \code{print(...)}.
\end{enumerate}

\section{Sum types and the shape of UI data}

Before we tour every variant in \maya{}'s element tree, it is
worth pausing on \emph{why} variants are the right tool for UI
trees. The decision was not aesthetic; it was forced by what UI
trees look like.

\subsection*{Product vs sum types}

The two basic shapes of compound data are products and sums:

\begin{itemize}[leftmargin=*]
  \item A \textbf{product type} carries one of each member. A
    \code{struct Point \{ int x; int y; \};} is a product: it has
    both an x \emph{and} a y.
  \item A \textbf{sum type} carries one of several possibilities.
    A \code{std::variant<int, double, std::string>} is a sum: it
    is an int \emph{or} a double \emph{or} a string, never more
    than one at a time.
\end{itemize}

A UI tree node is fundamentally a sum: a node is \emph{either} a
box \emph{or} a piece of text \emph{or} a list of children. It is
never simultaneously a box and a text. The right C++ tool for ``X
or Y or Z'' is \code{std::variant}.

\subsection*{Why not inheritance?}

For decades the conventional answer was a class hierarchy:

\begin{lstlisting}
class Element { /* virtual ~Element() = 0; */ };
class Box  : public Element { ... };
class Text : public Element { ... };
// hand around Element* everywhere
\end{lstlisting}

This works, but it pays for flexibility you do not need:

\begin{itemize}[leftmargin=*]
  \item \textbf{Heap allocation.} Polymorphism through base
    pointers forces heap storage; an \code{Element*} cannot be a
    stack value of unknown size.
  \item \textbf{Indirection.} Every method call is a virtual
    dispatch through a vtable. Cache-cold for typical UI trees.
  \item \textbf{Open set.} Anyone can derive a new
    \code{Element}. That sounds like a feature; for a library
    that wants to optimise diff and layout, it is a curse. The
    set of node kinds must be closed for the compiler to
    exhaustively check switches.
  \item \textbf{Object lifetime.} \code{std::unique\_ptr<Element>}
    works, but every operation that wants to copy a tree has to
    deep-clone via a virtual \code{clone()}.
\end{itemize}

A closed sum type fixes all four. The set of kinds is known at
compile time; storage is inline; dispatch is a small switch the
optimiser inlines; copies are bytewise.

\subsection*{std::variant: how it actually works}

A \code{std::variant<A, B, C>} is, internally, a tagged union: a
small integer (the \emph{index}) plus enough storage for the
largest alternative, aligned for any of them.

\begin{lstlisting}
// Simplified internal layout.
struct VariantABC {
    std::size_t index;        // 0, 1, or 2
    union {
        A a;
        B b;
        C c;
    };
};
\end{lstlisting}

When you write \code{std::get<A>(v)}, it checks \code{index == 0}
and returns a reference to the active member. When you write
\code{std::visit(f, v)}, the implementation generates a switch
on \code{index} that calls \code{f} with the active member.

In release builds the optimiser usually inlines the switch into
the caller. The result is dispatch as fast as a virtual call
\emph{and} no indirection through the heap.

\begin{funfact}
\code{std::variant} was added in C++17. Before that, the
C++ ecosystem had Boost.Variant (since 2002) and a handful of
hand-rolled tagged unions in every large codebase. The standard
finally caught up to ML, which had had tagged unions since 1973.
\end{funfact}

\subsection*{std::visit and the overload trick}

\code{std::visit(f, v)} requires \code{f} to be callable on every
alternative. The cleanest way to write \code{f} is the
\emph{overload pattern}:

\begin{lstlisting}
template <class... Fs> struct overload : Fs... { using Fs::operator()...; };
template <class... Fs> overload(Fs...) -> overload<Fs...>;

// Usage:
auto result = std::visit(overload{
    [](Box b)  { return "box"; },
    [](Text t) { return "text"; },
    [](auto)   { return "?"; }
}, element);
\end{lstlisting}

The \code{overload} struct inherits from a pack of lambdas and
uses \code{using Fs::operator()...} to bring them all into scope
as overloads of a single \code{operator()}. The compiler picks
the one whose parameter matches the active alternative.

This pattern is so common that C++26 is considering builtin
support. For now, every \maya{} translation unit has one
\code{overload} struct.

\begin{recap}
A tagged union is a value-typed sum: index + storage. Visiting it
is a switch the compiler optimises into a jump table. Inheritance
is the wrong tool for UI nodes because the set of kinds is
closed; sums make the closure explicit and let the compiler
check exhaustiveness.
\end{recap}

\section{Every Element variant, by kind}

With the theory settled, here is the actual variant in
\file{maya/include/maya/element/element.hpp}:

\begin{lstlisting}
using Element = std::variant<BoxElement, TextElement, ElementList>;
\end{lstlisting}

Three alternatives. Two represent a node, one represents a
container (a flat list of nodes, used where layout does not
benefit from a box). Each is a struct with public fields.

\subsection*{TextElement}

A single run of text with one style:

\begin{lstlisting}
struct TextElement {
    std::variant<std::string_view, std::string> text;
    Style style;
    std::optional<std::string> hyperlink;
};
\end{lstlisting}

The \code{text} field is itself a small variant: a non-owning view
(for compile-time literals) or an owning string (for runtime
content). This is the runtime echo of the DSL's TextNode /
RuntimeTextNode split.

\code{style} carries colours, attributes, and the foreground/
background tags. \code{hyperlink} is the OSC 8 target URL, if any.

\subsection*{BoxElement}

A container with optional decoration:

\begin{lstlisting}
struct BoxElement {
    std::vector<Element> children;
    FlexStyle layout;
    std::optional<Border> border;
    Padding padding;
    Style style;
    Visibility visibility;
};
\end{lstlisting}

This is the workhorse. \code{children} is a flat list of any
element kind. \code{layout} contains the flexbox properties
(direction, justify, align, grow, shrink) covered in
Chapter~\ref{ch:layout}. \code{border} is an optional decoration
struct; \code{padding} is per-side integer cells; \code{style}
applies to the entire box content.

\code{visibility} is a small enum: \code{Shown}, \code{Hidden}
(layout reserves space, nothing drawn), \code{Collapsed} (no
space, no drawing). The distinction matters when you want a
toggle that does not jiggle the layout.

\subsection*{ElementList}

A list of elements without a box around them:

\begin{lstlisting}
struct ElementList {
    std::vector<Element> items;
};
\end{lstlisting}

Useful when you want to return ``several elements'' from a helper
function without forcing them into a box. The layout pass treats
an \code{ElementList} as if its items had been spliced in place.

\section{Forward declarations: the chicken-and-egg dance}

A \code{BoxElement} contains \code{std::vector<Element>}. An
\code{Element} is a \code{std::variant<BoxElement, \ldots>}.
This is mutually recursive: \code{Element} depends on
\code{BoxElement}, and \code{BoxElement} depends on
\code{Element}.

C++ does not allow circular type definitions. You cannot write:

\begin{lstlisting}
struct BoxElement { std::vector<Element> children; };  // error: Element undefined
using Element = std::variant<BoxElement, TextElement>;
\end{lstlisting}

The standard library solves it in two steps:

\begin{enumerate}[leftmargin=*]
  \item \code{std::vector<T>} requires only that \code{T} be a
    forward-declared type, not a complete one. (This is unusual
    for the standard library --- most containers require complete
    types --- but vector got the relaxed rule in C++17.)
  \item Therefore you can declare \code{BoxElement} containing
    \code{std::vector<Element>} as long as \code{Element} is at
    least forward-declared at that point.
\end{enumerate}

The pattern in \maya{}:

\begin{lstlisting}
// 1. Forward-declare Element.
struct Element;

// 2. Define the alternatives. They may contain vectors of Element.
struct TextElement { /* ... */ };
struct BoxElement  { std::vector<Element> children; /* ... */ };
struct ElementList { std::vector<Element> items; };

// 3. Now define Element itself.
struct Element {
    std::variant<BoxElement, TextElement, ElementList> v;
};
\end{lstlisting}

In practice \maya{} wraps the variant in a struct (rather than
using a \code{using Element = std::variant<\ldots>;}) so it can
add convenience methods like \code{is\_text()} and
\code{as\_box()}.

\begin{aside}
This dance generalises. Any recursive data structure in C++ uses
the same trick: forward-declare the outer type, define the inner
types that contain a vector or pointer to the outer type, then
define the outer type. The classic example is a JSON value, where
\code{Object} maps strings to values and \code{Array} is a vector
of values.
\end{aside}

\section{Constructing an Element by hand}

The DSL builds elements for you, but sometimes you need to
assemble one without the compile-time machinery (loaded layouts,
user-themable trees, debugging). C++26's designated initialisers
make this readable:

\begin{lstlisting}
Element greeting{
    .v = TextElement{
        .text  = std::string_view{"Hello, maya"},
        .style = Style{
            .fg = Color::rgb(220, 220, 240),
            .attributes = Attribute::Bold,
        },
    },
};
\end{lstlisting}

No allocation: \code{string\_view} into the literal, small POD
style. The whole construction is on the stack.

For a box:

\begin{lstlisting}
Element panel{
    .v = BoxElement{
        .children = {
            Element{.v = TextElement{.text = std::string_view{"Top"}}},
            Element{.v = TextElement{.text = std::string_view{"Mid"}}},
            Element{.v = TextElement{.text = std::string_view{"Bot"}}},
        },
        .layout = FlexStyle{.direction = Direction::Column},
        .border = Border{.style = BorderStyle::Rounded},
        .padding = Padding::all(1),
    },
};
\end{lstlisting}

The \code{children} list contains nested \code{Element} values.
One allocation for the vector buffer; the children themselves are
inline within their slots.

\subsection*{The box() builder helper}

For tests and one-off scripts, \code{box(...)} bundles common
defaults:

\begin{lstlisting}
auto e = box(text("hi"))
             .border(BorderStyle::Single)
             .padding(1)
             .build();
\end{lstlisting}

The builder is a small fluent API that mutates an internal
\code{BoxElement} and returns the wrapped \code{Element} on
\code{build()}. Use it when designated initialisers feel verbose.

\section{The colour model, up close}

\code{Color} is a sum of three representations:

\begin{lstlisting}
struct Color {
    enum class Kind { Default, Ansi16, Ansi256, Rgb };
    Kind kind = Kind::Default;
    uint8_t r = 0, g = 0, b = 0;  // payload for Rgb
    uint8_t index = 0;            // payload for Ansi16/Ansi256
    bool is_default() const { return kind == Kind::Default; }
};
\end{lstlisting}

The four kinds:

\begin{itemize}[leftmargin=*]
  \item \textbf{Default.} ``No colour set; inherit from terminal.''
    The ANSI emitter writes nothing for this case.
  \item \textbf{Ansi16.} Indexes 0--15 from the classic palette
    (black, red, green, \ldots, plus bright variants). Maps to
    SGR codes 30--37 / 40--47 and 90--97 / 100--107.
  \item \textbf{Ansi256.} Indexes 0--255 from the 256-colour
    palette (16 ANSI + 6$\times$6$\times$6 colour cube + 24-step
    grey ramp). Maps to SGR 38;5;$n$ / 48;5;$n$.
  \item \textbf{Rgb.} 24-bit truecolour. Maps to
    SGR 38;2;$r$;$g$;$b$ / 48;2;$r$;$g$;$b$.
\end{itemize}

\subsection*{Why Default exists}

``No colour'' is not the same as ``black''. If the user has a
dark-on-light terminal theme, drawing ``black'' over their
background produces invisible text or a dark blob. Drawing ``no
colour'' lets the terminal supply the foreground, which adapts
to the user's theme.

The DSL exposes \code{Color::default\_()} (or, in the runtime
builders, leaves the \code{Color} field default-constructed) to
opt out of overriding the user's choice.

\subsection*{8/256/truecolour: a small history}

\begin{itemize}[leftmargin=*]
  \item The original ANSI standard (X3.64, 1979) defined eight
    colours: SGR 30--37.
  \item The bright variants (90--97) were added by \emph{aixterm}
    around 1995 and adopted widely.
  \item The 256-colour palette was popularised by xterm in 1999
    via the SGR 38;5;$n$ syntax; the cube and grey ramp were
    xterm's invention, not a standard.
  \item Truecolour (24-bit RGB) was first implemented by
    Konsole and iTerm2 in the early 2010s; ITU-T T.416 had
    described the syntax \code{38:2::r:g:b} since 1993, but
    everyone uses \code{38;2;r;g;b} now.
\end{itemize}

\begin{funfact}
If you read T.416 carefully, the separators between the RGB
values should be colons (\code{:}), not semicolons (\code{;}).
Because every implementation used semicolons, the de-facto
standard split from the de-jure one. \maya{} emits semicolons,
like everyone else.
\end{funfact}

\subsection*{Choosing a colour kind}

When you write \code{Fg<R, G, B>} in the DSL, you get
\code{Color::Kind::Rgb}. When you write \code{Fg<"red">} you get
\code{Color::Kind::Ansi16}. The runtime emitter degrades
automatically: if the user's terminal does not advertise
truecolour (\code{TERM} or \code{COLORTERM} environment), the
emitter quantises RGB to the nearest 256-colour index.

The quantisation uses a simple Euclidean-distance match against
the 256-colour palette. The result is usually visually close;
for brand-perfect colour, set \code{COLORTERM=truecolor} in your
shell config.

\section{The border gallery}

\code{BorderStyle} is an enum. Each value picks a set of six
byte sequences (corners, horizontal, vertical) drawn around the
box. The Unicode block-drawing range
U+2500--U+257F supplies all the glyphs.

\subsection*{Single (U+250C--U+2518, U+2500, U+2502)}

The canonical thin-line border:

\begin{terminalbox}
+--Border-------+\\
|\, body line  \, |\\
+--------------+
\end{terminalbox}

Corners use U+250C (top-left), U+2510 (top-right), U+2514 (bottom-
left), U+2518 (bottom-right). Horizontals use U+2500; verticals
use U+2502.

\subsection*{Rounded (U+256D--U+2570)}

The friendly variant with curved corners:

\begin{terminalbox}
.--Rounded----.\\
|\, body line\, |\\
`-------------'
\end{terminalbox}

Corners are U+256D / U+256E / U+2570 / U+256F. Lines reuse
U+2500 and U+2502.

\subsection*{Double (U+2554, U+2557, U+255A, U+255D, U+2550, U+2551)}

Line-drawing for headers and important regions:

\begin{terminalbox}
=====================\\
||\, body line\, ||\\
=====================
\end{terminalbox}

(Real characters: U+2554 / U+2557 / U+255A / U+255D corners,
U+2550 horizontal, U+2551 vertical.)

\subsection*{Heavy (U+250F, U+2513, U+2517, U+251B, U+2501, U+2503)}

A bold thick-line border for emphasis.

\subsection*{ASCII}

A portability fallback using only ASCII: hyphens, pipes, and
plus signs. Use this on terminals that cannot render Unicode
block-drawing characters (rare, but they exist).

\subsection*{Custom}

A \code{Border} struct can also carry six \code{char32\_t}
values directly, bypassing the enum. Useful for half-block fills
and custom themes.

\subsection*{Partial-side borders}

\code{Border} carries a \code{sides} bitmask (top, right, bottom,
left). Set only some bits to get a border on a subset of sides:

\begin{lstlisting}
Border b{
    .style = BorderStyle::Single,
    .sides = BorderSide::Bottom | BorderSide::Top,
};
\end{lstlisting}

A box with only top and bottom borders looks like a horizontal
rule sandwich --- useful for table rows and section dividers.

\section{Style: merging rules}

A \code{Style} struct carries:

\begin{itemize}[leftmargin=*]
  \item \code{fg}, \code{bg}: \code{Color} values (or default).
  \item \code{attributes}: a bitset of
    Bold/Dim/Italic/Underline/Blink/Reverse/Strike.
  \item \code{hyperlink}: optional URL for OSC 8.
\end{itemize}

When styles compose (a child inherits from a parent, or two
pipes attach styles), \maya{} applies a merge rule:

\begin{enumerate}[leftmargin=*]
  \item Attributes union: if either has Bold, the result has
    Bold.
  \item Colours override: a non-default colour replaces a
    default. Two non-default colours: child wins.
  \item Hyperlink: child wins if set.
\end{enumerate}

\begin{lstlisting}
Style merge(Style outer, Style inner) {
    Style r = outer;
    r.attributes |= inner.attributes;
    if (!inner.fg.is_default()) r.fg = inner.fg;
    if (!inner.bg.is_default()) r.bg = inner.bg;
    if (inner.hyperlink)        r.hyperlink = inner.hyperlink;
    return r;
}
\end{lstlisting}

The merge happens during the layout pass, when the renderer walks
the tree and assembles per-cell style. You do not call it
yourself; \maya{} does.

\begin{idea}
The merge rule is what makes
\code{box(text("x")|Italic)|Bold} produce text that is both bold
and italic, without you writing \code{Bold | Italic} explicitly.
Styles inherit through the tree the way CSS does, just with a
cleaner rule set.
\end{idea}

\section{Tree-walking patterns}

Once you have an \code{Element}, every interesting operation is a
recursive visit. Five patterns appear over and over.

\subsection*{Counting nodes of a kind}

\begin{lstlisting}
int count_text(const Element& e) {
    return std::visit(overload{
        [](const TextElement&)      { return 1; },
        [](const BoxElement& b) {
            int n = 0;
            for (auto& c : b.children) n += count_text(c);
            return n;
        },
        [](const ElementList& l) {
            int n = 0;
            for (auto& c : l.items) n += count_text(c);
            return n;
        }
    }, e.v);
}
\end{lstlisting}

Visit at the root, recurse into containers. The visit dispatches
on kind; the recursion threads through the children.

\subsection*{Finding the first text matching a predicate}

\begin{lstlisting}
std::optional<std::string_view>
find_text(const Element& e,
          std::function<bool(std::string_view)> pred) {
    return std::visit(overload{
        [&](const TextElement& t) -> std::optional<std::string_view> {
            auto sv = std::visit(
                [](auto& x) { return std::string_view{x}; }, t.text);
            return pred(sv) ? std::optional{sv} : std::nullopt;
        },
        [&](const BoxElement& b) -> std::optional<std::string_view> {
            for (auto& c : b.children)
                if (auto r = find_text(c, pred)) return r;
            return std::nullopt;
        },
        [&](const ElementList& l) -> std::optional<std::string_view> {
            for (auto& c : l.items)
                if (auto r = find_text(c, pred)) return r;
            return std::nullopt;
        }
    }, e.v);
}
\end{lstlisting}

Classic depth-first search. The optional return is the standard
way to express ``found / not found'' without exceptions.

\subsection*{Transforming text in place}

\begin{lstlisting}
void uppercase_text(Element& e) {
    std::visit(overload{
        [](TextElement& t) {
            std::string s;
            std::visit([&](auto& x) { s = std::string{x}; }, t.text);
            for (auto& c : s) c = std::toupper(c);
            t.text = std::move(s);
        },
        [](BoxElement& b) {
            for (auto& c : b.children) uppercase_text(c);
        },
        [](ElementList& l) {
            for (auto& c : l.items) uppercase_text(c);
        }
    }, e.v);
}
\end{lstlisting}

Mutating visitor: the lambdas take non-const references and
modify in place. The variant remains the same kind; only its
payload changes.

\subsection*{Hit-testing (mouse coordinate to element)}

Given a screen coordinate \code{(x, y)}, find the deepest
element under that point. Layout has annotated each box with its
bounding rectangle:

\begin{lstlisting}
const Element* hit_test(const Element& e, int x, int y) {
    return std::visit(overload{
        [&](const TextElement& t) -> const Element* {
            // text knows its rect; check membership
            return t.rect.contains(x, y) ? &e : nullptr;
        },
        [&](const BoxElement& b) -> const Element* {
            if (!b.rect.contains(x, y)) return nullptr;
            // Recurse to find a deeper hit, else return self.
            for (auto& c : b.children)
                if (auto h = hit_test(c, x, y)) return h;
            return &e;
        },
        [&](const ElementList& l) -> const Element* {
            for (auto& c : l.items)
                if (auto h = hit_test(c, x, y)) return h;
            return nullptr;
        }
    }, e.v);
}
\end{lstlisting}

Returns the deepest hit, or \code{nullptr} if the coordinate is
outside the entire tree.

\subsection*{Pretty-printing for debugging}

\begin{lstlisting}
void dump(const Element& e, int depth = 0) {
    auto pad = std::string(depth * 2, ' ');
    std::visit(overload{
        [&](const TextElement& t) {
            std::visit([&](auto& s) {
                std::cout << pad << "text: " << s << "\n";
            }, t.text);
        },
        [&](const BoxElement& b) {
            std::cout << pad << "box (" << b.children.size() << " children)\n";
            for (auto& c : b.children) dump(c, depth + 1);
        },
        [&](const ElementList& l) {
            std::cout << pad << "list (" << l.items.size() << ")\n";
            for (auto& c : l.items) dump(c, depth + 1);
        }
    }, e.v);
}
\end{lstlisting}

Useful while you are learning the tree shape. Drop a
\code{dump(my\_element)} in a test and you can see what the DSL
produced.

\begin{recap}
Every operation on an \code{Element} is a recursive
\code{std::visit}. The pattern is uniform: handle each
alternative, recurse into containers. There is no virtual
dispatch; the compiler turns the visit into a switch.
\end{recap}

\section{Hyperlinks: OSC 8}

Xterm and most modern terminals support clickable hyperlinks via
the OSC 8 escape sequence:

\begin{lstlisting}
\e]8;;https://example.com\e\\Visit\e]8;;\e\\
\end{lstlisting}

The text between the start (\code{ESC ]8;;URL ESC \textbackslash})
and end (\code{ESC ]8;; ESC \textbackslash}) sequences is
rendered as a clickable region.

\maya{}'s \code{TextElement::hyperlink} field, if set, causes the
emitter to wrap the cell in OSC 8 brackets. The text itself is
unchanged; only the metadata differs.

\begin{lstlisting}
Element link{
    .v = TextElement{
        .text = std::string_view{"the docs"},
        .style = Style{.attributes = Attribute::Underline},
        .hyperlink = std::string{"https://example.com"},
    },
};
\end{lstlisting}

Rendering this in a terminal that supports OSC 8 produces a
clickable ``the docs''. On a terminal that does not, the
emitter falls back to plain text (the hyperlink sequence is
ignored).

\begin{aside}
OSC 8 was first implemented by VTE (the terminal widget in
GNOME Terminal) in 2017 and is now supported by iTerm2, Kitty,
WezTerm, Alacritty, and modern xterm. Tmux passes it through
as of version 3.3. If your terminal is older than
2019-ish, expect plain text.
\end{aside}

\section{Pitfalls}

A few traps to watch for when working with elements directly.

\begin{warning}
\textbf{Mixing string\_view and string carelessly.} If a
\code{TextElement::text} holds a \code{string\_view} that points
into a temporary, the view will dangle. Use \code{string\_view}
only for literals and other long-lived storage; use
\code{std::string} for anything computed.
\end{warning}

\begin{warning}
\textbf{Copying large element trees.} A deep tree's copy is
proportional to its depth and width. If you find yourself
cloning elements every frame, you are paying for moves that
should be moves and copies that should be views. Profile
before you optimise --- the cost is rarely the bottleneck ---
but watch the allocator.
\end{warning}

\begin{warning}
\textbf{Setting Visibility::Collapsed and expecting layout to
stay stable.} Collapsed removes the element from layout
entirely; its siblings will reflow. If you want a stable layout,
use \code{Hidden} instead, which reserves space.
\end{warning}

\begin{warning}
\textbf{Custom Border characters with combining marks.} A single
Unicode codepoint can take zero, one, or two cells. \maya{}'s
border drawing assumes one cell per character. If you supply a
combining mark (e.g. U+0301), the border will misalign. Stick to
the block-drawing range.
\end{warning}

\begin{warning}
\textbf{Hyperlinks in text that wraps.} If a hyperlink text
wraps across two visual lines, some terminals split the link
into two regions; some keep them unified. The behaviour is
terminal-dependent; do not rely on either.
\end{warning}

\section{The bigger picture}

The element tree is what flows between \code{view} and the
renderer. Every frame, \code{view} returns a fresh \code{Element};
the renderer walks it, lays it out, diffs it, and emits bytes.
The tree is the lingua franca between your code and \maya{}'s
machinery.

Understanding the tree pays off in three places:

\begin{itemize}[leftmargin=*]
  \item Debugging a misrendered frame: \code{dump} the tree and
    spot the wrong child.
  \item Writing widget libraries: a widget is a function from
    state to \code{Element}, so you must build trees by hand.
  \item Reading \maya{}'s source: the layout, the diff, the
    ANSI emitter --- all walk the same variant.
\end{itemize}

In the next chapter we follow the tree into the layout engine
and see how the abstract description becomes a grid of cells.

\begin{recap}
An \code{Element} is a closed sum: \code{BoxElement},
\code{TextElement}, \code{ElementList}. Three small structs, one
variant. Every \maya{} algorithm visits it. There is no
inheritance; the set of node kinds is fixed and exhaustively
checkable. Constructing trees by hand is one designated
initialiser at a time.
\end{recap}

\chapter{Flexbox in integers}
\label{ch:layout}

Given an element tree, \maya{} must decide where every text run, every
box, every border goes --- in cells, not pixels. The algorithm it uses
is a tailored re-implementation of CSS Flexbox. This chapter explains
flexbox from scratch (no web background assumed) and then walks
\maya{}'s implementation.

\begin{aside}
Layout is the unglamorous middle of every UI library. Nobody talks
about it; the algorithm runs hundreds of times a second, makes
thousands of small decisions, and disappears the moment it gets the
rectangle right. Bad layout is instantly noticeable; good layout is
invisible. This chapter is a guided tour of \maya{}'s ten kilobytes
of layout code, because once you understand it you can stop guessing
at why a panel is one cell off and start reading the cause out of
the configuration.
\end{aside}

\section{The problem layout solves}

You have a rectangular area --- the terminal --- and a tree of
elements each with constraints (\emph{this one wants to be 20 cells
wide; this row should pack its children with a 1-cell gap; this panel
must grow to fill leftover space}). Layout is the function that
assigns concrete \code{(x, y, w, h)} rectangles to each node such that
all the constraints are satisfied.

There are many possible algorithms (table layout, constraint solving,
grid layout). Flexbox is the most appropriate for terminal UIs:
predictable, single-pass, no surprises, no solver.

\subsection{Why not the obvious algorithms?}

A brief tour of the alternatives \maya{} did not choose:

\paragraph{Table layout (HTML \code{<table>}).} A rectangular grid
of rows and columns. Each cell's size is determined by the largest
content in its row and column. Beautiful for tabular data; awkward
for anything else --- two unrelated cells affecting each other's
size is a constant surprise. Used in early web design; abandoned
as a general-purpose layout in the mid-2000s.

\paragraph{Constraint solving (Cassowary, Auto Layout).} Each
position is a variable; each layout rule is a linear constraint.
A simplex-like solver finds an assignment. Powerful (it can express
``these two are always equal'', ``this is between 10 and 30
cells''), but expensive and prone to producing surprising answers
when constraints conflict. Apple's Auto Layout works this way.

\paragraph{Grid (CSS Grid).} A two-dimensional placement model:
you declare a grid of named rows and columns, then place children
in specific cells. Excellent for page-level layout; overkill for
the nested-row-and-column shape most TUIs take.

\paragraph{Flow / line-box (HTML inline text).} Children are
packed left-to-right, top-to-bottom, like words in a paragraph.
Good for prose; bad for everything else.

Flexbox is the simplest model that handles the common cases ---
rows of widgets, columns of panels, a sidebar plus content, a
status bar pinned to the bottom --- without ceremony. It is also
linear-time and predictable, both crucial for a sixty-frames-per-second
renderer.

\begin{funfact}
Flexbox went through several names and standards before it stuck.
The original CSS3 ``flexible box'' module was proposed in 2009;
the syntax changed twice (the \code{display: box} from 2009, then
\code{display: flexbox} in 2011, then \code{display: flex} in 2012)
before browsers settled. Mosaic's table-based layouts and
Netscape's \code{<center>} tag ran the web for a decade because
nothing nicer was standardised. The model \maya{} uses --- main
axis, grow, shrink, basis, gap --- is the 2012 design, late enough
that the rough edges had been filed off, simple enough to fit in
800 lines of \cpp{}.
\end{funfact}

\section{Flexbox in one paragraph}

Flexbox arranges children along a \textbf{main axis} (row or column).
Children may \emph{grow} to fill leftover space and \emph{shrink} when
oversized; they may \emph{wrap} onto new lines when they don't fit;
they may be \emph{aligned} along the cross axis and \emph{justified}
along the main axis. Each container picks an axis (row or column);
each child carries a few numbers (basis, grow, shrink, gap, alignment
override).

\section{Vocabulary}

A worked example will be clearest. Take a horizontal row of three
children inside a 30-cell-wide box:

\begin{lstlisting}
h(
    panel_a,   // wants 10 cells
    panel_b,   // grow=1
    panel_c    // wants 8 cells
) | gap_<1>
\end{lstlisting}

\maya{} reasons about it like this:

\begin{enumerate}[leftmargin=*]
  \item \textbf{Sum the basis sizes.} A=10, B=0 (no fixed size,
    grows), C=8. Plus 2 gap cells. Total committed: 20.
  \item \textbf{Leftover space.} 30 - 20 = 10 cells.
  \item \textbf{Distribute by grow.} Only B has \code{grow > 0}, so
    B receives all 10 cells. New size: B = 10.
  \item \textbf{Place.} A at column 0, gap, B at column 11, gap, C
    at column 22.
\end{enumerate}

The cross axis (vertical, in this example) is handled by
\code{Align}: children stretch to the parent height by default, but
they can be \code{Start}, \code{Center}, \code{End}, or
\code{Stretch}.

\section{The flex properties}

A node in \maya{}'s layout engine carries the following style:

\begin{lstlisting}
struct FlexStyle {
    FlexDirection direction = FlexDirection::Column;
    FlexWrap      wrap      = FlexWrap::NoWrap;

    Align   align_items   = Align::Auto;
    Align   align_self    = Align::Auto;
    Justify justify       = Justify::Auto;

    int     gap           = 0;

    int     pad_top    = 0, pad_right = 0,
            pad_bottom = 0, pad_left  = 0;
    int     mar_top    = 0, mar_right = 0,
            mar_bottom = 0, mar_left  = 0;

    int     width       = -1;   // -1 = auto
    int     height      = -1;
    int     min_width   = 0, max_width  = INT_MAX;
    int     min_height  = 0, max_height = INT_MAX;

    int     basis       = -1;   // preferred main-axis size, -1 = auto
    float   grow        = 0.0f;
    float   shrink      = 1.0f;

    Display display     = Display::Flex;
    Overflow overflow   = Overflow::Visible;
    bool    has_border  = false;
};
\end{lstlisting}

These mirror CSS flexbox almost exactly. The differences:

\begin{itemize}[leftmargin=*]
  \item Every quantity is an \emph{integer} number of cells. No
    \code{em}, no \code{px}, no fractional pixels. Sub-cell positions
    don't exist in a terminal.
  \item \code{display: Flex} and \code{display: None} are the only
    options. There is no block, no inline.
  \item Percentages are limited; most sizes are explicit integers.
\end{itemize}

\section{The data structure}

A layout tree is a flat \code{std::vector<LayoutNode>} indexed by
\code{int}. Each node has a list of child indices.

\begin{lstlisting}
struct LayoutNode {
    FlexStyle           style;
    std::vector<int>    children;     // indices into the same vector
    layout::Rect        computed;     // output: x, y, w, h
};
\end{lstlisting}

Why a flat vector and not a pointer tree? Three reasons:

\begin{itemize}[leftmargin=*]
  \item \textbf{Cache locality.} Walking siblings is a stride of
    \code{sizeof(LayoutNode)}; the prefetcher loves it.
  \item \textbf{No allocation in the hot path.} A
    \code{vector<int>} for children is cheaper than a \code{vector
    <unique\_ptr<LayoutNode>>}.
  \item \textbf{Easier reasoning.} Indices cannot dangle; an index
    is either in range or out of range. Pointer trees create
    aliasing puzzles.
\end{itemize}

\begin{idea}
``Make it an index, not a pointer'' is a recurring optimisation in
\maya{}. The same idea drives the style pool (a 16-bit ID into a
table) and the hyperlink interning (a 32-bit ID).
\end{idea}

\section{The algorithm, step by step}

\code{layout::compute(nodes, root, available\_width)} runs in three
phases per box:

\begin{enumerate}[leftmargin=*]
  \item \textbf{Measure.} Walk children recursively and compute their
    intrinsic main-axis size (the basis if specified, else the
    content's measured size).
  \item \textbf{Distribute.} Add up basis sizes plus gaps. Compare
    to the parent's main-axis size. If there is leftover space and
    any child has \code{grow > 0}, distribute the leftover
    proportionally; if there is overflow and any child has
    \code{shrink > 0}, shrink proportionally.
  \item \textbf{Place.} Compute final \code{(x, y, w, h)} for each
    child given the resolved main-axis sizes and the cross-axis
    alignment.
\end{enumerate}

Wrapping (when \code{wrap = Wrap}) re-enters the algorithm per line.
Justify and align use simple integer division: \code{SpaceBetween}
puts the leftover between children, \code{SpaceAround} puts it on
both sides, \code{Center} halves it.

\subsection{Padding, border, margin}

Three rectangles concentric:

\begin{center}
\small
\begin{tabular}{ll}
\textbf{margin} & cells \emph{outside} the box (sibling spacing) \\
\textbf{border} & one cell thick if present (drawn glyphs) \\
\textbf{padding} & cells \emph{inside} the border (content inset) \\
\textbf{content} & where children are laid out \\
\end{tabular}
\end{center}

The available width for children is \code{w - left\_pad - right\_pad -
border\_l - border\_r}. The available height is the same on the
vertical. \maya{}'s renderer paints the border in the one-cell ring
around the padded content.

\section{Worked example: a status panel}

\begin{lstlisting}
auto panel = v(
    t<"Status"> | Bold,
    sep,
    h(t<"CPU:"> | Dim, space, t<"42%"> | Bold) | gap_<1>,
    h(t<"Mem:"> | Dim, space, t<"8.2G"> | Bold) | gap_<1>
) | border_<Round> | pad<1, 2> | gap_<1>;
\end{lstlisting}

Suppose the parent box gives us 24 cells wide. The layout proceeds:

\begin{enumerate}[leftmargin=*]
  \item Outer box: border (1 cell each side) + padding (2 on each
    side) leaves 24 - 2 - 4 = \textbf{18 cells} of inner width.
  \item Inner column has four children: title (1 row), separator (1
    row), CPU line (1 row), Mem line (1 row), with gap 1 between
    them. Total height 4 + 3 = 7 rows of content. Plus the border
    pads (1 + 1) and padding (1 + 1) makes the outer box 11 rows.
  \item Inside each \code{h(...)} line, the children are: a static
    label, a \code{space} (grow=1), and a value text. Total fixed
    width: maybe 4 + 3 + gap 1 = 8 cells. Leftover 10 goes entirely
    to the spacer. The value text is pushed to column 8 + 10 = 18,
    i.e. the right edge.
\end{enumerate}

The renderer then walks the resulting rectangles and writes glyphs.

\section{When layout disagrees with you}

Two situations confuse newcomers:

\begin{itemize}[leftmargin=*]
  \item \textbf{Nothing happens when I set \code{grow}.} The parent
    must have a fixed (or grow-controlled) size for the leftover
    space to be meaningful. A \code{grow=1} child of a parent that
    sizes itself to its content has nothing to grow into.
  \item \textbf{Text wraps inside a column unexpectedly.} A
    \code{TextElement}'s width is its content. If you put a long
    string inside a column, it tries to take its natural width; the
    parent box constrains it to its own width, which causes wrap.
    Set \code{TextWrap::Truncate} or \code{Ellipsize} if that is
    wrong.
\end{itemize}

\section{Where the code lives}

\file{maya/include/maya/layout/yoga.hpp} and
\file{maya/src/layout/yoga.cpp}. Despite the name, this is not a
binding to Facebook's Yoga library --- it is an independent
implementation that follows the same flexbox spec. ``Yoga'' here is a
nickname for the algorithm. The implementation is around 800 lines
and has no external dependencies.

\section{Justify and Align, every value}
\label{sec:layout:justify-align}

The two enums \code{Justify} and \code{Align} control how leftover
space is distributed and how children line up across the cross axis.
Both share a value list: \code{Start}, \code{End}, \code{Center},
\code{Stretch}, \code{SpaceBetween}, \code{SpaceAround},
\code{SpaceEvenly}, \code{Auto}. Worth walking each.

Imagine a row 30 cells wide, holding three children each 4 cells
wide. Leftover space: $30 - 12 = 18$ cells.

\paragraph{\code{Start} (default).} Pack children to the start.
No space inserted between them; all 18 leftover cells go to the
end.

\begin{terminalbox}
AAAA BBBB CCCC..................
\end{terminalbox}

\paragraph{\code{End}.} Pack children to the end. All leftover
space at the start.

\begin{terminalbox}
..................AAAA BBBB CCCC
\end{terminalbox}

\paragraph{\code{Center}.} Half the leftover before, half after.
With odd leftover (rare in cells but possible), the extra cell
goes to the start.

\begin{terminalbox}
.........AAAA BBBB CCCC.........
\end{terminalbox}

\paragraph{\code{SpaceBetween}.} Leftover divided into
$n - 1$ equal gaps between children. First child at start, last at
end. With 18 cells divided into 2 gaps: 9 each.

\begin{terminalbox}
AAAA.........BBBB.........CCCC
\end{terminalbox}

\paragraph{\code{SpaceAround}.} Leftover divided into $n$ equal
segments, half a segment on each end, full segments between.
With 18 cells over 3 children: 6 per segment, 3 on each end.

\begin{terminalbox}
...AAAA......BBBB......CCCC...
\end{terminalbox}

\paragraph{\code{SpaceEvenly}.} Leftover divided into $n + 1$
equal gaps. Same gap on each end as between children. With 18
cells over 4 gaps: 4 or 5 each (integer division leaves a remainder).

\begin{terminalbox}
....AAAA....BBBB....CCCC...
\end{terminalbox}

\paragraph{\code{Stretch}.} Cross-axis only. Children are sized to
the parent's cross-axis size. The default for cross-axis alignment
on boxes (so children fill the height of a row by default).

\paragraph{\code{Auto}.} The default. For \code{align\_items}
resolves to \code{Stretch}; for \code{align\_self} means ``inherit
from parent's \code{align\_items}''.

\subsection{Integer division and the remainder}

A wrinkle of cell-based layout: leftover space rarely divides
evenly. With \code{SpaceBetween} of three children in 30 cells
(leftover 18, two gaps), each gap is 9. With four children in 30
cells (leftover 14, three gaps), each gap is 4 and the integer
division leaves 2 cells unaccounted for.

\maya{}'s strategy: the first $k$ gaps get one extra cell, where
$k$ is the remainder. Children later in the list see a slightly
tighter gap. The alternative would be to round each gap and
accept that the total drifts by a cell or two, but accumulating
rounding errors quickly produces visible misalignment.

\begin{idea}
In integer cell-based layout, leftover space rarely divides
evenly. The convention is to give the remainder to the early gaps,
so the off-by-one shows up at one edge rather than accumulating.
This is a recurring theme in pixel-free graphics; the same
problem appears in audio sample rate conversion and frame-rate
remapping.
\end{idea}

\section{Grow and shrink, mechanically}
\label{sec:layout:grow-shrink}

The \code{grow} and \code{shrink} floats are the most subtle part
of flexbox. Worth walking through their full mechanics, with
worked examples.

\subsection{The setup}

Given children with basis sizes $b_1, b_2, \dots, b_n$ and an
available main-axis size $W$:

\begin{enumerate}[leftmargin=*]
  \item Sum the basis sizes and gaps: $B = \sum b_i + \text{gaps}$.
  \item If $B = W$, you are done. Place children at their basis
    sizes.
  \item If $B < W$: there is leftover space, and \textbf{grow}
    distributes it. If no child has \code{grow > 0}, leftover stays
    leftover (handled by \code{justify}).
  \item If $B > W$: children overflow, and \textbf{shrink}
    distributes the negative. If no child has \code{shrink > 0},
    the overflow stays (visible as content escaping the box, or
    clipped depending on \code{overflow}).
\end{enumerate}

\subsection{Distributing leftover by \code{grow}}

Let $G = \sum g_i$ be the total grow factor of all children. If
$G > 0$ and $L = W - B > 0$, each child $i$ with grow $g_i$ receives
an additional $\lfloor L \cdot g_i / G \rfloor$ cells.

A worked example. Three children, basis sizes 10, 5, 5 (total 20),
grow factors 1, 0, 2 (total 3), available 32 cells. Leftover
$L = 12$. Distribute:

\begin{itemize}[leftmargin=*]
  \item Child 1: $\lfloor 12 \cdot 1 / 3 \rfloor = 4$ extra cells.
    Final size 14.
  \item Child 2: grow 0, no extra. Final size 5.
  \item Child 3: $\lfloor 12 \cdot 2 / 3 \rfloor = 8$ extra cells.
    Final size 13.
\end{itemize}

Total: $14 + 5 + 13 = 32$, exactly $W$. The remainders from
division sum to zero in this case; when they don't, the first child
in the list absorbs the difference (so the total comes out exact).

\subsection{Distributing overflow by \code{shrink}}

The symmetric case. If $B > W$ and total shrink $S > 0$, each child
loses $\lfloor (B - W) \cdot s_i / S \rfloor$ cells. The defaults
are important: \code{shrink = 1.0f} means every child shrinks
equally; \code{shrink = 0.0f} means the child refuses to shrink and
overflow goes elsewhere.

A child whose basis is already at its \code{min\_width} cannot
shrink further; the algorithm marks it ineligible and redistributes
over the remaining children. This is one of two places \maya{}'s
flexbox implementation does a second pass.

\subsection{Common patterns}

The spinner of common grow/shrink configurations:

\begin{center}
\small
\begin{tabular}{ll}
\textbf{Effect} & \textbf{Setting} \\
\hline
Fixed size, never resizes        & \code{grow=0, shrink=0, basis=N} \\
Fills leftover space              & \code{grow=1, shrink=1} (defaults) \\
Fills, never overflows            & \code{grow=1, shrink=1, min\_width=0} \\
Fills proportionally (2:1 split)  & two siblings with \code{grow=2} and \code{grow=1} \\
Never shrinks below content       & \code{shrink=0} \\
Doesn't grow, will shrink         & \code{grow=0, shrink=1} \\
\end{tabular}
\end{center}

\code{space} in the DSL is a node with \code{grow=1, basis=0,
shrink=0}: it eagerly absorbs every leftover cell and refuses to be
shrunk if there isn't enough space.

\section{Why integer cells, not floats}
\label{sec:layout:integer-math}

Web flexbox uses floating-point pixels because the visual unit ---
a pixel --- can effectively be subdivided in CSS (sub-pixel
antialiasing, fractional positions, transforms). The terminal has
no such unit: a cell is a cell, you cannot draw half a character.
Layout has to terminate on whole cells.

Using \code{int} for every quantity has consequences worth
being explicit about:

\paragraph{No rounding surprises during composition.} Two layouts
computed independently and then combined have exactly the same
geometry as if computed together. With floats, rounding error
compounds: a panel that is 12.499... cells wide at one level
becomes 12 at the next, then 11.998... in the parent, and so on.
Integer math avoids the question.

\paragraph{Diff is bitwise comparable.} Frame-to-frame layout
diffs (Chapter~\ref{ch:simd}) can compare integer rectangles
with a single \code{==}; with floats the question of ``did this
position change'' becomes ``did this position change by more than
epsilon'', and the epsilon has to be chosen.

\paragraph{Constraints are exact.} ``This box is at least 10 cells
wide and at most 30'' is bit-exact; the solver never needs to
worry about a width that is mathematically 10.0001 cells.

\paragraph{The cost: integer division does not preserve totals.}
The technique from the previous section --- give remainders to the
earliest gaps --- is what we pay for using integers. It is much
less work than the floating-point sorrows we avoid.

\section{The three phases, expanded}
\label{sec:layout:phases}

The summary said ``measure, distribute, place''. The full version
for a single box:

\paragraph{Phase 1: measure.}
For each child, recursively compute its \emph{intrinsic} main-axis
size --- the size it wants to be if given no extra space.
\code{TextElement}'s intrinsic main-axis size is the width of its
content (for a row container) or the height after wrapping (for a
column). \code{BoxElement}'s intrinsic is its own measured size,
computed by recursing into its children.

The result is each child's \textbf{basis} --- either the explicit
\code{basis} field if set, or the measured intrinsic.

\paragraph{Phase 2: distribute.}
Compare $B = \sum b_i + \text{gaps}$ to $W$, the available
main-axis size. Distribute leftover or overflow per
Section~\ref{sec:layout:grow-shrink}. After this phase every
child has a final main-axis size.

\paragraph{Phase 3: place.}
Walk children, assigning each one an $(x, y)$. The cursor advances
by the child's final main-axis size plus the gap; the cross-axis
position is determined by \code{align}. If \code{justify} is
non-trivial, the gap is enlarged or padding inserted at the ends.

Each child's final \code{computed} rectangle is written into the
\code{LayoutNode} struct, and the algorithm recurses into the
child to lay out \emph{its} children, with the child's resolved
size as the new available area.

\subsection{Wrapping}

When \code{wrap = Wrap} and the children don't fit on one line,
phase 2 partitions them into lines greedily: pack until the next
child would overflow, start a new line. Each line is then
independently laid out and placed at increasing cross-axis offsets.

For a wrapping column (rare), the lines are columns and the cross
axis is horizontal. Symmetric.

\section{Debugging layout: the four-step ritual}
\label{sec:layout:debugging}

When something is one cell off, two cells over, or in the wrong
place entirely, the cause is almost always one of four things.
Worth checking each in order before opening a debugger.

\paragraph{1. Check the available size.} The terminal's columns
and rows are passed in at the top of the layout. Print them. A
resize between frames will sometimes lag a frame, and you might be
laying out into the wrong rectangle.

\paragraph{2. Check padding and border on every ancestor.} A box
with padding 1 and a border eats 4 cells of width before its
children see anything. Inside a $24$-column terminal, after a
bordered-padded panel and another bordered-padded inner box,
children get 24 - 4 - 4 = 16 cells. The defect is rarely the layout;
it is usually an extra padding nobody remembered.

\paragraph{3. Print the layout tree.} A small helper:

\begin{lstlisting}
void dump(const LayoutNode& n, int depth = 0) {
    std::string indent(depth * 2, ' ');
    std::println("{}{}: ({},{}) {}x{} basis={} grow={}",
        indent, kind_name(n),
        n.computed.x, n.computed.y,
        n.computed.w, n.computed.h,
        n.style.basis, n.style.grow);
    for (int c : n.children) dump(nodes[c], depth + 1);
}
\end{lstlisting}

The answer is almost always visible in the first three lines of
output.

\paragraph{4. Use the DSL's debug border.} Wrap the suspicious
node in a border (\code{| border\_<Single>}); the corner glyphs
show you where its rectangle actually is. If the border is one
cell off, the layout is one cell off.

\begin{tryit}
Run any \maya{} sample and resize your terminal slowly. The
layout reflows on every \code{SIGWINCH}; you should see boxes
rearrange in cell steps, never half a cell. Try narrowing until
your status bar wraps; the wrap point is where total basis +
gaps exceeds the new available width.
\end{tryit}

\section{Layout recipes}
\label{sec:layout:recipes}

A handful of common shapes and how to spell them in the DSL.

\paragraph{Sidebar + content.}

\begin{lstlisting}
h(
    sidebar  | w_<24>,
    content  | grow_<1>
)
\end{lstlisting}

Sidebar fixed at 24; content takes everything else.

\paragraph{Header + body + footer.}

\begin{lstlisting}
v(
    header,
    body | grow_<1>,
    footer
)
\end{lstlisting}

Header and footer take their natural height; body fills the
leftover vertical space.

\paragraph{Three-pane equal split.}

\begin{lstlisting}
h(
    pane_a | grow_<1>,
    pane_b | grow_<1>,
    pane_c | grow_<1>
) | gap_<1>
\end{lstlisting}

Three equal-width panes with one-cell gaps. Each gets
$(W - 2) / 3$ cells.

\paragraph{Status bar on the right.}

\begin{lstlisting}
h(
    text(left),
    space,                  // spacer (grow=1, basis=0)
    text(right) | Dim
)
\end{lstlisting}

The spacer absorbs all leftover space, pushing the right text to
the edge.

\paragraph{Centred dialog.}

\begin{lstlisting}
h(
    space,
    v(
        space,
        dialog | w_<60>,
        space
    ),
    space
)
\end{lstlisting}

Four spacers form a cross-shaped frame around a 60-column dialog,
centring it horizontally and vertically.

\paragraph{A list with a fixed header.}

\begin{lstlisting}
v(
    header_row | Bold,
    sep,
    map(items, render_row) | grow_<1>
)
\end{lstlisting}

Header takes its row, separator one row, the body takes the rest.
When the list overflows, the \code{| grow\_<1>} on the mapped
content makes its container scrollable (Chapter~\ref{ch:widgets}
covers scroll surfaces).

\section{Performance notes}
\label{sec:layout:perf}

Layout is $O(n)$ in the number of nodes for an unwrapped tree.
With wrap, the worst case is $O(n^2)$ (every child triggers a
re-evaluation of remaining line capacity), but typical trees do
not wrap and the constant factor is small.

Numbers from the maya benchmark suite, on a typical desktop CPU:

\begin{itemize}[leftmargin=*]
  \item A 50-node tree (typical small app): under 5\,\textmu{}s.
  \item A 500-node tree (a long list of items): around
    30\,\textmu{}s.
  \item A 5000-node tree (pathological --- a maximalist
    dashboard): around 300\,\textmu{}s.
\end{itemize}

Layout is rarely the bottleneck. If it shows up in your profiler,
the usual cause is rebuilding too many DSL nodes per frame, not
the solver itself. Flatten where you can; cache \code{constexpr}
subtrees that do not change.

\section{Pitfalls and anti-patterns}
\label{sec:layout:pitfalls}

\begin{warning}
\textbf{Negative space from over-applied padding.} If parent and
child both apply padding, the visible insets compound. Decide which
layer owns the padding and keep it there.
\end{warning}

\begin{warning}
\textbf{A grow child inside a non-growing parent.} The parent sizes
itself to content, so there is no leftover space to distribute, and
the child does not grow. To get a growing layout, every ancestor up
to the root must be growing or explicitly sized.
\end{warning}

\begin{warning}
\textbf{Mixing \code{shrink=0} children with insufficient width.}
If every child refuses to shrink and they collectively exceed the
parent's width, content overflows. The renderer clips; the user
sees text cut off mid-glyph. Allow at least the variable-content
child to shrink.
\end{warning}

\begin{warning}
\textbf{Forgetting that text has intrinsic width.} A
\code{TextElement} in a row sizes to its content width by default.
Long strings push siblings off the edge. Wrap them in a fixed-width
box (\code{| w\_<30>}) or set their wrap mode to
\code{TextWrap::Truncate}.
\end{warning}

\begin{recap}
Layout decides where things go. \maya{} uses an integer-cell flexbox
solver over a flat \code{vector<LayoutNode>} addressed by index.
Three phases: measure, distribute, place. The same model handles
rows, columns, wrapping, alignment, justification, and nested boxes.
\end{recap}

\section*{Workshop}

\begin{enumerate}[leftmargin=*]
  \item Write a row that displays your name on the left and the date
    on the right, with the space between them growing to fill the
    terminal width.
  \item Build a 3-column dashboard: three boxes, each with
    \code{grow=1}, separated by 1-cell gaps. Resize your terminal
    --- watch the boxes split the new width.
  \item Read \file{maya/src/layout/yoga.cpp}. Identify the three
    phases (measure, distribute, place). Note where the available
    width is subdivided.
\end{enumerate}

\chapter{Concepts: the contracts the DSL enforces}
\label{ch:concepts}

\begin{aside}
Templates can do almost anything. That is their power and their
curse. A library that uses templates freely will, eventually,
produce error messages thousands of lines long where the actual
problem is a single missing function. The C++20 concepts feature
is the cure: a way to say ``this template parameter must satisfy
\emph{this} contract'' and have the compiler enforce it directly.
\maya{} leans on concepts heavily. This chapter walks through the
ones you will encounter and the ones you will write.
\end{aside}

\section{What a concept is, mechanically}

A \emph{concept} is a named compile-time predicate on a type. You
declare it once; you use it as a constraint anywhere a template
parameter appears.

\begin{lstlisting}
template <class T>
concept Addable = requires (T a, T b) {
    { a + b } -> std::convertible_to<T>;
};

template <Addable T>           // T must satisfy Addable
T sum(T a, T b) { return a + b; }
\end{lstlisting}

The \code{requires} clause inside the concept lists expressions
that must compile. \code{a + b} must be a valid expression and
its result must convert to \code{T}. If you call \code{sum(1, 2)},
the compiler checks: does \code{int + int} exist? Yes. Does it
convert to \code{int}? Yes. Substitution proceeds.

If you call \code{sum(std::string\{"a"\}, std::string\{"b"\})}, the
check runs again: does \code{string + string} exist? Yes. Does it
convert to \code{string}? Yes. Proceed.

If you call \code{sum(std::ifstream\{\}, std::ifstream\{\})}, the
check fails: \code{ifstream + ifstream} is not a thing. The
compiler emits a focused error: ``constraint Addable<ifstream>
not satisfied because a + b is ill-formed''.

\subsection*{Before concepts: SFINAE}

Pre-C++20, the same effect was achieved with SFINAE (Substitution
Failure Is Not An Error):

\begin{lstlisting}
template <class T,
          class = std::enable_if_t<std::is_arithmetic_v<T>>>
T sum(T a, T b) { return a + b; }
\end{lstlisting}

This worked, but the error messages were brutal. The constraint
was hidden in a \code{std::enable\_if\_t} that produced no
explanation; failures looked like ``no matching function''.

Concepts fixed three things:

\begin{itemize}[leftmargin=*]
  \item The constraint has a \emph{name} you can see in errors.
  \item The constraint is \emph{first-class}: you can compose
    concepts with \code{\&\&}, \code{||}, \code{!}.
  \item The compiler checks the constraint \emph{before} trying
    to instantiate the function body, so errors point at the
    constraint, not at the failing internals.
\end{itemize}

\begin{funfact}
Concepts were proposed for C++0x and dropped in 2009 because the
design was too complex. A simplified version (Concepts Lite)
appeared in 2014 as TS and finally shipped in C++20. Bjarne
Stroustrup has been writing about concepts since the 1980s; the
wait was personal for him.
\end{funfact}

\section{The concepts in maya's core}

\maya{} defines a handful of concepts in
\file{maya/include/maya/core/concepts.hpp}. Each names a contract
that some other part of the library relies on. Walking them is
the fastest way to understand the library's surface.

\subsection*{Program}

The single most important concept. It states what an application
must provide to be runnable by \code{maya::run}.

\begin{lstlisting}
template <class P>
concept Program = requires {
    typename P::Model;
    typename P::Msg;
    { P::init() }
        -> std::same_as<std::pair<typename P::Model, Cmd<typename P::Msg>>>;
    requires requires (typename P::Model m, typename P::Msg msg) {
        { P::update(m, msg) }
            -> std::same_as<std::pair<typename P::Model, Cmd<typename P::Msg>>>;
        { P::view(m) }      -> std::same_as<Element>;
        { P::subscribe(m) } -> std::same_as<Sub<typename P::Msg>>;
    };
};
\end{lstlisting}

Six requirements:

\begin{enumerate}[leftmargin=*]
  \item \code{P::Model} is a nested type.
  \item \code{P::Msg} is a nested type.
  \item \code{P::init()} is callable and returns a pair of the
    initial model and an initial command.
  \item \code{P::update(model, msg)} returns the same pair shape.
  \item \code{P::view(model)} returns an \code{Element}.
  \item \code{P::subscribe(model)} returns a \code{Sub<Msg>}.
\end{enumerate}

The \code{run} function takes a \code{Program} as a template
parameter:

\begin{lstlisting}
template <Program P>
int run(Config cfg = {});
\end{lstlisting}

If you mistype \code{view} as \code{View}, the compile error names
\code{Program} as the unsatisfied concept and points at the
missing function. No twenty-deep template trace.

\begin{idea}
Concepts in \maya{} are not decorative. They are documentation,
contract, and error-message generator in one. When you read
\code{template <Program P>}, you know exactly what \code{P}
must look like, because the concept is right there.
\end{idea}

\subsection*{ElementLike}

The concept that gates the DSL's children:

\begin{lstlisting}
template <class T>
concept ElementLike = requires (T t) {
    { build(t) } -> std::same_as<Element>;
};
\end{lstlisting}

Any type that has a free function \code{build} returning an
\code{Element} satisfies it. \code{TextNode}, \code{BoxNode},
\code{Element} itself, and the runtime variants all qualify.
This is what lets \code{h(\ldots)} and \code{v(\ldots)} take
heterogeneous children.

\subsection*{Renderable}

Used by widgets that produce their own elements:

\begin{lstlisting}
template <class W>
concept Renderable = requires (const W& w) {
    { w.render() } -> std::same_as<Element>;
};
\end{lstlisting}

Widgets that satisfy \code{Renderable} can be dropped into the
DSL via a \code{widget(\ldots)} helper that calls
\code{w.render()}.

\subsection*{StyleSource}

A type that can supply a \code{Style} on demand. Used by the
theme system:

\begin{lstlisting}
template <class T>
concept StyleSource = requires (const T& t) {
    { t.style() } -> std::same_as<Style>;
};
\end{lstlisting}

A \code{Theme} struct, a styled text node, and a tagged builder
all satisfy this. The renderer asks any \code{StyleSource} for
its style during the merge pass.

\subsection*{Subscribable}

Some pieces of state can be observed for changes:

\begin{lstlisting}
template <class S>
concept Subscribable = requires (S s, std::function<void()> cb) {
    { s.subscribe(cb) } -> std::convertible_to<std::function<void()>>;
};
\end{lstlisting}

The signal system (Chapter~\ref{ch:signals}) uses this to
distinguish reactive sources from plain values. The subscribe
function takes a callback and returns an unsubscribe handle.

\section{Composing concepts}

Concepts compose with the usual logical operators:

\begin{lstlisting}
template <class T>
concept Drawable = ElementLike<T> && Subscribable<T>;

template <class T>
concept Numeric = std::integral<T> || std::floating_point<T>;
\end{lstlisting}

A type satisfies \code{Drawable} if and only if it satisfies both
\code{ElementLike} \emph{and} \code{Subscribable}. The standard
library's \code{std::integral} and \code{std::floating\_point}
are themselves concepts; their union is \code{Numeric}.

You can also negate:

\begin{lstlisting}
template <class T>
concept NotEmpty = !std::same_as<T, std::monostate>;
\end{lstlisting}

\begin{recap}
Concepts are named compile-time predicates. They give templates a
contract the compiler can check up front. \maya{}'s
\code{Program}, \code{ElementLike}, \code{Renderable},
\code{StyleSource}, and \code{Subscribable} concepts describe the
shape of every extensible point in the library.
\end{recap}

\section{Writing your own concept}

Suppose you have a family of components that all expose a
\code{focus()} method returning \code{void} and a \code{focused()}
method returning \code{bool}. You want a function that takes any
such component.

\begin{lstlisting}
template <class C>
concept Focusable = requires (C& c, const C& cc) {
    { c.focus() }      -> std::same_as<void>;
    { cc.focused() }   -> std::same_as<bool>;
};

template <Focusable C>
void cycle_focus(std::vector<C>& cs) {
    for (auto& c : cs)
        if (c.focused()) { c.focus(); return; }
}
\end{lstlisting}

Notice the const-correctness: \code{focused()} is called on a
const reference, so it must be a const method. The concept
catches the mistake of declaring \code{focused()} non-const.

\subsection*{Multiple requires clauses}

You can use multiple \code{requires} clauses for organisation:

\begin{lstlisting}
template <class C>
concept Container = requires (C& c, C const& cc) {
    typename C::value_type;
    typename C::iterator;
    { c.begin() }    -> std::same_as<typename C::iterator>;
    { c.end() }      -> std::same_as<typename C::iterator>;
    { cc.size() }    -> std::convertible_to<std::size_t>;
    { cc.empty() }   -> std::same_as<bool>;
};
\end{lstlisting}

Type aliases (\code{typename C::value\_type}) require the nested
name to exist. Expression checks (\code{c.begin()}) require the
expression to compile, optionally with a return type constraint.

\section{Where concepts appear in maya}

A non-exhaustive list:

\begin{itemize}[leftmargin=*]
  \item \code{template <Program P> int run(Config)} ---
    \file{maya/include/maya/app/app.hpp}.
  \item \code{template <ElementLike\ldots Cs> auto h(Cs\&\& \ldots)}
    --- \file{maya/include/maya/dsl.hpp}.
  \item \code{template <ElementLike\ldots Cs> auto v(Cs\&\& \ldots)}
    --- ditto.
  \item \code{template <Renderable W> auto widget(const W\&)} ---
    DSL helper.
  \item \code{template <class M> Cmd<M> Cmd::task(callable)}, where
    the callable must satisfy a small concept ensuring it returns
    \code{M}.
  \item Platform abstractions: \code{template <Terminal T> class
    Renderer\{\}}, where \code{Terminal} is the concept that the
    POSIX, macOS, and Win32 backends all satisfy.
\end{itemize}

Reading the header, you can see the contract before you see the
implementation. This is what makes \maya{}'s headers approachable.

\begin{funfact}
GCC, Clang, and MSVC implement concepts differently in two
places: the error message format and the order in which concept
checks happen during overload resolution. The standard pins down
the \emph{behaviour}, not the \emph{output}. If you see slightly
different messages on different compilers, that is expected.
\end{funfact}

\section{Pitfalls}

\begin{warning}
\textbf{Concepts can hide failure if you let them.} A
\code{template <class T>} with no constraint and a
\code{requires} clause inside the function body will still
produce ugly errors. Put the constraint on the parameter, not in
the body, so the compiler can reject mismatched callers before
substituting the body.
\end{warning}

\begin{warning}
\textbf{Constrained overloads need partial ordering.} If you have
two function templates with overlapping but distinct concepts,
the compiler picks the most-constrained match. If neither
subsumes the other, the call is ambiguous --- and the error is
not friendly. Compose carefully.
\end{warning}

\begin{warning}
\textbf{Concepts and forward declarations.} You cannot constrain
a forward-declared function template with a concept that
references the function's parameter types if those types depend
on the function itself. The cycle is real; break it by
introducing a helper type.
\end{warning}

\begin{warning}
\textbf{Slow compile times.} Heavy concept use slows the
front-end. \maya{} keeps concepts shallow (a handful of
requirements each) and uses them sparingly outside the public
API. If your build is slow, look for nested concepts more than
two levels deep.
\end{warning}

\section{Concepts as documentation}

The deepest payoff of concepts is not the error message; it is
the documentation. A header file that says

\begin{lstlisting}
template <Program P> int run(Config cfg = {});
\end{lstlisting}

tells you everything you need to know about \code{run}: it takes
a \code{Program} (whose contract is defined elsewhere) and a
config. If you want to know what \code{Program} requires, you go
to its definition --- one place, all the requirements listed.

Compare with the pre-concept C++ idiom:

\begin{lstlisting}
template <class P>  // P must have init(), update(), view(), subscribe()
int run(Config cfg = {});
\end{lstlisting}

The comment is wishful. Nothing enforces it. With concepts, the
comment becomes machine-checked.

\begin{idea}
A concept-driven header is a contract: the compiler is the
referee. When you read \maya{}'s headers, treat the concept
constraints as the most reliable documentation.
\end{idea}

\section*{Workshop}

\begin{enumerate}[leftmargin=*]
  \item Write a concept \code{Printable<T>} that requires
    \code{std::cout << t} to be valid. Use it on a function
    \code{void log(T)}.
  \item Read
    \file{maya/include/maya/core/concepts.hpp} and identify
    every concept declared. Note which other headers they show
    up in.
  \item Try compiling a \code{Counter} app that omits the
    \code{subscribe} method. Read the resulting error and find
    the line that names the failing concept.
\end{enumerate}

\begin{recap}
Concepts are the contracts that make \maya{}'s templates safe and
its error messages readable. Five core concepts shape the public
API. Composing concepts with \code{\&\&} and \code{||} lets you
express precise requirements. Treat them as machine-checked
documentation.
\end{recap}

\chapter{Expected: errors as values}
\label{ch:expected}

\begin{aside}
Most C++ codebases handle errors with exceptions. \maya{} mostly
does not. The reason is the same reason \code{update} is pure: an
exception is a non-local control transfer, and non-local control
transfers undermine the value-typed contract the architecture
depends on. Where \maya{} needs to fail, it returns the failure
as a value. The tool is \code{maya::Expected<T, E>}: a sum type
that holds either a success or an error.
\end{aside}

\section{Why not exceptions?}

C++ exceptions are powerful. They unwind the stack, they call
destructors, they let you write code that ignores the error path
entirely. For application code with deep call graphs, they are
often the right tool.

For a library, they are often the wrong one. Specific issues:

\begin{itemize}[leftmargin=*]
  \item \textbf{Cost.} Throwing is expensive. Modern
    implementations (Itanium ABI on x86-64, MSVC's
    \emph{Funclets}) have nearly-zero cost on the happy path,
    but the cost on the failure path is large (table walks,
    destructor calls, allocator pressure). For a render loop
    that wants tens of thousands of frames per second on the
    happy path \emph{and} a fast crash-and-recover on the rare
    failure path, the asymmetry hurts.
  \item \textbf{Invisibility.} An exception can come out of
    almost any function. The reader does not know without
    reading the body --- or the bodies of every callee.
  \item \textbf{Composability with sum-typed code.} An exception
    bypasses the type system: it does not appear in the return
    type. Two pieces of code that handle their own failures
    through \code{Expected} compose cleanly; mix one piece that
    throws and the type system is no longer the source of truth.
  \item \textbf{Coroutines, restricted environments.} Some C++
    contexts (kernel modules, freestanding, \code{-fno-except})
    do not allow exceptions. Code that uses \code{Expected}
    works there.
\end{itemize}

For \maya{}'s low-level pieces --- terminal I/O, signal handling,
SIMD selection, file watching --- errors are values. Exceptions
are reserved for genuinely exceptional, unrecoverable conditions
(out of memory, hardware failure).

\begin{funfact}
Rust's \code{Result<T, E>} type is the same idea with a different
spelling. Haskell's \code{Either e a} is older still. Swift uses
\code{throws} but its model is closer to an exhaustively-checked
sum than to C++ exceptions. The pattern is older than C++.
\end{funfact}

\section{The Expected type}

\code{maya::Expected<T, E>} is a value-typed sum of a success
\code{T} and an error \code{E}:

\begin{lstlisting}
template <class T, class E>
class Expected {
    std::variant<T, E> v;
public:
    bool has_value() const { return v.index() == 0; }
    explicit operator bool() const { return has_value(); }
    T&       value()        { return std::get<0>(v); }
    const T& value() const  { return std::get<0>(v); }
    E&       error()        { return std::get<1>(v); }
    const E& error() const  { return std::get<1>(v); }
    // monadic helpers below
};
\end{lstlisting}

The C++23 standard introduced \code{std::expected<T, E>}; \maya{}
predates that, so it ships its own. The interface is similar.

\subsection*{Construction}

You can construct from either side:

\begin{lstlisting}
Expected<int, std::string> ok{42};                       // success
Expected<int, std::string> bad{std::unexpect, "boom"};   // error
\end{lstlisting}

The \code{std::unexpect} tag distinguishes the error constructor
from the success one when both types are convertible.

\subsection*{Checking}

\begin{lstlisting}
if (auto r = parse(input); r) {
    use(r.value());
} else {
    log(r.error());
}
\end{lstlisting}

The contextual bool conversion lets \code{if (r)} mean ``has
value''. Pair with the structured \code{auto r = \ldots; if (r)}
pattern to keep the scope tight.

\subsection*{Monadic helpers}

The interesting power is the chain operations:

\begin{itemize}[leftmargin=*]
  \item \code{transform(f)}: if value, apply \code{f}; if error,
    pass through. Returns a new \code{Expected<U, E>}.
  \item \code{and\_then(f)}: like \code{transform} but \code{f}
    itself returns an \code{Expected<U, E>}. Flattens the result.
  \item \code{or\_else(f)}: if error, apply \code{f} to the error
    (which may itself return an \code{Expected}); if value, pass
    through.
  \item \code{value\_or(default)}: extract the value or a default.
\end{itemize}

These let you write a pipeline that handles errors implicitly:

\begin{lstlisting}
Expected<Config, std::string> result =
    read_file("config.toml")
        .and_then(parse_toml)
        .and_then(validate)
        .transform(apply_defaults);

if (!result) std::cerr << result.error() << "\n";
\end{lstlisting}

If any stage fails, the rest are skipped and the error flows to
the end. No exceptions, no manual short-circuit ladders.

\begin{idea}
Monadic error handling is a style choice. Some prefer the
ladder; some prefer the chain. \code{Expected} supports both
and lets you pick per call site.
\end{idea}

\section{Where Expected appears in maya}

The main producers of \code{Expected} in \maya{}:

\begin{itemize}[leftmargin=*]
  \item \code{Terminal::open()} returns
    \code{Expected<Terminal, IoError>}. Failure cases:
    permission denied, device gone, raw-mode setup failed.
  \item \code{ansi::parse(byte\_stream)} returns
    \code{Expected<Event, ParseError>}. Failure cases:
    malformed escape, unterminated sequence.
  \item \code{file::read(path)} returns
    \code{Expected<std::string, IoError>}.
  \item \code{platform::detect\_simd()} returns
    \code{Expected<SimdLevel, std::string>}. Failure case: CPU
    feature detection failed.
  \item Widget initialisation in \file{maya/include/maya/widget/}
    --- e.g. \code{Picker::open(items)} returns
    \code{Expected<Picker, std::string>}.
\end{itemize}

The recurring shape: a function that interacts with the outside
world (OS, hardware, files) returns \code{Expected}. Pure
functions return their value directly.

\section{Error types}

\code{Expected} is parameterised on the error type. \maya{} uses
a few:

\begin{itemize}[leftmargin=*]
  \item \code{IoError}: a struct with a \code{std::errc} code, a
    \code{std::string} context message, and the syscall name.
  \item \code{ParseError}: a struct with a source offset and a
    description.
  \item \code{std::string}: when nothing more structured is
    needed.
\end{itemize}

You can use whatever error type fits. Tip: keep the error type
small and easy to construct, so the failure path is cheap.

\subsection*{IoError, in detail}

\begin{lstlisting}
struct IoError {
    std::errc code;          // standard error category
    std::string context;     // free-form, e.g. "opening /dev/tty"
    std::string syscall;     // "open", "read", "tcsetattr"
};
\end{lstlisting}

A constructor takes the \code{errno} value via
\code{std::make\_error\_condition} and stashes everything in one
place. When you log it:

\begin{lstlisting}
std::cerr << e.syscall << " failed ("
          << std::make_error_code(e.code).message() << "): "
          << e.context << "\n";
\end{lstlisting}

You get a message like ``open failed (No such file or
directory): /dev/tty12''.

\section{Crossing the boundary into Cmd}

In Chapter~\ref{ch:elm} we said \code{Cmd::task} returns a message.
What if the task fails?

The pattern is to make the message a sum:

\begin{lstlisting}
struct LoadOk    { std::string content; };
struct LoadFail  { IoError err; };
using Msg = std::variant<LoadOk, LoadFail, /* ... */>;

static Cmd<Msg> load_cmd(std::string path) {
    return Cmd<Msg>::task([path]() -> Msg {
        auto r = file::read(path);
        if (r) return LoadOk{std::move(r.value())};
        return LoadFail{std::move(r.error())};
    });
}
\end{lstlisting}

The task takes an \code{Expected}, peels it open, and routes the
result through the message variant. \code{update} then handles
both cases as branches. The error is data flowing through the
same pipeline as success.

\begin{warning}
Do not unwrap an \code{Expected} with \code{.value()} without
checking. \code{.value()} on an error throws. The whole point of
\code{Expected} is to avoid exceptions; calling \code{.value()}
without an \code{if (r)} undoes that.
\end{warning}

\section{Patterns}

\subsection*{Early return ladder}

\begin{lstlisting}
Expected<Config, std::string> load_config(std::string_view path) {
    auto raw = file::read(path);
    if (!raw) return std::unexpect, "read failed: " + raw.error().context;

    auto parsed = parse_toml(raw.value());
    if (!parsed) return std::unexpect, "parse failed: " + parsed.error();

    auto valid = validate(parsed.value());
    if (!valid) return std::unexpect, "validate failed: " + valid.error();

    return apply_defaults(valid.value());
}
\end{lstlisting}

The C-style ``check at every step'' approach. Verbose but
explicit; everyone can read it.

\subsection*{Monadic chain}

The same function with \code{and\_then}:

\begin{lstlisting}
Expected<Config, std::string> load_config(std::string_view path) {
    return file::read(path)
        .transform_error([](auto& e) {
            return "read failed: " + e.context;
        })
        .and_then(parse_toml)
        .and_then(validate)
        .transform(apply_defaults);
}
\end{lstlisting}

Shorter but requires the reader to understand the chain. Both
styles are fine; pick by audience.

\subsection*{Try-helper macro}

Some codebases define a \code{TRY} macro that hides the ladder:

\begin{lstlisting}
#define TRY(x) ({ auto __r = (x); if (!__r) return std::unexpect, __r.error(); std::move(__r.value()); })

Expected<Config, std::string> load_config(std::string_view path) {
    auto raw    = TRY(file::read(path));
    auto parsed = TRY(parse_toml(raw));
    auto valid  = TRY(validate(parsed));
    return apply_defaults(valid);
}
\end{lstlisting}

\maya{} does not use \code{TRY} in its public API (it relies on
a GCC extension), but the pattern is common in private code.

\begin{aside}
Rust has \code{?} for exactly this pattern, built into the
language. C++ is exploring similar syntax (P0779, ``Notes on the
syntax of try operators''), but it is years away. The macro is
the workaround.
\end{aside}

\section{Performance}

\code{Expected<T, E>} costs a variant: a tag byte plus enough
storage for the larger of \code{T} and \code{E}. For \code{T =
int} and \code{E = std::string}, that is around 40 bytes.

The happy path is one branch (the tag check); a missed branch
prediction is the worst case, perhaps a dozen cycles. The error
path is a string move, which is cheap.

Compared to exceptions:

\begin{itemize}[leftmargin=*]
  \item Happy path: \code{Expected} is one branch. Exceptions are
    zero branches (the unwinding tables sit dormant).
    Effectively a tie; the exception edges ahead by a hair on
    the very hot path.
  \item Error path: \code{Expected} is a value construction and a
    branch. Exceptions are a table walk, destructor calls,
    allocator pressure. \code{Expected} wins by an order of
    magnitude.
  \item Throughput: \code{Expected} is uniformly fast. Exceptions
    are bimodal.
\end{itemize}

For a renderer that wants predictable latency, \code{Expected} is
the better trade.

\section{Pitfalls}

\begin{warning}
\textbf{Bare types in the variant.} If \code{T} and \code{E}
are the same type, you cannot distinguish them. Use a wrapper
or a tag. \code{Expected<int, int>} compiles but is a foot-gun.
\end{warning}

\begin{warning}
\textbf{Implicit conversions from E to T.} If \code{T} accepts
construction from \code{E}, calling
\code{Expected<T, E>\{some\_E\}} might construct a success
holding a converted \code{T}. Use \code{std::unexpect} to be
unambiguous.
\end{warning}

\begin{warning}
\textbf{Discarded errors.} If you call a function returning
\code{Expected} and ignore the result, you lose the error. GCC
warns with \code{-Wunused-result} if the return type is marked
\code{[[nodiscard]]}; \maya{}'s \code{Expected} carries that
attribute.
\end{warning}

\begin{warning}
\textbf{Nested Expected.} \code{Expected<Expected<int, E>, E>} is
legal but ugly. Use \code{and\_then} to flatten, or refactor the
inner call to return the same error type.
\end{warning}

\section{When to throw}

Not every failure is recoverable. For genuinely exceptional
conditions --- assertion failures, out-of-memory, undefined
behaviour detected at runtime --- exceptions or \code{abort()}
are appropriate. \maya{} uses:

\begin{itemize}[leftmargin=*]
  \item \code{std::bad\_alloc}: propagated from allocations, not
    caught.
  \item \code{std::logic\_error}: thrown from APIs whose
    preconditions were violated (a sign of a bug in the caller).
  \item \code{abort()}: from internal sanity checks that should
    never fail (paired with an assert).
\end{itemize}

The dividing line: if the failure is something the caller might
reasonably want to recover from, use \code{Expected}. If the
failure indicates a bug, throw or abort.

\section*{Workshop}

\begin{enumerate}[leftmargin=*]
  \item Write a function \code{Expected<int, std::string>
    parse\_int(std::string\_view)} that returns the parsed
    integer or an error message.
  \item Chain three calls of \code{parse\_int} on different
    strings with \code{and\_then}, summing the results. What
    happens when one fails?
  \item Refactor a function in your code that throws an
    exception into one that returns \code{Expected}. Note where
    the type signature change improves the call site.
\end{enumerate}

\begin{recap}
\code{Expected<T, E>} is a value-typed sum of success and error.
It lets functions return failures as data, making the call sites
explicit and the type system the source of truth. \maya{} uses
it for every I/O boundary and reserves exceptions for genuinely
exceptional conditions.
\end{recap}

\chapter{The render context and the pipeline}
\label{ch:pipeline}

\begin{aside}
We have talked about \code{view} returning an element, layout
assigning rects, and the renderer emitting bytes. The piece in
between --- the \emph{pipeline} that ties them together --- is
worth a chapter of its own. It is where \maya{} decides what to
recompute and what to reuse, and it is the boundary between the
pure world of \code{update}/\code{view} and the messy world of
file descriptors.
\end{aside}

\section{What the pipeline does}

Each frame, the pipeline:

\begin{enumerate}[leftmargin=*]
  \item Asks the program for a fresh \code{Element} via
    \code{view(model)}.
  \item Lays it out against the current terminal size.
  \item Paints it into a canvas of packed cells.
  \item Diffs the canvas against the previous one.
  \item Serialises the diff into ANSI bytes.
  \item Writes the bytes to the terminal.
\end{enumerate}

Six steps. Each step has a well-defined input and output. The
boundaries are stable; the implementations can change. That is
the discipline of a pipeline architecture.

\subsection*{Why a pipeline?}

Two alternatives existed:

\begin{itemize}[leftmargin=*]
  \item \textbf{Direct rendering.} \code{view} would return a
    sequence of ``draw'' calls that the terminal executes
    directly. Tight coupling. Tested only at the integration
    level. Imagine debugging \emph{any} part of it.
  \item \textbf{Single-pass rendering.} Merge layout and paint
    into one walk. Faster in principle. But you lose the ability
    to inspect intermediate results, and adding features
    (themes, hit testing, animations) requires touching the one
    monolithic pass.
\end{itemize}

A pipeline gives you intermediate values you can inspect, test,
and replace. Layout takes an element and produces a layout tree.
Paint takes a layout tree and produces a canvas. Diff takes two
canvases and produces a region list. Each is a function from one
value to another; each is testable in isolation.

\begin{idea}
A pipeline is what you get when you apply the Elm architecture's
``effects as data'' philosophy to the renderer itself. The
intermediate values are the data; the stages are the functions.
\end{idea}

\section{The RenderContext}

\file{maya/include/maya/core/render\_context.hpp} defines the
struct that flows through the pipeline:

\begin{lstlisting}
struct RenderContext {
    Element       tree;        // from view()
    LayoutTree    layout;      // after measure/distribute/place
    Canvas        canvas;      // after paint
    DiffResult    diff;        // diff against previous canvas
    SerializeResult bytes;     // ANSI bytes to write
    Stats         stats;       // timings, counts
};
\end{lstlisting}

Each field is populated by one stage. By the end of the frame,
all are filled. The context is then reused for the next frame
(only \code{canvas} and \code{stats} persist; the rest are
recomputed).

\subsection*{Why a single struct?}

We could have used six free functions chained by hand. Bundling
into a struct does three things:

\begin{itemize}[leftmargin=*]
  \item \textbf{Logging.} The context can be dumped at any
    point. ``Why is this frame slow?'' is answerable by
    inspecting \code{stats}; ``why is this cell wrong?'' by
    inspecting \code{canvas}.
  \item \textbf{Lifetimes.} The struct owns all its fields. Pass
    a \code{const RenderContext\&} into any stage and it can
    read whatever it needs without lifetime puzzles.
  \item \textbf{Threading.} If we ever want to run layout on a
    different thread from paint, the context is the unit of
    handoff.
\end{itemize}

\section{The six stages, expanded}

\subsection*{Stage 1: view}

The pipeline calls \code{P::view(model)}. The result is an
\code{Element}. This step is allocator-pressure heavy but
otherwise trivial: it is your code.

Cost: depends on \code{view}. Microseconds for a typical UI.

\subsection*{Stage 2: layout}

Convert the \code{Element} tree to a \code{LayoutTree}: a flat
\code{vector<LayoutNode>} where each node has a final
\code{Rect}. Layout is its own three-phase pass (measure,
distribute, place) detailed in Chapter~\ref{ch:layout}.

Cost: O(n) in tree size. Tens of microseconds for hundreds of
nodes.

\subsection*{Stage 3: paint}

Walk the layout tree and write packed cells into a \code{Canvas}.
A \code{PackedCell} is a 64-bit struct holding a codepoint, a
style ID (an index into a pool), and flags. The canvas is a flat
\code{vector<PackedCell>} of width $\times$ height entries.

Cost: O(cells) = O(width $\times$ height). For 80$\times$24,
that is 1920 entries. Microseconds.

\subsection*{Stage 4: diff}

Compare the current \code{Canvas} against the previous one. The
output is a list of \code{Row} differences: for each row that
changed, the leftmost and rightmost dirty cell.

The diff is SIMD-accelerated (Chapter~\ref{ch:simd}). For an
80$\times$24 grid, it is sub-microsecond on modern CPUs.

\subsection*{Stage 5: serialize}

Convert the diff into ANSI bytes. This stage minimises:

\begin{itemize}[leftmargin=*]
  \item Cursor moves: only emit when the cursor would not
    naturally arrive at the next dirty cell.
  \item SGR sequences: only emit when the style changes from the
    previous cell.
  \item Buffer reuse: the byte buffer carries over capacity from
    frame to frame.
\end{itemize}

Cost: O(dirty cells). For a typical frame with light changes, a
few thousand bytes.

\subsection*{Stage 6: write}

Write the bytes to the terminal via the writer abstraction
(\file{maya/include/maya/terminal/writer.hpp}). The writer
batches calls and flushes once per frame. Synchronous; the cost
is the kernel write syscall plus terminal latency.

Cost: dominated by kernel time. Tens of microseconds typically.

\section{Caching and the cache\_id}

Some pipeline stages can be skipped if their inputs have not
changed. \maya{} tracks input identity with \code{cache\_id} ---
a 64-bit fingerprint computed from the input.

\begin{lstlisting}
struct CacheId {
    uint64_t value;
    bool operator==(const CacheId&) const = default;
};
\end{lstlisting}

If the new \code{CacheId} for the element tree matches the
previous frame's, the entire pipeline can be skipped: no work,
no bytes. The frame is silent.

This is the optimisation that lets a 60 Hz idle app burn
essentially no CPU. \code{view} re-runs every frame (cheap), but
if it produces the same tree, no paint, no diff, no write.

\subsection*{Computing the fingerprint}

The fingerprint walks the tree once, mixing every leaf's content
into a hash. Two trees with identical content produce identical
fingerprints. Two trees that differ in even one cell produce
different fingerprints (with overwhelming probability; collisions
are extremely rare).

\maya{} uses xxHash3 (fast, non-cryptographic) for the mixing.
The walk is itself O(n) in tree size --- same as layout --- so the
trade-off is one extra walk to skip up to five.

\begin{aside}
This is the same trick React uses with \code{shouldComponentUpdate}
or the newer \code{React.memo}: skip work if the inputs are
unchanged. The difference: \maya{}'s cache key is computed
automatically from the tree, not declared by the user.
\end{aside}

\subsection*{Damage tracking}

Some changes do not affect the visual output. A model field that
\code{view} does not consult can change without producing a
different element. Cache fingerprints catch this for free: the
tree is the same, so the fingerprint is the same.

Other changes affect only a sub-region. If a single line changes
in a status bar, only that row's diff is non-empty. The
serializer emits a cursor move plus the changed cells; everything
else is silent.

\section{Stats and observability}

The \code{Stats} struct holds per-frame timings and counters:

\begin{lstlisting}
struct Stats {
    Duration view_time;
    Duration layout_time;
    Duration paint_time;
    Duration diff_time;
    Duration serialize_time;
    Duration write_time;
    Duration total;

    size_t   cells_total;
    size_t   cells_dirty;
    size_t   bytes_written;
    bool     full_redraw;
};
\end{lstlisting}

After each frame the pipeline emits the stats; if you wire a
debugging widget to read them, you can see exactly where time
goes. \file{maya/include/maya/widget/} ships a
\code{render\_stats} widget for this purpose.

\begin{idea}
The pipeline is observable because the context is a value.
Logging, tracing, profiling, replaying --- all become easy when
the data flows through a single typed channel.
\end{idea}

\section{Adding a stage}

Suppose you want to add a stage between paint and diff that
applies a CRT scanline effect (every other row darkened). The
pipeline accommodates this:

\begin{enumerate}[leftmargin=*]
  \item Add a new field to \code{RenderContext}:
    \code{Canvas filtered\_canvas;}
  \item Insert the stage:
    \code{filtered\_canvas = scanlines(canvas);}
  \item Change the diff input from \code{canvas} to
    \code{filtered\_canvas}.
\end{enumerate}

No other stage cares. The pipeline structure makes the change
local.

\maya{} does not ship a scanline filter, but the example shows
how the pipeline can be extended. Real extensions in flight:
gamma correction for image widgets, optional anti-aliased text.

\section{Threading model}

\maya{}'s pipeline is single-threaded on the main loop. The
reasons:

\begin{itemize}[leftmargin=*]
  \item \textbf{Determinism.} A single thread guarantees
    deterministic ordering, which simplifies replay and tests.
  \item \textbf{Latency.} The pipeline is fast enough that
    threading would add more overhead than it saves. A 60 Hz
    UI has 16ms per frame; \maya{}'s pipeline takes under 1ms.
  \item \textbf{Simplicity.} No locks, no thread-safety
    requirements on \code{view}, no race conditions.
\end{itemize}

\code{Cmd::task} and \code{Cmd::isolated\_task} use worker
threads, but they return their results as messages onto the main
loop's queue. The pipeline itself never runs on a worker thread.

\begin{aside}
If you ever need to render off-thread (e.g. a background
preview), you can run the pipeline as a free function on any
thread you like --- the stages are pure. \code{maya::run}'s
pipeline is single-threaded; the algorithm is not.
\end{aside}

\section{Pitfalls}

\begin{warning}
\textbf{Doing work in view that should be done in update.} If
your \code{view} performs an expensive computation (parsing JSON,
walking a graph), it pays that cost \emph{every frame}. Cache the
result in \code{Model} and update it only when the source data
changes.
\end{warning}

\begin{warning}
\textbf{Large element trees.} The cache fingerprint is fast but
not free. A tree with tens of thousands of nodes will spend
non-trivial time on the fingerprint walk. Use \code{when} and
\code{visible} to elide branches you do not need.
\end{warning}

\begin{warning}
\textbf{Misreading stats.} Frame total is dominated by
\code{write\_time} on most terminals; that is the terminal's
latency, not yours. Focus on \code{view\_time} +
\code{layout\_time} for things you can optimise.
\end{warning}

\begin{warning}
\textbf{Forcing redraw too often.} \code{Cmd::force\_redraw}
discards the cache and emits every cell. Useful after a
disturbance, abusive every frame. Use it only when you must.
\end{warning}

\section{What this means in practice}

You will rarely write code that touches the pipeline directly.
\code{maya::run} sets it up; you write \code{init}, \code{update},
\code{view}, \code{subscribe} and the pipeline handles the rest.

But knowing the pipeline exists lets you:

\begin{itemize}[leftmargin=*]
  \item Reason about performance. ``My frame is slow'' becomes
    ``which stage is slow?''.
  \item Reason about cost. ``Why is my CPU at 5\% on an idle
    app?'' is answered by ``the cache is not catching the
    no-change case''.
  \item Trust that the pipeline is testable. Each stage is a
    pure function from one value to another.
\end{itemize}

\section*{Workshop}

\begin{enumerate}[leftmargin=*]
  \item Find \file{maya/include/maya/render/pipeline.hpp} and
    identify the six stages. Note that the source closely
    mirrors the description in this chapter.
  \item Run a \maya{} app with the debug stats widget enabled
    (set \code{MAYA\_RENDER\_STATS=1}). Watch \code{view\_time},
    \code{layout\_time}, and \code{bytes\_written} change as
    you interact.
  \item Hypothesise: which stage will dominate if you render a
    full-screen text editor with 10\,000 lines? Try it and
    confirm.
\end{enumerate}

\begin{recap}
The pipeline is the six-stage transformation from \code{Model} to
ANSI bytes: view, layout, paint, diff, serialize, write. The
\code{RenderContext} struct carries the intermediates. Cache
fingerprints skip the entire pipeline on no-change. The result is
predictable latency, observable performance, and a clean place to
add new stages.
\end{recap}

\chapter{Style, theme, and the palette}
\label{ch:style}

\begin{aside}
Colours and fonts are the bit users see first. A library that
treats them carelessly produces UIs that look wrong on half the
terminals they run on. \maya{}'s style system is more elaborate
than the typing-friendly DSL might suggest: a palette of named
roles, a theme that maps roles to colours, a style struct that
overrides per-element, and a merger that combines them under a
defined precedence.
\end{aside}

\section{Three layers of style}

When the renderer paints a cell, the cell's style is the result
of three layers, applied in order:

\begin{enumerate}[leftmargin=*]
  \item \textbf{Theme.} The application-level palette. Maps named
    roles (\emph{foreground}, \emph{accent}, \emph{warning})
    to concrete colours.
  \item \textbf{Inherited style.} The style applied by parent
    boxes via the DSL pipe or runtime equivalent.
  \item \textbf{Local style.} The style attached to the
    text/box at the leaf.
\end{enumerate}

Each layer can override the previous. The merge rule
(Chapter~\ref{ch:elements}) decides per-attribute whether to
inherit or override.

\subsection*{Why three layers?}

Two layers (theme + override) would work for simple cases. The
third (inheritance) is what makes ``the whole panel is bold''
work without setting Bold on every leaf.

The trade-off: more layers means more bookkeeping. \maya{}'s
merge is a small constant-time operation per cell, so the cost
is acceptable.

\section{The palette}

A palette is a fixed-size array of \code{Color} values, each
indexed by a named role. \maya{}'s default palette has 24 slots:

\begin{itemize}[leftmargin=*]
  \item \code{Foreground}, \code{Background}: the default
    foreground and background.
  \item \code{Muted}: dim text (timestamps, labels).
  \item \code{Subtle}: secondary text (placeholders).
  \item \code{Primary}: the brand colour.
  \item \code{Accent}: highlight colour.
  \item \code{Success}, \code{Warning}, \code{Danger},
    \code{Info}: semantic colours.
  \item \code{Border}: default border colour.
  \item \code{Selection}, \code{Cursor}: text-input UI.
  \item Ten ``user'' slots for application-specific roles.
\end{itemize}

A palette is just a \code{std::array<Color, 24>}, plus a small
enum for the indices. Looking up is constant time.

\subsection*{Roles vs colours}

The point of the palette is to decouple ``what the cell means''
from ``what colour to draw''. Code that wants a warning highlight
asks for \code{Palette::Warning}; the theme decides whether that
is amber, orange, or yellow.

This is the same indirection CSS uses (via custom properties) and
that design systems (Tailwind, Material Design) build whole
ecosystems around.

\section{The Theme struct}

A theme is a palette plus metadata:

\begin{lstlisting}
struct Theme {
    Palette  palette;
    std::string name;        // "dark", "light", "solarized"
    bool     is_dark;        // light-mode UIs flip choices
    Border   default_border; // BorderStyle::Rounded by default
};
\end{lstlisting}

\maya{} ships three themes out of the box: \code{Dark} (default),
\code{Light}, and \code{HighContrast}. Loading a theme replaces
the palette wholesale.

\subsection*{Defining a theme}

\begin{lstlisting}
constexpr Theme my_theme = {
    .palette = {
        [Role::Foreground] = Color::rgb(220, 220, 240),
        [Role::Background] = Color::default_(),
        [Role::Muted]      = Color::rgb(120, 130, 150),
        [Role::Primary]    = Color::rgb(120, 180, 250),
        // ...
    },
    .name = "midnight",
    .is_dark = true,
};
\end{lstlisting}

Designated initialisers (C++20) let you skip slots you do not
care about; they retain their default values.

\section{The Style struct, revisited}

\begin{lstlisting}
struct Style {
    Color fg = Color::default_();
    Color bg = Color::default_();
    Attribute attributes = Attribute::None;
    std::optional<std::string> hyperlink;
};
\end{lstlisting}

We saw this in Chapter~\ref{ch:elements}. The colour fields can
be a literal RGB, an Ansi16 index, an Ansi256 index, or
``default'' (let the terminal pick).

A style created from a role:

\begin{lstlisting}
Style accent_style = Style::from_role(Role::Accent);
\end{lstlisting}

This is a thin helper that looks up the colour in the current
theme. The lookup is deferred until paint --- if the theme
changes between frames, the style picks up the new value.

\subsection*{Attribute, in detail}

\code{Attribute} is a bitset of typographic attributes:

\begin{itemize}[leftmargin=*]
  \item \code{Bold}: SGR 1.
  \item \code{Dim}: SGR 2.
  \item \code{Italic}: SGR 3.
  \item \code{Underline}: SGR 4.
  \item \code{Blink}: SGR 5. (Rarely supported well; avoid.)
  \item \code{Reverse}: SGR 7. Swaps fg and bg.
  \item \code{Strike}: SGR 9.
  \item \code{DoubleUnderline}: SGR 21.
  \item \code{Overline}: SGR 53. (Even more rarely supported.)
\end{itemize}

Stored as a small bitmask. Two attributes combine with bitwise
OR; this is what makes \code{Bold | Italic} mean ``both''.

\section{Style merge: the precedence rule}

The merge rule, again, with the formal precedence:

\begin{enumerate}[leftmargin=*]
  \item \textbf{Attributes.} Result is the union of both. Bold
    in either input means Bold in the output.
  \item \textbf{Colours.} If the inner style sets a non-default
    colour, the inner wins. Otherwise, the outer's colour
    propagates.
  \item \textbf{Hyperlink.} Inner wins if set.
\end{enumerate}

In code:

\begin{lstlisting}
Style merge(const Style& outer, const Style& inner) {
    Style r = outer;
    r.attributes |= inner.attributes;
    if (!inner.fg.is_default()) r.fg = inner.fg;
    if (!inner.bg.is_default()) r.bg = inner.bg;
    if (inner.hyperlink)        r.hyperlink = inner.hyperlink;
    return r;
}
\end{lstlisting}

The asymmetry (attribute union, colour override) is what makes
the system intuitive. ``Inherit the Bold from the parent and the
colour from the child'' is what most users mean.

\begin{idea}
This is similar to CSS cascade rules, simplified. CSS has six
levels of specificity; \maya{} has two (outer, inner) and one
rule (attribute-union, colour-override). The simpler rule wins
in clarity.
\end{idea}

\section{Themes at runtime}

To change theme:

\begin{lstlisting}
maya::set_theme(Theme::Light());  // or any custom theme
\end{lstlisting}

The pipeline notices on the next frame and reflows. Any style
that referenced a role updates automatically.

\subsection*{Persisting theme preferences}

If you store the theme name in a config file, you can restore it
on startup:

\begin{lstlisting}
auto config = load_config();
maya::set_theme(themes::by_name(config.theme).value_or(Theme::Dark()));
\end{lstlisting}

The pipeline does not care where the theme came from; it just
applies whatever is current.

\section{High-contrast and accessibility}

The \code{HighContrast} theme is designed for users who need
sharp distinctions: pure black background, pure white foreground,
saturated accent colours. The palette is calibrated to a minimum
WCAG-AA-equivalent contrast ratio.

If your app respects \code{NO\_COLOR=1} (an environment variable
that means ``do not use colour''), \maya{} will use a theme that
relies on attributes (Bold, Underline) instead of colour:

\begin{lstlisting}
if (std::getenv("NO_COLOR")) {
    maya::set_theme(themes::Monochrome());
}
\end{lstlisting}

\begin{aside}
The NO\_COLOR convention was proposed by Jason Boice in 2017 and
is now widely supported (ls, ripgrep, fd, bat, hyperfine,
\maya{}). When a user sets it, they mean it: drop all colour.
\end{aside}

\section{Detecting terminal capabilities}

Not every terminal supports truecolour, or italics, or OSC 8.
\maya{} detects capabilities at startup:

\begin{itemize}[leftmargin=*]
  \item \code{COLORTERM=truecolor} or \code{=24bit}: emit
    full RGB.
  \item \code{TERM} starts with one of a known set (xterm-256,
    \emph{etc.}): use 256-colour palette.
  \item Otherwise: 8/16-colour palette.
\end{itemize}

The detection happens once in
\file{maya/include/maya/terminal/terminal.hpp} and the result
threads through the pipeline as a \code{TerminalCaps} struct.

When the renderer emits a colour, it consults caps:

\begin{itemize}[leftmargin=*]
  \item RGB asked but only 256 supported: quantise.
  \item RGB asked but only 16 supported: nearest-neighbour to
    16-colour palette.
  \item Default-fg asked: emit nothing (terminal's choice).
\end{itemize}

\section{Building your own theme}

If you want to ship a theme with your app:

\begin{lstlisting}
namespace mythemes {
    inline constexpr Theme Sunset = {
        .palette = {
            [Role::Foreground] = Color::rgb(245, 230, 200),
            [Role::Background] = Color::rgb(40, 25, 35),
            [Role::Primary]    = Color::rgb(220, 140, 80),
            [Role::Accent]     = Color::rgb(255, 180, 120),
            [Role::Warning]    = Color::rgb(255, 200, 100),
            [Role::Danger]     = Color::rgb(220, 80, 80),
            [Role::Success]    = Color::rgb(140, 200, 130),
            [Role::Muted]      = Color::rgb(160, 140, 130),
            [Role::Border]     = Color::rgb(120, 90, 80),
        },
        .name = "sunset",
        .is_dark = true,
    };
}
\end{lstlisting}

Then set it:

\begin{lstlisting}
maya::set_theme(mythemes::Sunset);
\end{lstlisting}

Every styled element using \code{Role::Primary} will use
RGB(220, 140, 80). No code change required.

\section{Style preset gallery}

Reusable style preset patterns:

\begin{lstlisting}
constexpr auto H1 = Bold | Underline | Fg<255, 240, 200>;
constexpr auto H2 = Bold | Fg<200, 220, 240>;
constexpr auto Body = Style{};
constexpr auto Code = Italic | Fg<200, 200, 160>;
constexpr auto Subtle = Dim | Fg<140, 140, 150>;
constexpr auto Error = Bold | Fg<255, 100, 100>;
constexpr auto Success = Bold | Fg<100, 200, 130>;
\end{lstlisting}

Each is a \code{StyTag} at the type level. Reusing them keeps
your app's appearance consistent and lets you tweak in one
place.

\section{Pitfalls}

\begin{warning}
\textbf{Hard-coding colours instead of using roles.} If you write
\code{Fg<255, 200, 100>} everywhere, you cannot theme the app.
Use \code{role::Warning} (or the equivalent helper) so a theme
swap actually changes the appearance.
\end{warning}

\begin{warning}
\textbf{Setting BG to black on a non-default terminal.} Some
users have dark grey or off-black backgrounds. Setting your bg
to RGB(0,0,0) creates a visual line. Use \code{Color::default\_()}
unless you really need a specific bg.
\end{warning}

\begin{warning}
\textbf{Forgetting NO\_COLOR.} If your app relies on colour to
convey meaning (red = bad, green = good), users with
\code{NO\_COLOR=1} will lose the cue. Pair colour with text
or attributes (``[ERR]'', Underline).
\end{warning}

\begin{warning}
\textbf{Truecolour quantisation surprises.} A theme that uses
adjacent RGB values may collapse onto the same Ansi256 cell on
old terminals. Test your theme on a 256-colour terminal before
shipping.
\end{warning}

\section*{Workshop}

\begin{enumerate}[leftmargin=*]
  \item Define a \code{Theme} of your own with three named roles.
    Use it with \code{maya::set\_theme}. Toggle between yours
    and the default at runtime with a hotkey.
  \item Inspect what your terminal advertises: print
    \code{getenv("TERM")} and \code{getenv("COLORTERM")} in a
    counter app.
  \item Run with \code{NO\_COLOR=1} set and confirm the app
    still conveys distinctions through attributes.
\end{enumerate}

\begin{recap}
The style system has three layers (theme, inheritance, local) and
merges with a defined precedence (attribute-union,
colour-override). The palette decouples roles from colours;
themes swap palettes wholesale. Truecolour quantisation,
NO\_COLOR, and capability detection happen at the edges so your
code can think in roles, not raw RGB.
\end{recap}

\chapter{Builders and the runtime tree}
\label{ch:builders}

\begin{aside}
The DSL is the path most users take to construct UI. There is a
second path that does not appear in the chapters so far: the
\emph{builders}. A builder is a small fluent API that produces an
\code{Element} programmatically, without the compile-time
machinery. Builders are what widgets use internally, what
runtime-driven UIs use, and what you reach for when the DSL's
type-state machine is in the way.
\end{aside}

\section{Why builders exist}

The DSL is excellent when the structure is known at compile
time. When the structure depends on data, the DSL fights you:

\begin{itemize}[leftmargin=*]
  \item A list of items whose length depends on a vector.
  \item A box whose border style is chosen by a config setting.
  \item A widget whose internal layout depends on its current
    state.
\end{itemize}

You can do these things in the DSL with \code{when}, the runtime
spellings of pipes, and ranges --- but it gets noisy. A builder
sidesteps the compile-time pipeline entirely:

\begin{lstlisting}
auto e = box()
    .direction(Direction::Column)
    .border(BorderStyle::Single)
    .padding(1)
    .add(text("Title").bold())
    .add(text("Body line"))
    .build();
\end{lstlisting}

Each method returns a reference to the builder, so calls chain.
\code{build()} produces the final \code{Element}.

\subsection*{Builders are not anti-DSL}

The two work together. A widget can have its own builder
internally and the DSL call site is unchanged:

\begin{lstlisting}
auto bar = widget(StatusBar()
    .mode(\"NORMAL\")
    .file(\"~/notes.md\")
    .cursor(line, col));
\end{lstlisting}

Here \code{StatusBar} is built with a builder, but the result is
hosted in the DSL by \code{widget(\ldots)}.

\section{The BoxBuilder}

\file{maya/include/maya/element/builder.hpp} defines
\code{BoxBuilder}:

\begin{lstlisting}
class BoxBuilder {
public:
    BoxBuilder& direction(Direction);
    BoxBuilder& justify(Justify);
    BoxBuilder& align(Align);
    BoxBuilder& wrap(Wrap);
    BoxBuilder& gap(int);
    BoxBuilder& padding(Padding);
    BoxBuilder& padding(int);             // all sides
    BoxBuilder& padding(int v, int h);    // vert + horiz
    BoxBuilder& border(BorderStyle);
    BoxBuilder& border_color(Color);
    BoxBuilder& width(int);
    BoxBuilder& height(int);
    BoxBuilder& grow(int);
    BoxBuilder& shrink(int);
    BoxBuilder& style(Style);
    BoxBuilder& add(Element child);
    template <class... Cs> BoxBuilder& add_all(Cs&&... cs);

    Element build() &&;
};

BoxBuilder box();
\end{lstlisting}

Each setter returns a reference to itself, so calls compose.
\code{build()} is rvalue-qualified --- you can only call it on a
temporary (i.e. at the end of the chain), which discourages
accidentally reusing a half-built builder.

\subsection*{Why ampersand-ampersand on build()?}

The double-ampersand (\code{\&\&}) qualifier means ``this method
can only be called on an rvalue''. \code{build()} consumes the
builder. If you wrote:

\begin{lstlisting}
BoxBuilder b;
b.direction(Direction::Row);
auto e = b.build();   // error: build() is not callable on lvalues
\end{lstlisting}

\ldots{} the compiler would refuse. The fix:

\begin{lstlisting}
auto e = std::move(b).build();
\end{lstlisting}

This is a deliberate design: a builder is a one-shot value. The
\code{\&\&} qualifier enforces that.

\begin{funfact}
Ref-qualifiers (\code{\&} and \code{\&\&} on methods) were added
in C++11 specifically for builder patterns and string-builder
optimisations. \code{std::string::operator+=} can be optimised
when it knows it is being called on a temporary; the same idea
applies to fluent builders.
\end{funfact}

\section{The TextBuilder}

For text leaves:

\begin{lstlisting}
class TextBuilder {
public:
    TextBuilder& style(Style);
    TextBuilder& fg(Color);
    TextBuilder& bg(Color);
    TextBuilder& bold();
    TextBuilder& italic();
    TextBuilder& underline();
    TextBuilder& dim();
    TextBuilder& hyperlink(std::string);

    Element build() &&;
    operator Element() && { return std::move(*this).build(); }
};

TextBuilder text(std::string s);
TextBuilder text(std::string_view sv);
\end{lstlisting}

Most users never call \code{build()} on a \code{TextBuilder}
directly; the conversion operator does it implicitly when the
builder is consumed by another builder's \code{add()}:

\begin{lstlisting}
auto e = box()
    .add(text(\"hi\").bold())     // TextBuilder -> Element
    .build();
\end{lstlisting}

The \code{operator Element()} is rvalue-qualified for the same
reason as \code{build()}: it consumes the builder.

\section{ListBuilder for dynamic children}

When you have a runtime list of items:

\begin{lstlisting}
auto e = list()
    .for_each(items, [](auto& item) {
        return text(item.label).bold();
    })
    .with_separator(text(\" | \").dim())
    .build();
\end{lstlisting}

\code{ListBuilder::for\_each} takes a range and a function from
item to element. It eats the runtime cost (one allocation for
the children vector) so the DSL does not have to.

\subsection*{Implementation sketch}

\begin{lstlisting}
template <class R, class F>
ListBuilder& ListBuilder::for_each(R&& range, F&& f) {
    for (auto& item : range) {
        children_.push_back(Element{std::forward<F>(f)(item)});
    }
    return *this;
}
\end{lstlisting}

Generic over the range type and the projection. The result is
sequential children in the order of iteration.

\section{Widget builders}

Every widget in \file{maya/include/maya/widget/} is a builder
whose \code{build()} returns an \code{Element}. The pattern:

\begin{lstlisting}
class Progress {
public:
    Progress& value(double v);                 // 0..1
    Progress& total(int t);
    Progress& current(int c);
    Progress& label(std::string);
    Progress& width(int);
    Progress& style(Style);

    Element build() &&;
};
\end{lstlisting}

The widget knows how to render itself; the caller picks the
parameters. The DSL hosts the result:

\begin{lstlisting}
auto bar = widget(Progress{}
    .total(100)
    .current(progress_now)
    .label(\"Loading\\ldots\")
    .width(40));
\end{lstlisting}

\subsection*{Why fluent, not designated init?}

You can also write widgets with designated initialisers:

\begin{lstlisting}
auto bar = widget(Progress{
    .total = 100,
    .current = progress_now,
    .label = \"Loading\\ldots\",
    .width = 40,
});
\end{lstlisting}

Both work. The fluent form composes nicely with conditional
chains:

\begin{lstlisting}
auto bar = Progress{}.total(100).current(p);
if (show_label) bar.label(\"Loading\");
auto e = widget(std::move(bar));
\end{lstlisting}

The designated-init form requires you to commit to all fields at
once. Pick by ergonomics.

\section{Builders for tree construction}

A common pattern: build a tree by recursion.

\begin{lstlisting}
Element render_node(const Node& n) {
    auto b = box()
        .direction(Direction::Column)
        .add(text(n.label).bold());
    for (auto& child : n.children) {
        b.add(render_node(child));
    }
    return std::move(b).build();
}
\end{lstlisting}

Each call to \code{render\_node} produces an \code{Element}. The
recursion handles arbitrarily deep trees. No special API; just
\code{add} in a loop.

\subsection*{Composing with the DSL}

The output of a builder is an \code{Element}, so it composes with
the DSL:

\begin{lstlisting}
auto root = v(
    text(\"Header\") | Bold,
    render_tree(my_tree),       // a builder-produced Element
    text(\"Footer\") | Dim
);
\end{lstlisting}

You can switch between styles as the structure demands.

\section{Tradeoffs}

Builders are not free:

\begin{itemize}[leftmargin=*]
  \item \textbf{Allocation.} Each builder allocates the
    \code{vector<Element>} for its children eagerly. The DSL's
    compile-time path avoids this for fixed-size containers.
  \item \textbf{Type erasure.} The builder takes
    \code{Element} as the child type; the DSL's pipe overloads
    preserve more type information. If you need
    \code{constexpr}-correctness, the DSL is the path.
  \item \textbf{Error messages.} A typo in a setter
    (\code{borderr(\\ldots)}) is a normal C++ error: ``no member
    function''. The DSL's concept-based errors are sometimes
    clearer for structural mistakes.
\end{itemize}

The performance difference at frame time is usually negligible
(one vs zero allocations on a tree with hundreds of nodes is
microseconds). Pick by ergonomics, not by performance.

\begin{aside}
Internally, the DSL's \code{build()} produces the same
\code{Element} that a builder would. The DSL is a thin
compile-time wrapper that happens to skip some allocations; the
runtime model is the same.
\end{aside}

\section{Pitfalls}

\begin{warning}
\textbf{Forgetting to consume the builder.} If you do not call
\code{build()} (or pass it to \code{add()} / \code{widget()}),
the builder goes out of scope and its data is dropped.
\code{-Wunused-value} will warn if your builder is
\code{[[nodiscard]]}.
\end{warning}

\begin{warning}
\textbf{Reusing a moved-from builder.} After
\code{std::move(b).build()}, \code{b} is moved-from. Touching it
is technically allowed (the destructor runs) but reading it is
nonsense. The rvalue qualifier prevents this in practice.
\end{warning}

\begin{warning}
\textbf{Mixing units across setters.} \code{padding(1)} means
``all sides 1 cell''. \code{padding(1, 2)} means ``vert 1, horiz
2''. \code{padding(Padding\{.top=1, \\ldots\})} is the explicit
form. Read your code aloud; the unit should be obvious.
\end{warning}

\begin{warning}
\textbf{Reading layout from a half-built tree.} The builder does
not run layout; the rects are uninitialised until paint. If a
helper inspects \code{rect.w} on a builder-produced element
before painting, it will read zero.
\end{warning}

\section*{Workshop}

\begin{enumerate}[leftmargin=*]
  \item Write a function \code{Element render\_table(const
    std::vector<std::vector<std::string>>\& rows)} that uses
    builders to assemble a table.
  \item Read
    \file{maya/include/maya/element/builder.hpp} and identify
    every setter on \code{BoxBuilder}. Map each to a property
    on \code{FlexStyle}.
  \item Convert a small DSL expression into the equivalent
    builder code, then back. Note which form is more readable
    for each shape.
\end{enumerate}

\begin{recap}
Builders are the runtime spelling of the DSL. They produce the
same \code{Element} that the DSL produces; the difference is
where the work happens. Use builders when structure depends on
data; use the DSL when structure is fixed. Both compose at the
\code{Element} boundary.
\end{recap}

\chapter{Text, graphemes, and width}
\label{ch:text}

\begin{aside}
A terminal puts characters in cells. Two facts make this harder
than it sounds: not every Unicode codepoint is one character, and
not every character is one cell wide. \maya{} treats text as a
grapheme stream and resolves visual width at the boundary between
the element tree and the canvas. This chapter explains the model.
\end{aside}

\section{Code units, codepoints, graphemes}

Three terms for three different things:

\begin{itemize}[leftmargin=*]
  \item A \textbf{code unit} is the smallest piece of an encoded
    string. For UTF-8, a code unit is a byte. For UTF-16, two
    bytes. \code{std::string::size()} returns code units.
  \item A \textbf{codepoint} is a Unicode scalar value
    (U+0000--U+10FFFF, minus the surrogate range). One codepoint
    may take one to four UTF-8 bytes.
  \item A \textbf{grapheme cluster} is what a human would call ``a
    character'' --- a base character plus any combining marks,
    joined emoji, and similar. One grapheme may contain several
    codepoints.
\end{itemize}

The terminal renders graphemes, but it does not actually know what
a grapheme is. It writes codepoints sequentially and lets the
font shaper combine them. If you put two codepoints in two cells
when they should have combined into one cell, the cursor and the
rest of the row drift by one.

\subsection*{An example}

The visible character \emph{é} can be encoded two ways:

\begin{itemize}[leftmargin=*]
  \item U+00E9 (precomposed): one codepoint, one cell.
  \item U+0065 U+0301 (e + combining acute): two codepoints, one
    cell.
\end{itemize}

\maya{} normalises text to a grapheme stream when measuring width
and emitting cells. \code{std::string} is opaque to all this; the
text module wraps it.

\begin{funfact}
Unicode has at least three different forms of \emph{é} when you
count the precomposed/decomposed variants. Normalisation forms
(NFC, NFD, NFKC, NFKD) reduce them to canonical representations.
\maya{} does not normalise inputs; it preserves the user's
bytes and walks them as graphemes.
\end{funfact}

\section{East Asian Width}

A second wrinkle: some codepoints take two cells. Chinese,
Japanese, Korean characters typically do; emoji do; many symbols
do. Unicode's \emph{East Asian Width} property classifies each
codepoint as Narrow, Wide, Ambiguous, Fullwidth, Halfwidth, or
Neutral.

\maya{} ships with a precomputed table in
\file{maya/include/maya/text/unicode\_width\_table.hpp} mapping
codepoint to cell width:

\begin{lstlisting}
int width_of(char32_t cp);
// 0 for combining marks, zero-width joiners, control chars
// 1 for narrow / neutral / halfwidth
// 2 for wide / fullwidth
\end{lstlisting}

The table is generated from the Unicode 15.x EastAsianWidth
data. Lookup is a binary search over the table; sub-microsecond.

\subsection*{Why the table is precomputed}

Computing width at runtime would mean walking the Unicode
character database, which is large. \maya{} ships the table as a
constexpr array; the lookup is fast and the binary stays
self-contained.

The trade-off: when Unicode releases a new version with new
codepoints, \maya{} needs a regenerate. The generator script
lives at \file{maya/tools/gen\_unicode\_width.py}.

\section{The grapheme walker}

\file{maya/include/maya/text/unicode\_width.hpp} exposes:

\begin{lstlisting}
class GraphemeWalker {
public:
    explicit GraphemeWalker(std::string_view text);

    struct Grapheme {
        std::string_view text;   // bytes for this grapheme
        int width;               // 0, 1, or 2
    };

    std::optional<Grapheme> next();
    void reset();
};
\end{lstlisting}

The walker steps through the input grapheme by grapheme. For
each grapheme it returns the byte span and the cell width.

\subsection*{Implementation sketch}

The walker is a small state machine driven by Unicode's grapheme
cluster boundaries:

\begin{itemize}[leftmargin=*]
  \item Skip combining marks (they belong to the previous
    grapheme).
  \item Handle regional indicators (flag emoji are pairs).
  \item Handle ZWJ sequences (joiner-joined emoji are one
    grapheme).
  \item Handle Hangul syllables (\emph{L V T} sequences are one).
\end{itemize}

The result is a stream of cleanly-bounded graphemes that the
canvas can lay out cell by cell.

\section{Measuring width}

The most common use:

\begin{lstlisting}
int display_width(std::string_view text) {
    GraphemeWalker w{text};
    int total = 0;
    while (auto g = w.next()) total += g->width;
    return total;
}
\end{lstlisting}

The layout phase calls this for each text node to get its
intrinsic width. Cached on the text node so successive frames do
not repeat the walk.

\subsection*{Ambiguous-width chars}

Some characters have ``ambiguous'' width: 1 cell in Western
terminals, 2 cells in CJK terminals. \maya{} defaults to 1
(Western), with a config switch for users on CJK terminals:

\begin{lstlisting}
maya::set_ambiguous_width(2);  // CJK terminals
\end{lstlisting}

The cached widths invalidate when the setting changes.

\section{Truncation and wrapping}

When a text node must fit a fixed width, \maya{} provides:

\begin{itemize}[leftmargin=*]
  \item \code{TextWrap::Truncate}: drop characters from the end,
    optionally with an ellipsis.
  \item \code{TextWrap::Wrap}: spill onto the next line at a
    word boundary.
  \item \code{TextWrap::WrapAnywhere}: spill at any grapheme
    boundary, including mid-word.
\end{itemize}

\subsection*{Truncation with ellipsis}

\begin{lstlisting}
std::string truncate(std::string_view text, int max_width) {
    GraphemeWalker w{text};
    std::string out;
    int used = 0;
    while (auto g = w.next()) {
        if (used + g->width > max_width - 1) {
            out += \"...\";
            return out;
        }
        out += g->text;
        used += g->width;
    }
    return out;
}
\end{lstlisting}

Reserve one cell for the ellipsis itself (well, three for the
ASCII ``...''; \maya{} can also use U+2026 for one cell).

\subsection*{Word wrapping}

Word wrap is harder. The algorithm:

\begin{enumerate}[leftmargin=*]
  \item Walk graphemes accumulating into a current line.
  \item When the line reaches \code{max\_width}, scan backward
    for the last whitespace and break there. The trailing word
    becomes the start of the next line.
  \item If a single word exceeds \code{max\_width}, break it
    at a grapheme boundary (degrade to WrapAnywhere).
\end{enumerate}

The result is a vector of lines. The canvas paints each line at
its computed $y$.

\section{Special characters}

Several Unicode categories deserve special handling.

\subsection*{Tabs}

A tab character (U+0009) advances the cursor to the next tab
stop, typically every 8 cells. \maya{} expands tabs to spaces
before measuring; the tab stops are configurable.

\subsection*{Newlines}

\code{\\n} (U+000A) starts a new line. The canvas treats it as a
forced break: the current row ends; the next grapheme goes to
column 0 of the next row.

\subsection*{Control characters}

Most control characters (U+0001--U+0008, U+000B--U+001F) are
zero-width and invisible. \maya{} either drops them or renders a
caret notation (\code{\textasciicircum A} for U+0001) depending
on the \code{ControlPolicy} config.

\subsection*{Variation selectors}

U+FE0E (text presentation) and U+FE0F (emoji presentation) modify
the preceding codepoint. \maya{} honours them: a base character
followed by U+FE0F becomes the emoji form.

\subsection*{Right-to-left text}

\maya{} does not perform full BiDi (bidirectional) reordering.
RTL text (Arabic, Hebrew) is rendered in logical order; the
terminal handles visual order based on its own BiDi
implementation. This means RTL text inside boxes may align
unexpectedly. Full BiDi is a planned feature; see the issue
tracker.

\begin{warning}
Mixing LTR and RTL in the same text node is a known weak point.
If you need this, render the RTL portion as its own text node
and accept the visual quirks.
\end{warning}

\section{Performance}

Width measurement is fast:

\begin{itemize}[leftmargin=*]
  \item Grapheme walking: roughly 10ns per grapheme.
  \item Width lookup: roughly 20ns per codepoint (binary search
    in the precomputed table).
  \item Per-line measurement: microseconds for 80-character
    lines.
\end{itemize}

For an 80x24 screen of text, total measurement is well under a
millisecond. \maya{} caches per-text-node width so re-measurement
across frames is free unless the text changes.

\subsection*{Cache invalidation}

The cache key is the text content plus the ambiguous-width
setting. If either changes, the cache is dropped.

\section{Pitfalls}

\begin{warning}
\textbf{Using std::string::size() for visual width.} It is the
byte count, not the cell count. ``hello'' is 5 bytes and 5
cells, but ``\verb|privet|'' (transliterating six Cyrillic
letters) is 12 bytes and 6 cells, and a two-character CJK string
like ``ni hao'' encoded in Chinese is 6 bytes and 4 cells. Use
\code{display\_width} for any
visual sizing.
\end{warning}

\begin{warning}
\textbf{Calling display\_width every frame.} For text in a tight
inner loop, the cost adds up. Pre-compute and cache, or trust
the layout phase to do it once.
\end{warning}

\begin{warning}
\textbf{Mixing widths in a single text node.} If a text node
contains both narrow and wide characters, the layout phase
handles it. But if you try to truncate at a fixed
character count instead of cell count, you will get visual
misalignment. Always truncate by cell count.
\end{warning}

\begin{warning}
\textbf{Trailing combining marks.} A grapheme can extend past
the visual edge of a fixed-width region. The walker handles
this, but if you slice raw bytes you may split a grapheme. Always
slice at grapheme boundaries.
\end{warning}

\section*{Workshop}

\begin{enumerate}[leftmargin=*]
  \item Write a function that takes a \code{std::string\_view}
    and returns its display width using \code{GraphemeWalker}.
    Test it on plain ASCII, on emoji, and on CJK text.
  \item Write a truncator that respects grapheme boundaries.
    Test it on a string ending in a multi-codepoint emoji.
  \item Read
    \file{maya/include/maya/text/unicode\_width\_table.hpp}.
    Note the size of the table and the structure of its
    entries.
\end{enumerate}

\begin{recap}
Text is a stream of bytes, codepoints, and graphemes; the cell
grid wants graphemes. \maya{} ships a precomputed width table
and a grapheme walker that converts. Truncation and wrapping
respect grapheme boundaries. Tabs, newlines, control characters,
variation selectors, and ambiguous-width chars have explicit
policies.
\end{recap}

\chapter{The platform abstraction}
\label{ch:platform}

\begin{aside}
A TUI library that supports POSIX, macOS, and Windows lives or
dies by its platform abstraction. The OS differences are real:
POSIX uses \code{termios} and \code{epoll}; macOS uses
\code{kqueue}; Windows uses the Console API and
\code{WaitForMultipleObjects}. \maya{} hides all of this behind
a small set of platform classes that satisfy the same concept.
This chapter walks the shape of that abstraction.
\end{aside}

\section{The shape of the problem}

To run a TUI on three platforms, \maya{} must do five things
portably:

\begin{enumerate}[leftmargin=*]
  \item Open the terminal and put it in raw mode.
  \item Read keystrokes and mouse events as they arrive.
  \item Receive signals (window resize, interrupt, shutdown).
  \item Wake the main loop from a worker thread.
  \item Restore the terminal when the program exits.
\end{enumerate}

POSIX, macOS, and Windows each spell these five things
differently. \maya{} pins down the spelling with a concept and
provides three implementations.

\begin{funfact}
The first portable terminal library was \code{termcap} (1978),
which abstracted differences between physical terminal models
(VT100 vs ADM-3a). Modern terminal emulators ironed out most of
that diversity. The remaining portability story is between
\emph{operating systems}, not terminals --- a different axis.
\end{funfact}

\section{The Terminal concept}

\file{maya/include/maya/platform/concepts.hpp} declares:

\begin{lstlisting}
template <class T>
concept Terminal = requires (T& t, std::span<std::byte> buf) {
    { t.enter_raw_mode() }   -> std::same_as<Expected<void, IoError>>;
    { t.leave_raw_mode() }   -> std::same_as<void>;
    { t.read(buf) }          -> std::same_as<Expected<size_t, IoError>>;
    { t.write(buf) }         -> std::same_as<Expected<size_t, IoError>>;
    { t.size() }             -> std::same_as<Expected<TermSize, IoError>>;
    { t.caps() }             -> std::same_as<const TerminalCaps&>;
};
\end{lstlisting}

Six methods. Two for setup/teardown, two for I/O, one for size
queries, one for capabilities. Each platform implementation
satisfies the concept.

\subsection*{TermSize and TerminalCaps}

\begin{lstlisting}
struct TermSize {
    int rows, cols;
    int pixel_width = 0, pixel_height = 0;  // optional
};

struct TerminalCaps {
    bool truecolor = false;
    bool osc8_hyperlinks = false;
    bool kitty_keyboard = false;
    bool sync_update = false;
    int max_colors = 16;
};
\end{lstlisting}

\code{TermSize} comes from \code{ioctl(TIOCGWINSZ)} on POSIX,
\code{GetConsoleScreenBufferInfo} on Win32. \code{TerminalCaps}
is detected at startup by sniffing environment variables and
sometimes issuing terminal queries.

\section{The POSIX backend}

\file{maya/include/maya/platform/posix/} defines the POSIX
\code{Terminal}. The core is wrapping \code{termios}.

\subsection*{Raw mode, in detail}

\begin{lstlisting}
Expected<void, IoError> PosixTerminal::enter_raw_mode() {
    if (tcgetattr(fd_, &saved_) != 0)
        return std::unexpect, IoError{from_errno(errno), "tcgetattr"};

    termios raw = saved_;
    raw.c_iflag &= ~(IGNBRK | BRKINT | PARMRK | ISTRIP
                   | INLCR  | IGNCR  | ICRNL  | IXON);
    raw.c_oflag &= ~OPOST;
    raw.c_lflag &= ~(ECHO | ECHONL | ICANON | ISIG | IEXTEN);
    raw.c_cflag &= ~(CSIZE | PARENB);
    raw.c_cflag |= CS8;
    raw.c_cc[VMIN]  = 1;
    raw.c_cc[VTIME] = 0;

    if (tcsetattr(fd_, TCSAFLUSH, &raw) != 0)
        return std::unexpect, IoError{from_errno(errno), "tcsetattr"};

    raw_active_ = true;
    return {};
}
\end{lstlisting}

Each flag mask disables one piece of cooked behaviour. The
incantations come from W. Richard Stevens' \emph{Advanced
Programming in the UNIX Environment}; \maya{} just copies the
recipe.

\subsection*{Event source: epoll}

The POSIX backend uses \code{epoll} on Linux:

\begin{lstlisting}
PosixEventSource::PosixEventSource() {
    epfd_ = epoll_create1(EPOLL_CLOEXEC);
    epoll_ctl(epfd_, EPOLL_CTL_ADD, stdin_fd_,
              &epoll_event{.events = EPOLLIN, .data = {.fd = stdin_fd_}});
    epoll_ctl(epfd_, EPOLL_CTL_ADD, signal_fd_,
              &epoll_event{.events = EPOLLIN, .data = {.fd = signal_fd_}});
    epoll_ctl(epfd_, EPOLL_CTL_ADD, wake_fd_,
              &epoll_event{.events = EPOLLIN, .data = {.fd = wake_fd_}});
}

Event PosixEventSource::wait(Duration timeout) {
    epoll_event ev;
    int n = epoll_wait(epfd_, &ev, 1, to_ms(timeout));
    if (n == 0) return Event::Timeout{};
    if (ev.data.fd == stdin_fd_)   return drain_stdin();
    if (ev.data.fd == signal_fd_)  return drain_signal_fd();
    if (ev.data.fd == wake_fd_)    return Event::Wakeup{};
    return Event::Unknown{};
}
\end{lstlisting}

Three descriptors wait at once: stdin, signalfd, and a self-pipe
for cross-thread wakeups. Whichever fires first wins.

\subsection*{Signal handling}

POSIX signals are delivered out-of-band. \maya{} uses
\code{signalfd} on Linux to turn signals into file-descriptor
events:

\begin{lstlisting}
sigset_t mask;
sigemptyset(&mask);
sigaddset(&mask, SIGINT);
sigaddset(&mask, SIGTERM);
sigaddset(&mask, SIGWINCH);
sigprocmask(SIG_BLOCK, &mask, nullptr);
int sfd = signalfd(-1, &mask, SFD_NONBLOCK | SFD_CLOEXEC);
\end{lstlisting}

The signals are blocked from default delivery and routed
through \code{sfd}. When \code{epoll\_wait} reports activity on
\code{sfd}, a \code{read} drains the pending signals.

On BSD/macOS, \code{kqueue} has native \code{EVFILT\_SIGNAL}
support that does the same thing without the file-descriptor
indirection.

\subsection*{Wake fd}

\code{wake\_fd\_} is a self-pipe used by background threads to
wake the main loop. A worker thread that finished a task writes
one byte to \code{wake\_fd\_}; the main loop's epoll fires; the
loop drains the pipe and dispatches the queued message.

\begin{aside}
The self-pipe trick is one of the oldest patterns in Unix
network programming. It works because reading from an fd is
\emph{the} primitive the OS gives you for ``wake up when
something happens''.
\end{aside}

\section{macOS backend}

\file{maya/include/maya/platform/macos/} reuses much of the
POSIX backend (the \code{termios} flags are the same) but
overrides the event source.

macOS uses \code{kqueue} natively. The event source registers
stdin, signal filters, and a Mach port for cross-thread wakeups.

There is a quirk: macOS's \code{kqueue} delivers
\code{EVFILT\_SIGNAL} events with edge semantics, so the count
of pending signals must be tracked separately. \maya{}
maintains a small atomic counter for each watched signal.

\section{Win32 backend}

Windows is the outlier.

\subsection*{The Console API}

There is no \code{termios}. Raw mode is set via
\code{SetConsoleMode}:

\begin{lstlisting}
DWORD mode;
GetConsoleMode(hStdin, &mode);
mode |= ENABLE_VIRTUAL_TERMINAL_INPUT;
mode &= ~(ENABLE_PROCESSED_INPUT | ENABLE_LINE_INPUT | ENABLE_ECHO_INPUT);
SetConsoleMode(hStdin, mode);

GetConsoleMode(hStdout, &mode);
mode |= ENABLE_VIRTUAL_TERMINAL_PROCESSING;
mode |= DISABLE_NEWLINE_AUTO_RETURN;
SetConsoleMode(hStdout, mode);
\end{lstlisting}

The \code{ENABLE\_VIRTUAL\_TERMINAL\_INPUT} flag is what makes
Windows 10+ consoles emit VT escape sequences for keystrokes
instead of \code{INPUT\_RECORD}s. \maya{} relies on this; older
Windows versions are not supported.

\subsection*{Event source: WaitForMultipleObjects}

Win32's I/O completion ports are very different from epoll. The
\maya{} backend uses a hybrid: \code{WaitForMultipleObjects} on
the console input handle plus a manual-reset event for
cross-thread wakeup:

\begin{lstlisting}
HANDLE handles[] = { hStdin_, hWake_ };
DWORD r = WaitForMultipleObjects(2, handles, FALSE, timeout_ms);
if (r == WAIT_OBJECT_0) return read_stdin();
if (r == WAIT_OBJECT_0 + 1) {
    ResetEvent(hWake_);
    return Event::Wakeup{};
}
\end{lstlisting}

Simpler than epoll, less scalable, but the console input is the
only ``socket'' we need.

\subsection*{Signal handling}

Windows has fewer signals than POSIX. The main ones \maya{}
cares about:

\begin{itemize}[leftmargin=*]
  \item \code{CTRL\_C\_EVENT} and \code{CTRL\_BREAK\_EVENT},
    delivered via a console control handler.
  \item Window resize events, delivered as
    \code{WINDOW\_BUFFER\_SIZE\_EVENT} on the console input
    handle.
\end{itemize}

The control handler runs on a separate thread, so it writes to
the wake event and the main loop picks up the message.

\section{The select header}

\file{maya/include/maya/platform/select.hpp} picks the right
backend at compile time:

\begin{lstlisting}
#if defined(_WIN32)
using Terminal      = Win32Terminal;
using EventSource   = Win32EventSource;
#elif defined(__APPLE__)
using Terminal      = MacTerminal;
using EventSource   = MacEventSource;
#else
using Terminal      = PosixTerminal;
using EventSource   = PosixEventSource;
#endif
\end{lstlisting}

The rest of \maya{} writes \code{maya::platform::Terminal} and
gets the right type.

\subsection*{Why not virtual functions?}

We considered defining \code{ITerminal} with virtual methods and
selecting at runtime. We rejected it because:

\begin{itemize}[leftmargin=*]
  \item The choice is fixed at compile time anyway --- you do
    not target multiple platforms in one binary.
  \item Virtual dispatch on every read/write adds overhead.
  \item Concepts give better error messages when a method is
    missing.
\end{itemize}

\section{Cross-platform pitfalls}

\begin{warning}
\textbf{Line endings.} POSIX terminals send LF (U+000A) for
Enter; Windows sends CRLF (U+000D U+000A). \maya{} normalises in
the input parser, but if you bypass the parser you will see the
raw bytes.
\end{warning}

\begin{warning}
\textbf{Path separators in errors.} POSIX uses forward slashes,
Win32 uses backslashes. The \code{IoError::context} string
preserves whatever the OS gave us. If you display paths to users,
normalise.
\end{warning}

\begin{warning}
\textbf{Console code pages.} Older Windows consoles default to
code page 437 or 850, not UTF-8. \maya{} calls
\code{SetConsoleOutputCP(CP\_UTF8)} on startup. If you bypass
this, your Unicode output will be mojibake.
\end{warning}

\begin{warning}
\textbf{ConPTY vs raw console.} Windows Terminal hosts a
ConPTY (pseudo-terminal) layer that translates between VT
escapes and the legacy console. Most behaviour is identical to
Linux/macOS, but a handful of escapes (cursor shape, mouse
modes) behave subtly differently. If you see a glitch only on
Windows Terminal, check the escape sequence semantics.
\end{warning}

\section*{Workshop}

\begin{enumerate}[leftmargin=*]
  \item Open \file{maya/include/maya/platform/posix/terminal.hpp}
    and find the \code{enter\_raw\_mode} function. Match each
    flag mask to a line in \emph{stty -a} output.
  \item Read \file{maya/include/maya/platform/select.hpp} and
    confirm which backend is selected on your machine.
  \item On Linux, run \code{strace -e epoll\_wait,read,write}
    on a \maya{} app for one second. Note that almost all time
    is spent blocked in \code{epoll\_wait}.
\end{enumerate}

\begin{recap}
\maya{}'s platform layer is a concept (\code{Terminal}) and
three implementations. The POSIX backend uses \code{termios} and
\code{epoll}; macOS uses \code{kqueue}; Win32 uses
\code{SetConsoleMode} and \code{WaitForMultipleObjects}. The
header \code{platform/select.hpp} picks the right backend by
\code{\#ifdef}, so the rest of \maya{} sees one type.
\end{recap}

\chapter{Focus and keyboard navigation}
\label{ch:focus}

\begin{aside}
A UI with more than one focusable region needs a way to say
``input goes \emph{here} now''. \maya{} models focus as state in
your \code{Model}, not as a global property of the runtime.
That sounds restrictive; it is liberating. Focus becomes data,
flows through \code{update} like any other state, and is
testable without a runtime.
\end{aside}

\section{Why focus is not a global}

Most GUI frameworks store focus as a property of the widget
tree: a pointer to the ``focused widget'', set by mouse clicks
and Tab presses, read by event dispatchers. The pattern works
but it has costs:

\begin{itemize}[leftmargin=*]
  \item \textbf{Hidden state.} Focus is not visible in your
    model; it is owned by the framework. Diagnosing ``why did
    keystrokes go to the wrong place?'' means reading framework
    internals.
  \item \textbf{Time-travel breaks.} Restoring an old model does
    not restore the old focus; the framework's pointer survives.
  \item \textbf{Concurrency hazards.} Mouse and keyboard handlers
    might race on the focus pointer.
\end{itemize}

The Elm fix: focus is just a field in your \code{Model}. The
runtime never touches it; you decide who has focus by setting
the field in \code{update}.

\subsection*{The simplest form}

\begin{lstlisting}
struct Model {
    int focus;    // 0 = left, 1 = right
    LeftPanel  left;
    RightPanel right;
};
\end{lstlisting}

The \code{view} renders both panels but highlights only the
focused one. The keyboard handler routes input based on
\code{focus}. \code{Tab} flips \code{focus}.

That is the whole pattern. There is no \code{focus\_widget()}
call.

\section{The FocusManager helper}

For larger apps, naming focus targets by integers gets old.
\file{maya/include/maya/core/focus.hpp} ships a small helper:

\begin{lstlisting}
class FocusManager {
public:
    using Id = uint32_t;

    Id  register_target();
    void focus(Id);
    void focus_next();
    void focus_prev();
    bool is_focused(Id) const;
    Id   current() const;

private:
    std::vector<Id> order_;
    Id current_ = 0;
};
\end{lstlisting}

Each focusable region calls \code{register\_target()} at
construction to get a stable ID. \code{focus\_next} cycles
forward, \code{focus\_prev} cycles back.

The manager itself is plain data --- a vector and an integer.
You can store one in your \code{Model} and read it from
\code{view}.

\subsection*{Putting focus in the view}

\begin{lstlisting}
Element render_panel(const Panel& p, bool focused) {
    return box(
        text(p.label) | (focused ? Bold : Dim),
        text(p.body))
        | border<focused ? Heavy : Single>;
}
\end{lstlisting}

The same panel renders differently when focused. The decision is
the caller's; no global ``am I focused?'' query.

\subsection*{Keyboard routing}

\begin{lstlisting}
static std::pair<Model, Cmd<Msg>>
update(Model m, Msg msg) {
    return std::visit(overload{
        [&](KeyPress k) {
            if (k.key == Tab) {
                m.focus.focus_next();
                return std::pair{m, Cmd<Msg>{}};
            }
            // Route to the focused panel.
            if (m.focus.is_focused(m.left.id))
                return route_to_left(m, k);
            return route_to_right(m, k);
        },
        // ...
    }, msg);
}
\end{lstlisting}

The handler decides who gets the key based on
\code{focus.current()}. The focused panel's own \code{update}
processes it.

\section{Tab order}

\code{FocusManager} maintains the registration order. By default
\code{focus\_next} cycles through it. You can override:

\begin{lstlisting}
mgr.set_order({sidebar_id, main_id, search_id});
\end{lstlisting}

Or skip-disabled-elements:

\begin{lstlisting}
void focus_next_enabled(FocusManager& m, const Model& model) {
    do {
        m.focus_next();
    } while (!model.is_target_enabled(m.current())
             && m.current() != initial_id);
}
\end{lstlisting}

The pattern: focus order is a function of your model, not the
manager. The manager is a bookkeeping helper; the rules are
yours.

\section{Mouse focus}

A click can change focus. The hit test
(Chapter~\ref{ch:elements}) returns the deepest element under
the cursor; if that element has an associated focus ID, set it:

\begin{lstlisting}
[&](MouseEvent me) {
    if (me.kind == MouseKind::Press && me.button == 0) {
        if (auto id = hit_to_focus_id(m, me.x, me.y)) {
            m.focus.focus(*id);
        }
    }
    return std::pair{m, Cmd<Msg>{}};
}
\end{lstlisting}

The mapping from rect to focus ID lives in your model. The
runtime does not know which boxes are focusable; you do.

\section{Modal focus}

A modal (a dialog that captures all input) is just a focus stack:

\begin{lstlisting}
struct Model {
    FocusManager focus;
    std::vector<FocusManager> modal_stack;
};
\end{lstlisting}

When a modal opens, push a fresh \code{FocusManager}. While the
stack is non-empty, route keys to the top of the stack. When the
modal closes, pop.

The visible UI also follows the stack: only the top modal's
elements are interactive; everything below is dimmed.

\begin{idea}
Modal focus has been a recurring source of bugs in
retained-mode GUIs (``why doesn't Tab work in this dialog?'').
Modelling it as a stack of focus managers is a couple of
\code{push\_back} / \code{pop\_back} calls; the runtime stays
unaware.
\end{idea}

\section{Pitfalls}

\begin{warning}
\textbf{Focus on a non-rendered element.} If \code{view} hides
the focused element (visibility, conditional), keystrokes go
nowhere. Always reconcile after the model changes: if the
focused id is no longer rendered, advance focus.
\end{warning}

\begin{warning}
\textbf{Stable IDs.} If you regenerate focus IDs every frame,
focus moves randomly. \code{register\_target} should be called
once, in \code{init} or when the focusable region is created.
\end{warning}

\begin{warning}
\textbf{Capturing focus during scroll.} If a focused element
scrolls out of view, the user can no longer see what they are
typing into. Either keep the focused element on-screen by
auto-scrolling, or visually indicate that focus is off-screen.
\end{warning}

\begin{warning}
\textbf{Forgetting to update tab order.} Adding a new focusable
region without inserting it into the order is a frequent
oversight. Centralise the tab order in one place rather than
sprinkling \code{register\_target} calls.
\end{warning}

\section*{Workshop}

\begin{enumerate}[leftmargin=*]
  \item Add a third focusable region to the two-counter app from
    Chapter~\ref{ch:elm}. Make Tab cycle through all three.
  \item Add a modal that opens on \code{?} and shows a help
    screen. While open, all keys should go to the modal.
  \item Read \file{maya/include/maya/core/focus.hpp}. Identify
    the data members of \code{FocusManager} and the operations
    that mutate them.
\end{enumerate}

\begin{recap}
Focus is state, not a framework property. A \code{FocusManager}
tracks IDs and order; \code{view} reads focus to choose styling;
\code{update} routes input by inspecting focus. Modals stack
focus managers. Mouse clicks update focus through a hit-test
mapping. No magic, no global pointer.
\end{recap}

\chapter{Scroll, virtualisation, and viewports}
\label{ch:scroll}

\begin{aside}
Sooner or later your UI needs to show more content than fits the
terminal. The naive answer (render everything, clip what does
not fit) is wasteful for a list of ten thousand items.
\maya{}'s scroll model is a small data structure
(\code{ScrollState}) plus a viewport-aware rendering pattern.
This chapter walks it.
\end{aside}

\section{The ScrollState struct}

\file{maya/include/maya/core/scroll\_state.hpp} defines:

\begin{lstlisting}
struct ScrollState {
    int offset = 0;          // first visible item index
    int viewport_size = 0;   // number of rows visible
    int total_items = 0;     // total number of items in the source

    bool can_scroll_up()   const { return offset > 0; }
    bool can_scroll_down() const {
        return offset + viewport_size < total_items;
    }
    void scroll_to(int item_index, bool center = false);
    void scroll_by(int delta);
    void clamp();
};
\end{lstlisting}

Three integers (offset, viewport size, total). Every scrollable
widget stores one of these in its model. The widget's
\code{view} renders only items \code{[offset, offset+viewport\_size)}.

\subsection*{Why integers, not pointers}

A scroll position represented as ``a pointer to the visible row''
breaks the moment the underlying data changes. An integer offset
into a vector is stable: insertions and deletions can adjust it
deterministically.

\code{scroll\_to(index, center)} centres or pins the item at
\code{index}; \code{scroll\_by(delta)} moves by a relative
amount and clamps; \code{clamp()} ensures the offset is in
\code{[0, total\_items - viewport\_size]}.

\section{Virtualisation: rendering only what is visible}

A list of 10\,000 items, rendered naively, produces an element
tree with 10\,000 children. Layout walks them all. The canvas
diff walks them. ANSI emit walks them. Everything is slow.

The fix: render only the visible slice.

\begin{lstlisting}
Element render_list(const std::vector<Item>& items,
                    const ScrollState& s) {
    std::vector<Element> visible;
    for (int i = 0; i < s.viewport_size; ++i) {
        int idx = s.offset + i;
        if (idx >= s.total_items) break;
        visible.push_back(render_item(items[idx]));
    }
    return v(std::move(visible));
}
\end{lstlisting}

The element tree has at most \code{viewport\_size} children, no
matter how many items the source has. Layout, diff, emit ---
all proportional to the screen, not the data.

\begin{idea}
Virtualisation is one of the few optimisations that scales
\emph{better} than the underlying data structure. A list with a
million items renders as fast as one with twenty, because only
the twenty visible are ever materialised as elements.
\end{idea}

\subsection*{Handling tall items}

If items have variable height, the simple ``visible =
[offset, offset+viewport\_size)'' breaks: the slice's row count
depends on which items fall in it. Two approaches:

\begin{itemize}[leftmargin=*]
  \item \textbf{Constant-height items.} Easy: every item is one
    row. The slice is exactly \code{viewport\_size} items. This
    is what tables and pickers use.
  \item \textbf{Variable-height items.} Maintain a side table of
    per-item heights. To translate ``which item is on row $r$?''
    use a prefix sum or a Fenwick tree.
\end{itemize}

Most \maya{} widgets fix the height to one row. When you need
multi-row items, the widget either pre-measures or accepts the
$O($items$)$ pre-pass.

\section{Scroll bars}

The \code{scrollbar} widget reads a \code{ScrollState} and
renders a thumb:

\begin{lstlisting}
Element scrollbar(const ScrollState& s, int height) {
    if (s.total_items <= s.viewport_size) return text(\" \");
    int thumb_height = std::max(1, height * s.viewport_size / s.total_items);
    int thumb_y      = height * s.offset / s.total_items;
    // ... draw a vertical bar with thumb at thumb_y, thumb_height tall
}
\end{lstlisting}

The thumb size scales with the fraction of visible content; the
thumb position with the scroll offset. Standard math.

The scrollbar is decorative; the actual scrolling happens via
\code{ScrollState} mutations in \code{update}.

\section{Following the cursor}

A common pattern: a list with a current selection. When the
selection moves, the viewport should follow it.

\begin{lstlisting}
[&](KeyDown) {
    if (m.selected < m.items.size() - 1) m.selected++;
    if (m.selected >= m.scroll.offset + m.scroll.viewport_size)
        m.scroll.offset = m.selected - m.scroll.viewport_size + 1;
    return std::pair{m, Cmd<Msg>{}};
}
\end{lstlisting}

If the selection moves off the bottom of the viewport, scroll
down to keep it visible. The same logic applied symmetrically for
up movement keeps the cursor on-screen.

\subsection*{PageUp/PageDown}

Move the offset by \code{viewport\_size}:

\begin{lstlisting}
[&](PageDown) {
    m.scroll.scroll_by(m.scroll.viewport_size);
    m.selected = std::min(m.items.size() - 1,
                          m.selected + m.scroll.viewport_size);
    return std::pair{m, Cmd<Msg>{}};
}
\end{lstlisting}

Whether the cursor moves with the page or stays at a fixed
viewport offset is a UX choice; both are common.

\section{Horizontal scroll}

Some content (wide tables, log lines) extends past the right
edge. \code{ScrollState} can be reused on the horizontal axis:

\begin{lstlisting}
struct WideListModel {
    std::vector<std::string> lines;
    ScrollState vscroll;  // which rows are visible
    ScrollState hscroll;  // which columns
};
\end{lstlisting}

The render function clips each visible line to the horizontal
window:

\begin{lstlisting}
std::string clip_horizontal(std::string_view line, const ScrollState& h) {
    GraphemeWalker w{line};
    std::string out;
    int x = 0;
    while (auto g = w.next()) {
        if (x >= h.offset && x < h.offset + h.viewport_size) {
            out += g->text;
        }
        x += g->width;
        if (x >= h.offset + h.viewport_size) break;
    }
    return out;
}
\end{lstlisting}

Grapheme-aware, width-aware, allocation-conscious. The
performance is linear in visible cells, not in total line
width.

\section{Smooth scrolling}

For animated scroll (the viewport eases to a new position over
several frames), keep an animation state in the model:

\begin{lstlisting}
struct Scroll {
    int target = 0;       // where we want to go
    double current = 0;   // where we are now (fractional)
};
\end{lstlisting}

Each frame, advance \code{current} toward \code{target} by a
fraction of the gap. Round to integer for \code{offset}. Stop
when the difference is below 1.

The animation is driven by \code{Sub::animation\_frame}: as long
as \code{current != target}, the subscription is active.

\begin{warning}
Smooth scrolling on a terminal can look jittery because the
cell grid is integer. Snap to integer cells; smooth animations
of fractional offsets render as the same integer for multiple
frames in a row, which the user perceives as ``stuck''.
\end{warning}

\section{Sticky elements}

A header that stays visible while content scrolls under it is a
stylistic feature with one rule: layout the header separately
and never include it in the scrolled region.

\begin{lstlisting}
v(
  header() | height<1>,            // never scrolls
  scrolled_list(items, s) | grow<1>
)
\end{lstlisting}

The list grows to fill the rest of the column; the header stays
put.

For columns (a left rail that stays put while the main view
scrolls), the same pattern works on the horizontal axis.

\section{Snap-to-cell pitfalls}

\begin{warning}
\textbf{Offsets in graphemes, not bytes.} Scrolling a long line
horizontally must count graphemes, not bytes, or wide
characters will partially render.
\end{warning}

\begin{warning}
\textbf{Resize and ScrollState.} When the terminal resizes, the
viewport size changes. Always call \code{clamp()} to keep the
offset in bounds.
\end{warning}

\begin{warning}
\textbf{Dynamic total\_items.} If the source list grows or
shrinks (e.g. new log lines arrive), update
\code{total\_items} and clamp. Otherwise the user can scroll
past the end into empty space.
\end{warning}

\begin{warning}
\textbf{Offset becoming stale.} If items are inserted before
\code{offset}, the user's position drifts. Decide whether the
user's intent is ``stick with this item'' (adjust offset) or
``stay at this viewport position'' (keep offset). \maya{}'s
default is the former; sometimes you want the latter.
\end{warning}

\section{Performance numbers}

A virtualised list with a million items:

\begin{itemize}[leftmargin=*]
  \item Element tree size: bounded by viewport (e.g. 24
    children).
  \item Layout time: tens of microseconds.
  \item Diff time: sub-microsecond (only the viewport diffs).
  \item Bytes emitted per frame: a few hundred for incremental
    scroll.
\end{itemize}

The scaling is by viewport, not by source size. This is the
single optimisation that makes large data feasible in a TUI.

\section*{Workshop}

\begin{enumerate}[leftmargin=*]
  \item Build a list widget that virtualises 10\,000 items.
    Verify with a profiler that frame time is independent of
    item count.
  \item Add a scroll bar to the list using the \code{scrollbar}
    widget. Verify the thumb size and position match the
    \code{ScrollState}.
  \item Implement PageUp/PageDown that keep the selection
    centred when possible (only push it to the edge near the
    list bounds).
\end{enumerate}

\begin{recap}
\code{ScrollState} is three integers: offset, viewport size,
total. Virtualisation renders only the visible slice; layout,
diff, and emit cost are O(viewport). Sticky elements layout
outside the scrolled region. Animations are driven by signals
or \code{AnimationFrame}. The model is data; the runtime stays
unaware.
\end{recap}

\chapter{Error boundaries and recovery}
\label{ch:errorboundary}

\begin{aside}
The Elm architecture treats expected failures as messages. But
what about unexpected ones --- a stray null deref, an assertion
fired, a third-party library throwing an exception you did not
plan for? \maya{} provides \emph{error boundaries}: a way to
isolate a subtree of the UI so that when it fails, only that
subtree breaks, not the whole app. This chapter explains the
mechanism and how to use it.
\end{aside}

\section{What an error boundary is}

An error boundary is a wrapper component that catches exceptions
thrown by its child's \code{view} or \code{update}. When a
catch happens, the boundary swaps in a fallback view (typically a
\emph{something went wrong} message with a retry button) and
isolates the failure.

The pattern comes from React's \emph{error boundary} concept.
The mechanics in \maya{} are different (no class hierarchy, no
\code{componentDidCatch}) but the goal is the same: contain
failure, do not crash the whole UI.

\subsection*{Without an error boundary}

If \code{view} throws, the runtime catches the exception, prints
a diagnostic, restores the terminal, and exits. Your user sees
their app vanish. The reproducer is buried in stdout.

\subsection*{With an error boundary}

If \code{view} of a sub-tree throws, the boundary catches and
renders a small ``error'' panel in place of that sub-tree. The
rest of the UI continues. The user can keep working; the developer
can inspect the error message.

\section{The ErrorBoundary type}

\file{maya/include/maya/app/error\_boundary.hpp} defines:

\begin{lstlisting}
template <ElementLike Child>
class ErrorBoundary {
public:
    using Fallback = std::function<Element(std::exception_ptr)>;

    ErrorBoundary(Child child, Fallback fallback);

    Element build() &&;   // catches exceptions inside child.build()
};

template <ElementLike Child, class F>
auto error_boundary(Child&& child, F&& fallback);
\end{lstlisting}

The wrapper takes a child and a fallback. \code{build} wraps the
child's own \code{build} in a \code{try/catch} block:

\begin{lstlisting}
template <ElementLike Child>
Element ErrorBoundary<Child>::build() && {
    try {
        return std::move(child_).build();
    } catch (...) {
        return fallback_(std::current_exception());
    }
}
\end{lstlisting}

If the child's build (which recursively builds the entire
sub-tree) succeeds, the wrapper passes it through. If it throws
anywhere, the fallback is rendered instead.

\subsection*{The fallback function}

The fallback receives a \code{std::exception\_ptr} so it can
extract a message if desired:

\begin{lstlisting}
auto safe_panel = error_boundary(
    risky_panel(state),
    [](std::exception_ptr ep) {
        std::string msg;
        try { std::rethrow_exception(ep); }
        catch (const std::exception& e) { msg = e.what(); }
        catch (...)                       { msg = \"unknown error\"; }
        return box(
            text(\"This panel crashed:\") | Bold | Fg<255, 120, 120>,
            text(msg) | Dim
        ) | border<Single> | pad<1>;
    }
);
\end{lstlisting}

A user-friendly fallback hides the exception details by default
and reveals them on demand. A developer-mode fallback dumps the
full \code{what()} immediately.

\section{Where boundaries belong}

Not every sub-tree needs a boundary. Adding boundaries
everywhere defeats the purpose (the runtime becomes lossy). Add
one where:

\begin{itemize}[leftmargin=*]
  \item The sub-tree renders user-supplied content (a markdown
    preview, a code highlighter) where input may be malformed.
  \item The sub-tree calls third-party code you do not control.
  \item The sub-tree is independent of the rest of the UI ---
    losing it does not lose work in progress.
\end{itemize}

A reasonable rule: the panels of a multi-panel app each get a
boundary. The whole-app root does not (because if rendering the
whole app fails, there is nothing useful to show).

\section{Catching during update}

\code{ErrorBoundary} only catches in \code{build}/\code{view}.
Exceptions in \code{update} are different: they corrupt the
model.

\maya{} handles update exceptions at the runtime level:

\begin{lstlisting}
try {
    auto [m2, c2] = P::update(model, msg);
    model = std::move(m2);
    interpret(c2, loop);
} catch (...) {
    log_update_exception(std::current_exception(), model, msg);
    // The model is unchanged from before update was called.
    // No effects were interpreted.
}
\end{lstlisting}

If \code{update} throws, the model is not mutated and no
commands are issued. The next frame renders the previous model.
The user might not notice anything happened; the log records it.

This is intentionally conservative: an exception in \code{update}
indicates a programming bug, not an expected condition. The
right fix is for the developer to add a branch in \code{update}
that returns a message, not for the runtime to silently retry.

\begin{warning}
Do not rely on the runtime's update-exception catching as a
defence in depth. It exists to keep the program alive while the
bug is diagnosed, not as a feature.
\end{warning}

\section{Effect failures}

Failures inside \code{Cmd::task} are wrapped at the boundary:

\begin{lstlisting}
return Cmd<Msg>::task([url]() -> Msg {
    try {
        return FetchOk{do_fetch(url)};
    } catch (const std::exception& e) {
        return FetchFailed{e.what()};
    }
});
\end{lstlisting}

The task callable catches its own exceptions and routes them
through the message type. \code{update} branches on the result.
No exception ever crosses the task boundary.

If you fail to wrap, \maya{}'s task scheduler catches the
exception and logs it, but the corresponding message is never
delivered. Your \code{update} will not see a result. This is a
silent failure --- bad enough that you should always wrap.

\section{Recovery patterns}

\subsection*{Retry on a timer}

After a failure, schedule a retry:

\begin{lstlisting}
[&](FetchFailed e) {
    m.error = e.reason;
    m.retry_count++;
    if (m.retry_count < 3) {
        return std::pair{m,
            Cmd<Msg>::after(1s * (1 << m.retry_count),
                            Retry{})};
    }
    return std::pair{m, Cmd<Msg>{}};
}
\end{lstlisting}

Exponential backoff in three lines.

\subsection*{Reset on user action}

Provide a ``retry'' button in the error UI:

\begin{lstlisting}
auto error_panel = v(
    text(error_message) | Fg<255, 120, 120>,
    text(\"[r] retry  [q] quit\") | Dim
);
\end{lstlisting}

Pressing \code{r} fires a \code{Retry} message; \code{update}
re-issues the original \code{Cmd}.

\subsection*{Degraded mode}

If a feature fails repeatedly, disable it and keep going:

\begin{lstlisting}
if (m.feature_x_failures >= 3) {
    m.feature_x_disabled = true;
}
\end{lstlisting}

The view checks the flag and renders a placeholder. The rest of
the app stays alive.

\section{Logging the failure}

A boundary that catches but never reports is worse than no
boundary at all (you do not even know there was a bug). Pair the
boundary with a log:

\begin{lstlisting}
auto safe_panel = error_boundary(
    risky_panel(state),
    [&](std::exception_ptr ep) {
        log_exception(ep, \"risky_panel\");
        return fallback_ui(ep);
    }
);
\end{lstlisting}

Where \code{log\_exception} writes to a file or sends to a
telemetry service. For development, \code{std::cerr} is fine.
For production, write to a debug ring buffer that the user can
toggle on demand.

\begin{aside}
A debug ring buffer is a small in-memory log that survives
crashes. \maya{}'s built-in one keeps the last 1024 messages and
exposes them via a \code{maya::dump\_log()} call --- useful when
the user reports ``something looked weird''.
\end{aside}

\section{Testing error paths}

The boundary itself is testable:

\begin{lstlisting}
TEST_CASE(\"boundary catches throwing build\") {
    auto throwing = test::throw_on_build();
    auto wrapped = error_boundary(throwing, [](auto) {
        return text(\"fallback\");
    });
    Element e = std::move(wrapped).build();
    REQUIRE(is_text(e, \"fallback\"));
}
\end{lstlisting}

A test helper that throws on \code{build}, wrapped in a
boundary; the result should be the fallback. The runtime is
not involved.

\section{Pitfalls}

\begin{warning}
\textbf{Catching too broadly.} A boundary that catches all
exceptions can mask bugs you want to surface. Catch specific
exception types where possible, and re-throw the ones you do not
understand.
\end{warning}

\begin{warning}
\textbf{Side effects in the fallback.} The fallback is a render
function; it should not have side effects (file I/O, network).
If you need to log, do it before returning the fallback element.
\end{warning}

\begin{warning}
\textbf{Boundary on a leaf.} Wrapping a single text node in a
boundary is wasteful (text nodes do not throw). Boundaries pay
off on subtrees that aggregate work.
\end{warning}

\begin{warning}
\textbf{Boundary that itself throws.} If your fallback throws,
the runtime catches it and you crash anyway. Keep fallbacks
defensively simple.
\end{warning}

\section*{Workshop}

\begin{enumerate}[leftmargin=*]
  \item Wrap your app's largest panel in an
    \code{error\_boundary}. Force an exception inside it and
    verify only that panel is replaced.
  \item Add a debug overlay that shows the last logged
    exception when the user presses \code{F12}.
  \item Read
    \file{maya/include/maya/app/error\_boundary.hpp}. Identify
    where the exception is caught and how the fallback is
    invoked.
\end{enumerate}

\begin{recap}
Error boundaries isolate failure. A boundary catches exceptions
during \code{build}/\code{view} and renders a fallback in place
of the broken sub-tree. \code{update} exceptions are caught at
the runtime level and the model is left unchanged. Task
exceptions must be converted to messages at the task boundary.
Together these mechanisms keep \maya{} apps alive in the face of
bugs.
\end{recap}

\chapter{The reactive primitives, previewed}
\label{ch:reactivity-preview}

\begin{aside}
Most of \maya{} is plain functions. A handful of places ---
animations, watched files, derived state --- want \emph{reactive
primitives}: signals that notify subscribers when their value
changes. \maya{}'s signal system lives in
\file{maya/include/maya/core/signal.hpp} and gets a full
treatment in Chapter~\ref{ch:signals}. This chapter is the
architectural preview: what they are, why they exist, where the
Elm architecture leaves room for them.
\end{aside}

\section{Signals in one paragraph}

A \code{Signal<T>} is a value that notifies subscribers when it
changes. A \code{Computed<T>} is a derived signal whose value is
a function of other signals. An \code{Effect} is a callable that
re-runs whenever one of the signals it reads changes.

If you have used SolidJS, MobX, Vue, Knockout, or Svelte 5, this
is the same machinery. The mechanics differ; the model is the
same.

\section{Why signals at all?}

Pure \code{update} works for most of an app. But two patterns
are awkward without reactive helpers:

\begin{itemize}[leftmargin=*]
  \item \textbf{Animations.} A counter that ticks once per
    frame, a loading spinner, a smooth scroll. The state is not
    business logic; it is presentation.
  \item \textbf{Derived state.} A search box that filters a
    list. The filter result is a function of the list and the
    query; recomputing it inside \code{view} every frame is
    wasteful.
\end{itemize}

You could express both with \code{Cmd} and \code{Sub}, and many
\maya{} apps do. Signals are an alternative for fine-grained
reactivity without piping every change through the message
variant.

\subsection*{The trade}

\code{update} is pure and testable. Signals introduce mutable
shared state. That is a real trade --- but it is bounded: the
signal graph is small, its subscribers are explicit, and the
mutation is observable.

For the parts of an app where the signal model is a better fit
(animations, computed columns, file watches), the cost is worth
it. For business logic, prefer \code{update}.

\begin{warning}
The two models can mix in one app, but in any given subsystem
pick one. Mixing them inside a single widget is a recipe for
``why did this not re-render?'' bugs.
\end{warning}

\section{The basic shape}

\begin{lstlisting}
auto count = Signal{0};
auto doubled = Computed{[&] { return count() * 2; }};

Effect{[&] {
    std::cout << \"count=\" << count() << \" doubled=\" << doubled() << \"\\n\";
}};

count.set(5);   // prints \"count=5 doubled=10\"
count.set(7);   // prints \"count=7 doubled=14\"
\end{lstlisting}

Three pieces:

\begin{itemize}[leftmargin=*]
  \item A \code{Signal<int>} holds a value and tracks subscribers.
  \item A \code{Computed<int>} is a function over signals. When
    its inputs change, it recomputes; subscribers of the computed
    are notified.
  \item An \code{Effect} runs once at registration, then again
    whenever any signal it read changes.
\end{itemize}

Subscriber tracking is automatic. \code{count()} inside the
effect registers a dependency; \code{count.set(\ldots)} triggers
re-run.

\section{Why this works without explicit subscribe}

The key trick: while an effect or computed is evaluating,
\maya{} sets a thread-local current node. Any \code{Signal::get}
records the current node as a subscriber. When the effect
finishes, the dependency set is fixed until the next run.

\begin{lstlisting}
template <class T>
T Signal<T>::get() {
    if (auto* node = detail::current_node) {
        subscribers_.insert(node);
        node->dependencies.insert(this);
    }
    return value_;
}
\end{lstlisting}

The dependency graph builds itself. No \code{addEventListener},
no \code{watch} declarations. The function calls themselves
register the wiring.

\subsection*{What gets tracked}

Only reads through \code{Signal::get()} (which the call operator
forwards to). Plain variable reads, function calls, and
non-signal state are invisible to the tracker.

This means: if you compute a value from a signal and a plain
variable, only the signal is tracked. Changing the plain variable
will not re-run the effect; you must read it through a signal
or wrap it in one.

\section{When to use signals in maya}

\subsection*{Animations}

A spinner is the canonical case:

\begin{lstlisting}
auto frame = Signal{0};

// somewhere: a 60Hz timer increments frame.
schedule_every(16ms, [&] { frame.set(frame() + 1); });

// in view:
Element view(const Model& m) {
    return text(spinner_frames[frame() % spinner_frames.size()]);
}
\end{lstlisting}

The view reads \code{frame()} so it auto-subscribes. The signal
changes; \maya{}'s render pipeline knows the affected sub-tree
must re-render. The model is not involved.

\subsection*{Computed columns}

A list with a filter:

\begin{lstlisting}
auto items = Signal<std::vector<Item>>{};
auto query = Signal<std::string>{};
auto filtered = Computed{[&] {
    std::vector<Item> out;
    for (auto& i : items()) {
        if (i.label.contains(query())) out.push_back(i);
    }
    return out;
}};
\end{lstlisting}

\code{filtered} recomputes only when \code{items} or
\code{query} changes. The view reads \code{filtered()};
intermediate frames see the same vector without recomputing.

\subsection*{File watches}

A signal whose value is the contents of a file, automatically
updated when the file is modified:

\begin{lstlisting}
auto config = file_signal(\"~/.config/myapp.toml\");
auto parsed = Computed{[&] { return parse_toml(config()); }};
\end{lstlisting}

Behind the scenes, \code{file\_signal} uses inotify (Linux),
kqueue (macOS), or ReadDirectoryChangesW (Windows) to watch the
file. When it changes, \code{config} updates, \code{parsed}
recomputes, every subscriber re-renders.

\section{What signals do not replace}

The Elm core stays. \code{update} is still the place for
business logic. Signals supplement it for the cases above; they
do not subsume it.

A useful rule: if the state is something a user might want to
\emph{see in a debugger}, put it in \code{Model}. If the state
is a synthetic computation (animation frame, derived index),
\code{Signal} is fine.

\section{The cost}

Signals add:

\begin{itemize}[leftmargin=*]
  \item Heap allocations: each \code{Signal} stores a
    \code{std::vector} of subscribers.
  \item Thread-local state: the dependency tracker.
  \item Mutable state: signals are by definition not pure.
  \item Glitches if you do not understand the model: re-running
    an effect can trigger another effect, which can trigger
    another\ldots
\end{itemize}

Used sparingly (a handful of signals per app), the cost is
negligible. Used as the primary state-management tool, the cost
adds up --- and you have lost the testability of pure
\code{update}.

\section*{Workshop}

\begin{enumerate}[leftmargin=*]
  \item Create a \code{Signal<int>} and an \code{Effect} that
    prints the value. Set the signal three times; observe the
    effect runs four times (once on registration, three on
    changes).
  \item Add a \code{Computed} that doubles the signal. Confirm
    the effect prints both the signal and the computed.
  \item Read \file{maya/include/maya/core/signal.hpp}. Identify
    the data members of \code{Signal<T>} and the locks (if any)
    that protect them.
\end{enumerate}

\begin{recap}
Signals, computeds, and effects form a fine-grained reactive
system. They build a dependency graph automatically through
\code{get} calls, re-run effects when inputs change, and
complement (rather than replace) the Elm architecture.
Chapter~\ref{ch:signals} treats them in full; this chapter
explains where they fit.
\end{recap}

\chapter{Runtime internals: a guided tour}
\label{ch:runtime}

\begin{aside}
We have described \code{maya::run} as ``the runtime''. Plural,
really --- there are several cooperating subsystems behind the
single \code{run} call. This chapter is a guided tour of the
internals: what each piece does, where it lives in the source,
and how they fit together. By the end you should be able to
open any header in \file{maya/include/maya/app/} and know what
you are reading.
\end{aside}

\section{The roadmap}

A \maya{} runtime is six cooperating pieces:

\begin{enumerate}[leftmargin=*]
  \item \textbf{Context.} The bundle of platform handles,
    config, capabilities, and stats.
  \item \textbf{Event loop.} The driver that waits for events,
    dispatches them, and runs the pipeline.
  \item \textbf{Event source.} The platform abstraction (epoll,
    kqueue, WaitForMultipleObjects) that produces events.
  \item \textbf{Subscription manager.} The diff of declared
    subscriptions versus active OS-level wires.
  \item \textbf{Command interpreter.} The interpreter that walks
    \code{Cmd<Msg>} variants and effects them.
  \item \textbf{Renderer.} The pipeline that turns elements into
    ANSI bytes.
\end{enumerate}

Each lives behind its own header. \code{maya::run} composes them.

\section{Context}

\file{maya/include/maya/app/context.hpp} defines:

\begin{lstlisting}
struct Context {
    Config            config;
    platform::Terminal& terminal;
    TerminalCaps      caps;
    SignalDispatcher& signals;
    WakeFd&           wake;
    Stats             stats;
    std::optional<Logger> logger;
};
\end{lstlisting}

The context is constructed at startup and passed by reference
into every subsystem. Subsystems read what they need; the
context owns nothing they own.

\subsection*{Why a context, not globals?}

Globals would work, but:

\begin{itemize}[leftmargin=*]
  \item Two simultaneous \maya{} runtimes in one process (a test
    runner, for example) would clobber each other.
  \item Threading is harder to reason about with hidden state.
  \item Tests cannot vary the context without resetting globals.
\end{itemize}

A context is a small ceremony for a real benefit.

\section{Event loop}

The driver in \file{maya/include/maya/app/app.hpp}:

\begin{lstlisting}
template <Program P>
int run(Config cfg) {
    auto ctx = make_context(cfg);
    auto loop = EventLoop{ctx};
    auto renderer = Renderer{ctx};
    auto interpreter = Interpreter{ctx, loop};

    auto [model, init_cmd] = P::init();
    interpreter.interpret(init_cmd);

    while (loop.running()) {
        renderer.frame(P::view(model));
        loop.update_subscriptions(P::subscribe(model));

        Event ev = loop.wait();
        if (auto msg = to_message<typename P::Msg>(ev)) {
            try {
                auto [m2, c2] = P::update(model, *msg);
                model = std::move(m2);
                interpreter.interpret(c2);
            } catch (...) {
                log_update_exception(ctx, std::current_exception());
            }
        }
    }
    return 0;
}
\end{lstlisting}

The pattern is the loop we walked in Chapter~\ref{ch:elm}, now
with the supporting classes named.

\subsection*{Composition over inheritance}

\code{EventLoop}, \code{Renderer}, and \code{Interpreter} are
plain classes that take a \code{Context\&} in their constructor.
There is no \code{class App} hierarchy. The runtime is the
\code{run} function plus the helpers it instantiates.

\begin{idea}
Composition keeps the runtime debuggable. Each subsystem is a
file; each file is a class; each class has a clear contract.
You can read the runtime by reading six files in order, without
following virtual dispatches.
\end{idea}

\section{Event source}

The OS-level pump (Chapter~\ref{ch:platform}). It produces
\code{Event} values: \code{KeyPress}, \code{MouseEvent},
\code{Resize}, \code{Wakeup}, \code{Timeout}.

The event loop calls \code{source.wait(timeout)} and translates
the result into a message via the active subscription set.

\subsection*{The Event variant}

\begin{lstlisting}
using Event = std::variant<
    KeyPress, MouseEvent, Resize, Paste,
    Timeout, Wakeup, Unknown>;
\end{lstlisting}

The same overload pattern from Chapter~\ref{ch:elements} matches
each kind. The subscription manager (next section) handles the
match.

\section{Subscription manager}

The subscriber tracks which \code{Sub} alternatives are active
and how to translate events into messages:

\begin{lstlisting}
class SubscriptionManager {
public:
    void update(const Sub<Msg>& new_sub);
    std::optional<Msg> dispatch(const Event& ev);
};
\end{lstlisting}

\code{update} takes the new \code{Sub} from
\code{P::subscribe(model)} and diffs it against the previous one.
Subscriptions added are installed; subscriptions removed are
torn down. Subscriptions present in both are kept.

\code{dispatch} takes an event and returns a message if any
subscription matches. The matching is by event kind and any
subscription-specific filters.

\subsection*{Diffing subscriptions}

The diff uses structural equality on a flattened list of
(kind, parameters) tuples. Two \code{Every(50ms, Tick\{\})}
subscriptions are equal regardless of how they were constructed.
The fingerprint comparison is O(n) in subscription count, but
subscription counts are tiny (typically under ten).

\section{Command interpreter}

\file{maya/include/maya/app/} also houses the interpreter:

\begin{lstlisting}
class Interpreter {
public:
    void interpret(const Cmd<Msg>& cmd);

private:
    void interpret_none(None);
    void interpret_quit(Quit);
    void interpret_after(After<Msg>);
    void interpret_task(Task<Msg>);
    void interpret_batch(Batch<Msg>);
    // ... one per alternative
};
\end{lstlisting}

Each \code{interpret\_X} performs the real-world side effect
corresponding to that alternative. \code{interpret\_quit} sets a
flag on the event loop. \code{interpret\_after} schedules a
timer on the event source. \code{interpret\_task} hands the
callable to the worker pool.

The dispatch is a single \code{std::visit}:

\begin{lstlisting}
void Interpreter::interpret(const Cmd<Msg>& cmd) {
    std::visit([this](auto& alt) { this->do_interpret(alt); }, cmd.v);
}
\end{lstlisting}

Adding a new \code{Cmd} alternative is two steps: add to the
variant, add an \code{interpret} overload. No existing callers
change.

\subsection*{The worker pool}

\code{Task} and \code{IsolatedTask} run on threads. \maya{}
ships with a small fixed-size pool plus the option for isolated
threads. The pool size is configurable (default: thread::hardware
concurrency divided by two; minimum one).

A pool worker, when its callable returns, writes the resulting
message to a lock-free queue and signals the wake fd. The main
loop's event source picks up the wake; \code{dispatch} pulls the
message from the queue.

\begin{aside}
The wake-fd pattern is what makes \maya{} able to deliver
worker-thread messages without polling. The cost is one byte
written and one byte read per cross-thread message.
\end{aside}

\section{The renderer}

\file{maya/include/maya/render/renderer.hpp} houses:

\begin{lstlisting}
class Renderer {
public:
    explicit Renderer(Context&);
    void frame(Element root);
    void force_redraw();
    void set_caps(TerminalCaps);
};
\end{lstlisting}

\code{frame(Element)} runs the six-stage pipeline from
Chapter~\ref{ch:pipeline}: layout, paint, diff, serialise,
write. \code{force\_redraw} discards the cache. \code{set\_caps}
updates capability detection (e.g. on terminal-type change).

The renderer owns the previous canvas (for the diff) and the
output buffer. Each frame is amortised; the buffer's capacity
persists across frames.

\section{Shutdown}

When \code{Cmd::Quit} is interpreted, the event loop sets its
running flag to false. The main \code{while} exits on the next
iteration. Destructors restore the terminal:

\begin{enumerate}[leftmargin=*]
  \item \code{Renderer}'s destructor flushes any pending bytes.
  \item \code{Interpreter}'s destructor joins worker threads.
  \item \code{EventLoop}'s destructor closes the event source.
  \item \code{Terminal}'s destructor calls \code{leave\_raw\_mode}.
\end{enumerate}

The order matters: flush bytes \emph{before} restoring raw mode
(so the final frame appears), join workers \emph{before}
destroying the queue (so writes do not race), and so on. RAII
takes care of this if the construction order is right.

\subsection*{Signal-driven shutdown}

If the user presses Ctrl-C, the signal handler writes a special
\code{InterruptedSignal} message to the queue. \code{update}
sees it (most apps handle it as ``quit''); the rest is the
normal shutdown path.

If the signal handler cannot reach \code{update} (e.g. the
runtime is stuck in a long computation), the OS will eventually
kill the process. The terminal restore happens in
\code{atexit}-registered handlers as a last resort.

\section{Concurrency model summary}

\begin{itemize}[leftmargin=*]
  \item Main thread: event loop, renderer, interpreter,
    \code{update}, \code{view}, \code{subscribe}.
  \item Worker threads: \code{Task} callables.
  \item Isolated threads: \code{IsolatedTask} callables.
  \item Signal-disposition thread (on Win32): the console
    control handler.
\end{itemize}

The handoff between worker and main is via the wake fd and the
message queue. The main thread is the only one that touches
\code{Model}, \code{view}, or the renderer.

\begin{warning}
Do not share mutable state across the worker/main boundary
without synchronisation. The whole architecture is built on the
idea that \code{Model} is owned by the main thread; if you
violate that, races are likely.
\end{warning}

\section{Reading the source}

A suggested order for reading \maya{}'s runtime:

\begin{enumerate}[leftmargin=*]
  \item \file{maya/include/maya/app/app.hpp} (the \code{run}
    function and the main loop).
  \item \file{maya/include/maya/app/context.hpp}.
  \item \file{maya/include/maya/app/sub.hpp}.
  \item \file{maya/include/maya/core/cmd.hpp}.
  \item \file{maya/include/maya/app/events.hpp}.
  \item \file{maya/include/maya/render/renderer.hpp}.
\end{enumerate}

Six files. Combined, around 2000 lines of (heavily commented)
source. That is the entire runtime.

\section*{Workshop}

\begin{enumerate}[leftmargin=*]
  \item Open \file{maya/include/maya/app/app.hpp} and identify
    every subsystem class instantiated in \code{run}.
  \item Trace what happens when the user presses Ctrl-C from
    the terminal to the program's exit. Note each thread
    boundary crossed.
  \item Try implementing a simple\code{trace\_all} mode that
    logs every event, message, and command. Decide which
    subsystem you would hook into.
\end{enumerate}

\begin{recap}
The runtime is a composition of six subsystems: context, event
loop, event source, subscription manager, command interpreter,
renderer. Each lives behind its own header; \code{run} composes
them. The main thread owns the model; worker threads communicate
via a wake-fd and message queue. RAII handles shutdown order.
The whole thing is around 2000 lines of source.
\end{recap}

\chapter{Testing maya applications}
\label{ch:testing}

\begin{aside}
A library architecture is only worth the cost it imposes if the
benefits show up in tests. \maya{}'s purity discipline pays off
here: most of an app is testable without a terminal, a renderer,
or even a runtime. This chapter is a playbook of testing
patterns that work with the architecture.
\end{aside}

\section{What is and is not testable directly}

\code{update} is a pure function. So is \code{view}. So is the
fingerprint computation, the layout pass, and the diff. These
are testable with plain function calls and assertions; no setup.

The things that are \emph{not} directly testable are anything
that touches the OS:

\begin{itemize}[leftmargin=*]
  \item The event loop blocking on epoll.
  \item The renderer writing bytes to a terminal.
  \item Tasks spawning worker threads.
  \item Signals delivered to signalfd.
\end{itemize}

For these, you mock the boundary. \maya{} ships a test harness
that wraps the runtime with a controllable event source and
captured-byte renderer.

\section{Unit testing update}

The simplest test pattern:

\begin{lstlisting}
TEST_CASE(\"counter increments\") {
    auto m0 = Counter::init().first;
    auto [m1, cmd] = Counter::update(m0, Counter::Inc{});
    REQUIRE(m1.count == 1);
    REQUIRE(std::holds_alternative<None>(cmd.v));
}
\end{lstlisting}

Direct function call. No runtime, no terminal, no setup. The
test is the specification.

\subsection*{Testing a sequence}

A useful helper:

\begin{lstlisting}
template <Program P>
auto fold(std::initializer_list<typename P::Msg> msgs) {
    auto m = P::init().first;
    std::vector<Cmd<typename P::Msg>> cmds;
    for (auto& msg : msgs) {
        auto [m2, c] = P::update(m, msg);
        m = std::move(m2);
        cmds.push_back(c);
    }
    return std::pair{std::move(m), std::move(cmds)};
}

TEST_CASE(\"three increments and a decrement\") {
    auto [m, _] = fold<Counter>({
        Counter::Inc{}, Counter::Inc{}, Counter::Inc{},
        Counter::Dec{}
    });
    REQUIRE(m.count == 2);
}
\end{lstlisting}

The fold is the playback. The final model is the assertion target.

\subsection*{Testing emitted commands}

When \code{update} returns a non-trivial \code{Cmd}, assert
against the variant:

\begin{lstlisting}
TEST_CASE(\"save schedules a clipboard write\") {
    auto [m, cmd] = MyApp::update(
        MyApp::Model{.body = \"hello\"}, MyApp::Save{});
    auto* wc = std::get_if<WriteClipboard<MyApp::Msg>>(&cmd.v);
    REQUIRE(wc != nullptr);
    REQUIRE(wc->text == \"hello\");
}
\end{lstlisting}

\code{std::get\_if} returns a pointer or null; testing for null
is the cleanest pattern.

\section{Testing view}

\code{view} returns an \code{Element}. You can assert against
its structure with a visitor:

\begin{lstlisting}
TEST_CASE(\"counter view contains the count\") {
    auto e = Counter::view(Counter::Model{.count = 42});
    bool found = false;
    walk(e, [&](const TextElement& t) {
        auto s = text_view(t);
        if (s.find(\"42\") != std::string_view::npos) found = true;
    });
    REQUIRE(found);
}
\end{lstlisting}

\code{walk} is a tree-walker helper. You assert on the presence
or absence of text, not on the exact rendering. The exact bytes
depend on the terminal capabilities.

\subsection*{Snapshot tests}

For UIs that are mostly stable, snapshot tests work well:

\begin{lstlisting}
TEST_CASE(\"counter view snapshot\") {
    auto e = Counter::view(Counter::Model{.count = 42});
    auto dump = dump_element(e);
    REQUIRE(dump == load_snapshot(\"counter_view_42.txt\"));
}
\end{lstlisting}

\code{dump\_element} produces a textual representation; the
snapshot file is checked into source control. Differences are
visible in code review.

The risk: snapshot tests catch every change, including
intentional ones. Update them deliberately, not reflexively.

\section{Property tests}

Pure functions are gold for property-based testing:

\begin{lstlisting}
PROPERTY(\"increment then decrement is identity\", Counter::Model m) {
    auto [m1, _] = Counter::update(m, Counter::Inc{});
    auto [m2, __] = Counter::update(m1, Counter::Dec{});
    REQUIRE(m2.count == m.count);
}

PROPERTY(\"any sequence preserves count >= 0\", std::vector<Counter::Msg> seq) {
    auto m = Counter::init().first;
    for (auto& msg : seq) {
        auto [m2, _] = Counter::update(m, msg);
        m = std::move(m2);
    }
    REQUIRE(m.count >= 0);
}
\end{lstlisting}

The framework (rapidcheck, Catch2 generators) supplies the
inputs. Counterexamples are minimised automatically. Bugs that
would never appear in hand-written tests show up.

\subsection*{Generating Msg variants}

For a \code{Msg = std::variant<\ldots>}, the generator needs to
sample each alternative. Custom generators are straightforward:

\begin{lstlisting}
template <>
struct Arbitrary<Counter::Msg> {
    static Gen<Counter::Msg> arbitrary() {
        return gen::oneOf<Counter::Msg>(
            gen::just(Counter::Inc{}),
            gen::just(Counter::Dec{}),
            gen::just(Counter::Quit{})
        );
    }
};
\end{lstlisting}

For variants with payload, generate the payload with
\code{gen::build}.

\section{Testing the layout phase}

The layout pass is a pure function from
\code{(Element, Size)} to a layout tree:

\begin{lstlisting}
TEST_CASE(\"flex row distributes width\") {
    auto e = h(
        text(\"abc\") | grow<1>,
        text(\"def\") | grow<2>
    );
    auto layout = compute_layout(e, Size{30, 1});
    REQUIRE(rect_of_child(layout, 0).w == 10);
    REQUIRE(rect_of_child(layout, 1).w == 20);
}
\end{lstlisting}

The layout helper exposes the post-layout rects so you can
assert against them.

\section{Testing the diff and serialise stages}

These are pure too, but assertions become byte-counting:

\begin{lstlisting}
TEST_CASE(\"changing one cell emits a minimal diff\") {
    Canvas a = canvas_of(\"hello world\");
    Canvas b = canvas_of(\"hello WORLD\");
    auto diff = compute_diff(a, b);
    auto bytes = serialise(diff);
    REQUIRE(bytes.size() < 30);  // cursor move + 5 cells + reset
}
\end{lstlisting}

The exact byte count depends on the serialiser strategy; the
contract is ``less than full redraw''.

\section{End-to-end with a fake terminal}

For tests that exercise the runtime end-to-end:

\begin{lstlisting}
TEST_CASE(\"app handles quit message\") {
    test::FakeTerminal term;
    test::EventQueue events;
    events.push(Event::KeyPress{'q'});

    auto rc = test::run<Counter>(term, events);
    REQUIRE(rc == 0);
    REQUIRE(term.last_frame_contains(\"Count: 0\"));
}
\end{lstlisting}

The fake terminal records every byte written. The event queue
delivers events in order. The test runner exits when the program
returns. The terminal's \code{last\_frame\_contains} is a
helper for fuzzy assertions.

\subsection*{Time control}

Real-time-dependent tests use a virtual clock:

\begin{lstlisting}
test::FakeClock clock;
events.push_at(clock.now() + 1s, Event::KeyPress{'a'});

auto rc = test::run<MyApp>(term, events, clock);
\end{lstlisting}

The virtual clock advances only when the event queue asks for
a wait. Tests that need to verify ``300ms later, the toast
expired'' run instantly.

\section{Testing animation}

Animation tests use \code{AnimationFrame} subscriptions and
the virtual clock:

\begin{lstlisting}
TEST_CASE(\"spinner advances each frame\") {
    test::FakeTerminal term;
    test::EventQueue events;
    test::FakeClock clock;

    // Advance 10 frames.
    for (int i = 0; i < 10; ++i) {
        events.push_animation_frame();
    }

    test::run<App>(term, events, clock, /*max_frames=*/10);
    REQUIRE(term.frame_count() == 10);
    // Each frame should show a different spinner glyph.
}
\end{lstlisting}

The \code{max\_frames} parameter prevents infinite loops in
tests that would otherwise run forever.

\section{Testing concurrency}

\code{Cmd::task} runs on a thread. To test deterministically,
use a synchronous task pool:

\begin{lstlisting}
auto ctx = test::Context{.task_pool = test::SyncTaskPool{}};
\end{lstlisting}

Synchronous tasks run inline on the main thread; the assertion
about their result happens after the dispatch.

\subsection*{Testing race conditions}

Race-condition tests are the hardest. \maya{} ships a
\code{ThreadSanitizer}-aware build mode that flags races at
runtime; pair it with the test suite by running tests with
\code{-fsanitize=thread}.

\begin{warning}
Race tests are necessarily probabilistic. A test that passes
once does not prove the absence of races; it only proves the
race did not happen on that run. Loop the test a few thousand
times for confidence.
\end{warning}

\section{Coverage and confidence}

For a \maya{} app, aim for:

\begin{itemize}[leftmargin=*]
  \item 100\% statement coverage on \code{update}.
  \item Property tests for any invariant the model maintains.
  \item A handful of view snapshot tests for visual regressions.
  \item One or two end-to-end tests per major user flow.
\end{itemize}

The first two are cheap because \code{update} is pure. The last
two are the expensive ones; budget them deliberately.

\section{Pitfalls}

\begin{warning}
\textbf{Testing the renderer's exact output.} Different
terminals produce different bytes for the same content. Test
the content, not the bytes.
\end{warning}

\begin{warning}
\textbf{Snapshot fragility.} A snapshot that changes on every
unrelated UI tweak teaches the team to update snapshots without
reading them. Keep snapshots focused; cover one component at a
time.
\end{warning}

\begin{warning}
\textbf{Flaky timing tests.} If your end-to-end test sleeps for
real wall-clock time, it will be flaky on slow machines. Use
the virtual clock for any timing-dependent assertion.
\end{warning}

\begin{warning}
\textbf{Tests that depend on terminal capabilities.} A test
that checks ``truecolour bytes were emitted'' will fail on a
fake terminal with caps disabled. Either configure caps for
the test or assert against the more general behaviour.
\end{warning}

\section*{Workshop}

\begin{enumerate}[leftmargin=*]
  \item Write five unit tests for the two-counter app from
    Chapter~\ref{ch:elm}: each verifies one branch of
    \code{update}.
  \item Add a property test asserting that the counter never
    overflows under any sequence of messages.
  \item Configure your CI to fail on snapshot mismatches and to
    update snapshots only on a deliberate flag.
\end{enumerate}

\begin{recap}
The Elm architecture pays for itself most in tests.
\code{update} and \code{view} are pure and test directly with
function calls. The layout and serialise stages are pure too.
End-to-end tests use a fake terminal, an event queue, and a
virtual clock. Property tests cover the invariants the model
maintains. Concurrency tests need patience and a synchronous
pool.
\end{recap}

\chapter{The widget catalogue, by architecture}
\label{ch:widget-arch}

\begin{aside}
\maya{} ships dozens of widgets in \file{maya/include/maya/widget/}.
Each is a small builder that produces an \code{Element}. The
catalogue is too large to enumerate here (that is
Chapter~\ref{ch:widgets}), but the \emph{architecture} of a
widget --- the patterns every widget follows --- is small enough
to learn in one sitting. This chapter is the architectural
template.
\end{aside}

\section{What a widget is, exactly}

A widget is one of three things:

\begin{itemize}[leftmargin=*]
  \item \textbf{A pure renderer.} Takes parameters, returns an
    \code{Element}. No state, no input handling. The
    \code{badge}, \code{divider}, \code{spinner} widgets.
  \item \textbf{A stateful component.} Has its own \code{Model},
    \code{Msg}, \code{update}, \code{view}. Composes via
    \code{Cmd::map}. The \code{textarea}, \code{picker},
    \code{table} widgets.
  \item \textbf{A reactive widget.} Uses signals to drive
    self-contained animations or watched state. The
    \code{spinner} (running variant), \code{progress},
    \code{streaming\_cursor}.
\end{itemize}

The three layer cleanly: a pure renderer can be wrapped by a
stateful component, which can be wrapped by a reactive widget.
The \code{Element} type at the boundary lets them compose.

\section{The pure-renderer pattern}

Smallest possible widget:

\begin{lstlisting}
struct Badge {
    std::string label;
    Color fg = Color::default_();
    Color bg = Color::default_();

    Badge& with_label(std::string s) { label = std::move(s); return *this; }
    Badge& with_fg(Color c)         { fg = c; return *this; }
    Badge& with_bg(Color c)         { bg = c; return *this; }

    Element build() && {
        return text(\"[\" + label + \"]\")
            | Style{.fg = fg, .bg = bg, .attributes = Attribute::Bold};
    }
    operator Element() && { return std::move(*this).build(); }
};

Badge badge(std::string label = \"\") {
    return Badge{.label = std::move(label)};
}
\end{lstlisting}

Three setters, one builder, one conversion. The \code{build}
produces a styled text element. Trivially composable in the DSL:

\begin{lstlisting}
auto bar = h(
    widget(badge(\"NORMAL\").with_bg(Color::rgb(0,80,150))),
    text(\" file.txt\")
);
\end{lstlisting}

\subsection*{Pure renderer characteristics}

\begin{itemize}[leftmargin=*]
  \item Zero allocations beyond the produced element.
  \item No subscription requirements.
  \item Trivially testable: call \code{build}, inspect the result.
  \item No memory between frames; each \code{build} is fresh.
\end{itemize}

Use the pure-renderer pattern for any widget that is just a
visual arrangement. Badges, dividers, callouts, tags, labels,
chips, breadcrumbs.

\section{The stateful-component pattern}

A widget with its own state, message, update:

\begin{lstlisting}
namespace textarea {

struct Model {
    std::vector<std::string> lines;
    size_t cursor_row = 0, cursor_col = 0;
    size_t scroll_top = 0;
    bool focused = false;
};

struct InsertChar { char c; };
struct Backspace {};
struct CursorMove { int dr, dc; };
struct GainFocus {};
struct LoseFocus {};
using Msg = std::variant<InsertChar, Backspace, CursorMove,
                         GainFocus, LoseFocus>;

std::pair<Model, Cmd<Msg>> update(Model m, Msg msg);
Element view(const Model& m, int width, int height);
Sub<Msg> subscribe(const Model& m);

}
\end{lstlisting}

This is the same shape as a \maya{} application! That is the
point: widgets are sub-apps. Composing them is exactly the
parent/child pattern from Chapter~\ref{ch:elm}.

\subsection*{Parent integration}

\begin{lstlisting}
struct App {
    struct Model { textarea::Model editor; };
    struct ToEditor { textarea::Msg m; };
    using Msg = std::variant<ToEditor, Quit>;

    static std::pair<Model, Cmd<Msg>>
    update(Model m, Msg msg) {
        return std::visit(overload{
            [&](ToEditor te) {
                auto [e2, ecmd] = textarea::update(m.editor, te.m);
                m.editor = std::move(e2);
                return std::pair{m,
                    map(ecmd, [](auto x) -> Msg { return ToEditor{x}; })};
            },
            [&](Quit) { return std::pair{m, Cmd<Msg>::quit()}; }
        }, msg);
    }

    static Element view(const Model& m) {
        return textarea::view(m.editor, 80, 24);
    }

    static Sub<Msg> subscribe(const Model& m) {
        return Sub<Msg>::batch({
            Sub<Msg>::map(textarea::subscribe(m.editor),
                          [](auto x) -> Msg { return ToEditor{x}; }),
            // ... app-level subscriptions
        });
    }
};
\end{lstlisting}

The familiar \code{map} threading. Widget messages stay in the
widget's namespace; the parent's \code{Msg} variant has one
alternative per widget instance.

\section{The reactive widget pattern}

A widget that owns a signal for self-driven updates:

\begin{lstlisting}
class Spinner {
public:
    Spinner();
    Element build() &&;

private:
    Signal<int> frame_{0};
    // The signal advances on each animation frame.
};

Spinner::Spinner() {
    // Set up an animation-frame subscriber.
    AnimationFrame{[this] { frame_.set(frame_() + 1); }};
}

Element Spinner::build() && {
    constexpr std::array glyphs = {\"|\", \"/\", \"-\", \"\\\\\"};
    return text(glyphs[frame_() % glyphs.size()]);
}
\end{lstlisting}

The widget owns its signal; the build reads it; the animation
runs automatically. No parent integration needed for the
animation itself.

\subsection*{When reactive widgets shine}

Reactive widgets are best for:

\begin{itemize}[leftmargin=*]
  \item Decorative elements (spinners, pulsing cursors).
  \item Independent live values (clock, frame counter).
  \item Watched filesystem or network state.
\end{itemize}

For business logic, prefer the stateful-component pattern; the
reactive pattern's hidden state hurts testability.

\section{Widget composition}

Widgets compose three ways:

\subsection*{Vertical composition}

A widget contains another widget:

\begin{lstlisting}
class Modal {
    Element content;
    std::string title;

    Element build() && {
        return box(
            v(text(title) | Bold,
              std::move(content)))
            | border<Heavy> | pad<2>;
    }
};
\end{lstlisting}

The outer widget takes the inner as a parameter and embeds it.

\subsection*{Horizontal composition}

Several widgets share a frame:

\begin{lstlisting}
auto bar = h(
    widget(Badge{}.with_label(\"NORMAL\")),
    text(\" \") | grow<1>,
    widget(Clock{})
);
\end{lstlisting}

The DSL composes them; each widget contributes a sub-tree.

\subsection*{Decoration}

A widget wraps another with style:

\begin{lstlisting}
class Panel {
    Element content;
    std::string title;

    Element build() && {
        return box(
            v(text(title) | Bold,
              std::move(content)))
            | border<Rounded> | pad<1>;
    }
};
\end{lstlisting}

Same structural shape as Modal; different decoration. The pattern
recurs.

\section{Theming widgets}

A widget should use roles, not raw colours. The badge example
above takes \code{fg} and \code{bg} as parameters; a
theme-aware variant would default to roles:

\begin{lstlisting}
Element Badge::build() && {
    auto& theme = current_theme();
    Color f = fg.is_default() ? theme.palette[Role::Accent] : fg;
    Color b = bg.is_default() ? theme.palette[Role::Background] : bg;
    return text(\"[\" + label + \"]\")
        | Style{.fg = f, .bg = b, .attributes = Attribute::Bold};
}
\end{lstlisting}

The widget asks for the current theme's accent colour if the
caller did not specify. Theme swaps update the widget for free.

\section{Widget testing}

\subsection*{Pure renderers}

\begin{lstlisting}
TEST_CASE(\"badge renders with brackets\") {
    auto e = badge(\"hi\").with_fg(Color::rgb(0,0,0)).build();
    REQUIRE(text_view(e) == \"[hi]\");
}
\end{lstlisting}

Direct. No setup.

\subsection*{Stateful components}

Same patterns as Chapter~\ref{ch:testing}. The widget's
\code{update} is pure; test it with \code{fold}.

\subsection*{Reactive widgets}

Use \code{test::FakeClock} and a controlled signal driver.
Advance the clock, assert against \code{build()} output.

\section{When to write a widget}

Not every visual snippet needs a widget. A rule:

\begin{itemize}[leftmargin=*]
  \item One-off layout: write it inline.
  \item Reused in three places: extract a function.
  \item Reused in five places, or has state: write a widget.
  \item Public API for other apps: write a widget with docs.
\end{itemize}

Premature widget abstraction costs more than it saves. Wait
until you see the duplication.

\begin{warning}
A widget that takes too many parameters (more than five or
six) is a sign that it is doing two jobs. Split it.
\end{warning}

\section*{Workshop}

\begin{enumerate}[leftmargin=*]
  \item Write a \code{kbd} widget that renders a keyboard
    shortcut (\verb|[Ctrl-X]|) with consistent styling.
  \item Write a \code{toast} stateful component with an
    \code{auto-dismiss} timer (\code{Cmd::after(3s, Dismiss)}).
  \item Read three widgets in \file{maya/include/maya/widget/}
    and classify each as pure, stateful, or reactive.
\end{enumerate}

\begin{recap}
A widget is a pure renderer, a stateful component, or a
reactive widget. Each pattern has a small ceremony; the
\code{Element} type at the boundary lets them compose.
Stateful widgets are sub-apps and integrate via \code{Cmd::map}.
Reactive widgets own signals for self-driven updates. Theme-
awareness is opt-in but cheap. Premature abstraction is the
real enemy.
\end{recap}

\chapter{Patterns: the architecture in practice}
\label{ch:patterns}

\begin{aside}
The architecture chapters showed primitives. Real apps combine
the primitives in patterns that recur across many projects. This
chapter is a tour: ten common patterns, each with the
architectural shape, the trade-offs, and the code skeleton.
\end{aside}

\section{Pattern 1: master-detail}

\textbf{Shape.} A list on the left, a detail view on the right.
Selecting an item updates the detail.

\textbf{Model.}

\begin{lstlisting}
struct Model {
    std::vector<Item> items;
    int selected = 0;
};
\end{lstlisting}

\textbf{View.}

\begin{lstlisting}
Element view(const Model& m) {
    auto list = render_list(m.items, m.selected);
    auto detail = m.items.empty()
        ? text(\"(no items)\") | Dim
        : render_detail(m.items[m.selected]);
    return h(list | width<30>, detail | grow<1>);
}
\end{lstlisting}

\textbf{Update.} Arrow keys move \code{selected}; Enter triggers
a detail-specific action.

\textbf{Trade-offs.}

\begin{itemize}[leftmargin=*]
  \item Simple to write; works for small datasets.
  \item Re-renders the detail on every selection change; OK as
    long as detail rendering is cheap.
\end{itemize}

\section{Pattern 2: tabbed pages}

\textbf{Shape.} Several distinct pages, switched by Tab or by
keystrokes. Each page has its own state.

\textbf{Model.}

\begin{lstlisting}
struct Model {
    int active_tab = 0;
    HomePage home;
    InboxPage inbox;
    SettingsPage settings;
};
\end{lstlisting}

\textbf{View.}

\begin{lstlisting}
Element view(const Model& m) {
    return v(
        render_tab_bar(m.active_tab),
        switch_tab(m.active_tab,
            home_view(m.home),
            inbox_view(m.inbox),
            settings_view(m.settings))
    );
}
\end{lstlisting}

\textbf{Trade-offs.}

\begin{itemize}[leftmargin=*]
  \item Each page keeps state across tab switches.
  \item Memory cost: all pages' models live even when not
    visible.
  \item Pages can be tested independently.
\end{itemize}

\section{Pattern 3: wizard / stepper}

\textbf{Shape.} Several screens in sequence; each collects input,
the final screen submits.

\textbf{Model.}

\begin{lstlisting}
enum class Step { Welcome, Name, Email, Confirm, Done };
struct Model {
    Step step = Step::Welcome;
    std::string name, email;
};
\end{lstlisting}

\textbf{Update.} Next and Back transitions step the enum;
typing fills fields.

\textbf{Trade-offs.}

\begin{itemize}[leftmargin=*]
  \item Clean linear flow; easy to test (one path per step).
  \item Adding a step is local: enum value, view branch, update
    branch.
\end{itemize}

\section{Pattern 4: undo/redo}

\textbf{Shape.} Every state-changing message creates a new model
that is pushed onto an undo stack.

\textbf{Model.}

\begin{lstlisting}
struct Doc { /* the actual document state */ };
struct Model {
    Doc current;
    std::vector<Doc> undo_stack;
    std::vector<Doc> redo_stack;
};
\end{lstlisting}

\textbf{Update.} Each mutation pushes the previous \code{current}
to \code{undo\_stack}; an \code{Undo} message pops back, pushing
the now-current state to \code{redo\_stack}.

\textbf{Trade-offs.}

\begin{itemize}[leftmargin=*]
  \item Memory grows with the undo history. Cap the stack length
    to bound memory.
  \item Trivially correct because models are values --- no
    fiddly ``do I have a deep copy?''.
\end{itemize}

\section{Pattern 5: incremental search}

\textbf{Shape.} A search box and a filtered result list; results
update on every keystroke without blocking.

\textbf{Model.}

\begin{lstlisting}
struct Model {
    std::string query;
    std::vector<Item> all_items;
    std::vector<size_t> filtered_indices;
    int debounce_token = 0;
};
\end{lstlisting}

\textbf{Update.} A keystroke updates \code{query} and issues
\code{Cmd::after(50ms, Filter\{debounce++\})}. The filter, when
it fires, checks whether \code{debounce\_token} matches; if not,
the keystroke fired again and the filter is stale.

\textbf{Trade-offs.}

\begin{itemize}[leftmargin=*]
  \item Debouncing avoids re-filtering on every key.
  \item Stale-result rejection avoids race-y UI flashes.
  \item More complex \code{update} logic; tests are essential.
\end{itemize}

\section{Pattern 6: side-effecting buttons}

\textbf{Shape.} A button that, when clicked, performs an effect
(network call, file write) and updates UI based on the result.

\textbf{Model.}

\begin{lstlisting}
enum class Status { Idle, Pending, Ok, Failed };
struct Model { Status status = Status::Idle; std::string error; };
\end{lstlisting}

\textbf{Update.} On click, set status to \code{Pending}, issue
\code{Cmd::task}. On result, transition to \code{Ok} or
\code{Failed}.

\textbf{View.} The button shows ``Save'', ``Saving\ldots'',
``Saved!'', or ``Failed (retry)'' depending on status.

\textbf{Trade-offs.}

\begin{itemize}[leftmargin=*]
  \item User sees feedback immediately.
  \item No blocking on the main thread.
  \item Status enum makes the four states explicit; no
    \code{is\_loading} + \code{is\_error} bool combinations.
\end{itemize}

\section{Pattern 7: long-running stream}

\textbf{Shape.} A widget consumes a stream of events (log lines,
tokens from an LLM, network packets) and appends them to a
display.

\textbf{Model.}

\begin{lstlisting}
struct Model {
    std::vector<std::string> lines;
    ScrollState scroll;
    bool autoscroll = true;
};
\end{lstlisting}

\textbf{Update.} Each incoming line via \code{Cmd::task}
(reading from a pipe) is appended. If \code{autoscroll}, adjust
the offset to follow.

\textbf{Trade-offs.}

\begin{itemize}[leftmargin=*]
  \item Use virtualisation (Chapter~\ref{ch:scroll}) for
    large logs.
  \item Cap the vector length to bound memory.
  \item Autoscroll-while-following is a UX pattern users love;
    pair with ``stop following on manual scroll''.
\end{itemize}

\section{Pattern 8: keyboard-driven palette}

\textbf{Shape.} A command palette (Ctrl-P style): a text field, a
filtered list, Enter to execute the selected command.

\textbf{Model.}

\begin{lstlisting}
struct Command { std::string label; std::function<Cmd<Msg>()> action; };
struct Model {
    bool open = false;
    std::string query;
    int selected = 0;
    std::vector<Command> commands;
};
\end{lstlisting}

\textbf{Update.} Open/close on Ctrl-P; arrow keys move
selection; Enter dispatches the command's action.

\textbf{Trade-offs.}

\begin{itemize}[leftmargin=*]
  \item Commands are first-class values; new ones add as a
    push\_back.
  \item Filter is a simple substring match; full-text fuzzy
    match is a deeper rabbit hole.
\end{itemize}

\section{Pattern 9: split-pane editor}

\textbf{Shape.} Two editors side by side; focus moves between
them with a hotkey.

\textbf{Model.}

\begin{lstlisting}
struct Model {
    textarea::Model left;
    textarea::Model right;
    bool focus_right = false;
};
\end{lstlisting}

\textbf{Update.} Keystrokes route to the focused editor; Tab
toggles focus.

\textbf{Trade-offs.}

\begin{itemize}[leftmargin=*]
  \item Two editors compose by message-tagging
    (\code{ToLeft}/\code{ToRight}).
  \item Resizing splits the width evenly; flex handles it.
\end{itemize}

\section{Pattern 10: themed multi-screen app}

\textbf{Shape.} A full app combining several of the patterns
above: tabs at the top, a sidebar, a main area with master-
detail, a status bar, a palette overlay.

\textbf{Composition.}

\begin{lstlisting}
struct App {
    struct Model {
        Theme theme;
        int active_tab;
        SidebarModel sidebar;
        MasterDetailModel main;
        StatusModel status;
        std::optional<PaletteModel> palette;
    };
    // ... Msg with one alt per subsystem
    // ... update routes by alt
    // ... view composes the layout
    // ... subscribe collects every subsystem's subs
};
\end{lstlisting}

\textbf{Trade-offs.} The model is a struct of structs; the message
is a variant of variants; the update is a visitor that routes
to subsystems. The complexity is real but bounded: each
subsystem stays small and testable, and the wiring is
mechanical.

\section{When patterns conflict}

Two patterns might pull in opposite directions. Examples:

\begin{itemize}[leftmargin=*]
  \item \textbf{Undo + long-running stream.} A growing log
    polluting the undo stack is a memory leak. Skip
    appended-log messages from the undo recording.
  \item \textbf{Tabs + palette.} The palette is a modal; while
    open, tab switching is disabled. Handle modality at the
    top-level update.
  \item \textbf{Wizard + autosave.} A wizard step that goes
    back must restore the prior step's input. Save on every
    transition.
\end{itemize}

Patterns are tools, not rules. Read the situation; pick the
right combination.

\section{Anti-patterns}

\begin{warning}
\textbf{God model.} A single \code{Model} with hundreds of
fields. Symptom: \code{update} grows to thousands of lines.
Refactor by splitting into sub-models, one per concern.
\end{warning}

\begin{warning}
\textbf{Message explosion.} A \code{Msg} variant with hundreds
of alternatives. Symptom: every keystroke routes through a
massive visitor. Use sub-component patterns to push detail into
sub-modules.
\end{warning}

\begin{warning}
\textbf{Imperative inside view.} \code{view} calls a function
that performs file I/O. Symptom: the program hangs on render.
Move the I/O into a \code{Cmd::task}.
\end{warning}

\begin{warning}
\textbf{Subscribe of doom.} A monolithic \code{subscribe} with
dozens of always-on subscriptions. Symptom: removing one event
source causes mysterious freezes. Make subscriptions conditional
on relevant model state.
\end{warning}

\section*{Workshop}

\begin{enumerate}[leftmargin=*]
  \item Pick one of the ten patterns and build a small app
    around it.
  \item Combine two patterns (tabs + master-detail) in one
    app. Note where the patterns share state and where they
    are independent.
  \item Take a god-model from an old codebase and refactor it
    into three sub-models. The new \code{update} routes by
    message variant.
\end{enumerate}

\begin{recap}
Patterns are recurring combinations of the architectural
primitives. Master-detail, tabbed pages, wizards, undo, search,
side-effecting buttons, log streams, palettes, split panes,
multi-screen apps --- each combines a model shape, an update
strategy, and a view. Pick patterns intentionally; combine them
carefully; avoid the anti-patterns.
\end{recap}

\chapter{Performance, profiling, and the budget}
\label{ch:perf}

\begin{aside}
A TUI feels good when frames land in time. A 60Hz cycle gives
16.67ms; a 30Hz cycle 33.3ms. Most of that budget is the
terminal's own latency. \maya{}'s work fits in microseconds, but
only if you stay in the lanes the architecture invites. This
chapter is the budget: where the cycles go, how to measure,
how to fix what you find.
\end{aside}

\section{The frame budget}

A modern terminal commits a frame when:

\begin{enumerate}[leftmargin=*]
  \item Your program writes bytes to stdout.
  \item The terminal's render loop wakes (typically every
    frame).
  \item The GPU draws.
\end{enumerate}

The terminal contributes most of the latency: 5--15 ms is
normal. \maya{}'s pipeline runs in well under 1 ms. If your
frame total exceeds 16 ms, the work is almost certainly yours
(in \code{update} or \code{view}), not the library's.

\subsection*{Where the budget breaks down}

For a typical UI:

\begin{itemize}[leftmargin=*]
  \item \code{view}: 100--500 \textmu s.
  \item Layout: 20--100 \textmu s.
  \item Paint: 5--50 \textmu s.
  \item Diff: 1--10 \textmu s.
  \item Serialise: 5--50 \textmu s.
  \item Write to stdout: 50--200 \textmu s.
  \item Terminal latency: 5--15 ms.
\end{itemize}

Your work is the 100--500 \textmu s in \code{view}. The rest is
fixed cost.

\begin{idea}
The biggest single optimisation for most apps is moving
expensive computation out of \code{view} and into \code{update}
(cached in the model). \code{view} runs every frame; \code{update}
runs once per relevant message.
\end{idea}

\section{Profiling the right way}

Three measurement layers, from coarse to fine.

\subsection*{The pipeline stats}

\maya{}'s pipeline exposes per-frame timings via the \code{Stats}
struct. Enable the overlay:

\begin{lstlisting}
maya::config().show_render_stats = true;
\end{lstlisting}

A small widget appears in a corner with \code{view\_time},
\code{layout\_time}, \code{diff\_time}, \code{serialize\_time},
and frame total. Watch which number changes when you do
something.

\subsection*{Sampling profiler}

For deeper analysis, run under \code{perf} (Linux) or
\code{Instruments} (macOS):

\begin{lstlisting}
perf record -g ./my_maya_app
perf report
\end{lstlisting}

The flamegraph reveals which functions consume CPU. \code{view}
calls dominate; below them, the cost lives in your business
logic or in \code{element::build}.

\subsection*{Per-frame log}

For occasional spikes, log per-frame timings to a file:

\begin{lstlisting}
maya::config().log_frame_times = \"/tmp/frames.log\";
\end{lstlisting}

The log writes one line per frame with the timings. Grep for
outliers; correlate with what the user was doing.

\section{Optimisation playbook}

When you find a slow frame, the playbook in order:

\subsection*{1. Move work to update}

If \code{view} parses a string, formats a number, or walks a
big list, do it once in \code{update} and cache the result:

\begin{lstlisting}
// Bad: parses every frame.
Element view(const Model& m) {
    auto parsed = parse_huge_thing(m.input);
    return render(parsed);
}

// Good: cached in the model.
struct Model {
    std::string input;
    ParsedThing cached;  // refreshed in update when input changes
};
\end{lstlisting}

\subsection*{2. Virtualise long lists}

If \code{view} renders thousands of items, only render the
visible slice (Chapter~\ref{ch:scroll}). The layout pass scales
with the rendered count, not the source size.

\subsection*{3. Use compile-time DSL where possible}

A static structure built at compile time has fewer allocations
than the same structure built at runtime. For panels with fixed
shape, the DSL pays for itself.

\subsection*{4. Reuse buffers}

If a widget allocates a large vector every frame, capture the
vector in the widget's state and reset (\code{vec.clear()})
instead of reconstructing. Capacity persists; allocations
disappear.

\subsection*{5. Reduce allocations in update}

Update is hot. Every \code{std::string} concatenation allocates;
every \code{std::vector} push that exceeds capacity allocates.
Profile, then move construction outside hot loops.

\subsection*{6. Cache derived state}

Computed values that depend on slow inputs deserve a cache.
\code{Signal} + \code{Computed} (Chapter~\ref{ch:reactivity-preview})
gives you automatic recomputation only on change.

\subsection*{7. Tune the cache fingerprint}

If your tree is huge but mostly stable, the fingerprint walk is
the bottleneck. Two options: split the tree into smaller
sub-trees (each fingerprinted independently) or mark stable
sub-trees as \emph{frozen} (skipped in the walk).

\section{Memory budget}

A \maya{} app's resident memory:

\begin{itemize}[leftmargin=*]
  \item Model: usually small (KB). Watch for large vectors
    that grow unbounded (logs, history).
  \item Element tree: KB to low MB. Freshly allocated each
    frame; garbage-collected by RAII at frame end.
  \item Canvas: width $\times$ height $\times$ cell\_size. For
    80x24, that is 1920 cells $\times$ 8--16 bytes = a few KB
    each (current + previous = two copies).
  \item Style pool: a few KB.
  \item Renderer buffers: KB.
\end{itemize}

A small \maya{} app sits at 5--20 MB resident, dominated by the
binary itself. If you see hundreds of MB, look for unbounded
vectors in the model.

\subsection*{Allocations per frame}

A well-tuned app makes a handful of allocations per frame:

\begin{itemize}[leftmargin=*]
  \item Element tree allocations: a vector per box,
    \code{string\_view}s for literals.
  \item Layout vector resize (cap reused).
  \item Render buffer resize (cap reused).
\end{itemize}

\code{strace -c -e brk,mmap} on a running app shows the rate.
If it climbs into the thousands per second, hunt the
allocation sources.

\section{Latency vs throughput}

The pipeline optimises for \emph{latency} (frame-to-frame
time), not throughput (frames per second). The trade-offs:

\begin{itemize}[leftmargin=*]
  \item Low-priority diff: skip if last frame was less than
    16ms ago. Reduces CPU on busy apps; adds latency.
  \item Triple buffering: render frame N+2 while frame N is on
    the wire. Adds latency; increases throughput.
\end{itemize}

\maya{}'s defaults favour latency. If you have a soft-realtime
requirement (audio visualiser, game), tune the renderer
explicitly.

\section{Common slowdowns}

\subsection*{Reading large files in view}

The classic. \code{view} opens a file, reads it, formats it,
returns the element. Every frame. The file is re-read 60 times
a second. Cache the contents in the model and refresh only when
the file changes (use \code{Cmd::task} with a file-watch).

\subsection*{Heavy regex in view}

Regex compilation and matching are slow. Cache the compiled
pattern; cache the match results.

\subsection*{Recursive view functions}

A \code{view} that calls itself deeply on a recursive data
structure can blow stack and time. Convert to iterative;
materialise levels separately.

\subsection*{Mouse events firing on every move}

\code{OnMouse} fires for every motion event, which can be
thousands per second on a fast trackpad. Throttle the message
delivery (\code{Sub::throttle(30ms, OnMouse)}) or filter at the
handler.

\subsection*{Style pool exhaustion}

The style pool interns styles to small IDs. If you generate
thousands of distinct styles (e.g. one per character of a
gradient), the pool grows and dispatches slow down. Reuse a
few common styles where possible.

\section{When to give up tuning}

If you have:

\begin{itemize}[leftmargin=*]
  \item Frame time under 5 ms.
  \item No visible jank to the user.
  \item No CPU usage during idle.
\end{itemize}

\ldots{} you are done. Further optimisation is anti-economic.

\maya{}'s architecture lets you stop tuning earlier than other
frameworks because the defaults are aggressive: cache
fingerprints, SIMD diff, batched serialise. Profile only the
remaining hot spots in your own code.

\begin{aside}
Donald Knuth's ``premature optimisation is the root of all evil''
is often misquoted to argue against any optimisation. The full
quote ends ``yet we should not pass up our opportunities in that
critical 3\%''. Identify the 3\%; tune it; leave the rest.
\end{aside}

\section{Benchmarking widgets}

For widgets you ship to others, include a benchmark:

\begin{lstlisting}
BENCHMARK(\"render textarea with 1000 lines\") {
    Model m = make_model_with_lines(1000);
    return view(m);
}
\end{lstlisting}

The benchmark suite measures across compilers, terminal sizes,
and content. Track regressions over time.

\section*{Workshop}

\begin{enumerate}[leftmargin=*]
  \item Enable \code{show\_render\_stats} on a small app.
    Note the per-stage timings. Identify which stage dominates.
  \item Add a deliberately slow function to \code{view} (e.g.
    \code{std::this\_thread::sleep\_for(10ms)}). Watch the
    frame time rise.
  \item Profile a real app under \code{perf}. Identify the top
    three functions; fix the slowest one.
\end{enumerate}

\begin{recap}
The pipeline runs in microseconds; user work is the variable
cost. Move expensive computation from \code{view} to
\code{update} with caching. Virtualise long lists. Reuse
buffers. Avoid recomputing what has not changed. Memory budget
is small. The defaults are aggressive; tune only the residual
hot spots.
\end{recap}

\chapter{The DSL's static guarantees, in detail}
\label{ch:dsl-guarantees}

\begin{aside}
Chapter~\ref{ch:dsl} introduced the DSL's type-state machine.
This chapter zooms in on the specific guarantees: what the DSL
prevents at compile time, what it allows at runtime, and how
each guarantee maps to a constraint or concept. Treat it as a
reference for ``what can I assume about my code if it
compiled?''.
\end{aside}

\section{Guarantee 1: borderless boxes cannot have border colour}

\begin{lstlisting}
auto a = box(text(\"x\"));                       // no border
auto b = a | bcol<60, 65, 80>;                  // compile error
\end{lstlisting}

Mechanism: the \code{operator|(BoxNode<Cfg>, bcol\_t)} has a
\code{requires (Cfg.has\_border)} clause. Without a border, the
clause fails.

Why it matters: a border colour without a border produces no
visible effect at runtime --- a silent bug. The compile error
makes the intent explicit.

\section{Guarantee 2: padding is non-negative}

\begin{lstlisting}
auto a = text(\"x\") | pad<-1>;                  // compile error
\end{lstlisting}

Mechanism: the \code{pad} template uses an NTTP constrained to
non-negative integers via a concept.

Why it matters: negative padding has no defined meaning; a
runtime crash or visual glitch would be the alternative.

\section{Guarantee 3: each pipe attribute applies at most once}

\begin{lstlisting}
auto a = text(\"x\") | width<10> | width<20>;    // compile error
\end{lstlisting}

Mechanism: the second \code{| width<20>} has a \code{requires
(!Cfg.has\_width)} clause; the first pipe set the flag, so the
second fails.

Why it matters: ``which width wins?'' would be ambiguous;
forbidding the double-pipe makes the intent explicit. If you
really want to override, build a fresh node.

\section{Guarantee 4: children must be elements}

\begin{lstlisting}
auto a = h(text(\"x\"), 42);                     // compile error
\end{lstlisting}

Mechanism: \code{h(\ldots)} requires all arguments to satisfy
\code{ElementLike}. An \code{int} does not.

Why it matters: silent conversion would produce a runtime crash
or empty element.

\section{Guarantee 5: NTTPs are constant expressions}

\begin{lstlisting}
std::string s = compute();
auto a = t<s>;                                 // compile error
\end{lstlisting}

Mechanism: \code{t<S>} takes a \code{Str} NTTP; only literals
and constexpr values qualify.

Why it matters: the compile-time path is exclusive to compile-
time strings. For runtime strings, use \code{text(s)} explicitly.

\section{Guarantee 6: style composition is associative}

\begin{lstlisting}
auto a = Bold | Italic | Underline;
auto b = (Bold | Italic) | Underline;
auto c = Bold | (Italic | Underline);
// a, b, c are the same type.
\end{lstlisting}

Mechanism: \code{StyTag<A> | StyTag<B>} returns \code{StyTag<merge(A,
B)>}, and \code{merge} is associative on the underlying flag
struct.

Why it matters: refactoring style expressions is safe; the
order of pipes does not affect the result.

\section{Guarantee 7: TextNode preserves literal storage}

\begin{lstlisting}
auto a = t<\"hello\">;                           // TextNode<Str{\"hello\"}>
// At runtime, the bytes \"hello\" are in the static string table.
// The runtime element holds a string_view into them.
\end{lstlisting}

Mechanism: the NTTP carries the literal; the runtime element
holds a \code{string\_view} that points into the static segment.

Why it matters: no allocation for compile-time text; no
copy-on-construction.

\section{Guarantee 8: Style merge precedence is deterministic}

\begin{lstlisting}
auto a = text(\"x\") | Bold;                     // Bold
auto b = box(a) | Italic;                       // Bold + Italic
// The order of Bold and Italic across the parent/child is unambiguous.
\end{lstlisting}

Mechanism: the merge function is pure and deterministic; same
inputs, same output.

Why it matters: behavioural reasoning is local; reading a pipe
chain tells you the resulting style without running the program.

\section{Guarantee 9: width-clipped text never exceeds width}

\begin{lstlisting}
auto a = text(very_long) | width<10>;
// At runtime, the rendered text is exactly 10 cells wide.
\end{lstlisting}

Mechanism: the layout phase respects the width as a hard cap.
The text renderer truncates (with optional ellipsis) when
content exceeds.

Why it matters: layout stability; widgets that allocate 10
cells get exactly 10.

\section{Guarantee 10: flex direction is one-way}

\begin{lstlisting}
auto a = h(text(\"x\")) | direction<Column>;     // ?? not allowed
\end{lstlisting}

Mechanism: \code{h(\ldots)} creates a row; the direction is
baked in. Overriding via pipe is not supported; if you want a
column, write \code{v(\ldots)}.

Why it matters: ambiguity in direction would force a runtime
choice; making it explicit at the DSL level avoids the
confusion.

\section{Guarantees the DSL does not give}

For completeness, the things the DSL does \emph{not} guarantee:

\begin{itemize}[leftmargin=*]
  \item Runtime content fits the layout. A long runtime string
    might overflow its width container.
  \item Terminal capabilities. The DSL emits truecolour even
    if the terminal does not support it; quantisation happens
    later.
  \item Sub-tree count. A vector of children can be of any
    length; the DSL does not constrain.
  \item Memory budget. A deeply nested tree can blow the
    stack; the DSL does not bound recursion.
\end{itemize}

These are runtime concerns that the DSL leaves to the
runtime. The discipline is: use the DSL for shape, use the
runtime for content.

\section{Cumulative effect}

The ten guarantees, plus a dozen more that fall out of concept
constraints, mean a \maya{} UI that compiles is structurally
sound. You cannot ship a panel with a border colour and no
border; you cannot pass an \code{int} where an element is
expected; you cannot accidentally set width twice.

This is what the type-state machine buys you: the compiler
caught the bugs before you ran the program.

\begin{idea}
A compiler that catches structural bugs is the best code review
you will ever get. The DSL turns the C++ compiler into a UI
linter, automatically and continuously. Every successful build is
a small certificate of correctness.
\end{idea}

\section{Mapping guarantees to error messages}

When a guarantee fires, the compiler emits an error. The shape
is consistent:

\begin{lstlisting}
error: no match for 'operator|' (operand types are X and Y)
note: candidate: 'requires Z ...'
note:   constraint not satisfied: 'Z evaluated to false'
\end{lstlisting}

Read top-down: ``no match'' is the call site; ``candidate'' is
the overload that almost matched; ``constraint not satisfied''
is the specific rule. Map the rule to one of the guarantees in
this chapter; the fix follows.

\section{Guarantees in widget code}

Widget builders can declare their own constraints. If a widget
requires certain configuration, encode it as a concept:

\begin{lstlisting}
template <class T>
concept HasLabel = requires (T t) { { t.label } -> std::convertible_to<std::string_view>; };

template <HasLabel T>
Element labelled(const T& t, std::string_view extra);
\end{lstlisting}

Now callers without a \code{label} field get a clear error,
not an SFINAE wall.

\section*{Workshop}

\begin{enumerate}[leftmargin=*]
  \item Pick three of the ten guarantees. For each, write code
    that violates it, read the resulting error, identify which
    \code{requires} clause caught it.
  \item Write a small widget with a concept-constrained
    parameter. Verify the error message is clear when the
    concept is violated.
  \item Read the source of a DSL pipe overload. Confirm the
    \code{requires} clause is human-readable.
\end{enumerate}

\begin{recap}
The DSL gives ten structural guarantees, all enforced at compile
time. They eliminate whole classes of bugs (silent decorations,
ambiguous widths, wrong child types). Errors map back to a
specific \code{requires} clause; the fix is mechanical. The
trade-off is some DSL complexity; the payoff is a UI that
compiles is a UI that is structurally sound.
\end{recap}

\chapter{The C++26 features maya leans on}
\label{ch:cpp26-features}

\begin{aside}
\maya{} is named ``C++26 for the terminal'' for a reason: a
handful of newer language and library features make the
architecture practical. This chapter is a tour of those features
in isolation: what each adds, how \maya{} uses it, and what the
pre-C++26 alternative would look like.
\end{aside}

\section{Concepts (C++20)}

Already covered in Chapter~\ref{ch:concepts}. Without concepts,
template error messages would be wall-of-text SFINAE diagnostics.
With concepts, the constraint is named; the failure point is
specific.

\maya{} uses concepts for: \code{Program}, \code{ElementLike},
\code{Renderable}, \code{StyleSource}, \code{Subscribable},
\code{Terminal}, \code{Addable}, etc.

\section{Designated initialisers (C++20)}

\begin{lstlisting}
Element panel{
    .v = BoxElement{
        .children = {...},
        .layout = FlexStyle{.direction = Direction::Column},
        .border = Border{.style = BorderStyle::Rounded},
    },
};
\end{lstlisting}

Designated initialisers let you construct aggregates by name,
not by position. The pre-C++20 alternative was a fluent builder
(verbose) or careful brace-init (fragile).

\maya{}'s use: \code{Style}, \code{Color}, \code{Border},
\code{FlexStyle}, \code{Padding} are all aggregates and are
typically constructed with designated init.

\section{Structural NTTPs (C++20, refined C++26)}

The DSL's headline feature. \code{Str}, \code{CTStyle},
\code{BoxCfg} are structural types that can appear as template
parameters. Without them, the DSL would have to encode every
parameter as a separate template specialisation by hand.

\maya{}'s use: every compile-time configuration --- text
literals, styles, padding values, border styles, widths, heights,
flex factors --- is carried in an NTTP.

\section{std::variant (C++17)}

The fundamental sum type. \code{Element}, \code{Cmd<Msg>},
\code{Sub<Msg>}, \code{Event} are all variants. Without them,
the same shape would require an OOP hierarchy plus virtual
dispatch.

\section{std::visit and overload (C++17)}

\code{std::visit} matches a variant against a callable. The
overload pattern (multiple lambdas inheriting into one struct)
makes the visitor readable.

\maya{}'s \code{maya::overload} (in
\file{maya/include/maya/core/overload.hpp}) is the pattern as a
named struct.

\section{std::expected (C++23)}

Value-typed error handling (Chapter~\ref{ch:expected}). Before
C++23, you used Boost.Outcome, the absl variant, or a hand-rolled
expected. The standard version simplifies portability.

\maya{} ships its own \code{Expected} that predates the standard
addition; the API is compatible.

\section{Coroutines (C++20)}

\maya{} does not currently use coroutines, but the
\code{Cmd::task} interface is coroutine-shaped: you could
\code{co\_await} on a task. Future versions may take this on.

\section{std::format (C++20)}

\begin{lstlisting}
auto label = std::format(\"Count: {}\", n);
\end{lstlisting}

A type-safe \code{printf}. Used internally by widget code that
constructs labels and tooltips.

The C++23 \code{std::print} extends \code{format} to direct
stdout output; \maya{} prefers \code{print} where supported.

\section{std::span (C++20)}

A view over a contiguous range, like \code{string\_view} but
generic. Used heavily in \maya{}'s I/O code:

\begin{lstlisting}
Expected<size_t, IoError> read(std::span<std::byte> buf);
\end{lstlisting}

No ownership, no copy. The buffer is supplied by the caller.

\section{Ranges (C++20 and C++23)}

\maya{} uses ranges for view-and-transform pipelines:

\begin{lstlisting}
auto labels = items | std::views::transform([](auto& i) { return i.label; });
\end{lstlisting}

Lazy evaluation: the transform runs as the range is consumed,
not eagerly. Useful for building element trees from large
sources.

C++23 added \code{std::views::join\_with} which the DSL uses to
intersperse separators between list items.

\section{Concepts-aware overload resolution (C++20)}

Two templates with overlapping concepts: the more-specialised
one wins. The DSL relies on this heavily; multiple pipe overloads
target the same operator and the compiler picks by the most-
specific concept.

\section{constinit and consteval (C++20)}

\code{constinit} guarantees a variable is initialised at
compile time; \code{consteval} requires a function to run only
at compile time. The DSL's \code{merge}, \code{parse\_hex},
\code{compute\_box\_cfg} are \code{consteval}; the result is
that runtime code never sees the cost.

\section{Lambda extensions (C++23, C++26)}

C++23 added attributes on lambdas; C++26 is exploring
\code{static} lambdas (no captures, no per-call storage).
\maya{}'s tight loops use lambdas; the future-version
optimisations will improve performance without source changes.

\section{Class types in NTTPs (C++20)}

The key C++20 relaxation that made compile-time strings
possible. A structural type with all-public members can be an
NTTP. Without this, \maya{}'s entire DSL would be impossible.

\section{Class template argument deduction (C++17, refined C++20)}

\begin{lstlisting}
auto e = Element{TextElement{\"hi\"}};            // deduces Element type
\end{lstlisting}

The compiler infers the template arguments from constructor
arguments. \maya{} uses CTAD throughout to keep the syntax
clean.

\section{Three-way comparison operator (C++20)}

\begin{lstlisting}
auto cmp(const Style& a, const Style& b) = default;  // <=>
\end{lstlisting}

A single defaulted spaceship operator generates all six
comparison operators. \maya{}'s structural types use this for
NTTP equality.

\section{std::array (C++11) and constexpr containers (C++20)}

The Unicode width table is a \code{constexpr std::array}.
Lookup is binary search. The whole table sits in read-only data;
no startup cost.

\section{The if-with-init form (C++17)}

\begin{lstlisting}
if (auto r = parse(input); r) {
    use(r.value());
}
\end{lstlisting}

The init clause scopes the variable to the \code{if}/\code{else}
chain. Used in every error-handling site.

\section{P2300 senders/receivers (C++26)}

A long-awaited concurrency model. \maya{} does not depend on it
yet but the \code{Cmd::task} abstraction is sender-shaped; future
versions may integrate.

\section{Reflection (C++26)}

The reflection proposal (P2996) would let \maya{} generate
visitors and serialisers automatically. The current code uses
templates and constexpr to approximate; reflection would simplify.

\section{What it costs to use C++26}

The features above require GCC 14+, Clang 17+, or MSVC v19.40+.
Older compilers cannot build \maya{}. The trade-off:

\begin{itemize}[leftmargin=*]
  \item Library is shorter (fewer workarounds).
  \item Code is easier to read (named constraints, designated
    init).
  \item Distribution is harder (newer compiler required).
\end{itemize}

For a library targeting modern terminals on modern OSes,
requiring a modern compiler is acceptable.

\section{What the pre-C++26 version would look like}

If we backported \maya{} to C++17:

\begin{itemize}[leftmargin=*]
  \item Concepts become SFINAE. Error messages explode.
  \item NTTP-based DSL becomes type-tag-based with
    template specialisations per literal.
  \item Designated init becomes builders or positional init.
  \item Variants exist but visit-with-overload is uglier.
\end{itemize}

The result would be 2--3 times as much code for the same
functionality. C++26 is the right floor.

\section*{Workshop}

\begin{enumerate}[leftmargin=*]
  \item Read \file{maya/include/maya/core/concepts.hpp} and
    list every concept by feature it depends on.
  \item Find a single function in \maya{}'s source that uses
    designated initialisers, NTTPs, ranges, and visit in one
    expression. (There are many candidates.)
  \item Try to compile \maya{} with \code{-std=c++17}. Count the
    errors. (They will be many.)
\end{enumerate}

\begin{recap}
\maya{} leans on concepts, designated initialisers, structural
NTTPs, variants, visit, expected, format, span, ranges, and
\code{consteval}. Each of these features removes a workaround
that pre-C++26 codebases live with. The cost is a modern
compiler requirement; the payoff is a small, readable library.
\end{recap}

\chapter{Comparing maya to other TUI libraries}
\label{ch:comparison}

\begin{aside}
\maya{} is one of many TUI libraries. Each has its own
opinions; \maya{}'s opinions only make sense in context. This
chapter compares \maya{} to the major alternatives: ncurses,
notcurses, tui-rs / ratatui, Bubble Tea, Textual, FTXUI, Brick.
The point is not to declare a winner but to make the design
trade-offs explicit.
\end{aside}

\section{ncurses (C, 1980s)}

The grandparent of TUI libraries. Origin: BSD's curses, extended
by GNU. The model:

\begin{itemize}[leftmargin=*]
  \item Imperative: \code{move(y, x); addstr(s); refresh();}
  \item Window-oriented: every region is a \code{WINDOW*}.
  \item Manual layout: you compute coordinates.
  \item Manual input loop: \code{wgetch} in a \code{while}.
\end{itemize}

\textbf{Strengths.} Ubiquitous, stable, available everywhere. The
``every Unix machine has it'' guarantee.

\textbf{Weaknesses.} Imperative model fights composition. No
flexbox, no themes, no rich text. Truecolour is bolted on. Modern
keyboard handling (KKP, bracketed paste) requires manual escape
parsing.

\textbf{\maya{} contrast.} Declarative \code{view}, flexbox, type-
safe DSL. The trade: \maya{} needs a modern compiler; ncurses
runs on anything.

\section{notcurses (C, 2020)}

A modern reimagining of curses by Nick Black. Targets the
absolute highest fidelity of terminal output.

\textbf{Strengths.} Pixel-perfect rendering (cell-by-cell control).
First-class support for images via sixel and the Kitty graphics
protocol. Truecolour everywhere.

\textbf{Weaknesses.} Imperative API (similar to ncurses but
modernised). Steep learning curve. C API; the C++ wrapper is
thin.

\textbf{\maya{} contrast.} \maya{} is declarative and offers
fewer rendering capabilities (no images, no sixel). The trade:
\maya{}'s code is shorter and more testable; notcurses output is
visually richer.

\section{tui-rs / ratatui (Rust, 2019-)}

The Rust ecosystem's standard. Originally by Florian Dehau as
tui-rs; the community fork ratatui continues development.

\textbf{Strengths.} Declarative widget API. Strong Rust type
safety. Active community, dozens of community widgets.

\textbf{Weaknesses.} Rust-only. Widget composition uses
\code{Box<dyn Widget>} for dynamism; \maya{}'s sum-type model
avoids the indirection. Layout is constraint-list-based, not
flexbox-based, and feels less ergonomic for nested cases.

\textbf{\maya{} contrast.} Different languages, similar
declarative spirit. \maya{}'s DSL pushes more into compile time;
ratatui pushes more into runtime composition.

\section{Bubble Tea (Go, 2020)}

By Charm. The library that introduced the Elm architecture to
many programmers via Go.

\textbf{Strengths.} The Elm architecture made approachable.
Beautiful default styles via the Lip Gloss style library.
Excellent docs and examples.

\textbf{Weaknesses.} Go's lack of generics (until recently)
forced \code{interface\{\}} in many APIs. Rendering is line-based
(reasoned about as strings), not cell-based, so subtle Unicode
width bugs creep in.

\textbf{\maya{} contrast.} Same architectural family. \maya{}
is cell-based (graphemes + widths handled at the canvas level),
type-safer, and lower-level. Bubble Tea is faster to prototype in.

\section{Textual (Python, 2021)}

By the Rich author Will McGugan. Brings web-app ergonomics to
the terminal.

\textbf{Strengths.} CSS-like styling, including selectors and
themes. Reactive properties via Python descriptors. Beautiful
widget catalogue (Rich tables, syntax highlighting).

\textbf{Weaknesses.} Python overhead is real: a Textual app can
hit 5--10\% CPU at idle. The reactive system is implicit and
sometimes surprising.

\textbf{\maya{} contrast.} Different language ecosystems. Textual
appeals to web devs; \maya{} appeals to systems devs. Textual's
themes and styles are richer; \maya{}'s performance ceiling is
higher.

\section{FTXUI (C++, 2018-)}

The closest C++ analogue to \maya{}. By Arthur Sonzogni.

\textbf{Strengths.} Header-only. Functional-style API: components
are functions returning elements. Good widget catalogue.

\textbf{Weaknesses.} Uses inheritance for some component
hierarchies. Layout is less expressive than flexbox. The
component model is hookier than the Elm architecture (state lives
in closures, not in a plain model).

\textbf{\maya{} contrast.} Same language, different design
choices. FTXUI: closure-based state, classes for some
abstractions. \maya{}: value-based model, sum types throughout.
FTXUI is easier to learn; \maya{} scales further.

\section{Brick (Haskell, 2015)}

By Jonathan Daugherty. Haskell's Elm-inspired TUI library.

\textbf{Strengths.} Haskell's type system shines: the Elm
architecture is the natural shape. Pure functions everywhere.
Layout combinators are elegant.

\textbf{Weaknesses.} Haskell ecosystem; less mainstream than C++,
Go, or Python.

\textbf{\maya{} contrast.} Same architectural family. Brick
benefits from Haskell's type inference and pattern matching;
\maya{} works around C++'s limitations with concepts and visit.

\section{Where maya fits}

\maya{} occupies a specific niche:

\begin{itemize}[leftmargin=*]
  \item C++26 modern: NTTP-based DSL, concepts everywhere.
  \item Elm architecture: pure update, side effects as data.
  \item Cell-based renderer: SIMD diff, packed cells, integer
    layout.
  \item Compile-time guarantees: shape correctness checked at
    build, not runtime.
\end{itemize}

If you are a C++ shop that wants modern declarative TUIs with
strong static guarantees and high performance, \maya{} is in
your shortlist. If you need to deploy on legacy systems, want
to use Python or Go, or need image rendering, look at the
alternatives.

\section{Migration paths}

\subsection*{From ncurses}

The architectural shift is the hard part. Stop thinking in
``move cursor here, write text''; start thinking in ``what is
the model''. The rendering details (raw mode, escape sequences)
disappear into the runtime; you keep the business logic.

\subsection*{From Bubble Tea or Brick}

Mechanical: the model/message/update/view shape is the same.
You will rewrite Go's structs as C++ structs, channels as
\code{Cmd::task}, and the type system goes from inferred to
explicit. Most of your logic translates line by line.

\subsection*{From FTXUI}

Less mechanical: closures-as-state become explicit fields in
\code{Model}. The flow of data is more visible; the trade is
some up-front design work for testability later.

\subsection*{From Textual}

A bigger jump: dynamic Python to static C++. Worth it if you
need the performance; do not bother if Textual is fast enough.

\section{Choosing wisely}

A handful of questions to ask:

\begin{itemize}[leftmargin=*]
  \item \textbf{Language?} The library should match your
    ecosystem; you will not enjoy fighting bindings.
  \item \textbf{Performance?} If you need sub-millisecond
    frame budget, prefer C++ or Rust libraries. Python and
    Go have their own ceilings.
  \item \textbf{Type safety?} If you value compile-time
    catching of structural bugs, \maya{}, ratatui, and
    Brick are good fits.
  \item \textbf{Visual fidelity?} If you need images and
    fancy rendering, notcurses leads.
  \item \textbf{Familiarity?} If your team already knows one
    library, the switching cost may dominate the benefits.
\end{itemize}

There is no objectively best library. The best is the one
whose trade-offs match your project.

\section*{Workshop}

\begin{enumerate}[leftmargin=*]
  \item Pick two libraries from this chapter. Build a
    \emph{Hello World} in each. Compare the code.
  \item Identify three patterns from
    Chapter~\ref{ch:patterns} and find how each library
    expresses them.
  \item Choose a real project. List the requirements, then
    pick the library best suited to them. Defend the choice.
\end{enumerate}

\begin{recap}
\maya{} is one TUI library among many. ncurses is universal but
imperative; notcurses is high-fidelity; ratatui is the Rust
standard; Bubble Tea is the Elm story in Go; Textual is the
Python web-ish option; FTXUI is the C++ closure-based option;
Brick is the Haskell idiom. \maya{}'s position: modern C++,
Elm architecture, compile-time guarantees, high performance.
Pick by trade-off.
\end{recap}

\chapter{Putting it all together: an end-to-end app}
\label{ch:end-to-end}

\begin{aside}
Eighteen chapters of architecture. Now we synthesise. This
chapter builds an end-to-end app --- a small task manager ---
that exercises every concept from Part II. It is the test of
whether the model holds together. Spoiler: it does.
\end{aside}

\section{Specification}

The app: a terminal task manager that:

\begin{itemize}[leftmargin=*]
  \item Loads tasks from a JSON file on startup.
  \item Lets the user add, complete, delete, and edit tasks.
  \item Shows a list on the left, a detail panel on the right.
  \item Filters by status (all / open / done).
  \item Auto-saves to disk after each change.
  \item Supports undo with Ctrl-Z.
  \item Has a help overlay on \code{?}.
\end{itemize}

Six features. Each maps onto a Part II concept.

\section{The model}

\begin{lstlisting}
struct Task {
    int id;
    std::string title;
    std::string notes;
    bool done = false;
    std::chrono::system_clock::time_point created;
};

enum class Filter { All, Open, Done };

enum class Mode { Browsing, Editing, Modal };

struct Model {
    std::vector<Task> tasks;
    int selected = 0;
    Filter filter = Filter::All;
    Mode mode = Mode::Browsing;
    std::optional<int> editing_id;
    std::string edit_buffer;
    std::vector<Model> undo_stack;     // pattern 4
    std::string error;
    bool dirty = false;                 // unsaved
};
\end{lstlisting}

The model is data: vectors, enums, primitives. No pointers, no
shared state, no hidden caches.

\section{The message type}

\begin{lstlisting}
struct Loaded { std::vector<Task> tasks; };
struct LoadFailed { std::string error; };
struct Saved {};
struct SaveFailed { std::string error; };
struct AddTask {};
struct CompleteSelected {};
struct DeleteSelected {};
struct StartEditing {};
struct EditChar { char c; };
struct EditBackspace {};
struct EditCommit {};
struct EditCancel {};
struct CursorUp {};
struct CursorDown {};
struct ChangeFilter { Filter f; };
struct ToggleHelp {};
struct Undo {};
struct Quit {};
using Msg = std::variant<
    Loaded, LoadFailed, Saved, SaveFailed,
    AddTask, CompleteSelected, DeleteSelected,
    StartEditing, EditChar, EditBackspace, EditCommit, EditCancel,
    CursorUp, CursorDown, ChangeFilter, ToggleHelp, Undo, Quit>;
\end{lstlisting}

Eighteen alternatives. Each captures one user intent or system
event. The visitor in \code{update} will have one branch per
alternative; the compiler enforces exhaustiveness.

\section{init}

\begin{lstlisting}
static std::pair<Model, Cmd<Msg>> init() {
    return {
        Model{},
        Cmd<Msg>::task([]() -> Msg {
            auto r = file::read(\"tasks.json\");
            if (!r) return LoadFailed{r.error().context};
            try {
                return Loaded{parse_tasks(r.value())};
            } catch (const std::exception& e) {
                return LoadFailed{e.what()};
            }
        })
    };
}
\end{lstlisting}

The model starts empty; the initial \code{Cmd} kicks off a
file read. The result becomes a \code{Loaded} or \code{LoadFailed}
message. Errors as values (Chapter~\ref{ch:expected}).

\section{update: a tour of the architecture}

A few representative branches.

\subsection*{Loaded: replacing the model}

\begin{lstlisting}
[&](Loaded l) {
    m.tasks = std::move(l.tasks);
    m.selected = 0;
    m.error = \"\";
    return std::pair{m, Cmd<Msg>{}};
}
\end{lstlisting}

Pure replacement. No effect.

\subsection*{AddTask: undo-recorded mutation}

\begin{lstlisting}
[&](AddTask) {
    m.undo_stack.push_back(m);       // pattern 4
    m.tasks.push_back(Task{
        .id = next_id(m.tasks),
        .title = \"New task\",
        .created = std::chrono::system_clock::now()
    });
    m.selected = m.tasks.size() - 1;
    m.dirty = true;
    return std::pair{m, save_cmd(m)};
}
\end{lstlisting}

Three steps: snapshot for undo, mutate, queue a save. The save
command is a \code{Cmd::task} that returns \code{Saved} or
\code{SaveFailed}.

\subsection*{Undo: pattern 4 in action}

\begin{lstlisting}
[&](Undo) {
    if (m.undo_stack.empty()) return std::pair{m, Cmd<Msg>{}};
    auto prev = std::move(m.undo_stack.back());
    m.undo_stack.pop_back();
    prev.undo_stack = std::move(m.undo_stack);  // preserve history
    prev.dirty = true;
    return std::pair{std::move(prev), save_cmd(prev)};
}
\end{lstlisting}

Pop the stack; restore the previous model; queue a save.

\subsection*{EditChar: mode-aware routing}

\begin{lstlisting}
[&](EditChar e) {
    if (m.mode != Mode::Editing) return std::pair{m, Cmd<Msg>{}};
    m.edit_buffer += e.c;
    return std::pair{m, Cmd<Msg>{}};
}
\end{lstlisting}

Mode gates the message: a typed character is meaningful only in
edit mode. Outside of it, ignored.

\subsection*{ToggleHelp: modal stack}

\begin{lstlisting}
[&](ToggleHelp) {
    m.mode = (m.mode == Mode::Modal) ? Mode::Browsing : Mode::Modal;
    return std::pair{m, Cmd<Msg>{}};
}
\end{lstlisting}

A modal is a mode (sorry). The view consults \code{mode} to
draw the help overlay.

\section{view: composing the patterns}

\begin{lstlisting}
static Element view(const Model& m) {
    auto list = render_task_list(m);
    auto detail = render_detail(m);
    auto status = render_status_bar(m);

    auto main = h(list | width<40>, detail | grow<1>);
    auto root = v(main | grow<1>, status | height<1>);

    if (m.mode == Mode::Modal) {
        return overlay(root, render_help());
    }
    return root;
}
\end{lstlisting}

The structure: master-detail (pattern 1) inside a tabbed layout.
The status bar is a footer. The modal overlays everything when
help is open.

\subsection*{render\_task\_list}

\begin{lstlisting}
Element render_task_list(const Model& m) {
    std::vector<Element> rows;
    int idx = 0;
    for (auto& t : m.tasks) {
        if (!filter_match(t, m.filter)) continue;
        bool selected = (idx == m.selected);
        rows.push_back(h(
            text(t.done ? \"[x] \" : \"[ ] \"),
            text(t.title) | (selected ? Bold : Style{})
        ) | (selected ? Bg<60, 65, 80> : Style{}));
        idx++;
    }
    return box(v(rows)) | border<Single>;
}
\end{lstlisting}

A runtime loop builds the rows; the DSL wraps each. The
selected row is highlighted.

\subsection*{Modal overlay}

\begin{lstlisting}
Element overlay(Element bg, Element fg) {
    return z(
        bg | Dim,
        fg | center<Vertical> | center<Horizontal>
    );
}
\end{lstlisting}

\code{z} is a Z-axis stack: children render at the same
position; the last one wins on conflict. Used for tooltips,
modals, dropdowns.

\section{subscribe: pattern by mode}

\begin{lstlisting}
static Sub<Msg> subscribe(const Model& m) {
    auto keys = Sub<Msg>::on_key([&](Key k) -> std::optional<Msg> {
        if (m.mode == Mode::Editing) return edit_key(k);
        if (m.mode == Mode::Modal) return modal_key(k);
        return browse_key(k);
    });
    return keys;
}
\end{lstlisting}

The subscription is mode-aware: each mode dispatches keys
through its own handler. The model decides where input goes.

\section{What this exercises}

\begin{itemize}[leftmargin=*]
  \item Concepts and the \code{Program} contract.
  \item \code{Expected} for file I/O errors.
  \item \code{Cmd::task} for async load and save.
  \item Designated initialisers for tasks.
  \item Variants and \code{visit} for messages.
  \item The DSL for static layout.
  \item Builders for runtime list rows.
  \item Style merge for selection highlighting.
  \item Focus and mode for input routing.
  \item Undo (pattern 4) for history.
  \item Master-detail (pattern 1) for layout.
  \item Modal (overlay) for help.
\end{itemize}

A 200-line app. Every Part II chapter showed up.

\section{Where to take it next}

The app is small; the path forward is large:

\begin{itemize}[leftmargin=*]
  \item Add a filter palette (pattern 8).
  \item Persist undo to disk (so it survives restart).
  \item Add a tag system (pattern 5: incremental search by
    tag).
  \item Sync to a remote backend (pattern 6: side-effecting
    save, optimistic UI).
\end{itemize}

Each is an evening's work because the architecture absorbs the
features without refactoring. That is the win.

\section{The reading order, in retrospect}

If you read Part II in order, you saw the architecture from the
bottom up:

\begin{enumerate}[leftmargin=*]
  \item The Elm loop, conceptually (Ch.~\ref{ch:elm}).
  \item The compile-time DSL that produces elements
    (Ch.~\ref{ch:dsl}).
  \item Elements, the data the renderer consumes
    (Ch.~\ref{ch:elements}).
  \item Layout, the algorithm that places elements
    (Ch.~\ref{ch:layout}).
  \item Concepts, the contracts everything uses
    (Ch.~\ref{ch:concepts}).
  \item Expected, the error-as-value type (Ch.~\ref{ch:expected}).
  \item The pipeline (Ch.~\ref{ch:pipeline}).
  \item Style, theme, palette (Ch.~\ref{ch:style}).
  \item Builders for runtime trees (Ch.~\ref{ch:builders}).
  \item Text and graphemes (Ch.~\ref{ch:text}).
  \item Platform abstraction (Ch.~\ref{ch:platform}).
  \item Focus (Ch.~\ref{ch:focus}).
  \item Scroll (Ch.~\ref{ch:scroll}).
  \item Error boundaries (Ch.~\ref{ch:errorboundary}).
  \item Reactivity preview (Ch.~\ref{ch:reactivity-preview}).
  \item Runtime internals (Ch.~\ref{ch:runtime}).
  \item Testing (Ch.~\ref{ch:testing}).
  \item Widget patterns (Ch.~\ref{ch:widget-arch}).
  \item Compositional patterns (Ch.~\ref{ch:patterns}).
  \item Performance (Ch.~\ref{ch:perf}).
  \item DSL guarantees (Ch.~\ref{ch:dsl-guarantees}).
  \item C++26 features (Ch.~\ref{ch:cpp26-features}).
  \item Comparison to other libraries (Ch.~\ref{ch:comparison}).
\end{enumerate}

Twenty-three chapters. Each is a piece of the architecture. The
end-to-end app above ties them together.

\section{What is next}

Part III walks the rendering pipeline: how the element tree
becomes a canvas of packed cells, how SIMD makes the diff
sub-microsecond, how ANSI is emitted. Part IV walks the
reactive layer in depth and the widget catalogue.

But before turning the page, build something. Even a small
counter. The architecture only feels right once you have lived
in it.

\begin{recap}
A 200-line task manager exercises every architectural concept
from Part II. The pattern: model is data, messages drive state
transitions, the view is a pure function, effects are described
as values. Every concept earned its keep. Part III walks the
rendering details that make the model fast.
\end{recap}

\chapter{Message routing and dispatch internals}
\label{ch:dispatch}

\begin{aside}
Behind \code{std::visit} and \code{Cmd::map}, there is a real
dispatch machinery that turns variant types into jump tables and
threaded function pointers. This chapter opens the hood and walks
through what the compiler actually emits, how dispatch costs
compare across alternatives, and why \maya{}'s sum-type approach
beats both inheritance and \code{std::function} in practice.
\end{aside}

\section{The dispatch options C++ gives us}

Three mainstream ways to dispatch a message to a handler:

\begin{enumerate}[leftmargin=*]
  \item \textbf{Virtual functions.} Classic OOP. A vtable pointer
    per object, a vtable per class, one indirect call.
  \item \textbf{std::function or function pointer.} A type-erased
    callable; storage is heap-allocated for non-trivial
    closures.
  \item \textbf{Sum type + visit.} A discriminator integer plus
    inline storage; visit compiles to a switch.
\end{enumerate}

\maya{} uses option 3 throughout. The trade-offs are subtle but
real.

\subsection*{Virtual functions: the cost model}

\begin{lstlisting}
class Handler {
public:
    virtual void handle(Msg) = 0;
};

class Counter : public Handler {
public:
    void handle(Msg) override;
};

Handler* h = new Counter();
h->handle(msg);
\end{lstlisting}

The call \code{h->handle(msg)}:

\begin{itemize}[leftmargin=*]
  \item Load \code{h} (likely cached).
  \item Load vtable pointer at \code{*h} (one extra cache line).
  \item Load \code{handle} slot from vtable.
  \item Indirect call.
\end{itemize}

Costs: typically 2--5 cycles in the best case (vtable cached),
20+ cycles in the worst (vtable cold). Plus the heap allocation
for \code{new Counter()}.

\subsection*{std::function: the cost model}

\begin{lstlisting}
std::function<void(Msg)> handler = [](Msg m) { /* ... */ };
handler(msg);
\end{lstlisting}

Costs:

\begin{itemize}[leftmargin=*]
  \item If the captured state is small-object-optimised, two
    indirect calls (the \code{function} dispatches via vtable
    too).
  \item If the captured state exceeds SBO, one heap allocation
    plus the indirection.
  \item Inlining is rare; the compiler cannot prove
    \code{handler} points to a specific function.
\end{itemize}

Plus the heap allocation cost amortised across calls.

\subsection*{Sum type + visit: the cost model}

\begin{lstlisting}
std::variant<MsgA, MsgB, MsgC> msg = MsgA{};
std::visit([](auto& m) { /* ... */ }, msg);
\end{lstlisting}

Costs:

\begin{itemize}[leftmargin=*]
  \item Read the discriminator (one byte).
  \item Switch on the discriminator (jump table).
  \item Direct call to the matched branch.
  \item Inlining is common; the compiler sees every branch.
\end{itemize}

Costs: 1--3 cycles in the typical case. No heap allocation.

\section{What the compiler emits}

Let us look at the assembly.

\subsection*{Virtual call disassembly}

\begin{lstlisting}
; h->handle(msg)
mov rax, qword ptr [rdi]            ; load vtable ptr
mov rdx, qword ptr [rax]            ; load handle slot
jmp rdx                              ; indirect tail call
\end{lstlisting}

Three instructions, two memory loads, one indirect call. The
branch predictor must learn the target.

\subsection*{std::visit disassembly}

\begin{lstlisting}
; std::visit(f, v)
movzx eax, byte ptr [rdi + 16]      ; load discriminator
cmp eax, 3                           ; bounds check
jae default
jmp qword ptr [jump_table + rax*8]  ; jump-table dispatch

case_0:
    ; inlined body for MsgA
    ret
case_1:
    ; inlined body for MsgB
    ret
case_2:
    ; inlined body for MsgC
    ret
\end{lstlisting}

A jump table the size of the variant. Each case is inlined
where possible. The branch predictor learns the table entries.

\subsection*{The point}

The compiler can specialise the visit much more aggressively
than the virtual call. Inlining a variant case eliminates
function-call overhead; inlining a virtual call requires
devirtualisation, which is rare.

\section{Cmd::map: closure threading}

\code{Cmd::map} is more involved. It transforms a
\code{Cmd<ChildMsg>} into a \code{Cmd<ParentMsg>} via a
function. The implementation walks the variant and rebuilds it:

\begin{lstlisting}
template <class A, class B, class F>
Cmd<B> map(Cmd<A> cmd, F f) {
    return std::visit(overload{
        [&](None)      -> Cmd<B> { return Cmd<B>::none(); },
        [&](Quit)      -> Cmd<B> { return Cmd<B>::quit(); },
        [&](After<A> a) -> Cmd<B> {
            return Cmd<B>::after(a.delay, f(a.message));
        },
        [&](Task<A> t) -> Cmd<B> {
            auto fn = std::move(t.fn);
            auto f2 = std::move(f);
            return Cmd<B>::task([fn = std::move(fn), f = std::move(f2)]() mutable -> B {
                return f(fn());
            });
        },
        // ... one branch per alternative
    }, cmd.v);
}
\end{lstlisting}

\code{None}, \code{Quit}: no payload to transform; return the
same shape in the new type.

\code{After<A>}: transform the payload message via \code{f}.

\code{Task<A>}: build a new closure that runs the old task and
applies \code{f} to its result.

The cost: one allocation for \code{Task} (the closure storing
\code{f}); zero for the simpler alternatives.

\section{Why visit, not switch}

We could write \code{visit} by hand:

\begin{lstlisting}
switch (msg.index()) {
    case 0: handle(std::get<MsgA>(msg)); break;
    case 1: handle(std::get<MsgB>(msg)); break;
    case 2: handle(std::get<MsgC>(msg)); break;
}
\end{lstlisting}

This compiles to the same assembly as \code{std::visit} ---
but with two problems:

\begin{itemize}[leftmargin=*]
  \item Adding an alternative breaks the switch; the compiler
    does not warn unless you enable \code{-Wswitch}.
  \item \code{std::get<MsgA>} can throw if the variant is not
    \code{MsgA}; \code{visit}'s callable cannot.
\end{itemize}

\code{std::visit} is the type-safe spelling that the compiler
enforces. It is also the recipe the overload pattern targets.

\section{The exhaustiveness check}

When the variant grows, \code{std::visit} requires the visitor
to cover every alternative. If you forget one, the compiler
emits a clear error:

\begin{lstlisting}
error: no matching function for call to 'operator()' with 'MsgD'
\end{lstlisting}

For an overload-pattern visitor:

\begin{lstlisting}
error: 'overload<...>::operator()(MsgD)' is not callable
\end{lstlisting}

Either way: a new alternative without a corresponding handler
fails the build. You cannot ship a half-implemented message
variant.

\subsection*{When you want a catch-all}

Sometimes you want a default. An \code{auto} lambda matches any
alternative:

\begin{lstlisting}
std::visit(overload{
    [](MsgA a) { /* specific */ },
    [](MsgB b) { /* specific */ },
    [](auto)   { /* default */ }
}, msg);
\end{lstlisting}

The compiler picks the most-specialised overload for each
alternative; \code{auto} catches the rest. Use sparingly ---
it defeats exhaustiveness.

\section{Tagged union vs raw union}

\code{std::variant} is a tagged union: a discriminator plus
storage. A raw \code{union} has no discriminator; you must
track it yourself.

\begin{lstlisting}
union Msg {
    MsgA a;
    MsgB b;
    MsgC c;
};

// You must remember which one is active.
\end{lstlisting}

The trade-off:

\begin{itemize}[leftmargin=*]
  \item Variant: safer, slightly larger (one byte for tag plus
    alignment).
  \item Union: smaller, requires manual bookkeeping.
\end{itemize}

For one-off micro-optimisations a hand-rolled union can shave
bytes; for application code, the variant's safety wins every
time.

\section{Inlining and the optimiser}

Three factors determine how well the optimiser inlines visit:

\begin{enumerate}[leftmargin=*]
  \item \textbf{Visibility.} If the visitor and the visit call
    are in the same translation unit, the compiler can see all
    branches.
  \item \textbf{Constness.} A \code{const} variant is easier to
    reason about.
  \item \textbf{Branch count.} Above 10--12 branches, the
    compiler switches from inlined cases to a jump table; both
    are fast.
\end{enumerate}

In practice, \maya{}'s visit calls inline cleanly. The
generated code is competitive with hand-rolled switches.

\section{Pitfalls}

\begin{warning}
\textbf{Visit returning different types.} Each branch must
return a type the others convert to. If one returns
\code{int} and another \code{string}, the visit does not
compile. Wrap in \code{std::variant} or a base type.
\end{warning}

\begin{warning}
\textbf{Recursive visit.} Visiting from inside another visit
can be confusing. Capture the inner type clearly; write each
visit as its own named function and call them in sequence rather
than nesting visits at the call site.
\end{warning}

\begin{warning}
\textbf{Sharing state across visits.} The visitor lambda
captures by reference; the closure lives as long as the visit
returns. Do not store the visit-returned value beyond the
visitor's scope without copying.
\end{warning}

\begin{warning}
\textbf{Performance pessimism for small variants.}
\code{std::variant<int, double>} adds a byte. For tight inner
loops over millions of items, that byte can matter. Profile;
do not assume.
\end{warning}

\section{What the runtime emits per frame}

A typical \maya{} update loop has ten to twenty visit calls per
frame (one for the main message, one for each Cmd
interpretation, several for tree walks). The total dispatch
cost is microseconds, dominated by the actual work in each
branch, not by the dispatch itself.

If you ever profile a frame and the dispatch shows up high in
the flamegraph, the message variant is probably too large
(hundreds of alternatives) or the visit is in a tight loop
where it should not be. Move the visit outside the loop.

\section*{Workshop}

\begin{enumerate}[leftmargin=*]
  \item Write a small program with a 5-alternative variant and
    a \code{std::visit}. Compile with optimisation; inspect the
    assembly. Identify the jump table.
  \item Add a sixth alternative. Confirm the compiler errors
    on the visitor.
  \item Replace the visit with a \code{switch} on
    \code{msg.index()}. Confirm the assembly is identical.
\end{enumerate}

\begin{recap}
Sum types plus visit are the dispatch mechanism throughout
\maya{}. They compile to jump tables, inline cleanly, avoid
heap allocations, and enforce exhaustiveness. Virtual functions
and \code{std::function} have higher overhead and weaker static
guarantees. The choice is not stylistic; it is mechanical.
\end{recap}

\chapter{Constexpr deep dive: how the DSL computes at build time}
\label{ch:constexpr-deep}

\begin{aside}
\code{constexpr} and \code{consteval} are the engines of the
DSL. Most users see them as decorations; behind the scenes they
are full programming languages running inside the compiler. This
chapter takes a deep look at what is computable at compile time
in C++26, where the limits are, and how the DSL uses every
available feature.
\end{aside}

\section{What runs at compile time}

A \code{constexpr} function may run at compile time when its
arguments are constant expressions. A \code{consteval} function
\emph{must} run at compile time --- the compiler errors out if a
runtime invocation is attempted.

The set of operations allowed inside \code{constexpr} has
grown dramatically:

\begin{itemize}[leftmargin=*]
  \item C++11: arithmetic, simple branches, single return.
  \item C++14: loops, mutation of locals, multiple returns.
  \item C++17: lambda invocation, \code{if constexpr}.
  \item C++20: dynamic allocation (transient), virtual functions,
    \code{try/catch} (limited).
  \item C++23: \code{static\_assert} of non-literal types,
    \code{constexpr} std::optional, std::variant.
  \item C++26: even more relaxations; reflection in the pipeline.
\end{itemize}

The DSL uses C++20 and later. The compiler can run real
programs --- not just constant folding --- to compute its NTTPs.

\section{A worked example: parsing a hex colour at compile time}

\code{Fg<"\#FF6432">} requires the compiler to parse the hex
literal and convert it to three integers. The code:

\begin{lstlisting}
constexpr uint8_t hex_digit(char c) {
    if (c >= '0' && c <= '9') return c - '0';
    if (c >= 'a' && c <= 'f') return c - 'a' + 10;
    if (c >= 'A' && c <= 'F') return c - 'A' + 10;
    throw std::invalid_argument(\"not a hex digit\");
}

constexpr std::array<uint8_t, 3> parse_hex(std::string_view s) {
    if (s.size() != 7 || s[0] != '#')
        throw std::invalid_argument(\"expected #RRGGBB\");
    return {
        static_cast<uint8_t>(hex_digit(s[1]) * 16 + hex_digit(s[2])),
        static_cast<uint8_t>(hex_digit(s[3]) * 16 + hex_digit(s[4])),
        static_cast<uint8_t>(hex_digit(s[5]) * 16 + hex_digit(s[6]))
    };
}

template <Str s>
constexpr auto fg = []() {
    constexpr auto rgb = parse_hex(s.view());
    return Fg<rgb[0], rgb[1], rgb[2]>;
}();
\end{lstlisting}

Several things happen here:

\begin{itemize}[leftmargin=*]
  \item \code{parse\_hex} is \code{constexpr}. It throws on
    malformed input.
  \item The thrown exception, in a \code{constexpr} context, is
    a compile-time error. The compiler reports it as
    ``call to non-constexpr function'' or ``thrown exception''.
  \item \code{template <Str s>} takes the literal as an NTTP;
    the immediately-invoked lambda runs at compile time and
    instantiates \code{Fg} with the parsed values.
\end{itemize}

\subsection*{What the compiler sees}

When you write \code{fg<"\#FF6432">}, the compiler:

\begin{enumerate}[leftmargin=*]
  \item Builds a \code{Str<8>} from the literal.
  \item Substitutes into the template.
  \item Evaluates the lambda body: calls \code{parse\_hex}.
  \item \code{parse\_hex} returns \code{[255, 100, 50]}.
  \item Instantiates \code{Fg<255, 100, 50>}.
  \item Caches the result; \code{fg<"\#FF6432">} is now this
    instance.
\end{enumerate}

All of that happens before any runtime code runs.

\section{Throwing in constexpr}

A throw inside \code{constexpr} is interesting:

\begin{itemize}[leftmargin=*]
  \item If the function is called at runtime, throwing is
    normal: the exception propagates.
  \item If the function is called at compile time, the throw
    is a compile error: ``a thrown exception is not a constant
    expression''.
\end{itemize}

This is the mechanism by which \code{parse\_hex} validates input.
You get runtime exceptions in runtime context and compile errors
in constexpr context, from the same source.

\begin{idea}
This dual nature is what makes \code{constexpr} so powerful for
DSLs. A single function can validate its inputs at \emph{either}
compile time or runtime, depending on how it is called. The DSL
is just a function called at compile time with literal
arguments.
\end{idea}

\section{consteval: compile-time only}

When you want a function that must never run at runtime, use
\code{consteval}:

\begin{lstlisting}
consteval CTStyle merge_styles(CTStyle a, CTStyle b) {
    return CTStyle{
        .bold = a.bold || b.bold,
        .italic = a.italic || b.italic,
        // ...
    };
}
\end{lstlisting}

The compiler errors if you try to call \code{merge\_styles}
with a runtime argument. This is what gates the DSL's pipe
operators: they call \code{merge\_styles}, which insists on
constexpr inputs, which forces the pipe to be a compile-time
expression.

\section{constinit: ensuring static-init order}

\begin{lstlisting}
constinit Theme default_theme = make_dark_theme();
\end{lstlisting}

\code{constinit} guarantees the variable is initialised at
compile time, not via dynamic initialisation. This avoids the
static-init order fiasco: \code{constinit} variables are ready
before \code{main} runs.

\maya{} uses \code{constinit} for theme defaults, palette
tables, and other startup-critical values.

\section{constexpr std::vector and std::string}

C++20 made \code{std::vector} and \code{std::string} usable in
\code{constexpr}. The compiler implements a tiny heap inside the
constant evaluator:

\begin{lstlisting}
consteval std::vector<int> build_table() {
    std::vector<int> v;
    for (int i = 0; i < 100; ++i) v.push_back(i * i);
    return v;
}

constexpr auto squares = build_table();
\end{lstlisting}

The catch: the vector must not escape the constant evaluator.
You cannot return a \code{constexpr std::vector} to runtime;
the heap allocation is freed when constexpr evaluation ends.

The workaround: copy into a \code{std::array}:

\begin{lstlisting}
consteval auto fixed_squares() {
    auto v = build_table();
    std::array<int, 100> a{};
    for (size_t i = 0; i < 100; ++i) a[i] = v[i];
    return a;
}

constexpr auto squares = fixed_squares();
\end{lstlisting}

\maya{} uses this pattern for the Unicode width table: build at
compile time, store as \code{constexpr std::array}.

\section{if constexpr: compile-time branches}

\begin{lstlisting}
template <class T>
auto encode(T value) {
    if constexpr (std::is_integral_v<T>) {
        return std::to_string(value);
    } else if constexpr (std::is_floating_point_v<T>) {
        return std::format(\"{:f}\", value);
    } else {
        return std::string{value};
    }
}
\end{lstlisting}

Branches whose condition is \code{constexpr} are evaluated at
compile time; only the chosen branch is instantiated. The other
branches do not have to be well-formed for the unselected type.

\maya{} uses \code{if constexpr} for platform-specific code,
SIMD level selection, and template specialisation without
explicit overloading.

\section{requires expressions in constexpr}

Inside a function, you can use \code{requires} to check
properties:

\begin{lstlisting}
template <class T>
auto format(T value) {
    if constexpr (requires { value.to_string(); }) {
        return value.to_string();
    } else if constexpr (requires { std::to_string(value); }) {
        return std::to_string(value);
    } else {
        return std::string{\"(no format)\"};
    }
}
\end{lstlisting}

The requires-expression is a compile-time predicate. The
branches are selected based on what types support what
operations.

\section{The limits of compile-time computation}

What you cannot do at compile time:

\begin{itemize}[leftmargin=*]
  \item I/O. No file reads, no network, no stdout.
  \item Reading the system clock.
  \item Random numbers (other than constexpr-friendly PRNGs).
  \item Reflection on names (until C++26 reflection).
  \item Unbounded recursion (the compiler has a depth limit,
    typically 1024).
\end{itemize}

The DSL stays within these limits. Every literal it processes,
every style it merges, every box it configures --- all done
without I/O or randomness.

\section{Compile-time vs runtime: who wins?}

A computation that fits in either context has a choice.
Compile-time advantages:

\begin{itemize}[leftmargin=*]
  \item Zero runtime cost.
  \item Errors surface at build, not in production.
  \item Result is a known value; the optimiser sees it.
\end{itemize}

Runtime advantages:

\begin{itemize}[leftmargin=*]
  \item Faster builds: not every value needs to be a template
    parameter.
  \item More flexible: runtime input.
  \item Debuggable.
\end{itemize}

The DSL's principle: shape at compile time, content at runtime.
Box configurations, style flags, sizes --- compile time. Text
contents, list lengths, model values --- runtime.

\section{Build time as a cost}

A complex DSL expression compiles slowly. \maya{}'s typical
build:

\begin{itemize}[leftmargin=*]
  \item A small app (1000 lines): 2--5 seconds.
  \item A medium app (10\,000 lines): 30--60 seconds.
  \item A large app (100\,000 lines): minutes.
\end{itemize}

Most of the time is spent in template instantiation. If you
care about build times, use the runtime variants
(\code{text(\ldots)}, \code{box(\ldots)}) for the bulk of your
UI and reserve the compile-time DSL for hot or critical paths.

\section*{Workshop}

\begin{enumerate}[leftmargin=*]
  \item Write a \code{constexpr} function that computes the
    factorial. Call it at compile time and at runtime; observe
    both work.
  \item Add a \code{static\_assert} that fails. Read the
    compiler error. Now make it pass.
  \item Read \file{maya/include/maya/dsl.hpp} and find every
    \code{consteval} function. Each is a piece of the DSL's
    compile-time computation.
\end{enumerate}

\begin{recap}
\code{constexpr} and \code{consteval} are the engines of
compile-time computation. The DSL uses them to parse literals,
merge styles, validate configurations, and build static lookup
tables. The compiler runs real programs to compute the types
your runtime sees. The cost is build time; the payoff is zero
runtime overhead and compile-time error detection.
\end{recap}

\chapter{The element identity story}
\label{ch:identity}

\begin{aside}
Two element trees that look the same may or may not be the same.
The renderer needs to know which: identical trees mean ``skip
the work''; different trees mean ``diff and emit''. The story
of how \maya{} establishes identity --- and why it matters ---
is more subtle than ``the trees are equal''.
\end{aside}

\section{Three notions of identity}

For an element tree, identity could mean:

\begin{itemize}[leftmargin=*]
  \item \textbf{Pointer identity.} Same address.
  \item \textbf{Structural identity.} Same shape, same contents.
  \item \textbf{Semantic identity.} Same shape, same contents,
    same focus / cursor / scroll positions.
\end{itemize}

Different renderer stages care about different ones. The cache
fingerprint cares about (2). The hit-tester cares about (1).
The animation system cares about (3).

\section{Pointer identity in the DSL}

The DSL builds a fresh tree every frame. The root pointer
changes; the children pointers change. Pointer identity is
useless for cross-frame reasoning.

If you need to keep state across frames keyed by ``the same
button'', use a stable ID:

\begin{lstlisting}
struct Button {
    int id;            // stable across frames
    std::string label;
};
\end{lstlisting}

The ID lives in your model. The element produced by rendering
the button gets the ID as a property; the runtime uses it for
focus, hit-test, animation state.

\section{Structural identity: the cache fingerprint}

The cache fingerprint (Chapter~\ref{ch:pipeline}) is structural
identity. Two trees with identical contents produce identical
fingerprints; two different trees produce different ones (with
high probability).

The fingerprint is computed by walking the tree and mixing each
leaf's content into a hash:

\begin{lstlisting}
uint64_t fingerprint(const Element& e) {
    XXH64State state;
    walk(e, [&](const TextElement& t) {
        XXH64_update(&state, t.text.data(), t.text.size());
        XXH64_update(&state, &t.style, sizeof(t.style));
    });
    walk(e, [&](const BoxElement& b) {
        XXH64_update(&state, &b.layout, sizeof(b.layout));
        // children handled by recursion
    });
    return XXH64_digest(&state);
}
\end{lstlisting}

The walk is depth-first, left-to-right. Two trees that differ
only in child order produce different fingerprints.

\subsection*{Collision probability}

A 64-bit hash collides at roughly the square-root rate: two
different trees out of $2^{32}$ are expected to share a
fingerprint. For a \maya{} app that produces at most a few
thousand distinct trees, the probability of a false ``no-change''
verdict in any given session is negligible.

If collision matters (e.g. you ship a long-running server with
millions of users), use a 128-bit hash. \maya{}'s fingerprint
is parameterised; switching is a one-line change.

\section{Semantic identity: keys and reconciliation}

When trees change, the renderer wants to preserve cursor
positions, scroll offsets, animation states. Naively, ``the
third child is now in position four'' would lose this state.

The fix: attach \emph{keys} to children. A key is a stable
identifier on an element that the reconciler uses to match
across frames.

\begin{lstlisting}
v(
    keyed(\"item-42\", render_item(items[42])),
    keyed(\"item-43\", render_item(items[43])),
    keyed(\"item-44\", render_item(items[44]))
);
\end{lstlisting}

The reconciler matches by key, not position. Inserting a new
item at the top does not shift the existing items' state.

\subsection*{When keys matter most}

\begin{itemize}[leftmargin=*]
  \item Lists where items can reorder.
  \item Lists with per-item animation state (slide-in,
    fade-out).
  \item Forms where input focus must follow a moving row.
\end{itemize}

For static lists, position-based reconciliation works fine.

\section{Identity in the layout phase}

Layout assigns each node a final rect. The rects are derived
from the tree, not from any cross-frame identity. Two frames
with the same tree produce the same rects.

This is intentional: layout is a pure function. There is no
``layout cache'' keyed by node identity. If you want the layout
to be different from frame to frame, the tree must differ.

\section{Identity in the canvas}

The canvas is a flat grid of cells. Cells have no identity
across frames; each frame produces a fresh canvas, and the diff
compares them cell-by-cell.

This means: cell-level animations (a single character pulsing)
are tracked at the cell level, not the element level. The
canvas does not know about elements.

\section{Identity in the focus system}

Focus is keyed by ID, not by element pointer. A focused button
keeps focus across re-renders because its ID is stable. If you
generate fresh IDs each frame, focus moves randomly --- the
classic ``why is focus jumping?'' bug.

\section{Identity in the signal system}

Signals are objects with their own identity. Subscribing to a
signal is keyed by the signal's address. A signal whose address
changes across frames cannot be subscribed to reliably; keep
signals in long-lived containers (the model, a long-lived
struct).

\section{When identity breaks down}

A few patterns cause identity to misbehave:

\subsection*{Anonymous lambdas as keys}

\begin{lstlisting}
auto button = [&] { return text(\"click\"); };
// Each frame, button is a different lambda --- different type, different identity.
\end{lstlisting}

Capture the lambda's expression, not the lambda itself. Or
better, write a function.

\subsection*{Index-based keys for reorderable lists}

\begin{lstlisting}
for (int i = 0; i < items.size(); ++i) {
    rows.push_back(keyed(std::to_string(i), render(items[i])));
}
\end{lstlisting}

When an item is inserted at the front, every key shifts. Use the
item's own stable ID, not the position.

\subsection*{Mutable state in element fields}

If an element carries mutable state via a captured reference,
the fingerprint may not detect changes (the data the reference
points to is not part of the element's bytes). Either include
the data by value in the element or mark the element as
``do not cache''.

\section{Pitfalls}

\begin{warning}
\textbf{Forgetting keys in reorderable lists.} Symptoms:
focus jumps to a different row, animation restarts, scroll
position drifts. Fix: add keys.
\end{warning}

\begin{warning}
\textbf{Keys that are not unique.} Two children with the same
key under one parent break the reconciler. Use IDs from
your model.
\end{warning}

\begin{warning}
\textbf{Cache fingerprint false negatives.} If your tree
references external mutable state, the fingerprint may not
change when the rendered output should. Make the data part
of the model, not external.
\end{warning}

\begin{warning}
\textbf{Cache fingerprint false positives.} If two visually
distinct trees somehow produce the same fingerprint
(collision), the renderer skips a real update. Vanishingly
rare with 64-bit hashes for typical app sizes; switch to
128-bit if you ship long-running services.
\end{warning}

\section*{Workshop}

\begin{enumerate}[leftmargin=*]
  \item Build a list of three items. Add a keyed item at the
    top. Confirm the bottom two items keep their state (focus,
    scroll).
  \item Replace the keys with the position index. Add another
    item at the top. Observe the state drift.
  \item Compute the fingerprint of a tree by hand (small one).
    Confirm two different trees yield different fingerprints.
\end{enumerate}

\begin{recap}
Identity in \maya{} has three notions: pointer (per frame),
structural (cross-frame, by content hash), and semantic
(cross-frame, by key). The renderer uses each appropriately.
Keys preserve focus, scroll, and animation state across
re-renders. The fingerprint enables ``skip everything'' on no
change. Mistakes in identity usually surface as visible state
loss.
\end{recap}

\chapter{Configuration and persistence}
\label{ch:config}

\begin{aside}
Every non-trivial app reads configuration from somewhere ---
defaults, environment variables, command-line flags, a config
file, a database. Persisting state across runs follows the same
pattern in reverse. \maya{} does not impose a config system, but
the architecture invites a specific pattern: read at startup,
hold in the model, never re-read inside the loop.
\end{aside}

\section{The startup pipeline}

A typical \maya{} app does the following before \code{run}:

\begin{enumerate}[leftmargin=*]
  \item Parse command-line arguments.
  \item Load the config file.
  \item Read relevant environment variables.
  \item Apply precedence: CLI overrides env overrides file
    overrides defaults.
  \item Construct the initial \code{Model} including the merged
    config.
  \item Hand the model to \code{run}.
\end{enumerate}

The startup pipeline is the only place that touches the
filesystem for config. Inside the loop, the model carries
everything.

\subsection*{Why not reload on demand?}

A frequent temptation: ``read the config file every few seconds
so changes take effect''. The Elm-style approach: model the
file as state, watch it with a \code{Cmd::task}, and turn
changes into messages. The runtime is the same; the contract
stays clean.

\section{Defaults}

A small struct with sensible defaults:

\begin{lstlisting}
struct Config {
    bool   show_help      = true;
    int    history_size   = 100;
    Theme  theme          = themes::Dark();
    std::string log_path  = \"~/.cache/myapp/log\";
    std::chrono::milliseconds idle_timeout{5000};
};

constexpr Config default_config{};
\end{lstlisting}

The struct holds everything user-tunable. Designated init
allows partial override:

\begin{lstlisting}
Config cfg{
    .theme = themes::Light(),
    .history_size = 500,
};
\end{lstlisting}

Fields not mentioned retain their defaults.

\section{Config files}

\subsection*{Format}

TOML, YAML, JSON, INI --- pick one. \maya{} examples use TOML
because:

\begin{itemize}[leftmargin=*]
  \item Human-readable.
  \item Strict parser available (\code{toml++}).
  \item No indentation sensitivity (unlike YAML).
  \item Comments allowed.
\end{itemize}

\subsection*{Loading}

\begin{lstlisting}
Expected<Config, std::string> load_config(std::string_view path) {
    return file::read(path)
        .transform_error([](auto& e) { return e.context; })
        .and_then(parse_toml)
        .transform(toml_to_config);
}
\end{lstlisting}

A monadic chain (Chapter~\ref{ch:expected}): read, parse,
convert. Each step can fail; the chain handles propagation.

\subsection*{Default location}

XDG conventions on Linux:

\begin{lstlisting}
std::string config_path() {
    if (auto* x = std::getenv(\"XDG_CONFIG_HOME\")) {
        return std::string{x} + \"/myapp/config.toml\";
    }
    return std::string{std::getenv(\"HOME\")}
        + \"/.config/myapp/config.toml\";
}
\end{lstlisting}

On macOS, \code{\$HOME/Library/Application Support/myapp/}; on
Windows, \code{\%APPDATA\%\textbackslash myapp\textbackslash}.
The platform abstraction layer (Chapter~\ref{ch:platform}) can
provide the helper.

\section{Environment variables}

For machine-specific overrides, env vars are common:

\begin{lstlisting}
if (auto* t = std::getenv(\"MYAPP_THEME\")) {
    cfg.theme = themes::by_name(t).value_or(cfg.theme);
}
if (auto* l = std::getenv(\"MYAPP_LOG\")) {
    cfg.log_path = l;
}
\end{lstlisting}

Convention: prefix env vars with your app name to avoid
collisions.

\subsection*{Standard variables to honour}

A small list of variables most users expect:

\begin{itemize}[leftmargin=*]
  \item \code{NO\_COLOR}: disable colour output.
  \item \code{FORCE\_COLOR}: enable colour even on non-tty.
  \item \code{COLORTERM}: terminal colour capabilities.
  \item \code{TERM}: terminal type.
  \item \code{HOME}, \code{XDG\_CONFIG\_HOME}, etc.: paths.
  \item \code{LANG}, \code{LC\_*}: locale (affects width tables).
\end{itemize}

\section{Command-line flags}

The simplest parser is hand-written; more elaborate apps use
\code{argparse} or \code{cxxopts}:

\begin{lstlisting}
int main(int argc, char** argv) {
    Config cfg = default_config;
    for (int i = 1; i < argc; ++i) {
        std::string_view arg = argv[i];
        if (arg == \"--theme\" && i + 1 < argc) {
            cfg.theme = themes::by_name(argv[++i]).value_or(cfg.theme);
        } else if (arg == \"--help\") {
            print_help();
            return 0;
        }
    }
    return maya::run<App>(cfg);
}
\end{lstlisting}

For complex apps with subcommands, use a library. The patterns
are the same.

\section{Precedence}

CLI overrides env overrides file overrides defaults. In code:

\begin{lstlisting}
Config build_config(int argc, char** argv) {
    Config cfg = default_config;
    if (auto file_cfg = load_config(config_path())) {
        merge(cfg, file_cfg.value());
    }
    apply_env(cfg);
    apply_cli(cfg, argc, argv);
    return cfg;
}
\end{lstlisting}

Each step layered on top. \code{merge} is field-by-field
override; only fields explicitly set in the higher layer
replace the lower layer.

\section{Persisting state}

The flip side of config: state that survives across runs.
Examples:

\begin{itemize}[leftmargin=*]
  \item Window size (the app should reopen at the same size).
  \item Last opened file.
  \item Cursor position.
  \item Command history.
\end{itemize}

The pattern:

\begin{enumerate}[leftmargin=*]
  \item On startup, load a state file (often JSON or binary).
  \item Restore the state into the model.
  \item On meaningful changes, write the state back (debounced
    to avoid thrashing the filesystem).
  \item On quit, write a final state.
\end{enumerate}

\subsection*{Write strategy: debounced}

\begin{lstlisting}
[&](StateChanged) {
    m.dirty = true;
    return std::pair{m, Cmd<Msg>::after(2s, FlushState{})};
}

[&](FlushState) {
    if (!m.dirty) return std::pair{m, Cmd<Msg>{}};
    m.dirty = false;
    return std::pair{m, save_state_cmd(m.persistable_state())};
}
\end{lstlisting}

A change sets a flag and schedules a flush. The flush, when it
fires, checks the flag and writes. Multiple changes within the
window collapse to one write.

\subsection*{Write atomically}

To avoid corrupted state files (if the app crashes mid-write):

\begin{lstlisting}
Expected<void, IoError> atomic_write(std::string_view path, std::string_view data) {
    std::string tmp = std::string{path} + \".tmp\";
    auto w = file::write(tmp, data);
    if (!w) return std::unexpect, w.error();
    return file::rename(tmp, path);
}
\end{lstlisting}

Write to a temp file, then atomically rename. The rename is
atomic on POSIX; on Win32 you need \code{ReplaceFileW}.

\section{Schema migration}

State files persist across versions. If you add a field or
change a type, old state files must still load. Two strategies:

\subsection*{Versioned schemas}

\begin{lstlisting}
struct StateV1 { int cursor; };
struct StateV2 { int cursor; std::string filename; };

StateV2 migrate(StateV1 s) {
    return StateV2{.cursor = s.cursor, .filename = \"\"};
}
\end{lstlisting}

Read the version field; pick the right struct; migrate forward.

\subsection*{Tolerant readers}

Each field is optional; missing fields use defaults:

\begin{lstlisting}
struct State {
    int cursor = 0;
    std::string filename = \"\";
    // ... older versions had fewer fields
};
\end{lstlisting}

The deserialiser accepts any TOML object and populates known
fields. Unknown fields are ignored (or logged).

\section{Secrets}

Some config values are sensitive (API keys, tokens). The
patterns:

\begin{itemize}[leftmargin=*]
  \item Read from a file that only the user can read
    (mode 0600).
  \item Read from a system keychain (libsecret, Keychain
    Services).
  \item Read from an environment variable that the user sets
    in their shell config.
\end{itemize}

Do not write secrets to the state file. If the user wants to
persist them, that is a separate, deliberate choice.

\section{Pitfalls}

\begin{warning}
\textbf{Reading config inside view or update.} The config
should be loaded once at startup and stored in the model.
Re-reading mid-loop is slow and racy.
\end{warning}

\begin{warning}
\textbf{Tight coupling to file format.} If your model is a
TOML-shaped struct, swapping to JSON requires changing the
model. Keep the on-disk shape separate from the in-memory
model.
\end{warning}

\begin{warning}
\textbf{Saving on every keystroke.} The filesystem will hate
you. Debounce; consider write barriers.
\end{warning}

\begin{warning}
\textbf{Forgetting to validate.} A config file written by a
human can contain garbage. Validate on load; reject with a
clear error.
\end{warning}

\section*{Workshop}

\begin{enumerate}[leftmargin=*]
  \item Write a small config struct with five fields. Load
    defaults; override one from an env var; verify in the
    model.
  \item Add atomic writes for a state file. Force-kill the
    app mid-write; verify the file is not corrupted.
  \item Implement schema migration from a V1 state file to a
    V2. Write a unit test for the migration.
\end{enumerate}

\begin{recap}
Read config once at startup; carry it in the model; never re-
read inside the loop. Files in TOML or JSON, environment
variables for overrides, CLI flags for ad-hoc tweaks. Persist
state debounced and atomically. Migrate schemas across
versions. Secrets stay out of state files.
\end{recap}

\chapter{Logging, tracing, and observability}
\label{ch:observability}

\begin{aside}
An interactive program that does not log is opaque when things
go wrong. \maya{} ships a small observability toolkit: a ring
buffer for last-N messages, structured tracing for the render
pipeline, and a debug overlay you can toggle at runtime. None
of it is mandatory; all of it pays off the first time a user
files a bug.
\end{aside}

\section{Why not just std::cerr?}

A TUI app cannot use \code{std::cerr} freely because stderr is
the user's screen. Writing log messages to stderr corrupts the
display.

Three workable alternatives:

\begin{itemize}[leftmargin=*]
  \item Redirect stderr to a file at startup.
  \item Buffer log messages in memory; surface on demand.
  \item Write to a separate log file always.
\end{itemize}

\maya{} supports all three. The default is the ring buffer:
in-memory by default, dump on request.

\section{The Logger interface}

\file{maya/include/maya/app/context.hpp} exposes:

\begin{lstlisting}
class Logger {
public:
    void debug(std::string_view msg);
    void info(std::string_view msg);
    void warn(std::string_view msg);
    void error(std::string_view msg);

    template <class... Args>
    void info(std::string_view fmt, Args&&... args) {
        info(std::format(fmt, std::forward<Args>(args)...));
    }

    std::vector<LogEntry> recent(size_t n) const;
    void dump_to(std::ostream& os) const;
};
\end{lstlisting}

Four levels (debug, info, warn, error). Each takes a message;
the \code{std::format}-aware overloads support typed
arguments. \code{recent(n)} returns the last $n$ entries;
\code{dump\_to} writes them to a stream.

\subsection*{LogEntry}

\begin{lstlisting}
struct LogEntry {
    std::chrono::system_clock::time_point time;
    Level level;
    std::string category;
    std::string message;
};
\end{lstlisting}

Time, level, category (subsystem tag), message. Categories let
you filter logs by source.

\section{The ring buffer}

The default logger is a fixed-size ring:

\begin{lstlisting}
class RingLogger {
    std::vector<LogEntry> buffer_;
    size_t head_ = 0;
    size_t count_ = 0;
    mutable std::mutex mu_;

public:
    explicit RingLogger(size_t capacity);
    void append(LogEntry e);
    std::vector<LogEntry> snapshot() const;
};
\end{lstlisting}

Capacity is typically 1024 or 4096 entries. Older entries are
overwritten in place; the ring stays at constant memory.

The mutex is fine-grained: a single mutex protects the buffer.
Logging from multiple threads is supported (worker tasks, the
main thread).

\subsection*{Snapshot}

\begin{lstlisting}
std::vector<LogEntry> RingLogger::snapshot() const {
    std::lock_guard lk{mu_};
    std::vector<LogEntry> out;
    out.reserve(count_);
    for (size_t i = 0; i < count_; ++i) {
        out.push_back(buffer_[(head_ - count_ + i) % buffer_.size()]);
    }
    return out;
}
\end{lstlisting}

Returns the entries in chronological order. Threadsafe; can be
called while logging continues.

\section{Structured logging}

For tracing, \maya{} uses a slightly richer entry shape:

\begin{lstlisting}
struct TraceEntry {
    std::chrono::steady_clock::time_point time;
    std::string event;
    std::unordered_map<std::string, std::string> attributes;
};
\end{lstlisting}

The attributes map lets you attach context: \code{frame=42},
\code{node\_count=128}, \code{cell\_count=1920}. When you
later inspect the trace, you can filter and group by attribute.

\subsection*{Tracing the render pipeline}

The pipeline emits trace events at each stage:

\begin{lstlisting}
trace(\"view.begin\", {{\"frame\", std::to_string(frame_no)}});
auto e = P::view(model);
trace(\"view.end\", {{\"frame\", std::to_string(frame_no)},
                    {\"node_count\", std::to_string(count_nodes(e))}});
\end{lstlisting}

A trace pair: \code{begin}/\code{end} with the same key
attributes. The duration is the difference; the attribute set
lets you correlate.

\section{The debug overlay}

A widget that renders the most recent log entries:

\begin{lstlisting}
Element render_log_overlay(const Logger& l) {
    auto entries = l.recent(20);
    std::vector<Element> lines;
    for (auto& e : entries) {
        auto style = level_style(e.level);
        lines.push_back(h(
            text(format_time(e.time)) | Dim,
            text(\" \"),
            text(e.message) | style
        ));
    }
    return box(v(lines)) | border<Single> | pad<1>;
}
\end{lstlisting}

Toggled with a hotkey (typically F12 or Ctrl-Shift-L). The
overlay sits on top of the app via the z-stack from
Chapter~\ref{ch:elements}.

\section{Categories and filters}

Tag log messages with a category:

\begin{lstlisting}
logger.info(\"render\", \"frame {} ({} cells dirty)\", n, dirty);
logger.info(\"network\", \"GET {} -> 200\", url);
logger.info(\"input\", \"key {} pressed\", key.name());
\end{lstlisting}

Then filter:

\begin{lstlisting}
auto render_logs = logger.recent_filtered(
    100, [](auto& e) { return e.category == \"render\"; });
\end{lstlisting}

For development, set an environment variable to control which
categories are recorded:

\begin{lstlisting}
MAYA_LOG_CATEGORIES=render,network ./my_app
\end{lstlisting}

Only the listed categories are appended; others are dropped
cheaply.

\section{Telemetry to a remote sink}

For deployed apps, log to a file or service. The pattern:

\begin{lstlisting}
auto telemetry_cmd = Cmd<Msg>::task([entries] {
    upload(entries);
    return TelemetryDone{};
});
\end{lstlisting}

A background task uploads batches. The main loop is unaffected;
the user does not see the upload.

\subsection*{Privacy}

Be careful about what you log. PII (personally identifiable
information), file paths, secrets --- redact before recording.
A simple redactor:

\begin{lstlisting}
std::string redact(std::string_view msg) {
    // Replace anything that looks like an email, a token, a path.
    return std::regex_replace(std::string{msg},
        std::regex{R\"(\\w+@\\w+\\.\\w+)\"}, \"[REDACTED]\");
}
\end{lstlisting}

The list of patterns is app-specific. Err on the side of
redacting too much.

\section{Crash diagnostics}

If your app crashes, what survives?

\subsection*{The ring buffer is volatile}

A SIGSEGV does not flush. You need a signal handler:

\begin{lstlisting}
extern \"C\" void crash_handler(int sig) {
    logger.dump_to(std::ofstream{\"/tmp/myapp.crashlog\"});
    std::signal(sig, SIG_DFL);
    std::raise(sig);
}
\end{lstlisting}

Install with \code{std::signal(SIGSEGV, crash\_handler)}. The
handler must be async-signal-safe; in practice ``open a file
and write'' is OK but allocating memory is not. \maya{}'s
\code{dump\_to} is designed to be async-signal-safe.

\subsection*{Symbolicated stacks}

The crash log should include a stack trace. Use
\code{std::stacktrace} (C++23):

\begin{lstlisting}
auto trace = std::stacktrace::current();
for (auto& frame : trace) {
    log_line(std::format(\"{}\", frame));
}
\end{lstlisting}

Compile with debug symbols (\code{-g}); the trace will name
functions and line numbers.

\section{Performance considerations}

Logging is on the hot path. Cheap logging:

\begin{itemize}[leftmargin=*]
  \item String formatting is the slowest part; do it lazily.
  \item Disabled levels (e.g. \code{debug} in release) should
    cost zero --- one branch on a constant flag, not even a
    function call.
  \item Lock contention is real; if multiple threads log
    heavily, the buffer mutex becomes a bottleneck. Use
    thread-local buffers and a periodic merge.
\end{itemize}

The simplest lazy pattern:

\begin{lstlisting}
if (logger.level_enabled(Level::Debug)) {
    logger.debug(\"value = {}\", expensive_compute());
}
\end{lstlisting}

The expensive compute only runs if debug is on.

\section{Pitfalls}

\begin{warning}
\textbf{Logging in the inner loop.} Even cheap logging adds up
at 10\,000 calls per frame. Sample (log every 100th occurrence)
or gate by level.
\end{warning}

\begin{warning}
\textbf{Logging secrets.} A bug report with a config file
containing API keys is a security incident. Redact aggressively.
\end{warning}

\begin{warning}
\textbf{Lost messages on crash.} A buffer that lives only in
memory disappears with the process. Pair with a crash handler
that flushes.
\end{warning}

\begin{warning}
\textbf{Overlapping logs.} Multiple threads writing to the
same buffer can produce interleaved messages. The mutex
serialises them, but if you log a multi-line message, the
lines may not be adjacent. Log one message per call.
\end{warning}

\section*{Workshop}

\begin{enumerate}[leftmargin=*]
  \item Add a ring logger to a small app. Log on every key
    press. Toggle the debug overlay; verify entries appear.
  \item Install a crash handler that dumps the log on
    \code{SIGSEGV}. Force a crash (deref a null pointer);
    verify the log is written.
  \item Add structured tracing to the render pipeline (begin
    + end events with attributes). Compute frame durations
    from the trace.
\end{enumerate}

\begin{recap}
A TUI cannot use stderr freely; use a ring buffer or a log
file. Structured tracing captures context with attributes.
A debug overlay surfaces recent entries at runtime. A crash
handler dumps the buffer on SIGSEGV. Categories and levels
let you filter cheaply. Redact secrets before they leave the
process.
\end{recap}

\chapter{Internationalisation and locale}
\label{ch:i18n}

\begin{aside}
A TUI library used outside one language and one terminal soon
encounters every weird edge of internationalisation: non-ASCII
text, right-to-left scripts, locale-dependent number
formatting, time zones. \maya{} does not try to solve every i18n
problem (no library does), but it gives you primitives that
play well with the standard C++ approaches. This chapter is the
guide.
\end{aside}

\section{Three i18n axes}

International apps face three orthogonal concerns:

\begin{itemize}[leftmargin=*]
  \item \textbf{Translation.} The text shown to the user is in
    their language.
  \item \textbf{Locale formatting.} Numbers, dates, currencies
    follow local conventions.
  \item \textbf{Text rendering.} Non-Latin scripts, RTL, complex
    glyphs render correctly.
\end{itemize}

Each is a separate problem with its own tools.

\section{Translation}

The canonical pattern: wrap user-visible strings in a function
that looks up the translation:

\begin{lstlisting}
std::string_view _(const char* key);  // looks up by key

auto label = text(_(\"open_file\"));
\end{lstlisting}

\code{\_} (underscore) is the conventional name. The key is a
short identifier; the actual text comes from a translation
table indexed by current locale.

\subsection*{Loading translations}

\begin{lstlisting}
struct Translations {
    std::unordered_map<std::string, std::string> strings;

    static Expected<Translations, std::string>
    load(std::string_view lang_tag);
};
\end{lstlisting}

Translation files are typically JSON or PO format. Load once at
startup; hold in a global (or in the model).

\subsection*{Plurals}

English has two plural forms (one, many); some languages have
six. The standard solution is the CLDR plural rules:

\begin{lstlisting}
std::string_view ngettext(const char* key, int n);
\end{lstlisting}

The function picks the right form based on $n$ and the current
locale's plural rule. \maya{} does not ship one; use a library
(\code{boost::locale}, ICU) or roll your own.

\section{Locale formatting}

Numbers, dates, currencies follow locale conventions. C++26's
\code{std::format} respects the current locale by default:

\begin{lstlisting}
std::string s = std::format(\"{:L}\", 1234567);
// en-US: \"1,234,567\"
// de-DE: \"1.234.567\"
// fr-FR: \"1 234 567\"
\end{lstlisting}

The \code{L} format specifier enables locale awareness.

\subsection*{Setting the locale}

\begin{lstlisting}
std::locale::global(std::locale{\"de_DE.UTF-8\"});
\end{lstlisting}

Setting the global locale affects \code{format}, \code{ostream},
and the C functions that read the locale. Set it once at
startup based on the user's environment.

\subsection*{Time zones}

C++20 introduced \code{std::chrono::zoned\_time}:

\begin{lstlisting}
auto now = std::chrono::system_clock::now();
auto local = std::chrono::zoned_time{std::chrono::current_zone(), now};
std::cout << local << \"\\n\";
\end{lstlisting}

The output respects the system time zone. \maya{} apps that
display timestamps should always use \code{zoned\_time}
through the user's local zone, not the system clock directly.

\section{Text rendering}

\maya{}'s grapheme walker (Chapter~\ref{ch:text}) handles
non-Latin scripts correctly. The width table covers all
codepoints. But two complications remain.

\subsection*{Bidirectional text}

Arabic and Hebrew are read right-to-left. Modern terminals
(GNOME Terminal, Konsole, Alacritty) implement BiDi rendering;
\maya{} relies on this and emits text in logical order.

The catch: \emph{mixed} LTR/RTL inside a single text node is
where edge cases live. If the line contains both English and
Arabic, the visual order depends on the terminal's BiDi
implementation. \maya{} does not attempt to predict it.

\subsection*{Complex shaping}

Devanagari, Arabic ligatures, vowel marks --- these require
the font shaper to combine sequences into glyphs. The terminal's
font handles this, not \maya{}. The grapheme walker counts the
cells correctly; the terminal renders the glyphs.

\section{Identifying the locale}

\begin{lstlisting}
std::string detect_locale() {
    if (auto* l = std::getenv(\"LANG\")) return l;
    if (auto* l = std::getenv(\"LC_ALL\")) return l;
    return \"C\";
}
\end{lstlisting}

POSIX uses \code{LANG} and \code{LC\_*} variables. Win32 uses
\code{GetUserDefaultLocaleName}. The platform abstraction can
unify them.

\section{Number input}

When parsing user-typed numbers, the locale matters. \code{1,5}
is one and a half in German, fifteen in English. The cleanest
approach: pick a canonical form for storage, convert only at
display.

\begin{lstlisting}
double parse_locale_number(std::string_view s);  // tolerant
std::string format_locale_number(double n);       // locale-aware

// Internally: always double.
double value = parse_locale_number(user_input);
display = format_locale_number(value);
\end{lstlisting}

The model holds the unambiguous form; the view formats for
display.

\section{Plural-aware UIs}

For status messages with counts:

\begin{lstlisting}
auto msg = std::format(\"{} {}\",
                       std::format(\"{:L}\", count),
                       ngettext(\"file_singular\", \"file_plural\", count));
\end{lstlisting}

The translation table has separate entries for the singular and
plural forms. The format pulls the right one.

\section{The C locale}

The default C locale (\code{LC\_ALL=C}) is the lowest common
denominator: ASCII only, no comma separators, English. For
servers and scripts, this is intentional --- predictable.

For an interactive TUI, default to the user's locale. Honour
\code{LC\_ALL=C} as an explicit opt-out.

\section{Right-to-left UI shells}

If the entire UI is in Arabic or Hebrew, you may want the
\emph{layout} flipped (sidebar on the right). This is
deliberate; the flexbox model supports it via
\code{Direction::RowReverse}:

\begin{lstlisting}
auto root = h(...) | (rtl ? direction<RowReverse> : direction<Row>);
\end{lstlisting}

Each row places its children in reverse order; columns are
unaffected.

\section{Pitfalls}

\begin{warning}
\textbf{Hard-coded English in widget defaults.} If your
widget shows ``OK'' and ``Cancel'', those strings need
translation. Avoid hard-coded UX strings; route them through
\code{\_(\ldots)}.
\end{warning}

\begin{warning}
\textbf{Assuming character width.} A name field that allocates
20 characters works for English ``Alice'' but fails for the
six-cell ``\\verb|ya-ma-da|'' (a CJK transliteration). Allocate
by cell, not by character count.
\end{warning}

\begin{warning}
\textbf{Date formatting in strings.} Hard-coding ``MM/DD/YYYY''
breaks for users who expect ``DD.MM.YYYY''. Use
\code{std::format} with a locale-aware specifier.
\end{warning}

\begin{warning}
\textbf{Truncating at byte boundaries.} A UTF-8 truncation at
byte $n$ might split a multi-byte character. Always truncate at
grapheme boundaries (Chapter~\ref{ch:text}).
\end{warning}

\section*{Workshop}

\begin{enumerate}[leftmargin=*]
  \item Add a small translation table to a counter app. Switch
    locales at runtime; verify the strings change.
  \item Format a number with \code{std::format}'s \code{L}
    specifier in three locales. Note the differences.
  \item Type Arabic text in a text input widget. Verify the
    cells are counted correctly even if the visual order is
    terminal-dependent.
\end{enumerate}

\begin{recap}
i18n is three problems: translation, locale formatting, text
rendering. \code{std::format} handles locale numbers and dates;
the grapheme walker handles non-Latin text; the terminal handles
shaping and BiDi. \maya{} stays out of the way; you compose the
right tools. Default to the user's locale, opt out only when you
need predictability.
\end{recap}

\chapter{Lifecycle: from main to exit}
\label{ch:lifecycle}

\begin{aside}
A \maya{} app starts with \code{main} and ends with the
terminal restored. Between those two moments, a precise sequence
of events happens. This chapter traces them: who runs, in what
order, why each phase matters. The mental model lets you
predict where a hang, a crash, or a hot-reload should be hooked.
\end{aside}

\section{Phase 0: pre-main}

Before \code{main} runs, C++ does work:

\begin{enumerate}[leftmargin=*]
  \item Static-storage-duration objects are initialised in
    definition order within a translation unit; order between
    TUs is unspecified.
  \item \code{constinit} variables are guaranteed to be ready
    by the start of \code{main}.
  \item Function-local static objects are deferred to first use.
\end{enumerate}

\maya{} uses \code{constinit} for theme defaults and palette
tables to avoid the static-init order fiasco. Anything that
needs to be ready before \code{main} should be \code{constinit}.

\section{Phase 1: argument parsing and config}

\code{main} typically:

\begin{lstlisting}
int main(int argc, char** argv) {
    auto cfg = build_config(argc, argv);
    if (cfg.show_help) { print_help(); return 0; }
    return maya::run<App>(cfg);
}
\end{lstlisting}

Three steps:

\begin{itemize}[leftmargin=*]
  \item Parse args.
  \item Decide whether to run (or print help / version).
  \item Call \code{run<App>}.
\end{itemize}

The handoff to \code{run} is where \maya{}'s lifecycle starts.

\section{Phase 2: run setup}

\code{run<P>(cfg)} does the following before the loop starts:

\begin{enumerate}[leftmargin=*]
  \item Construct the \code{Context} (Chapter~\ref{ch:runtime}).
  \item Open the terminal and enter raw mode.
  \item Install signal handlers (SIGINT, SIGWINCH, SIGTERM).
  \item Start the worker pool.
  \item Construct the renderer, the event loop, the interpreter.
  \item Call \code{P::init()} to get the initial model and
    command.
  \item Interpret the initial command (e.g. fire a startup
    task).
\end{enumerate}

If any step fails, the runtime unwinds the prior steps and
returns an error code without entering the main loop. The
terminal is restored, the workers are joined, and resources are
released.

\subsection*{Failure modes during setup}

Common failures:

\begin{itemize}[leftmargin=*]
  \item Terminal does not exist (running under cron, no tty).
  \item Permission denied (the user's stdin is read-only).
  \item Raw mode setup fails (rare; usually permission).
  \item Worker pool cannot start (out of threads).
\end{itemize}

\maya{} propagates these as exit codes and writes a single
line to stderr explaining.

\section{Phase 3: the main loop}

The loop from Chapter~\ref{ch:runtime}:

\begin{enumerate}[leftmargin=*]
  \item Render the current model.
  \item Compute the current subscription set.
  \item Wait for an event.
  \item Translate the event into a message.
  \item Run \code{update}; replace the model.
  \item Interpret the resulting command.
  \item Repeat.
\end{enumerate}

The loop ends when an interpretation sets the running flag to
false. The only message that does this is \code{Cmd::quit}; the
typical trigger is the user pressing \code{q} or \code{Ctrl-C}.

\subsection*{What can interrupt the loop}

\begin{itemize}[leftmargin=*]
  \item Quit command (graceful exit).
  \item SIGINT / SIGTERM (signal handler sets running flag).
  \item An exception escaping the runtime (rare; logged and
    aborts).
  \item Out-of-memory (terminates).
\end{itemize}

\section{Phase 4: shutdown}

When the loop exits, destructors run in reverse construction
order:

\begin{enumerate}[leftmargin=*]
  \item \code{Interpreter}: drain pending tasks, signal workers
    to stop, join.
  \item \code{Renderer}: flush any pending bytes, reset cursor,
    restore default style.
  \item \code{EventLoop}: close epoll/kqueue/handle, clean up
    the wake fd.
  \item \code{Terminal}: leave raw mode, leave alt screen if
    entered.
  \item \code{Context}: close logger, release config.
\end{enumerate}

After destructors, \code{run} returns. The exit code is the
last value the application set (typically 0).

\section{The order matters}

If you reverse the order ---  for example, restore the terminal
before joining workers --- a worker thread might still be
writing bytes when the terminal is back in cooked mode, and the
user sees corrupted output.

The construction order in \code{run} sets the destruction
order. The pattern: construct in dependency order; let RAII
handle the reverse.

\subsection*{What about atexit?}

\code{atexit} handlers run after \code{main} returns. \maya{}
registers a fallback handler that calls \code{leave\_raw\_mode}
in case the normal destructor did not. This handles
\code{std::exit} and uncaught exceptions.

\section{Hot reload}

In development, a hot-reload pattern is sometimes useful: detect
a code change, rebuild, restart the app while preserving the
model.

The pattern:

\begin{enumerate}[leftmargin=*]
  \item The app writes its model to a file on \code{Quit}.
  \item A watcher process notices the executable changed.
  \item The watcher kills the app, rebuilds, restarts it.
  \item The new app reads the model file and restores state.
\end{enumerate}

\maya{}'s architecture makes this tractable: the model is a
plain struct, serialisable to JSON or another format.

\section{Daemonisation}

\maya{} apps are interactive; daemonising them does not make
sense. If you want background work, use \code{Cmd::task}; if
you want a long-running service, write a regular service, not
a TUI.

\section{Subprocess management}

A TUI app sometimes runs subprocesses (git, ripgrep, an LLM).
The cleanest pattern:

\begin{lstlisting}
[&](RunGitDiff) {
    return std::pair{m, Cmd<Msg>::task([] -> Msg {
        auto out = exec(\"git\", {\"diff\", \"--stat\"});
        return GitDiffResult{std::move(out)};
    })};
}
\end{lstlisting}

The subprocess runs on a worker thread; the result message
arrives back. The user does not see a frozen UI.

\subsection*{Pty subprocesses}

For interactive subprocesses (a shell, vim), you need a
pseudo-terminal. \maya{} does not ship a pty manager, but the
patterns are well-known (forkpty on POSIX, ConPTY on Win32).
Wrap the pty in a \code{Cmd::task} that pumps bytes back as
messages.

\section{Crash recovery}

If the app crashes during a long operation, the user wants:

\begin{itemize}[leftmargin=*]
  \item The terminal restored (so they can keep using it).
  \item A crash log written to disk.
  \item State restorable on restart.
\end{itemize}

The signal handler from Chapter~\ref{ch:observability} handles
the first two. The third requires periodic state saves
(Chapter~\ref{ch:config}).

\section{Exit codes}

Convention:

\begin{itemize}[leftmargin=*]
  \item \code{0}: normal exit.
  \item \code{1}: user-facing error.
  \item \code{2}: internal error (bug).
  \item \code{3}: configuration error.
  \item \code{130}: interrupted by SIGINT (128 + 2).
\end{itemize}

\maya{} returns the value from \code{run<P>(cfg)}; you control
the rest via \code{Cmd::quit\_with(code)}.

\section{Pitfalls}

\begin{warning}
\textbf{Static destructors after main.} If you have a static
object whose destructor writes to stdout, the writes happen
\emph{after} \code{main} returns. By then \maya{} has restored
the terminal but the program is still alive. The writes corrupt
the user's prompt. Avoid output-producing destructors in
statics.
\end{warning}

\begin{warning}
\textbf{Forgetting to set a quit hotkey.} An app with no
\code{q} or \code{Ctrl-C} handler can only be exited by the
user killing the process. Always provide a clean exit path.
\end{warning}

\begin{warning}
\textbf{Holding the terminal during a long task.} If
\code{update} takes seconds, the UI freezes and the user
suspects a hang. Move long tasks to \code{Cmd::task}; show a
spinner while they run.
\end{warning}

\begin{warning}
\textbf{Joining workers in main thread before they finish.}
The interpreter's destructor waits for workers to drain.
If a worker is stuck, the shutdown hangs. Use timeouts on
tasks that might block forever.
\end{warning}

\section*{Workshop}

\begin{enumerate}[leftmargin=*]
  \item Add startup logging to a small app: log each phase as
    it starts. Verify the order matches this chapter.
  \item Force a crash in \code{update}. Confirm the terminal
    is still usable after the process dies.
  \item Add a hot-reload cycle: persist model on quit; restore
    on startup; observe the seamless restart.
\end{enumerate}

\begin{recap}
A \maya{} app's lifecycle has four phases: pre-main static
init, run setup, main loop, shutdown. Each phase has clear
ordering and clear failure modes. RAII handles shutdown order
correctly. The architecture invites long-running tasks to be
\code{Cmd::task}s and signals to be subscriptions; the loop
itself is small and uniform.
\end{recap}

\chapter{Accessibility and inclusive design}
\label{ch:a11y}

\begin{aside}
A TUI is often the most accessible kind of UI: screen readers
work with it, low-vision users can configure terminal fonts and
contrast, and there is no mouse-only flow. But ``most accessible
by default'' is not ``always accessible''. A few choices in your
app design can welcome or exclude users with disabilities. This
chapter walks the considerations and the patterns.
\end{aside}

\section{Why TUIs start accessible}

A typical TUI:

\begin{itemize}[leftmargin=*]
  \item Renders text, which screen readers can read.
  \item Uses high-contrast cells (the user picks the colours).
  \item Has keyboard navigation by default.
  \item Avoids animation that distracts.
  \item Works with the user's chosen font size.
\end{itemize}

These are wins by default. The patterns to maintain them
require thought.

\section{Screen readers}

Modern screen readers (Orca on Linux, JAWS / NVDA on Windows,
VoiceOver on macOS) can read terminal output, but they only see
what is on the screen \emph{after} ANSI escapes are interpreted.

A few principles:

\begin{itemize}[leftmargin=*]
  \item Information conveyed by colour alone is invisible to
    screen reader users. Pair it with a text marker.
  \item Status indicators that flash (using SGR 5 blink) are
    not announced. Use static markers.
  \item Boxes drawn with Unicode borders read as random
    characters. Provide an option to hide borders or use the
    \code{ASCII} border style.
\end{itemize}

\subsection*{Live regions}

When the UI updates rapidly (a log streaming, a chat coming
in), screen readers can drop messages. Patterns:

\begin{itemize}[leftmargin=*]
  \item Only announce new lines (no re-reads of the whole
    panel).
  \item Provide a ``pause output'' toggle so the user can
    catch up.
  \item Offer a transcript view that does not auto-scroll.
\end{itemize}

\section{Colour and contrast}

\code{NO\_COLOR=1} is a hint, not a request. Honour it
strictly:

\begin{lstlisting}
if (std::getenv(\"NO_COLOR\")) {
    cfg.theme = themes::Monochrome();
}
\end{lstlisting}

The monochrome theme uses only attributes (Bold, Underline,
Dim) for distinction. Information that was conveyed by colour
is still distinguishable.

\subsection*{Contrast ratios}

WCAG 2.1 defines contrast ratios:

\begin{itemize}[leftmargin=*]
  \item AA Normal text: 4.5:1.
  \item AA Large text: 3:1.
  \item AAA Normal text: 7:1.
\end{itemize}

For a TUI, ``large text'' rarely applies (cells are uniform).
Aim for 4.5:1 minimum on default themes; provide a
high-contrast option at 7:1 for users who need it.

\subsection*{Colour blindness}

About 8\% of men and 0.5\% of women have some form of colour
blindness. Common patterns:

\begin{itemize}[leftmargin=*]
  \item Do not rely on red vs green alone. Pair with shapes
    (\code{[+]} vs \code{[-]}).
  \item Use distinct hues, not adjacent ones (red + blue, not
    red + orange).
  \item Test palettes in a simulator (Color Oracle, Sim Daltonism).
\end{itemize}

\section{Keyboard navigation}

A TUI cannot rely on a mouse; mouse support is a convenience.
The keyboard must reach every interactive element.

\subsection*{Tab order}

Document the tab order explicitly. The visible order should
match the logical order:

\begin{itemize}[leftmargin=*]
  \item Top to bottom.
  \item Left to right (or right to left in RTL languages).
  \item Skip non-interactive elements.
\end{itemize}

\maya{}'s \code{FocusManager} (Chapter~\ref{ch:focus}) stores
order explicitly. Set it to match the visual order; do not
rely on registration order alone.

\subsection*{Shortcut redundancy}

For every shortcut, ask: can the user discover it? Patterns:

\begin{itemize}[leftmargin=*]
  \item A help overlay listing all bindings.
  \item Status bar hints (e.g. \code{[q] quit  [/] search}).
  \item A command palette (Ctrl-P) so commands are discoverable
    by name even if the user does not know the binding.
\end{itemize}

\section{Animation and motion}

Animation can:

\begin{itemize}[leftmargin=*]
  \item Distract or trigger discomfort for users with motion
    sensitivity.
  \item Mask important information.
  \item Slow down task completion.
\end{itemize}

Honour \code{prefers-reduced-motion}:

\begin{lstlisting}
if (std::getenv(\"MAYA_REDUCED_MOTION\")) {
    cfg.animations_enabled = false;
}
\end{lstlisting}

When disabled, spinners become static markers, smooth scrolling
becomes instant, transitions are skipped.

\subsection*{Avoid blinking}

SGR 5 (blink) was deprecated in many terminals and is generally
considered hostile. Do not use it.

\section{Font size and readability}

The user controls font size at the terminal level. Your app
should:

\begin{itemize}[leftmargin=*]
  \item Adapt to any terminal width (no hard-coded 80 columns).
  \item Wrap or truncate gracefully.
  \item Avoid tiny indicator characters (a single \code{x}
    when a clearer indicator like \code{[done]} is available).
\end{itemize}

\section{Time-sensitive UI}

If something requires the user to act quickly:

\begin{itemize}[leftmargin=*]
  \item Default timeouts should be generous.
  \item Provide a way to extend.
  \item Show the remaining time clearly.
\end{itemize}

Patterns like a confirmation prompt that closes after 5 seconds
are hostile to users who read slowly.

\section{Error messages}

Error messages should:

\begin{itemize}[leftmargin=*]
  \item Be visible to screen readers (not just colour).
  \item Persist long enough to read.
  \item Suggest a remedy (\emph{press r to retry}).
  \item Not blame the user.
\end{itemize}

A pattern: errors in a status bar with a fixed duration plus a
\code{[?]} indicator that opens a detail panel.

\section{Cognitive load}

Reduce things to remember:

\begin{itemize}[leftmargin=*]
  \item Consistent terminology across screens.
  \item Visible state (don't hide what mode you are in).
  \item Confirmations for destructive actions.
  \item Undo for everything reversible.
\end{itemize}

The Elm architecture's explicit \code{Model} is a gift here:
the state is a value the user can be shown.

\section{Internationalisation revisited}

Translation (Chapter~\ref{ch:i18n}) is an accessibility
concern: users whose first language is not the app's default
need translations.

The pattern: wrap all UX strings; ship at least one alternate
locale to prove the wrapping works.

\section{Test with assistive tech}

A few practical tests:

\begin{itemize}[leftmargin=*]
  \item Try the app with a screen reader.
  \item Try with no colour (\code{NO\_COLOR=1}).
  \item Try with a 16-colour terminal.
  \item Try with a 40-column-wide terminal.
  \item Try with motion disabled.
\end{itemize}

If any of these break the app, the app is not accessible.

\section{Pitfalls}

\begin{warning}
\textbf{Colour-only status indicators.} The most common
mistake: a red dot to mean ``error''. Add text or a shape so
non-colour users can see it.
\end{warning}

\begin{warning}
\textbf{Hard-coded narrow widths.} Apps that assume 80 columns
or 120 columns break at other sizes. \maya{}'s flex layout
helps; do not undo it with \code{width<80>}.
\end{warning}

\begin{warning}
\textbf{Rapid auto-updates.} A log scrolling past faster than
the user can read is hostile. Provide a pause toggle.
\end{warning}

\begin{warning}
\textbf{Modal dialogs that grab focus.} A modal that prevents
\code{Ctrl-C} from quitting traps the user. Always allow the
quit hotkey, even inside modals.
\end{warning}

\section*{Workshop}

\begin{enumerate}[leftmargin=*]
  \item Run your app with \code{NO\_COLOR=1}. Identify every
    place colour-only conveyed meaning. Add text or shapes.
  \item Test the app with Orca (Linux) or NVDA (Windows).
    Note which controls are not announced.
  \item Add a help overlay that lists every keyboard shortcut.
    Trigger it on \code{F1} or \code{?}.
\end{enumerate}

\begin{recap}
TUIs start accessible by default. Maintain that with: text
markers paired with colour, monochrome theme on NO\_COLOR,
explicit tab order, discoverable shortcuts, generous timeouts,
motion controls, persistent and respectful error messages.
Test with assistive tech. The Elm architecture's explicit
state makes accessibility easier, not harder.
\end{recap}

\chapter{Security and untrusted input}
\label{ch:security}

\begin{aside}
A TUI app processes input from many sources: the user's
keyboard, the user's clipboard, files the user opens, network
data the app fetches, output from subprocesses. Any of these
can be hostile. This chapter is a short tour of the patterns
\maya{} encourages for keeping the app safe.
\end{aside}

\section{What \emph{can} go wrong}

A handful of failure modes are specific to terminal apps:

\begin{itemize}[leftmargin=*]
  \item \textbf{ANSI injection.} Content that contains escape
    codes can change the terminal's state in ways your code
    did not authorise. \code{\\e[2J} clears the screen;
    \code{\\e]52;c;\ldots\\e\\textbackslash} writes the
    clipboard.
  \item \textbf{Format string vulnerabilities.} If you pass
    untrusted text to \code{printf}-like functions as the
    format string, you have an information leak.
  \item \textbf{Buffer overflows in C interop.} If you call C
    code with the wrong size, you can corrupt memory.
  \item \textbf{Subprocess injection.} If you build a command
    line by concatenating user input, you create shell-
    injection paths.
  \item \textbf{Path traversal.} If you open a file by
    user-supplied path, the user can read anywhere your
    process has rights.
\end{itemize}

Each has a specific mitigation.

\section{ANSI injection}

The classic. A log line containing
\code{\\e[2J\\e[H} clears the screen and homes the cursor.
A pasted blob with \code{\\e]52;c;<base64>\\e\\textbackslash}
writes the clipboard.

\subsection*{Mitigation: strip on display}

\maya{} provides \code{sanitize\_text}:

\begin{lstlisting}
std::string sanitize_text(std::string_view input) {
    std::string out;
    out.reserve(input.size());
    for (size_t i = 0; i < input.size(); ++i) {
        unsigned char c = input[i];
        if (c == 0x1B) { i++; while (i < input.size() && !is_terminator(input[i])) i++; continue; }
        if (c < 0x20 && c != '\\t' && c != '\\n') continue;
        out += c;
    }
    return out;
}
\end{lstlisting}

Strip every escape sequence; strip every control character
except tab and newline. The result is safe to display.

\maya{} runs \code{sanitize\_text} on:

\begin{itemize}[leftmargin=*]
  \item Pasted text (from bracketed paste).
  \item File contents you display (when wrapped in
    \code{text\_safe}).
  \item Network-fetched strings via \code{Cmd::task}.
\end{itemize}

\subsection*{Mitigation: lock down the terminal}

Some terminals support a ``deny mode'' for clipboard writes
(tmux \code{set-clipboard off}). Document this for users who
display untrusted content frequently.

\section{Format string safety}

C++26's \code{std::format} is type-safe: the format string must
match the argument count and types, checked at compile time:

\begin{lstlisting}
std::format(\"{}\", value);                     // OK
std::format(input, value);                     // ERROR
\end{lstlisting}

The second form does not compile because \code{input} is not a
literal. Use \code{std::vformat} only when you genuinely need
a runtime format string; even then, validate it first.

\maya{} provides a wrapper that rejects suspicious format
strings (containing \code{\%n}, etc.). For ordinary use,
\code{std::format} is the safe path.

\section{Subprocess injection}

\code{std::system} and \code{popen} are unsafe with user input:

\begin{lstlisting}
std::system(\"grep \" + user_pattern + \" file\");  // injection!
\end{lstlisting}

If \code{user\_pattern} is \code{;rm -rf /}, the result is a
disaster.

\subsection*{Mitigation: argv-based exec}

\maya{} ships an exec helper that takes a vector, not a string:

\begin{lstlisting}
auto result = exec(\"grep\", {user_pattern, \"file\"});
\end{lstlisting}

The shell is never involved. Each element of the vector is one
argv slot. No quoting, no escaping, no injection.

On POSIX this maps to \code{execvp}; on Win32 to
\code{CreateProcessW} with the argv re-encoded.

\section{Path traversal}

\code{file::read(\"~/data/\" + filename)} where \code{filename}
is \code{../../etc/passwd} is a problem.

\subsection*{Mitigation: canonicalise and verify}

\begin{lstlisting}
std::filesystem::path requested = base / filename;
std::filesystem::path canonical = std::filesystem::weakly_canonical(requested);
if (!canonical.string().starts_with(base.string())) {
    return std::unexpect, IoError{\"path escape\"};
}
\end{lstlisting}

\code{weakly\_canonical} resolves \code{..} segments. After
resolution, check that the result lives under the intended
base directory.

\section{Untrusted serialised data}

A JSON or TOML file the user supplies could be malicious:

\begin{itemize}[leftmargin=*]
  \item Deeply nested objects can stack-overflow a recursive
    parser.
  \item Large arrays can exhaust memory.
  \item Specific patterns can trigger expensive regex matches
    (\emph{ReDoS}).
\end{itemize}

Use a strict parser with explicit limits:

\begin{lstlisting}
JsonOptions opt;
opt.max_depth = 64;
opt.max_array_size = 100'000;
opt.max_string_size = 1'000'000;

auto r = parse_json(input, opt);
\end{lstlisting}

The parser rejects content that exceeds limits with a clear
error.

\section{Memory safety}

C++ does not give you memory safety by default. Patterns that
help:

\begin{itemize}[leftmargin=*]
  \item Use \code{std::string\_view}, \code{std::span},
    \code{std::vector}, \code{std::array} instead of raw
    pointers and lengths.
  \item Enable sanitizers (\code{-fsanitize=address,undefined})
    in debug builds.
  \item Use \code{static\_assert} liberally to catch size
    mistakes at build.
  \item Bounds-check in places that touch external data
    (parsers, network, file readers).
\end{itemize}

\maya{} is internally written in this style. External callers
should follow suit.

\section{Cryptography}

\maya{} does not ship crypto. If your app needs it:

\begin{itemize}[leftmargin=*]
  \item Use \code{libsodium} for general-purpose crypto.
  \item Use the OS's TLS library for HTTPS (\code{libcurl}
    + system trust store).
  \item Do not roll your own. Ever.
\end{itemize}

\section{Privilege management}

If your app runs as root (unusual for a TUI), drop privileges
as early as possible:

\begin{itemize}[leftmargin=*]
  \item \code{setuid}/\code{setgid} to a non-privileged user.
  \item \code{chroot} into a restricted directory if you only
    need a subtree.
  \item Use \code{seccomp} or platform sandboxing to restrict
    system calls.
\end{itemize}

For ordinary apps, the user already runs you with their own
privileges. The risk is in what you do with those privileges.

\section{Logging secrets}

Already covered (Chapter~\ref{ch:observability}). Redact:

\begin{itemize}[leftmargin=*]
  \item API tokens.
  \item Passwords (never log these in any form).
  \item Personal data (emails, addresses).
  \item File contents in error messages (a path is fine; the
    body is not).
\end{itemize}

\section{Supply chain}

Every dependency is an attack surface. Mitigations:

\begin{itemize}[leftmargin=*]
  \item Pin versions in the build system; do not float to
    latest.
  \item Vendor critical dependencies (commit them to your
    repository).
  \item Audit periodically (CVE databases, \code{cargo audit}
    equivalents for C++).
  \item Limit the dependency surface --- prefer the standard
    library where possible.
\end{itemize}

\maya{} itself has a small dependency footprint: the standard
library, optional toml++ for config, optional xxHash for
fingerprints. Each is auditable.

\section{Pitfalls}

\begin{warning}
\textbf{Displaying raw subprocess output.} If \code{git diff}
returns content with ANSI codes (which it does for colour),
those codes will be interpreted by the terminal. Sanitize
first.
\end{warning}

\begin{warning}
\textbf{Writing user-controlled text to the clipboard.} If a
user pastes a malicious blob and you write it to the clipboard
unmodified, the next paste somewhere else can carry the
payload. Validate before writing.
\end{warning}

\begin{warning}
\textbf{Trusting file paths.} A symlink can point anywhere.
Resolve fully before opening.
\end{warning}

\begin{warning}
\textbf{Verbose error messages with paths.} A bug report
including the absolute path of a user's home directory leaks
information. Print relative paths or anonymised paths.
\end{warning}

\section*{Workshop}

\begin{enumerate}[leftmargin=*]
  \item Add \code{sanitize\_text} to a log viewer. Test with a
    file containing ANSI escapes; verify they are stripped.
  \item Replace a \code{std::system} call with the argv-based
    \code{exec} helper. Confirm the same behaviour without
    shell involvement.
  \item Audit your dependency list. For each, note the last
    update and any known CVEs.
\end{enumerate}

\begin{recap}
Untrusted input is everywhere: keyboard, clipboard, files,
network, subprocesses. Each has a mitigation: sanitize ANSI
escapes, use type-safe formatting, argv-based exec, canonical
path checks, strict parsers with limits, redacted logging. The
mitigations are small individually; together they keep the app
out of the news.
\end{recap}

\chapter{Cross-cutting concerns: logging, time, ID}
\label{ch:cross-cutting}

\begin{aside}
Some concerns do not belong to a single architectural layer.
Time, logging, ID generation, randomness --- each is touched
by many components, each is a source of bugs if treated
carelessly. This chapter pulls them together and shows the
patterns \maya{} uses for each.
\end{aside}

\section{Time}

A surprising amount of TUI code touches time:

\begin{itemize}[leftmargin=*]
  \item Animations advance over time.
  \item Timeouts fire after a duration.
  \item Timestamps display ``now''.
  \item Debounce windows compare ``now'' to previous.
  \item Profiling measures elapsed time.
\end{itemize}

C++26 gives several clocks:

\begin{itemize}[leftmargin=*]
  \item \code{std::chrono::system\_clock}: wall-clock time.
    Can jump backwards (NTP, user adjustment).
  \item \code{std::chrono::steady\_clock}: monotonically
    increasing. Cannot jump back. Use for durations.
  \item \code{std::chrono::high\_resolution\_clock}: alias for
    one of the above; system-dependent.
\end{itemize}

\subsection*{When to use which}

\begin{itemize}[leftmargin=*]
  \item Display: \code{system\_clock} (the user wants to see
    actual wall-clock).
  \item Animations, timeouts, debounce, profiling:
    \code{steady\_clock}.
\end{itemize}

Mixing them is a bug source. A debounce that uses
\code{system\_clock} can fire repeatedly or never if the clock
jumps.

\subsection*{Wrapping time for testability}

For test determinism, abstract the clock:

\begin{lstlisting}
class Clock {
public:
    virtual std::chrono::steady_clock::time_point now() const = 0;
};

class RealClock : public Clock {
    std::chrono::steady_clock::time_point now() const override {
        return std::chrono::steady_clock::now();
    }
};

class FakeClock : public Clock {
    std::chrono::steady_clock::time_point current_;
public:
    void advance(std::chrono::milliseconds d) { current_ += d; }
    std::chrono::steady_clock::time_point now() const override {
        return current_;
    }
};
\end{lstlisting}

Inject the clock into modules that care; use \code{RealClock}
in production, \code{FakeClock} in tests.

\section{Logging}

Already covered (Chapter~\ref{ch:observability}). The
cross-cutting point: every layer of your app may want to log,
not just the runtime. Make logging accessible:

\begin{itemize}[leftmargin=*]
  \item Pass a \code{Logger\&} via the context.
  \item Inside widgets, accept an optional \code{Logger*}.
  \item For tests, supply a \code{NullLogger}.
\end{itemize}

\subsection*{Avoiding global loggers}

A global \code{spdlog::get(\"app\")} pattern is convenient but
breaks test isolation. The injection pattern is better; one
logger per app instance, one app per test.

\section{IDs}

Generating unique IDs comes up everywhere:

\begin{itemize}[leftmargin=*]
  \item Focus IDs.
  \item Element keys for reconciliation.
  \item Database row IDs.
  \item Correlation IDs for log entries.
\end{itemize}

\subsection*{Patterns}

\begin{enumerate}[leftmargin=*]
  \item \textbf{Monotonic counter.} Simple, safe within one
    process. Not unique across runs.
  \item \textbf{UUID.} Unique across processes, machines, time.
    16 bytes; \code{libuuid} or \code{stduuid}.
  \item \textbf{Hash of content.} Deterministic; same input
    gives same ID. Useful for cache keys.
  \item \textbf{Timestamp + counter.} Sortable IDs.
    Snowflake-style.
\end{enumerate}

\subsection*{Monotonic counter implementation}

\begin{lstlisting}
class IdGenerator {
    std::atomic<uint64_t> next_{1};
public:
    uint64_t next() { return next_.fetch_add(1, std::memory_order_relaxed); }
};
\end{lstlisting}

Thread-safe via atomic. Zero is reserved for ``invalid'';
counts up from 1.

\section{Randomness}

\code{std::rand} is broken (not really random, not thread-safe).
Use \code{<random>}:

\begin{lstlisting}
thread_local std::mt19937 rng{std::random_device{}()};
\end{lstlisting}

Each thread gets its own seeded engine. Use
\code{std::uniform\_int\_distribution} for ranges.

\subsection*{Reproducibility}

For tests, seed the RNG deterministically:

\begin{lstlisting}
std::mt19937 rng{42};
\end{lstlisting}

Same seed, same sequence. The test fails deterministically when
it fails at all.

\subsection*{Cryptographic randomness}

\code{std::mt19937} is not cryptographically secure. For tokens
or session IDs, use the OS:

\begin{itemize}[leftmargin=*]
  \item Linux: \code{getrandom}.
  \item BSD: \code{arc4random}.
  \item Win32: \code{BCryptGenRandom}.
\end{itemize}

The platform abstraction layer should provide a thin wrapper.

\section{Errors}

Cross-cutting because every layer can produce them. Patterns:

\begin{itemize}[leftmargin=*]
  \item Use \code{Expected<T, E>} for recoverable failures.
  \item Use exceptions for unrecoverable bugs (assertion
    failures).
  \item Log every error at the boundary it crossed; do not
    re-log as it propagates.
  \item Surface to the user via the UI, not via stderr.
\end{itemize}

\section{Configuration access}

Already covered (Chapter~\ref{ch:config}). The cross-cutting
point: any layer might need to read config. The pattern:

\begin{itemize}[leftmargin=*]
  \item The model holds the config struct.
  \item \code{view} reads from \code{model.config.\ldots}.
  \item Widgets that need config take it as a parameter.
\end{itemize}

Do not reach into a global; thread the config through.

\section{Caches}

A cache is cross-cutting: it sits between a producer and a
consumer. Patterns:

\begin{itemize}[leftmargin=*]
  \item \textbf{Memoise} pure functions.
    \code{std::unordered\_map<Key, Result>}.
  \item \textbf{LRU caches} for bounded memory.
  \item \textbf{Signal-driven} caches that invalidate when
    inputs change (Chapter~\ref{ch:reactivity-preview}).
\end{itemize}

The trade-off is always memory vs CPU. Profile before adding a
cache; the framework's fingerprint cache often makes app-level
caches unnecessary.

\section{Locking and synchronisation}

\maya{}'s threading model (Chapter~\ref{ch:runtime}) keeps most
synchronisation invisible. Where you need to share state across
threads:

\begin{itemize}[leftmargin=*]
  \item \code{std::atomic} for single primitives.
  \item \code{std::mutex} + \code{std::lock\_guard} for short
    critical sections.
  \item \code{std::shared\_mutex} for read-heavy workloads.
  \item Lock-free queues for high-throughput producer/consumer.
\end{itemize}

The pattern: take the lock for the minimum scope; never call
out (especially not to user code) while holding a lock.

\section{Async cleanup}

Resources that need explicit cleanup (file descriptors, sockets,
GPU handles) should use RAII:

\begin{lstlisting}
class FdGuard {
    int fd_;
public:
    explicit FdGuard(int fd) : fd_{fd} {}
    ~FdGuard() { if (fd_ >= 0) ::close(fd_); }
    FdGuard(const FdGuard&) = delete;
    FdGuard(FdGuard&& o) noexcept : fd_{std::exchange(o.fd_, -1)} {}
};
\end{lstlisting}

The destructor handles cleanup; move-only semantics prevent
double-close. Use \code{std::unique\_ptr} with a custom deleter
for non-trivial resources.

\section{Validation}

Whenever data crosses a trust boundary (network, file, user
input), validate. Patterns:

\begin{itemize}[leftmargin=*]
  \item Validate at the boundary, not deep in the call stack.
  \item Reject early; return \code{Expected<\ldots>::error}.
  \item Log the rejection (so you can spot patterns).
  \item Never silently truncate or coerce; that hides bugs.
\end{itemize}

\section{Telemetry}

If your app reports anonymous usage stats, the patterns from
Chapter~\ref{ch:security} apply: redact, opt-in by default
in some jurisdictions, transparent about what is sent.

The simplest pattern: log a structured event every time a
meaningful action happens; upload batches in
\code{Cmd::isolated\_task}. Privacy and bandwidth both win.

\section{Memory pressure}

When the system is low on memory:

\begin{itemize}[leftmargin=*]
  \item \maya{}'s frame buffers can be shrunk on demand.
  \item Logs can be trimmed.
  \item Caches can be flushed.
\end{itemize}

The Linux kernel signals via cgroups and OOM thresholds; the
app can register a handler. For most apps, this is overkill;
mention it for the few that need it.

\section{Pitfalls}

\begin{warning}
\textbf{Time arithmetic.} Subtracting two
\code{std::chrono::time\_point}s of different clocks does not
compile (good). Subtracting two \code{system\_clock}
timestamps that span a clock jump gives nonsense (bad). Use
\code{steady\_clock} for durations.
\end{warning}

\begin{warning}
\textbf{Global RNG seeded once.} A single global RNG seeded
at startup gives every call the same sequence relative to
program order. Tests become flaky as soon as call order
changes. Use \code{thread\_local} engines or pass a seed.
\end{warning}

\begin{warning}
\textbf{Caching too eagerly.} A cache that never invalidates
holds stale data forever. Pair every cache with an
invalidation rule.
\end{warning}

\begin{warning}
\textbf{Validating in the wrong place.} Validating deep in a
function means callers have already done partial work that
must be undone. Validate at the entry point.
\end{warning}

\section*{Workshop}

\begin{enumerate}[leftmargin=*]
  \item Inject a fake clock into a debounce helper. Test that
    the debounce fires exactly when the clock crosses the
    threshold.
  \item Generate IDs with a thread-safe counter. Verify
    uniqueness under high concurrency.
  \item Wrap a file descriptor in a \code{FdGuard}. Test that
    closing happens deterministically even on early return.
\end{enumerate}

\begin{recap}
Time, logging, IDs, randomness, errors, config, caches, locks,
cleanup, validation, telemetry, memory pressure --- each is a
cross-cutting concern. None belongs to one layer. The patterns:
inject (do not globalise), validate at the boundary, RAII for
resources, atomic or thread-local for shared state. Test with
fakes; production uses the real thing.
\end{recap}

\chapter{Documentation and developer experience}
\label{ch:devx}

\begin{aside}
A library that is easy to use is not the same as a library that
is easy to learn. \maya{}'s architecture is small, but the
surface area --- DSL, builders, widgets, signals, themes ---
takes a while to absorb. This chapter is about \emph{developer
experience}: how to document a \maya{} library, how to write
examples, how to teach the architecture. Part of the book's job
is to teach by example; here we name the techniques.
\end{aside}

\section{Three audiences}

When you write docs for a \maya{}-based project, you are writing
for three audiences:

\begin{itemize}[leftmargin=*]
  \item \textbf{Users.} People who run your app. They need
    the keyboard shortcuts, the troubleshooting, the upgrade
    guide.
  \item \textbf{Integrators.} People who embed your library
    in their app. They need the API reference, the example
    code, the architectural overview.
  \item \textbf{Contributors.} People who modify your source.
    They need the design rationale, the testing setup, the
    build instructions.
\end{itemize}

A single \code{README.md} cannot serve all three. Split them.

\section{The README}

For \maya{} projects, the README answers ``what is this and
why should I care?'' in 60 seconds:

\begin{itemize}[leftmargin=*]
  \item One-line description.
  \item Screenshot or recorded animation.
  \item Install command.
  \item First-run instructions.
  \item Links to deeper docs.
\end{itemize}

That is it. Long READMEs lose readers. Short READMEs invite the
right next click.

\section{API reference}

For each public type or function, document:

\begin{itemize}[leftmargin=*]
  \item One-sentence purpose.
  \item Parameters: types, units, valid ranges.
  \item Return value.
  \item Errors / exceptions.
  \item Example of typical use.
\end{itemize}

Doxygen extracts most of this from comments. The discipline:
write the comment before the code, not after.

\subsection*{Comment style}

\begin{lstlisting}
/// Render an element to ANSI bytes.
///
/// Lays out the element against the current terminal size,
/// diffs against the previous canvas, and writes the
/// minimum-change byte sequence.
///
/// \\param e   the element to render.
/// \\returns  the number of bytes written.
/// \\throws   io_error on terminal write failure.
size_t render(const Element& e);
\end{lstlisting}

Brief, parameter, return, throws. The Doxygen tags are
optional; the structure matters.

\section{Examples}

The fastest way to learn an API is to see it used. For each
major feature, ship at least one example:

\begin{itemize}[leftmargin=*]
  \item A minimal sample that demonstrates the feature in
    isolation.
  \item A realistic sample showing typical composition.
  \item Tests that double as examples.
\end{itemize}

\maya{}'s \code{examples/} directory holds small one-file
programs: counter, todo, file-browser, dashboard. Each
exercises a slice of the library.

\subsection*{Runnable in seconds}

Examples should build and run without setup. A \code{cmake}
sub-target builds each example; \code{make examples} produces
binaries; the README shows the commands.

\section{Architecture documents}

Beyond per-function docs, your project needs a few
architecture-level documents:

\begin{itemize}[leftmargin=*]
  \item \textbf{Overview.} The big picture in one page.
  \item \textbf{Decision records.} Why we chose X over Y.
  \item \textbf{Concepts.} Vocabulary specific to your app.
  \item \textbf{Sequence diagrams.} For complex flows
    (startup, network sync, crash recovery).
\end{itemize}

The book you are reading is an architecture document for
\maya{} itself. For your own apps, a few-page overview is
usually enough.

\subsection*{Decision records (ADRs)}

A lightweight format for design decisions:

\begin{itemize}[leftmargin=*]
  \item Title.
  \item Context: what was the situation?
  \item Decision: what did we choose?
  \item Consequences: what does this mean for callers?
  \item Status: proposed, accepted, deprecated, superseded.
\end{itemize}

A six-paragraph ADR is enough. Keep them in
\code{docs/adr/}; number them sequentially.

\section{Tutorials}

Tutorials are different from reference. They take a learner
from zero to a working application via a guided sequence:

\begin{itemize}[leftmargin=*]
  \item State assumed knowledge upfront.
  \item Build one feature at a time.
  \item Show the diff at each step.
  \item Explain \emph{why}, not just \emph{how}.
\end{itemize}

A good tutorial leaves the learner with a working app and the
confidence to extend it.

\section{Error message quality}

User-facing errors are docs too. A bad error:

\begin{lstlisting}
error: parse failed
\end{lstlisting}

A good error:

\begin{lstlisting}
error: parse failed at config.toml:14:5
  expected: a string for key 'theme'
  found:    integer 42
  hint:     theme = \"dark\"
\end{lstlisting}

Three things distinguish them: location, expected vs found,
a hint. Build the helper once; use it for every parser, schema,
or validator in your app.

\subsection*{Compile errors}

For libraries that lean on templates, compile errors are
docs. \maya{}'s concept names appear in errors; the names are
chosen to be human-readable. Apply the same discipline to your
own templates.

\section{Changelog}

Every release needs a changelog. The format that wins:

\begin{lstlisting}
## [0.5.0] - 2025-03-15

### Added
- Mouse focus on click.
- New `picker` widget with fuzzy matching.

### Changed
- Default theme now uses 16-colour fallback on dumb terminals.

### Fixed
- Crash on resize during animation.
\end{lstlisting}

Sections by kind, dates in ISO format. The first thing users
read after upgrading; the first thing developers grep before
upgrading.

\section{Onboarding contributors}

The \code{CONTRIBUTING.md}:

\begin{itemize}[leftmargin=*]
  \item Build instructions.
  \item Test instructions.
  \item Style guide.
  \item How to file an issue / open a PR.
  \item Areas that need help.
\end{itemize}

Make the first PR easy. A ``good first issue'' label and a
sample issue with hand-holding goes a long way.

\section{Versioning}

Semantic versioning is the convention:

\begin{itemize}[leftmargin=*]
  \item \code{MAJOR}: incompatible API changes.
  \item \code{MINOR}: new functionality, backward-compatible.
  \item \code{PATCH}: bug fixes.
\end{itemize}

Pin dependencies to compatible versions; bump \code{MAJOR} only
when truly necessary.

\section{Deprecation policy}

Things you deprecate should:

\begin{itemize}[leftmargin=*]
  \item Be announced one minor version before removal.
  \item Print a warning when used.
  \item Have a documented replacement.
  \item Be removed only in a major version.
\end{itemize}

The C++ attribute \code{[[deprecated(\"use foo instead\")]]}
generates compile-time warnings.

\section{Building trust}

Documentation is also a trust signal. Users adopt libraries
that:

\begin{itemize}[leftmargin=*]
  \item Are well-tested (visible CI).
  \item Have a changelog.
  \item Respond to issues.
  \item Have other users (a few public projects using it).
\end{itemize}

You cannot fake any of these; build them over time.

\section{Pitfalls}

\begin{warning}
\textbf{Documenting only the happy path.} The error paths are
often where users get stuck. Document failure modes
explicitly.
\end{warning}

\begin{warning}
\textbf{Stale examples.} Examples that compiled in version
0.1 but not in 0.5 are worse than no examples. CI should
build every example on every commit.
\end{warning}

\begin{warning}
\textbf{Magical docs.} ``Just call \code{run}'' is not
documentation. Explain what \code{run} does, what it requires,
what it produces.
\end{warning}

\begin{warning}
\textbf{Code without comments.} For non-trivial logic, a
sentence explaining \emph{why} saves future-you hours. Avoid
restating \emph{what} the code already says.
\end{warning}

\section*{Workshop}

\begin{enumerate}[leftmargin=*]
  \item Take a function from your project. Write a 4-line
    doc comment: purpose, params, return, error.
  \item Write a 6-paragraph ADR for one design choice in your
    project.
  \item Add an example to your repo. Wire it into CI so it
    builds on every commit.
\end{enumerate}

\begin{recap}
A library is its docs plus its code. Write for three audiences:
users, integrators, contributors. Keep READMEs short; document
APIs with purpose, parameters, returns, errors; ship runnable
examples; record architecture decisions; produce a changelog.
The discipline pays back when someone else adopts your library
or when future-you returns to it.
\end{recap}

\chapter{The architecture, in three sentences}
\label{ch:summary}

\begin{aside}
You have just spent thirty chapters on Part II. Before turning
the page to Part III, it is worth distilling everything to its
shortest form. If you remember only three things, remember
these.
\end{aside}

\section{Sentence 1: the data flow}

\begin{quote}
A \maya{} program is four pure functions
(\code{init}, \code{update}, \code{view}, \code{subscribe}) over
two sum types (\code{Msg}, \code{Element}) plus two value-typed
effect languages (\code{Cmd}, \code{Sub}).
\end{quote}

Everything in Part II descended from that sentence. The Elm
loop, the DSL, elements, layout, concepts, expected, the
pipeline, style, builders, text, focus, scroll, error
boundaries, signals, the runtime, widgets, performance,
patterns --- all of them are either inside one of those types
or transformations between them.

\subsection*{Why the four functions}

Each one separates a concern that becomes a bug when entangled:

\begin{itemize}[leftmargin=*]
  \item \code{init}: how does the program start? Separated so
    boot-time state has its own home.
  \item \code{update}: how does state evolve? Separated so
    business logic is testable in isolation.
  \item \code{view}: how does state become pixels? Separated so
    the renderer is a pure function of state.
  \item \code{subscribe}: how does the program listen?
    Separated so external events are declarative.
\end{itemize}

Combine any two and you lose the property. Keeping them apart
is the discipline.

\subsection*{Why two sum types}

\code{Msg} is the alphabet of state transitions. Each
alternative is one kind of thing that can happen. The compiler
checks exhaustiveness; new alternatives must be handled.

\code{Element} is the alphabet of UI shapes. Each alternative is
one kind of thing the renderer can draw. The same exhaustiveness
property applies.

Together they form the program's two languages: one for
inputs, one for outputs. Both checked at compile time.

\subsection*{Why effects as data}

\code{Cmd} and \code{Sub} are descriptions of side effects,
not the effects themselves. The runtime interprets them. This
buys:

\begin{itemize}[leftmargin=*]
  \item Tests without the OS.
  \item Replay of any session.
  \item Composition without callbacks.
  \item A single place to add a new effect type.
\end{itemize}

\section{Sentence 2: the renderer}

\begin{quote}
The renderer is a six-stage pipeline (view, layout, paint,
diff, serialize, write) over a packed-cell canvas, fingerprinted
to skip whole frames when nothing changed.
\end{quote}

Part II touched the renderer at the edges; Part III walks
inside it. The architectural point: each stage is a pure
function from one value to another, the values are inspectable,
and the whole pipeline costs microseconds.

\subsection*{Why six stages, not one}

The same reason \code{init} and \code{update} are separate. Each
stage has a clear contract; bugs and optimisations are local.
Testing, profiling, replacement all happen one stage at a time.

\subsection*{Why packed cells}

A cell is a struct: codepoint + style id + flags. Eight or
sixteen bytes. The canvas is a flat array. Diff is SIMD over
the array. There is no widget hierarchy; only cells.

\subsection*{Why fingerprints}

A 64-bit hash of the element tree. Identical trees produce
identical hashes; the renderer skips the entire pipeline. An
idle app burns microseconds, not milliseconds.

\section{Sentence 3: the type system}

\begin{quote}
The DSL puts structural shape into the type system (via NTTPs
and concepts) and leaves content to runtime, so the compiler
catches shape errors and the cost of correctness is paid at
build time.
\end{quote}

Part II's longest chapter (the DSL) explored this. The
guarantees chapter listed ten of them. The C++26 chapter
explained how the language makes it possible.

\subsection*{Why structural in types}

The shape of a box (has-border, has-padding, has-width) does
not change at runtime. Encoding it in the type lets the
compiler check it. A box without a border that asks for a
border colour fails to build, not silently misrender at
runtime.

\subsection*{Why content at runtime}

Text strings, list lengths, model values: these depend on
input. Forcing them into types would balloon the type system
without buying safety. The DSL leaves these at runtime
(\code{text(s)}, runtime children) and only enforces shape.

\subsection*{Why both}

The DSL is the compile-time spelling; builders and runtime
constructors are the runtime spelling. They produce the same
\code{Element}. You pick by what your data looks like.

\section{What you should be able to do}

After Part II, you should:

\begin{itemize}[leftmargin=*]
  \item Write the four pure functions for a small app from
    scratch.
  \item Read \code{maya/include/maya/app/app.hpp} and
    understand it.
  \item Decode a typical DSL compile error and fix it.
  \item Pick the right pattern for a UI shape (master-detail,
    wizard, palette, log stream).
  \item Test \code{update} without a terminal.
  \item Reason about performance without profiling first.
\end{itemize}

If any of these is unclear, the relevant chapter is the place
to revisit.

\section{What Part III adds}

Part III opens the rendering pipeline. You will see:

\begin{itemize}[leftmargin=*]
  \item How a tree becomes a canvas (Chapter 14).
  \item How SIMD makes diffing sub-microsecond (Chapter 15).
  \item How ANSI bytes are serialised efficiently (Chapter 16).
  \item How inline mode coexists with the terminal scrollback
    (Chapter 17).
\end{itemize}

The architecture stays the same; Part III is the
implementation tour.

\section{What Part IV adds}

Part IV is reactivity and widgets:

\begin{itemize}[leftmargin=*]
  \item The signal/computed/effect system in depth.
  \item Events, focus, input parsing under the hood.
  \item \code{Cmd} and \code{Sub} reference.
  \item The widget catalogue.
  \item Building, packaging, shipping.
\end{itemize}

By the end of Part IV you will know enough to ship a non-
trivial \maya{} app in production.

\section{Where to take the book next}

Three suggestions:

\begin{itemize}[leftmargin=*]
  \item Build the task manager from Chapter~\ref{ch:end-to-end}.
    Force every concept to earn its place.
  \item Read \file{maya/include/maya/app/app.hpp} from top to
    bottom. The comments in the source extend the book.
  \item Ship something. Even a small tool. The architecture
    only feels right after the first 1000-line app you
    write.
\end{itemize}

The book is a guide; the apps you write are the test of
whether the guide worked.

\section*{Workshop}

\begin{enumerate}[leftmargin=*]
  \item Restate the three sentences in your own words. Compare
    with someone who has read the book.
  \item For each sentence, find a counter-example in another
    framework you have used.
  \item Imagine extending \maya{} with a new \code{Cmd}
    alternative. List every place you would need to change.
\end{enumerate}

\begin{recap}
Three sentences. Four functions, two sum types, two effect
languages, one pipeline, one canvas, types for shape and
runtime for content. Everything else in Part II is a
specialisation, a worked example, or a pattern. Internalising
the three sentences is the goal; the rest follows.
\end{recap}

\chapter{Glossary and reference}
\label{ch:glossary}

\begin{aside}
Part II has introduced a lot of vocabulary. This chapter is a
reference: every term in one place, with a one-paragraph
explanation and a cross-reference back to its detailed chapter.
Use it as a refresher before tackling Parts III and IV.
\end{aside}

\section{Architecture terms}

\textbf{Element.} The unit of UI tree. A sum of
\code{BoxElement}, \code{TextElement}, \code{ElementList}.
Chapter~\ref{ch:elements}.

\textbf{Model.} The application state. A struct holding every
piece of state the app needs. Pure data; no pointers, no
shared state. Chapter~\ref{ch:elm}.

\textbf{Msg.} The alphabet of state transitions. A
\code{std::variant} where each alternative is one kind of
event the program responds to. Chapter~\ref{ch:elm}.

\textbf{Cmd<Msg>.} A description of side effects. A sum type
whose alternatives include \code{None}, \code{Quit},
\code{Task}, \code{After}, etc. Interpreted by the runtime.
Chapter~\ref{ch:elm}.

\textbf{Sub<Msg>.} A description of event sources the program
listens to. A sum type whose alternatives include \code{OnKey},
\code{OnMouse}, \code{Every}, etc. Chapter~\ref{ch:elm}.

\textbf{view.} A pure function from \code{Model} to
\code{Element}. Called every frame. Chapter~\ref{ch:elm}.

\textbf{update.} A pure function from
$(\text{Model}, \text{Msg})$ to $(\text{Model}, \text{Cmd})$.
Drives state transitions. Chapter~\ref{ch:elm}.

\textbf{subscribe.} A pure function from \code{Model} to
\code{Sub}. Declares which event sources are active.
Chapter~\ref{ch:elm}.

\textbf{init.} A pure function returning the initial
\code{Model} and \code{Cmd}. Called once on startup.
Chapter~\ref{ch:elm}.

\textbf{Program.} A concept aggregating the four functions
and their types. \code{maya::run} takes a \code{Program}.
Chapter~\ref{ch:concepts}.

\section{Rendering terms}

\textbf{Pipeline.} The six-stage transformation from
\code{Element} to ANSI bytes: view, layout, paint, diff,
serialize, write. Chapter~\ref{ch:pipeline}.

\textbf{Layout.} The phase that assigns each node a final
\code{Rect}. Uses a flexbox-style algorithm with integer cells.
Chapter~\ref{ch:layout}.

\textbf{Canvas.} A flat grid of packed cells, sized to the
terminal. Painted from the layout result. Chapter~\ref{ch:pipeline}.

\textbf{Packed cell.} A 64-bit struct holding a codepoint, a
style ID, and flags. Chapter~\ref{ch:pipeline}.

\textbf{Diff.} The phase that compares the current canvas to
the previous one and produces a per-row dirty range. SIMD-
accelerated. Chapter~\ref{ch:pipeline}.

\textbf{Serialize.} The phase that turns the diff into a
minimal sequence of ANSI bytes: cursor moves, SGR sequences,
cell writes. Chapter~\ref{ch:pipeline}.

\textbf{Style pool.} An intern table mapping unique
\code{Style} structs to small integer IDs. Used to keep packed
cells small. Chapter~\ref{ch:pipeline}.

\textbf{Cache fingerprint.} A 64-bit hash of the element
tree. Equal fingerprints mean the renderer can skip the entire
pipeline. Chapter~\ref{ch:pipeline}.

\textbf{Force redraw.} A \code{Cmd} alternative that discards
the cache and emits every cell on the next frame. Used after
terminal disturbances. Chapter~\ref{ch:elm}.

\section{DSL terms}

\textbf{NTTP.} Non-Type Template Parameter. A value (not a
type) that appears in a template parameter list. C++20
relaxed the rules so structural types (\code{Str},
\code{CTStyle}) can be NTTPs. Chapter~\ref{ch:dsl}.

\textbf{consteval.} A function that must run at compile time.
The compiler errors on runtime invocation. The DSL's pipe
operators are typically \code{consteval}. Chapter~\ref{ch:constexpr-deep}.

\textbf{constexpr.} A function that may run at compile time if
its arguments are constant expressions. May also run at runtime.
Chapter~\ref{ch:constexpr-deep}.

\textbf{Str.} A structural type wrapping a fixed-size character
array. Lets string literals appear as NTTPs.
Chapter~\ref{ch:dsl}.

\textbf{CTStyle.} A compile-time style struct holding boolean
flags and colour fields. Forms the payload of \code{StyTag}.
Chapter~\ref{ch:dsl}.

\textbf{StyTag.} A wrapper template indexed by a \code{CTStyle}
NTTP. Two \code{StyTag} values with the same \code{CTStyle}
have the same type. Chapter~\ref{ch:dsl}.

\textbf{TextNode.} A compile-time element wrapping a literal
string and an optional style. Pipe operators on it remain
compile-time. Chapter~\ref{ch:dsl}.

\textbf{RuntimeTextNode.} The runtime counterpart of
\code{TextNode}. Holds an owning string and style.
Chapter~\ref{ch:dsl}.

\textbf{BoxNode.} A compile-time element representing a box
with children. Indexed by a \code{BoxCfg} NTTP.
Chapter~\ref{ch:dsl}.

\textbf{Type-state machine.} The pattern where the type of a
DSL node encodes which operations are legal on it. The pipe
overloads have \code{requires} clauses that enforce the
transitions. Chapter~\ref{ch:dsl}.

\textbf{Concept.} A named compile-time predicate on a type. Used
to constrain template parameters. Chapter~\ref{ch:concepts}.

\section{Effect and event terms}

\textbf{Task.} A \code{Cmd} alternative that runs a callable on
a worker thread; its return value becomes a message.
Chapter~\ref{ch:elm}.

\textbf{IsolatedTask.} Like \code{Task} but on a dedicated
thread, not the pool. For long blocking work.
Chapter~\ref{ch:elm}.

\textbf{Batch.} A \code{Cmd} or \code{Sub} alternative that
groups multiple effects or sources. Chapter~\ref{ch:elm}.

\textbf{After.} A \code{Cmd} alternative that delivers a
message after a duration. Used for debounce, retry backoff,
toast expiry. Chapter~\ref{ch:elm}.

\textbf{Every.} A \code{Sub} alternative that delivers a
periodic message. Used for animations and polling.
Chapter~\ref{ch:elm}.

\textbf{AnimationFrame.} A \code{Sub} alternative that delivers
a message after each render frame. Used for frame-synchronised
animation. Chapter~\ref{ch:elm}.

\textbf{OnKey, OnMouse, OnResize, OnPaste.} \code{Sub}
alternatives for the various input events. Chapter~\ref{ch:elm}.

\textbf{Cmd::map.} A function that lifts a
\code{Cmd<ChildMsg>} into a \code{Cmd<ParentMsg>} via a
mapping function. The mechanism for composing components.
Chapter~\ref{ch:elm}.

\section{Style and theme terms}

\textbf{Style.} A struct holding fg, bg, attributes, and
optional hyperlink. The per-cell colouring directive.
Chapter~\ref{ch:style}.

\textbf{Attribute.} A bitset of typographic flags (Bold, Dim,
Italic, etc.). Combine with bitwise OR. Chapter~\ref{ch:style}.

\textbf{Color.} A sum type representing default, Ansi16,
Ansi256, or Rgb. Quantises if the terminal does not support
the requested form. Chapter~\ref{ch:elements}.

\textbf{Theme.} A palette plus metadata. Maps named roles to
colours. Swappable at runtime. Chapter~\ref{ch:style}.

\textbf{Palette.} An array of \code{Color} values indexed by
\code{Role}. The application-level colour vocabulary.
Chapter~\ref{ch:style}.

\textbf{Role.} A named slot in the palette: Foreground,
Background, Primary, Accent, Warning, etc. Decouples ``what
the colour means'' from ``what the colour is''.
Chapter~\ref{ch:style}.

\section{Layout terms}

\textbf{FlexStyle.} The layout parameters on a box: direction,
justify, align, grow, shrink, basis, wrap, gap, padding,
border. Chapter~\ref{ch:layout}.

\textbf{Direction.} Row or Column. Determines main axis.
Chapter~\ref{ch:layout}.

\textbf{Justify.} Where children sit along the main axis:
Start, End, Center, SpaceBetween, SpaceAround, SpaceEvenly.
Chapter~\ref{ch:layout}.

\textbf{Align.} Where children sit along the cross axis:
Start, End, Center, Stretch, Baseline. Chapter~\ref{ch:layout}.

\textbf{Grow.} A non-negative number; children with positive
grow share extra space proportionally. Chapter~\ref{ch:layout}.

\textbf{Shrink.} A non-negative number; children with positive
shrink absorb deficit proportionally. Chapter~\ref{ch:layout}.

\textbf{Basis.} The starting size before grow/shrink.
Chapter~\ref{ch:layout}.

\textbf{Wrap.} What to do when children overflow: NoWrap, Wrap,
WrapReverse. Chapter~\ref{ch:layout}.

\textbf{Gap.} Fixed cell spacing between children.
Chapter~\ref{ch:layout}.

\textbf{Padding.} Per-side integer cells, inside the box.
Chapter~\ref{ch:layout}.

\textbf{Border.} A decoration around a box; takes one cell per
side. Chapter~\ref{ch:elements}.

\section{Runtime terms}

\textbf{Context.} A bundle of runtime resources (platform
handles, config, caps, logger, stats) passed to subsystems.
Chapter~\ref{ch:runtime}.

\textbf{EventLoop.} The driver that waits for events and drives
the pipeline. Chapter~\ref{ch:runtime}.

\textbf{EventSource.} The platform abstraction for the
underlying OS event mechanism (epoll, kqueue,
WaitForMultipleObjects). Chapter~\ref{ch:platform}.

\textbf{SubscriptionManager.} The component that diffs the
declared subscription set against the active OS-level wires.
Chapter~\ref{ch:runtime}.

\textbf{Interpreter.} The component that walks
\code{Cmd<Msg>} variants and performs the side effects.
Chapter~\ref{ch:runtime}.

\textbf{Renderer.} The component that runs the six-stage
pipeline and writes to stdout. Chapter~\ref{ch:runtime}.

\textbf{Wake fd.} A self-pipe used by worker threads to wake
the main loop. Chapter~\ref{ch:platform}.

\section{Reactivity terms}

\textbf{Signal<T>.} A value that notifies subscribers when it
changes. Chapter~\ref{ch:reactivity-preview}.

\textbf{Computed<T>.} A derived signal whose value is a
function of other signals. Recomputes on input change.
Chapter~\ref{ch:reactivity-preview}.

\textbf{Effect.} A callable that re-runs whenever any signal it
reads changes. Chapter~\ref{ch:reactivity-preview}.

\section{Other terms}

\textbf{Expected<T, E>.} A value-typed sum of success and
error. \maya{}'s value-typed error handling type.
Chapter~\ref{ch:expected}.

\textbf{IoError.} A struct describing an I/O failure: code,
context, syscall name. The default error type for boundary
operations. Chapter~\ref{ch:expected}.

\textbf{Error boundary.} A wrapper that catches exceptions in
\code{build}/\code{view} and renders a fallback element.
Chapter~\ref{ch:errorboundary}.

\textbf{FocusManager.} A small helper that tracks focus IDs
and order. Stored in the model.
Chapter~\ref{ch:focus}.

\textbf{ScrollState.} A small struct (offset, viewport,
total) representing scroll position. Used for virtualisation.
Chapter~\ref{ch:scroll}.

\textbf{GraphemeWalker.} An iterator over Unicode grapheme
clusters. Cell-width-aware. Chapter~\ref{ch:text}.

\textbf{Widget.} A reusable UI component. May be a pure
renderer, a stateful component, or a reactive widget.
Chapter~\ref{ch:widget-arch}.

\textbf{Z-stack.} A composition mode where children render at
the same position; the last one wins on conflict. Used for
modals, overlays. Chapter~\ref{ch:end-to-end}.

\textbf{Hot reload.} Restarting an app while preserving its
model. Chapter~\ref{ch:lifecycle}.

\textbf{NO\_COLOR.} An environment variable that, when set,
asks programs to disable colour. \maya{} honours it via the
monochrome theme. Chapter~\ref{ch:a11y}.

\textbf{OSC 8.} The ANSI escape for clickable hyperlinks.
\maya{}'s \code{TextElement::hyperlink} emits it.
Chapter~\ref{ch:elements}.

\textbf{Kitty Keyboard Protocol.} An extended keyboard reporting
mode supported by modern terminals. Provides better key
disambiguation. Foundations chapter (Part I).

\textbf{Bracketed paste.} A terminal mode where pasted text is
wrapped in escape sequences, distinguishable from typed input.
Chapter~\ref{ch:elm}.

\textbf{Inline mode.} A rendering mode where \maya{} draws to
the existing terminal scrollback without taking over the screen.
Treated fully in Chapter~\ref{ch:inline}.

\section{Index of cross-references}

Each chapter has a \code{\textbackslash label}; cross-references
use \code{\textbackslash ref}. To find where a term is treated
in depth, look up the term here and follow the chapter link.

For online reading, the PDF is hyperlinked; clicking a chapter
reference jumps to the target. For print reading, the page
numbers in the TOC are the navigation tool.

\section*{Workshop}

\begin{enumerate}[leftmargin=*]
  \item Pick five terms from this glossary you cannot define
    from memory. Re-read the relevant chapter sections.
  \item Map each term to its position in the architecture
    diagram on page xxx (the architecture overview).
  \item For each \code{Cmd} alternative listed in the glossary,
    write the smallest possible app that uses it.
\end{enumerate}

\begin{recap}
The glossary collects every term Part II introduced. Use it as
a reference when reading later parts. Three sentences from the
previous chapter summarise the architecture; this chapter is
the vocabulary supporting those sentences.
\end{recap}


% ---- Part III: Rendering -----------------------------------------------
\part{Rendering}
\include{chapters/08-canvas}
\include{chapters/09-simd}
\include{chapters/10-ansi}
\include{chapters/11-inline}

% ---- Part IV: Reactivity, runtime, widgets -----------------------------
\part{Reactivity, runtime, and widgets}
\include{chapters/12-signals}
\include{chapters/13-events}
\include{chapters/14-cmdsub}
\include{chapters/15-widgets}
\include{chapters/16-build}

\backmatter
\include{chapters/99-colophon}

\end{document}
